% Options for packages loaded elsewhere
\PassOptionsToPackage{unicode}{hyperref}
\PassOptionsToPackage{hyphens}{url}
\documentclass[
  british,
  a4paper,
]{article}
\usepackage{xcolor}
\usepackage[margin=1in]{geometry}
\usepackage{amsmath,amssymb}
\setcounter{secnumdepth}{-\maxdimen} % remove section numbering
\usepackage{iftex}
\ifPDFTeX
  \usepackage[T1]{fontenc}
  \usepackage[utf8]{inputenc}
  \usepackage{textcomp} % provide euro and other symbols
\else % if luatex or xetex
  \usepackage{unicode-math} % this also loads fontspec
  \defaultfontfeatures{Scale=MatchLowercase}
  \defaultfontfeatures[\rmfamily]{Ligatures=TeX,Scale=1}
\fi
\usepackage{lmodern}
\ifPDFTeX\else
  % xetex/luatex font selection
  \setmainfont[]{Noto Serif}
  \setmonofont[]{Noto Sans Mono}
\fi
% Use upquote if available, for straight quotes in verbatim environments
\IfFileExists{upquote.sty}{\usepackage{upquote}}{}
\IfFileExists{microtype.sty}{% use microtype if available
  \usepackage[]{microtype}
  \UseMicrotypeSet[protrusion]{basicmath} % disable protrusion for tt fonts
}{}
\makeatletter
\@ifundefined{KOMAClassName}{% if non-KOMA class
  \IfFileExists{parskip.sty}{%
    \usepackage{parskip}
  }{% else
    \setlength{\parindent}{0pt}
    \setlength{\parskip}{6pt plus 2pt minus 1pt}}
}{% if KOMA class
  \KOMAoptions{parskip=half}}
\makeatother
\usepackage{longtable,booktabs,array}
\newcounter{none} % for unnumbered tables
\usepackage{calc} % for calculating minipage widths
% Correct order of tables after \paragraph or \subparagraph
\usepackage{etoolbox}
\makeatletter
\patchcmd\longtable{\par}{\if@noskipsec\mbox{}\fi\par}{}{}
\makeatother
% Allow footnotes in longtable head/foot
\IfFileExists{footnotehyper.sty}{\usepackage{footnotehyper}}{\usepackage{footnote}}
\makesavenoteenv{longtable}
% definitions for citeproc citations
\NewDocumentCommand\citeproctext{}{}
\NewDocumentCommand\citeproc{mm}{%
  \begingroup\def\citeproctext{#2}\cite{#1}\endgroup}
\makeatletter
 % allow citations to break across lines
 \let\@cite@ofmt\@firstofone
 % avoid brackets around text for \cite:
 \def\@biblabel#1{}
 \def\@cite#1#2{{#1\if@tempswa , #2\fi}}
\makeatother
\newlength{\cslhangindent}
\setlength{\cslhangindent}{1.5em}
\newlength{\csllabelwidth}
\setlength{\csllabelwidth}{3em}
\newenvironment{CSLReferences}[2] % #1 hanging-indent, #2 entry-spacing
 {\begin{list}{}{%
  \setlength{\itemindent}{0pt}
  \setlength{\leftmargin}{0pt}
  \setlength{\parsep}{0pt}
  % turn on hanging indent if param 1 is 1
  \ifodd #1
   \setlength{\leftmargin}{\cslhangindent}
   \setlength{\itemindent}{-1\cslhangindent}
  \fi
  % set entry spacing
  \setlength{\itemsep}{#2\baselineskip}}}
 {\end{list}}
\usepackage{calc}
\newcommand{\CSLBlock}[1]{\hfill\break\parbox[t]{\linewidth}{\strut\ignorespaces#1\strut}}
\newcommand{\CSLLeftMargin}[1]{\parbox[t]{\csllabelwidth}{\strut#1\strut}}
\newcommand{\CSLRightInline}[1]{\parbox[t]{\linewidth - \csllabelwidth}{\strut#1\strut}}
\newcommand{\CSLIndent}[1]{\hspace{\cslhangindent}#1}
\ifLuaTeX
\usepackage[bidi=basic,shorthands=off]{babel}
\else
\usepackage[bidi=default,shorthands=off]{babel}
\fi
\ifPDFTeX
\else
\babelfont{rm}[]{Noto Serif}
\fi
\ifLuaTeX
  \usepackage{selnolig} % disable illegal ligatures
\fi
\setlength{\emergencystretch}{3em} % prevent overfull lines
\providecommand{\tightlist}{%
  \setlength{\itemsep}{0pt}\setlength{\parskip}{0pt}}
\usepackage{array}
\usepackage{tabularx}
\usepackage{bookmark}
\IfFileExists{xurl.sty}{\usepackage{xurl}}{} % add URL line breaks if available
\urlstyle{same}
\hypersetup{
  pdflang={en-GB},
  hidelinks,
  pdfcreator={LaTeX via pandoc}}

\author{}
\date{}

\begin{document}

\section{Germanic Lexeme Reports: Lexical Derivation
Volume}\label{germanic-lexeme-reports-lexical-derivation-volume}

\emph{Alpha 01 compact-regular assembly scaffold. This document
assembles the current lexeme-report corpus in manifest order. Regular
entries use the compact book-prose layer requested for rollout;
non-regular entries retain the current model-entry prose.}

\subsection{Introduction}\label{introduction}

The lexical catalogue is organized by derivation class rather than as a
single undifferentiated list. This makes the interpretive burden of each
entry explicit. Regular derivations establish the baseline relation
between selected input and Old English target under the current cascade.
The variant, analogy, reconstructed-comparator, and exception classes
then show where that baseline is not sufficient on its own.

This lexical catalogue is a word-centered volume. It shows how
individual selected inputs develop to individual Old English targets and
how each entry is classified. A later sound-change volume or report
should remain separate and rule-centered, covering chronology,
interactions among rules, and broader system-level exception handling.

\subsection{Data and sources}\label{data-and-sources}

This alpha assembles the current model-entry corpus against the live
manifest, compact derivation-trace report, and project bibliography. The
lexical data layer is the current aligned Germanic dataset as
represented by the model-entry metadata and the current compact trace
source; comparative dictionaries, Old English dictionaries, and
historical grammars remain in the entry prose exactly as already cited
there.

Broad citations are carried forward honestly where the citation-layer
audit already judged them mechanically acceptable for assembly. This
alpha therefore tests book structure and technical integration rather
than attempting a final source-polish pass.

\subsection{Transducer and derivation
method}\label{transducer-and-derivation-method}

Regular entries use the compact book-prose overlay for this assembly
pass, while the remaining derivation classes still reproduce the current
model-entry prose. All entries retain a generated derivation summary and
a boxed derivation trace split into Earlier Germanic changes and Old
English changes.

\subsection{Derivation classes}\label{derivation-classes}

The lexical catalogue is ordered by the seven live
\texttt{DERIVATION\_CLASS} values in the manifest. Counts in this alpha
are taken directly from \texttt{manifest\_all\_by\_class.tsv}:

\begin{itemize}
\tightlist
\item
  \texttt{regular}: \textbf{70}
\item
  \texttt{attested\_variant}: \textbf{4}
\item
  \texttt{early\_analogy}: \textbf{35}
\item
  \texttt{late\_analogy}: \textbf{28}
\item
  \texttt{reconstructed\_oe}: \textbf{3}
\item
  \texttt{known\_unmodelled}: \textbf{2}
\item
  \texttt{unexplained\_unmodelled}: \textbf{5}
\end{itemize}

\clearpage

\subsection{Part I. Regular
derivations}\label{part-i.-regular-derivations}

Regular derivations are entries where the selected transducer input and
the Old English target stand in a straightforward relation under the
current cascade. These entries form the baseline against which the
analogy and exception classes are interpreted.

\subsubsection{adder --- OE nǣdre}\label{adder-oe-nux1e3dre}

Derivation: \emph{*nḗdrōn} \textgreater{} \emph{nǣdre} (regular).

\begingroup
\setlength{\fboxsep}{6pt}

\noindent\fbox{%
\begin{minipage}{0.97\linewidth}
\small
\begin{tabularx}{\linewidth}{@{}>{\raggedright\arraybackslash}p{0.460\linewidth}>{\centering\arraybackslash}X>{\raggedright\arraybackslash}p{0.460\linewidth}@{}}
\begin{minipage}[t]{\linewidth}
\centering\textbf{Earlier Germanic changes}\par
\vspace{0.35em}
\raggedright
\centering\textbf{West Germanic}\par
\raggedright
\vspace{0.2em}
\raggedright [no change]\par
\vspace{0.6em}
\centering\textbf{Northwest Germanic}\par
\raggedright
\vspace{0.2em}
\begin{tabular}{@{}>{\raggedright\arraybackslash}p{0.74\linewidth}@{\hspace{0.55em}}>{\raggedright\arraybackslash}p{0.16\linewidth}@{\hspace{0.25em}}}
NWGmc N Stem N Loss & \emph{*nḗdrǭ} \\
\mbox{NWGmc Long E Lowering} & \emph{*nǣdrǭ} \\
\end{tabular}
\end{minipage}
&
&
\begin{minipage}[t]{\linewidth}
\centering\textbf{Old English changes}\par
\vspace{0.35em}
\raggedright
\begin{tabular}{@{}>{\raggedright\arraybackslash}p{0.74\linewidth}@{\hspace{0.55em}}>{\raggedright\arraybackslash}p{0.16\linewidth}@{\hspace{0.25em}}}
OE Unstressed Long Vowel Shortening & \emph{*nǣdræ} \\
OE Unstressed AE Merger & \emph{*nǣdre} \\
\end{tabular}
\end{minipage}
\\
\end{tabularx}
\end{minipage}%
} \endgroup

Kroonen distinguishes the masculine snake word \emph{*nadra-} from a
feminine ablauting formation \emph{*nēdrōn-}, and gives Old English
\emph{nǣdre}, \emph{næddre} under the latter
(\citeproc{ref-Kroonen2013}{Kroonen 2013, 426}). Orel likewise points
from the masculine entry to a feminine \emph{*nēdrōn} \textasciitilde{}
\emph{*nadrōn} type (\citeproc{ref-Orel2003}{Orel 2003, 325}). The Old
English word is securely represented by \emph{nǣdre}, with \emph{næddre}
as a secondary variant. From \emph{*nḗdrōn}, the stressed long mid vowel
develops to Old English \emph{ǣ}, and the weak feminine ending remains
as final \emph{-e}, giving \emph{nǣdre}.

\subsubsection{bake --- OE bacan}\label{bake-oe-bacan}

Derivation: \emph{*bákaną} \textgreater{} \emph{bacan} (regular).

\begingroup
\setlength{\fboxsep}{6pt}

\noindent\fbox{%
\begin{minipage}{0.97\linewidth}
\small
\begin{tabularx}{\linewidth}{@{}>{\raggedright\arraybackslash}p{0.300\linewidth}>{\centering\arraybackslash}X>{\raggedright\arraybackslash}p{0.540\linewidth}@{}}
\begin{minipage}[t]{\linewidth}
\centering\textbf{Earlier Germanic changes}\par
\vspace{0.35em}
\raggedright
\centering\textbf{West Germanic}\par
\raggedright
\vspace{0.2em}
\raggedright [no change]\par
\vspace{0.6em}
\centering\textbf{Northwest Germanic}\par
\raggedright
\vspace{0.2em}
\raggedright [no change]\par
\end{minipage}
&
&
\begin{minipage}[t]{\linewidth}
\centering\textbf{Old English changes}\par
\vspace{0.35em}
\raggedright
\begin{tabular}{@{}>{\raggedright\arraybackslash}p{0.74\linewidth}@{\hspace{0.55em}}>{\raggedright\arraybackslash}p{0.16\linewidth}@{\hspace{0.25em}}}
Anglo Frisian Brightening & \emph{*bækaną} \\
OE A Restoration & \emph{*bakaną} \\
OE Heavy Syllable Nasal Apocope & \emph{*bakan} \\
OE Secondary Nasalization & \emph{*bakąn} \\
OE Weak Tail Reduction & \emph{*bakan} \\
\end{tabular}
\end{minipage}
\\
\end{tabularx}
\end{minipage}%
} \endgroup

Orel reconstructs \emph{*bakanan}, and Campbell cites \emph{bacan} as a
standard example of A-restoration (\citeproc{ref-Orel2003}{Orel 2003,
33}; \citeproc{ref-Campbell1959}{Campbell 1959, 61};
\citeproc{ref-RingeTaylor2014}{Ringe and Taylor 2014, 270}).
Bosworth-Toller and Clark Hall record \emph{bacan} as the Old English
infinitive (\citeproc{ref-BosworthToller1898}{Bosworth and Toller 1898,
72}; \citeproc{ref-ClarkHall1960}{Clark Hall 1960, 19}). The selected
input \emph{*bákaną} represents the same weak verb in the accent
notation used here. From \emph{*bákaną}, Anglo-Frisian brightening first
gives \emph{*bækaną}. A-restoration before single \emph{k} returns the
stem vowel to \emph{a}, and later weak-tail reductions yield
\emph{bacan}.

\subsubsection{beech --- OE bōc}\label{beech-oe-bux14dc}

Derivation: \emph{*bōkō} \textgreater{} \emph{bōc} (regular).

\begingroup
\setlength{\fboxsep}{6pt}

\noindent\fbox{%
\begin{minipage}{0.97\linewidth}
\small
\begin{tabularx}{\linewidth}{@{}>{\raggedright\arraybackslash}p{0.460\linewidth}>{\centering\arraybackslash}X>{\raggedright\arraybackslash}p{0.460\linewidth}@{}}
\begin{minipage}[t]{\linewidth}
\centering\textbf{Earlier Germanic changes}\par
\vspace{0.35em}
\raggedright
\centering\textbf{West Germanic}\par
\raggedright
\vspace{0.2em}
\raggedright [no change]\par
\vspace{0.6em}
\centering\textbf{Northwest Germanic}\par
\raggedright
\vspace{0.2em}
\begin{tabular}{@{}>{\raggedright\arraybackslash}p{0.74\linewidth}@{\hspace{0.55em}}>{\raggedright\arraybackslash}p{0.16\linewidth}@{\hspace{0.25em}}}
\mbox{NWGmc Final Long O Raising} & \emph{*bōku} \\
\end{tabular}
\end{minipage}
&
&
\begin{minipage}[t]{\linewidth}
\centering\textbf{Old English changes}\par
\vspace{0.35em}
\raggedright
\begin{tabular}{@{}>{\raggedright\arraybackslash}p{0.74\linewidth}@{\hspace{0.55em}}>{\raggedright\arraybackslash}p{0.16\linewidth}@{\hspace{0.25em}}}
\mbox{OE High Vowel Apocope} & \emph{*bōk} \\
\end{tabular}
\end{minipage}
\\
\end{tabularx}
\end{minipage}%
} \endgroup

Kroonen gives \emph{*bōk(j)ō-} with Old English \emph{boc} and oblique
\emph{bēce} (\citeproc{ref-Kroonen2013}{Kroonen 2013, 81}). The selected
input \emph{*bōkō} is the nominative-singular shape of that family. The
selected form here is nominative \emph{bōc}; \emph{bēċe} belongs to the
same paradigm but is not the citation form used for this derivation.
From \emph{*bōkō}, final long \emph{ō} raises in Northwest Germanic, and
high-vowel apocope yields \emph{bōc}.

\subsubsection{begin --- OE beġinnan}\label{begin-oe-beux121innan}

Derivation: \emph{*bigínnaną} \textgreater{} \emph{beġinnan} (regular).

\begingroup
\setlength{\fboxsep}{6pt}

\noindent\fbox{%
\begin{minipage}{0.97\linewidth}
\small
\begin{tabularx}{\linewidth}{@{}>{\raggedright\arraybackslash}p{0.300\linewidth}>{\centering\arraybackslash}X>{\raggedright\arraybackslash}p{0.540\linewidth}@{}}
\begin{minipage}[t]{\linewidth}
\centering\textbf{Earlier Germanic changes}\par
\vspace{0.35em}
\raggedright
\centering\textbf{West Germanic}\par
\raggedright
\vspace{0.2em}
\raggedright [no change]\par
\vspace{0.6em}
\centering\textbf{Northwest Germanic}\par
\raggedright
\vspace{0.2em}
\raggedright [no change]\par
\end{minipage}
&
&
\begin{minipage}[t]{\linewidth}
\centering\textbf{Old English changes}\par
\vspace{0.35em}
\raggedright
\begin{tabular}{@{}>{\raggedright\arraybackslash}p{0.68\linewidth}@{\hspace{0.55em}}>{\raggedright\arraybackslash}p{0.22\linewidth}@{\hspace{0.25em}}}
OE Heavy Syllable Nasal Apocope & \emph{*bigínnan} \\
OE Secondary Nasalization & \emph{*bigínnąn} \\
OE Velar Palatalization & \emph{*biʤínnąn} \\
OE Prefix I Reduction & \emph{*bĕʤínnąn} \\
OE Weak Tail Reduction & \emph{*bĕʤínnan} \\
\end{tabular}
\end{minipage}
\\
\end{tabularx}
\end{minipage}%
} \endgroup

The verb is modeled here as inherited \emph{*bigínnaną}. Ringe and
Taylor state that intervocalic \emph{*g} is palatalized between front
vowels in Old English (\citeproc{ref-RingeTaylor2014}{Ringe and Taylor
2014, 218}), and Campbell lists \emph{ginnan} among familiar examples of
palatal \emph{g} in this verb family
(\citeproc{ref-Campbell1959}{Campbell 1959, 174}). Bosworth-Toller
lemmatizes the verb as \emph{be-ginnan}
(\citeproc{ref-BosworthToller1898}{Bosworth and Toller 1898, 84}). Clark
Hall likewise gives \emph{beginnan} (\citeproc{ref-ClarkHall1960}{Clark
Hall 1960, 34}). From \emph{*bigínnaną}, heavy-syllable nasal apocope
yields \emph{*bigínnan}. Ringe and Taylor explicitly cite \emph{bi-
\textgreater{} be-} as an Old English unstressed-prefix development
(\citeproc{ref-RingeTaylor2014}{Ringe and Taylor 2014, 350}).

\subsubsection{bier --- OE bǣr}\label{bier-oe-bux1e3r}

Derivation: \emph{*bḗrō} \textgreater{} \emph{bǣr} (regular).

\begingroup
\setlength{\fboxsep}{6pt}

\noindent\fbox{%
\begin{minipage}{0.97\linewidth}
\small
\begin{tabularx}{\linewidth}{@{}>{\raggedright\arraybackslash}p{0.460\linewidth}>{\centering\arraybackslash}X>{\raggedright\arraybackslash}p{0.460\linewidth}@{}}
\begin{minipage}[t]{\linewidth}
\centering\textbf{Earlier Germanic changes}\par
\vspace{0.35em}
\raggedright
\centering\textbf{West Germanic}\par
\raggedright
\vspace{0.2em}
\raggedright [no change]\par
\vspace{0.6em}
\centering\textbf{Northwest Germanic}\par
\raggedright
\vspace{0.2em}
\begin{tabular}{@{}>{\raggedright\arraybackslash}p{0.74\linewidth}@{\hspace{0.55em}}>{\raggedright\arraybackslash}p{0.16\linewidth}@{\hspace{0.25em}}}
\mbox{NWGmc Final Long O Raising} & \emph{*bḗru} \\
\mbox{NWGmc Long E Lowering} & \emph{*bǣru} \\
\end{tabular}
\end{minipage}
&
&
\begin{minipage}[t]{\linewidth}
\centering\textbf{Old English changes}\par
\vspace{0.35em}
\raggedright
\begin{tabular}{@{}>{\raggedright\arraybackslash}p{0.74\linewidth}@{\hspace{0.55em}}>{\raggedright\arraybackslash}p{0.16\linewidth}@{\hspace{0.25em}}}
\mbox{OE High Vowel Apocope} & \emph{*bǣr} \\
\end{tabular}
\end{minipage}
\\
\end{tabularx}
\end{minipage}%
} \endgroup

Kroonen reconstructs \emph{*bērō-} and cites Old English \emph{bar},
\emph{bær} among the reflexes (\citeproc{ref-Kroonen2013}{Kroonen 2013,
717}; \citeproc{ref-ClarkHall1960}{Clark Hall 1960, 32};
\citeproc{ref-BosworthToller1898}{Bosworth and Toller 1898, 73}). The
selected input \emph{*bḗrō} is the same noun in the accent notation used
here. Clark Hall and Bosworth-Toller give \emph{bær}, and Kroonen also
records \emph{bar}. The normalized form \emph{bǣr} marks the long vowel
explicitly. From \emph{*bḗrō}, final long \emph{ō} raises, long \emph{ē}
lowers to \emph{ǣ}, and apocope yields \emph{bǣr}.

\subsubsection{birth --- OE byrd}\label{birth-oe-byrd}

Derivation: \emph{*búrdiz} \textgreater{} \emph{byrd} (regular).

\begingroup
\setlength{\fboxsep}{6pt}

\noindent\fbox{%
\begin{minipage}{0.97\linewidth}
\small
\begin{tabularx}{\linewidth}{@{}>{\raggedright\arraybackslash}p{0.460\linewidth}>{\centering\arraybackslash}X>{\raggedright\arraybackslash}p{0.460\linewidth}@{}}
\begin{minipage}[t]{\linewidth}
\centering\textbf{Earlier Germanic changes}\par
\vspace{0.35em}
\raggedright
\centering\textbf{West Germanic}\par
\raggedright
\vspace{0.2em}
\raggedright [no change]\par
\vspace{0.6em}
\centering\textbf{Northwest Germanic}\par
\raggedright
\vspace{0.2em}
\begin{tabular}{@{}>{\raggedright\arraybackslash}p{0.74\linewidth}@{\hspace{0.55em}}>{\raggedright\arraybackslash}p{0.16\linewidth}@{\hspace{0.25em}}}
\mbox{PGmc Final Z Deletion} & \emph{*búrdi} \\
\end{tabular}
\end{minipage}
&
&
\begin{minipage}[t]{\linewidth}
\centering\textbf{Old English changes}\par
\vspace{0.35em}
\raggedright
\begin{tabular}{@{}>{\raggedright\arraybackslash}p{0.74\linewidth}@{\hspace{0.55em}}>{\raggedright\arraybackslash}p{0.16\linewidth}@{\hspace{0.25em}}}
OE I Umlaut & \emph{*byrdi} \\
\mbox{OE High Vowel Apocope} & \emph{*byrd} \\
\end{tabular}
\end{minipage}
\\
\end{tabularx}
\end{minipage}%
} \endgroup

Kroonen cites the noun under stem-level \emph{*burdi-} and gives Old
English \emph{(ge-)byrd} among the reflexes
(\citeproc{ref-Kroonen2013}{Kroonen 2013, 122}). The selected input
\emph{*búrdiz} is the nominative-style form that stands behind that stem
label. Clark Hall and Bosworth-Toller both attest simplex \emph{byrd} as
an Old English noun meaning `birth' (\citeproc{ref-ClarkHall1960}{Clark
Hall 1960, 70}; \citeproc{ref-BosworthToller1898}{Bosworth and Toller
1898, 125}). From \emph{*búrdiz}, loss of final \emph{z} gives
\emph{*búrdi}.

Form note. The relevant comparator here is the simplex noun \emph{byrd}.
The prefixed forms remain related attested material within the same
lexical family, and Hogg's discussion of deverbal feminines provides the
broader derivational setting (\citeproc{ref-Hogg1992}{Hogg 1992}).

\subsubsection{bone --- OE bān}\label{bone-oe-bux101n}

Derivation: \emph{*báiną} \textgreater{} \emph{bān} (regular).

\begingroup
\setlength{\fboxsep}{6pt}

\noindent\fbox{%
\begin{minipage}{0.97\linewidth}
\small
\begin{tabularx}{\linewidth}{@{}>{\raggedright\arraybackslash}p{0.300\linewidth}>{\centering\arraybackslash}X>{\raggedright\arraybackslash}p{0.540\linewidth}@{}}
\begin{minipage}[t]{\linewidth}
\centering\textbf{Earlier Germanic changes}\par
\vspace{0.35em}
\raggedright
\centering\textbf{West Germanic}\par
\raggedright
\vspace{0.2em}
\begin{tabular}{@{}>{\raggedright\arraybackslash}p{0.70\linewidth}@{\hspace{0.55em}}>{\raggedright\arraybackslash}p{0.16\linewidth}@{\hspace{0.25em}}}
PWGmc Ai Monophthongization & \emph{*bāną} \\
\end{tabular}
\vspace{0.6em}
\centering\textbf{Northwest Germanic}\par
\raggedright
\vspace{0.2em}
\raggedright [no change]\par
\end{minipage}
&
&
\begin{minipage}[t]{\linewidth}
\centering\textbf{Old English changes}\par
\vspace{0.35em}
\raggedright
\begin{tabular}{@{}>{\raggedright\arraybackslash}p{0.74\linewidth}@{\hspace{0.55em}}>{\raggedright\arraybackslash}p{0.16\linewidth}@{\hspace{0.25em}}}
OE Heavy Syllable Nasal Apocope & \emph{*bān} \\
\end{tabular}
\end{minipage}
\\
\end{tabularx}
\end{minipage}%
} \endgroup

Kroonen cites the noun as \emph{*baina-}, and Orel gives the same lexeme
under \emph{*bainan} (\citeproc{ref-Kroonen2013}{Kroonen 2013, 86};
\citeproc{ref-Orel2003}{Orel 2003, 71}). Both are comparative headword
conventions for the same neuter noun whose Old English reflex is
\emph{bān}. Clark Hall and Bosworth-Toller record \emph{bān} as the
ordinary Old English noun (\citeproc{ref-ClarkHall1960}{Clark Hall 1960,
44}; \citeproc{ref-BosworthToller1898}{Bosworth and Toller 1898}). West
Germanic monophthongization turns stressed \emph{*ai} into \emph{ā},
giving \emph{*bāną}; heavy-syllable nasal apocope then yields
\emph{bān}.

Source note. The comparative headwords \emph{*baina-} and \emph{*bainan}
provide lexeme background. The relevant comparison form here is the
nominative-accusative singular \emph{*báiną}.

\subsubsection{both --- OE bū}\label{both-oe-bux16b}

Derivation: \emph{*bō} \textgreater{} \emph{bū} (regular).

\begingroup
\setlength{\fboxsep}{6pt}

\noindent\fbox{%
\begin{minipage}{0.97\linewidth}
\small
\begin{tabularx}{\linewidth}{@{}>{\raggedright\arraybackslash}p{0.540\linewidth}>{\centering\arraybackslash}X>{\raggedright\arraybackslash}p{0.300\linewidth}@{}}
\begin{minipage}[t]{\linewidth}
\centering\textbf{Earlier Germanic changes}\par
\vspace{0.35em}
\raggedright
\centering\textbf{West Germanic}\par
\raggedright
\vspace{0.2em}
\raggedright [no change]\par
\vspace{0.6em}
\centering\textbf{Northwest Germanic}\par
\raggedright
\vspace{0.2em}
\begin{tabular}{@{}>{\raggedright\arraybackslash}p{0.74\linewidth}@{\hspace{0.55em}}>{\raggedright\arraybackslash}p{0.16\linewidth}@{\hspace{0.25em}}}
NWGmc Stressed Monosyllable O Raising & \emph{*bū} \\
\end{tabular}
\end{minipage}
&
&
\begin{minipage}[t]{\linewidth}
\centering\textbf{Old English changes}\par
\vspace{0.35em}
\raggedright
\raggedright [no change]\par
\end{minipage}
\\
\end{tabularx}
\end{minipage}%
} \endgroup

Kroonen treats the numeral under \emph{*ba-} and includes the inherited
\emph{*bō} beside the extended dual forms
(\citeproc{ref-Kroonen2013}{Kroonen 2013, 47}). Brunner, Campbell, and
Fulk preserve the Old English dual paradigm with masculine \emph{bēġen},
feminine \emph{bā}, and neuter \emph{bū}
(\citeproc{ref-SieversBrunner1965}{Sievers 1965, 324};
\citeproc{ref-Campbell1959}{Campbell 1959, 295};
\citeproc{ref-Fulk2018}{Fulk 2018, 223}). The selected form is neuter
\emph{bū}. Campbell and Brunner treat stressed monosyllabic \emph{ō
\textgreater{} ū} as a regular development in forms of this type, and
\emph{bū} belongs to that set (\citeproc{ref-Campbell1959}{Campbell
1959, 122}; \citeproc{ref-SieversBrunner1965}{Sievers 1965, 69}).

Comparison note. Orel cites an older \emph{*bō-j-} explanation for
\emph{bēġen}, while Fulk reports Seebold's preference for \emph{*bō-þ-}
(\citeproc{ref-Orel2003}{Orel 2003, 45}; \citeproc{ref-Fulk2018}{Fulk
2018, 223}). Those analyses concern \emph{bēġen} and the related
extended forms. The derivation here is the neuter \emph{*bō
\textgreater{} bū}.

\subsubsection{bow --- OE bīeġan}\label{bow-oe-bux12beux121an}

Derivation: \emph{*báugijaną} \textgreater{} \emph{bīeġan} (regular).

\begingroup
\setlength{\fboxsep}{6pt}

\noindent\fbox{%
\begin{minipage}{0.97\linewidth}
\footnotesize
\begin{tabularx}{\linewidth}{@{}>{\raggedright\arraybackslash}p{0.280\linewidth}>{\centering\arraybackslash}X>{\raggedright\arraybackslash}p{0.560\linewidth}@{}}
\begin{minipage}[t]{\linewidth}
\centering\textbf{Earlier Germanic changes}\par
\vspace{0.35em}
\raggedright
\centering\textbf{West Germanic}\par
\raggedright
\vspace{0.2em}
\raggedright [no change]\par
\vspace{0.6em}
\centering\textbf{Northwest Germanic}\par
\raggedright
\vspace{0.2em}
\raggedright [no change]\par
\end{minipage}
&
&
\begin{minipage}[t]{\linewidth}
\centering\textbf{Old English changes}\par
\vspace{0.35em}
\raggedright
\begin{tabular}{@{}>{\raggedright\arraybackslash}p{0.62\linewidth}@{\hspace{0.55em}}>{\raggedright\arraybackslash}p{0.28\linewidth}@{\hspace{0.25em}}}
OE Au Fronting & \emph{*báeugijaną} \\
OE Diphthong Leveling & \emph{*bēagijaną} \\
OE Heavy Syllable Nasal Apocope & \emph{*bēagijan} \\
OE Secondary Nasalization & \emph{*bēagijąn} \\
Sievers Law Syncope & \emph{*bēagjąn} \\
OE Velar Palatalization & \emph{*bēaʤjąn} \\
OE I Umlaut & \emph{*bīeʤjąn} \\
OE Weak Tail Reduction & \emph{*bīeʤjan} \\
OE J Loss After Heavy & \emph{*bīeʤan} \\
\end{tabular}
\end{minipage}
\\
\end{tabularx}
\end{minipage}%
} \endgroup

Kroonen reconstructs the weak causative \emph{*baugjan-}, and Ringe and
Taylor give the fuller development to West Saxon \emph{biegan}
(\citeproc{ref-Kroonen2013}{Kroonen 2013, 75};
\citeproc{ref-RingeTaylor2014}{Ringe and Taylor 2014, 94}). The verb
belongs to the weak causative member of the bend-family. Clark Hall and
Bosworth-Toller record the Old English weak verb
(\citeproc{ref-ClarkHall1960}{Clark Hall 1960, 34};
\citeproc{ref-BosworthToller1898}{Bosworth and Toller 1898, 102}). The
spelling \emph{bīeġan} used here normalizes that attested form. From
pre-Old-English \emph{*bēagjan}, palatalization of \emph{gj} and
i-mutation yield the West Saxon infinitive.

\subsubsection{breeches --- OE brēċ}\label{breeches-oe-brux113ux10b}

Derivation: \emph{*brōkiz} \textgreater{} \emph{brēċ} (regular).

\begingroup
\setlength{\fboxsep}{6pt}

\noindent\fbox{%
\begin{minipage}{0.97\linewidth}
\small
\begin{tabularx}{\linewidth}{@{}>{\raggedright\arraybackslash}p{0.460\linewidth}>{\centering\arraybackslash}X>{\raggedright\arraybackslash}p{0.460\linewidth}@{}}
\begin{minipage}[t]{\linewidth}
\centering\textbf{Earlier Germanic changes}\par
\vspace{0.35em}
\raggedright
\centering\textbf{West Germanic}\par
\raggedright
\vspace{0.2em}
\raggedright [no change]\par
\vspace{0.6em}
\centering\textbf{Northwest Germanic}\par
\raggedright
\vspace{0.2em}
\begin{tabular}{@{}>{\raggedright\arraybackslash}p{0.74\linewidth}@{\hspace{0.55em}}>{\raggedright\arraybackslash}p{0.16\linewidth}@{\hspace{0.25em}}}
\mbox{PGmc Final Z Deletion} & \emph{*brōki} \\
\end{tabular}
\end{minipage}
&
&
\begin{minipage}[t]{\linewidth}
\centering\textbf{Old English changes}\par
\vspace{0.35em}
\raggedright
\begin{tabular}{@{}>{\raggedright\arraybackslash}p{0.74\linewidth}@{\hspace{0.55em}}>{\raggedright\arraybackslash}p{0.16\linewidth}@{\hspace{0.25em}}}
OE Velar Palatalization & \emph{*brōʧi} \\
OE I Umlaut & \emph{*brēʧi} \\
\mbox{OE High Vowel Apocope} & \emph{*brēʧ} \\
\end{tabular}
\end{minipage}
\\
\end{tabularx}
\end{minipage}%
} \endgroup

Kroonen cites the noun under \emph{*brōk-}, with Old English \emph{brōc}
and plural `breeches' among its reflexes
(\citeproc{ref-Kroonen2013}{Kroonen 2013}). Ringe and Taylor give the
plural development directly as PNWGmc \emph{*brokiz} \textgreater{}
\emph{*breeci} \textgreater{} OE \emph{bréc}
(\citeproc{ref-RingeTaylor2014}{Ringe and Taylor 2014, 223}). Bright
notes \emph{brōc} with plural \emph{brēc}, and Clark Hall gives
\emph{bréc} fp. breeches while also listing \emph{broc} as a feminine
noun probably represented chiefly in the plural
(\citeproc{ref-BrightCassidyRingler1971}{Bright 1971, 38};
\citeproc{ref-ClarkHall1960}{Clark Hall 1960, 64}). After loss of final
\emph{-z}, the stem ends in \emph{-ki}, so the velar palatalizes and
\emph{ō} undergoes i-umlaut to \emph{ē}; final high-vowel apocope then
yields \emph{brēċ} (\citeproc{ref-RingeTaylor2014}{Ringe and Taylor
2014, 223}).

\subsubsection{calf --- OE ċealf}\label{calf-oe-ux10bealf}

Derivation: \emph{*kálbaz} \textgreater{} \emph{ċealf} (regular).

\begingroup
\setlength{\fboxsep}{6pt}

\noindent\fbox{%
\begin{minipage}{0.97\linewidth}
\small
\begin{tabularx}{\linewidth}{@{}>{\raggedright\arraybackslash}p{0.300\linewidth}>{\centering\arraybackslash}X>{\raggedright\arraybackslash}p{0.540\linewidth}@{}}
\begin{minipage}[t]{\linewidth}
\centering\textbf{Earlier Germanic changes}\par
\vspace{0.35em}
\raggedright
\centering\textbf{West Germanic}\par
\raggedright
\vspace{0.2em}
\raggedright [no change]\par
\vspace{0.6em}
\centering\textbf{Northwest Germanic}\par
\raggedright
\vspace{0.2em}
\begin{tabular}{@{}>{\raggedright\arraybackslash}p{0.74\linewidth}@{\hspace{0.55em}}>{\raggedright\arraybackslash}p{0.16\linewidth}@{\hspace{0.25em}}}
\mbox{PGmc Final Z Deletion} & \emph{*kálba} \\
\end{tabular}
\end{minipage}
&
&
\begin{minipage}[t]{\linewidth}
\centering\textbf{Old English changes}\par
\vspace{0.35em}
\raggedright
\begin{tabular}{@{}>{\raggedright\arraybackslash}p{0.70\linewidth}@{\hspace{0.55em}}>{\raggedright\arraybackslash}p{0.16\linewidth}@{\hspace{0.25em}}}
PWGmc Final Bare A Loss & \emph{*kálb} \\
Anglo Frisian Brightening & \emph{*kælb} \\
OE Breaking & \emph{*kealb} \\
PGmc B Allophony & \emph{*kealβ} \\
OE Velar Palatalization & \emph{*ʧealβ} \\
\end{tabular}
\end{minipage}
\\
\end{tabularx}
\end{minipage}%
} \endgroup

Kroonen treats the noun under \emph{*kalbiz-} and notes an older s-stem
\emph{*kalbaz}, pl. \emph{*kalbizō}, while Orel cites \emph{*kalbaz} as
the citation form and Ringe and Taylor derive West Saxon \emph{Cealf}
from \emph{*kalbaz}, \emph{*kalbiz-} (\citeproc{ref-Kroonen2013}{Kroonen
2013, 318}; \citeproc{ref-Orel2003}{Orel 2003, 248};
\citeproc{ref-RingeTaylor2014}{Ringe and Taylor 2014, 220}). Clark Hall
gives \emph{cealf}, and Bosworth-Toller likewise records \emph{Caelf /
Cealf} among the ordinary Old English spellings
(\citeproc{ref-ClarkHall1960}{Clark Hall 1960, 73};
\citeproc{ref-BosworthToller1898}{Bosworth and Toller 1898, 131}). After
loss of final \emph{-z} and bare \emph{-a}, Anglo-Frisian brightening
gives \emph{*kælb}, and breaking before \emph{l} plus consonant yields
\emph{*kealb}.

\subsubsection{corn --- OE corn}\label{corn-oe-corn}

Derivation: \emph{*kúrną} \textgreater{} \emph{corn} (regular).

\begingroup
\setlength{\fboxsep}{6pt}

\noindent\fbox{%
\begin{minipage}{0.97\linewidth}
\small
\begin{tabularx}{\linewidth}{@{}>{\raggedright\arraybackslash}p{0.460\linewidth}>{\centering\arraybackslash}X>{\raggedright\arraybackslash}p{0.460\linewidth}@{}}
\begin{minipage}[t]{\linewidth}
\centering\textbf{Earlier Germanic changes}\par
\vspace{0.35em}
\raggedright
\centering\textbf{West Germanic}\par
\raggedright
\vspace{0.2em}
\raggedright [no change]\par
\vspace{0.6em}
\centering\textbf{Northwest Germanic}\par
\raggedright
\vspace{0.2em}
\begin{tabular}{@{}>{\raggedright\arraybackslash}p{0.74\linewidth}@{\hspace{0.55em}}>{\raggedright\arraybackslash}p{0.16\linewidth}@{\hspace{0.25em}}}
\mbox{NWGmc U Lowering} & \emph{*kórną} \\
\end{tabular}
\end{minipage}
&
&
\begin{minipage}[t]{\linewidth}
\centering\textbf{Old English changes}\par
\vspace{0.35em}
\raggedright
\begin{tabular}{@{}>{\raggedright\arraybackslash}p{0.74\linewidth}@{\hspace{0.55em}}>{\raggedright\arraybackslash}p{0.16\linewidth}@{\hspace{0.25em}}}
OE Heavy Syllable Nasal Apocope & \emph{*kórn} \\
\end{tabular}
\end{minipage}
\\
\end{tabularx}
\end{minipage}%
} \endgroup

Kroonen cites the noun as \emph{*kurna-}, and Orel gives the citation
form \emph{*kurnan}, both with Old English \emph{corn} among the
reflexes (\citeproc{ref-Kroonen2013}{Kroonen 2013, 352};
\citeproc{ref-Orel2003}{Orel 2003, 264}). The singular form
\emph{*kúrną} is the nominative-accusative singular appropriate to the
citation noun. Clark Hall gives \emph{corn n.~`corn,' grain}, Bright's
glossary lists \emph{corn, n.} with genitive singular \emph{cornes}, and
Bosworth-Toller treats \emph{corn} as an ordinary noun headword
(\citeproc{ref-ClarkHall1960}{Clark Hall 1960, 79};
\citeproc{ref-BrightCassidyRingler1971}{Bright 1971, 347};
\citeproc{ref-BosworthToller1898}{Bosworth and Toller 1898, 144}). With
northwest Germanic lowering, \emph{*kúrną} becomes \emph{*kórną}, and
later loss of final nasal after a heavy syllable yields \emph{*kórn},
whence \emph{corn}.

\subsubsection{deed --- OE dǣd}\label{deed-oe-dux1e3d}

Derivation: \emph{*dḗdiz} \textgreater{} \emph{dǣd} (regular).

\begingroup
\setlength{\fboxsep}{6pt}

\noindent\fbox{%
\begin{minipage}{0.97\linewidth}
\small
\begin{tabularx}{\linewidth}{@{}>{\raggedright\arraybackslash}p{0.460\linewidth}>{\centering\arraybackslash}X>{\raggedright\arraybackslash}p{0.460\linewidth}@{}}
\begin{minipage}[t]{\linewidth}
\centering\textbf{Earlier Germanic changes}\par
\vspace{0.35em}
\raggedright
\centering\textbf{West Germanic}\par
\raggedright
\vspace{0.2em}
\raggedright [no change]\par
\vspace{0.6em}
\centering\textbf{Northwest Germanic}\par
\raggedright
\vspace{0.2em}
\begin{tabular}{@{}>{\raggedright\arraybackslash}p{0.74\linewidth}@{\hspace{0.55em}}>{\raggedright\arraybackslash}p{0.16\linewidth}@{\hspace{0.25em}}}
\mbox{PGmc Final Z Deletion} & \emph{*dḗdi} \\
\mbox{NWGmc Long E Lowering} & \emph{*dǣdi} \\
\end{tabular}
\end{minipage}
&
&
\begin{minipage}[t]{\linewidth}
\centering\textbf{Old English changes}\par
\vspace{0.35em}
\raggedright
\begin{tabular}{@{}>{\raggedright\arraybackslash}p{0.74\linewidth}@{\hspace{0.55em}}>{\raggedright\arraybackslash}p{0.16\linewidth}@{\hspace{0.25em}}}
\mbox{OE High Vowel Apocope} & \emph{*dǣd} \\
\end{tabular}
\end{minipage}
\\
\end{tabularx}
\end{minipage}%
} \endgroup

Orel reconstructs the noun as \emph{*dēdiz}, and Ringe and Taylor derive
the same inherited i-stem from Proto-Germanic \emph{*dédiz} through
northwest Germanic \emph{*dadiz} (\citeproc{ref-Orel2003}{Orel 2003};
\citeproc{ref-RingeTaylor2014}{Ringe and Taylor 2014, 27}). The
stress-marked form \emph{*dḗdiz} represents that same inherited noun in
a notation that keeps the stressed long vowel explicit. Campbell states
that Primitive Germanic long \emph{ē} appears as West Saxon \emph{ǣ} but
in other Old English dialects mostly as \emph{ē}
(\citeproc{ref-Campbell1959}{Campbell 1959, sect. 128}). Brunner gives
the contrast explicitly as West Saxon \emph{dǣd} beside non-West-Saxon
\emph{dēd} (\citeproc{ref-SieversBrunner1965}{Sievers 1965, sect. 98}).
From inherited \emph{*dēdiz}, loss of final \emph{-z} and the West Saxon
lowering of stressed long \emph{ē} yield \emph{dǣd}
(\citeproc{ref-Campbell1959}{Campbell 1959, sect. 128}); Anglian
\emph{dēd} preserves the non-West-Saxon outcome
(\citeproc{ref-SieversBrunner1965}{Sievers 1965, sect. 98}).

\subsubsection{door --- OE dor}\label{door-oe-dor}

Derivation: \emph{*dúrą} \textgreater{} \emph{dor} (regular).

\begingroup
\setlength{\fboxsep}{6pt}

\noindent\fbox{%
\begin{minipage}{0.97\linewidth}
\small
\begin{tabularx}{\linewidth}{@{}>{\raggedright\arraybackslash}p{0.460\linewidth}>{\centering\arraybackslash}X>{\raggedright\arraybackslash}p{0.460\linewidth}@{}}
\begin{minipage}[t]{\linewidth}
\centering\textbf{Earlier Germanic changes}\par
\vspace{0.35em}
\raggedright
\centering\textbf{West Germanic}\par
\raggedright
\vspace{0.2em}
\raggedright [no change]\par
\vspace{0.6em}
\centering\textbf{Northwest Germanic}\par
\raggedright
\vspace{0.2em}
\begin{tabular}{@{}>{\raggedright\arraybackslash}p{0.74\linewidth}@{\hspace{0.55em}}>{\raggedright\arraybackslash}p{0.16\linewidth}@{\hspace{0.25em}}}
\mbox{NWGmc U Lowering} & \emph{*dórą} \\
\end{tabular}
\end{minipage}
&
&
\begin{minipage}[t]{\linewidth}
\centering\textbf{Old English changes}\par
\vspace{0.35em}
\raggedright
\begin{tabular}{@{}>{\raggedright\arraybackslash}p{0.74\linewidth}@{\hspace{0.55em}}>{\raggedright\arraybackslash}p{0.16\linewidth}@{\hspace{0.25em}}}
OE Heavy Syllable Nasal Apocope & \emph{*dór} \\
\end{tabular}
\end{minipage}
\\
\end{tabularx}
\end{minipage}%
} \endgroup

Kroonen reconstructs a neuter \emph{*dura-} `gate, (single) door' and
cites Old English \emph{dor} among its reflexes. In the same entry he
separates Old English \emph{duru}, Old Frisian \emph{dore}, and Old High
German \emph{tura} as reflexes of \emph{*durō-} instead
(\citeproc{ref-Kroonen2013}{Kroonen 2013, 150}). Clark Hall records
\emph{dor} as a neuter noun and separately records feminine \emph{duru}
with its own inflection (\citeproc{ref-ClarkHall1960}{Clark Hall 1960,
92}). From \emph{*dúrą}, Northwest Germanic u-lowering gives
\emph{*dórą}, and heavy-syllable nasal apocope then yields \emph{dor}.

\subsubsection{fare --- OE faran}\label{fare-oe-faran}

Derivation: \emph{*fáraną} \textgreater{} \emph{faran} (regular).

\begingroup
\setlength{\fboxsep}{6pt}

\noindent\fbox{%
\begin{minipage}{0.97\linewidth}
\small
\begin{tabularx}{\linewidth}{@{}>{\raggedright\arraybackslash}p{0.300\linewidth}>{\centering\arraybackslash}X>{\raggedright\arraybackslash}p{0.540\linewidth}@{}}
\begin{minipage}[t]{\linewidth}
\centering\textbf{Earlier Germanic changes}\par
\vspace{0.35em}
\raggedright
\centering\textbf{West Germanic}\par
\raggedright
\vspace{0.2em}
\raggedright [no change]\par
\vspace{0.6em}
\centering\textbf{Northwest Germanic}\par
\raggedright
\vspace{0.2em}
\raggedright [no change]\par
\end{minipage}
&
&
\begin{minipage}[t]{\linewidth}
\centering\textbf{Old English changes}\par
\vspace{0.35em}
\raggedright
\begin{tabular}{@{}>{\raggedright\arraybackslash}p{0.74\linewidth}@{\hspace{0.55em}}>{\raggedright\arraybackslash}p{0.16\linewidth}@{\hspace{0.25em}}}
Anglo Frisian Brightening & \emph{*færaną} \\
OE A Restoration & \emph{*faraną} \\
OE Heavy Syllable Nasal Apocope & \emph{*faran} \\
OE Secondary Nasalization & \emph{*farąn} \\
OE Weak Tail Reduction & \emph{*faran} \\
\end{tabular}
\end{minipage}
\\
\end{tabularx}
\end{minipage}%
} \endgroup

Kroonen gives the inherited strong verb as \emph{*faran-}, and Orel
gives the same lexeme as \emph{*faranan}, both with Old English
\emph{faran} among the reflexes (\citeproc{ref-Kroonen2013}{Kroonen
2013, 149}; \citeproc{ref-Orel2003}{Orel 2003, 132};
\citeproc{ref-Campbell1959}{Campbell 1959, 61}). Clark Hall and
Bosworth-Toller distinguish strong \emph{faran} from weak \emph{færan},
while forms such as \emph{fære} and \emph{færð} belong to the present
tense (\citeproc{ref-ClarkHall1960}{Clark Hall 1960, 113};
\citeproc{ref-BosworthToller1898}{Bosworth and Toller 1898, 250}). From
\emph{*fáraną}, brightening first gives \emph{*færaną}; A-restoration
before single \emph{r} returns \emph{*faraną}, and later reductions
yield \emph{faran}.

\subsubsection{fell --- OE fell}\label{fell-oe-fell}

Derivation: \emph{*féllą} \textgreater{} \emph{fell} (regular).

\begingroup
\setlength{\fboxsep}{6pt}

\noindent\fbox{%
\begin{minipage}{0.97\linewidth}
\small
\begin{tabularx}{\linewidth}{@{}>{\raggedright\arraybackslash}p{0.300\linewidth}>{\centering\arraybackslash}X>{\raggedright\arraybackslash}p{0.540\linewidth}@{}}
\begin{minipage}[t]{\linewidth}
\centering\textbf{Earlier Germanic changes}\par
\vspace{0.35em}
\raggedright
\centering\textbf{West Germanic}\par
\raggedright
\vspace{0.2em}
\raggedright [no change]\par
\vspace{0.6em}
\centering\textbf{Northwest Germanic}\par
\raggedright
\vspace{0.2em}
\raggedright [no change]\par
\end{minipage}
&
&
\begin{minipage}[t]{\linewidth}
\centering\textbf{Old English changes}\par
\vspace{0.35em}
\raggedright
\begin{tabular}{@{}>{\raggedright\arraybackslash}p{0.74\linewidth}@{\hspace{0.55em}}>{\raggedright\arraybackslash}p{0.16\linewidth}@{\hspace{0.25em}}}
OE Heavy Syllable Nasal Apocope & \emph{*féll} \\
\end{tabular}
\end{minipage}
\\
\end{tabularx}
\end{minipage}%
} \endgroup

Kroonen reconstructs the noun as \emph{*fella-} `membrane, skin, hide'
and cites Old English \emph{fell} beside Dutch \emph{vel} and German
\emph{Fell} (\citeproc{ref-Kroonen2013}{Kroonen 2013}). The selected
input \emph{*féllą} is the derivable singular form of that same
inherited noun. Clark Hall records \emph{fell} as the noun `fell, skin,
hide' (\citeproc{ref-ClarkHall1960}{Clark Hall 1960, 115}). Bright's
glossary likewise gives \emph{fell, n.}, with inflected forms such as
accusative singular \emph{fel} and dative plural \emph{fellum}
(\citeproc{ref-BrightCassidyRingler1971}{Bright 1971, 277}). With
\emph{*féllą}, no special earlier reshaping is needed: heavy-syllable
nasal apocope yields \emph{*féll}, surfacing as \emph{fell}.

\subsubsection{fern --- OE fearn}\label{fern-oe-fearn}

Derivation: \emph{*fárnaz} \textgreater{} \emph{fearn} (regular).

\begingroup
\setlength{\fboxsep}{6pt}

\noindent\fbox{%
\begin{minipage}{0.97\linewidth}
\small
\begin{tabularx}{\linewidth}{@{}>{\raggedright\arraybackslash}p{0.300\linewidth}>{\centering\arraybackslash}X>{\raggedright\arraybackslash}p{0.540\linewidth}@{}}
\begin{minipage}[t]{\linewidth}
\centering\textbf{Earlier Germanic changes}\par
\vspace{0.35em}
\raggedright
\centering\textbf{West Germanic}\par
\raggedright
\vspace{0.2em}
\raggedright [no change]\par
\vspace{0.6em}
\centering\textbf{Northwest Germanic}\par
\raggedright
\vspace{0.2em}
\begin{tabular}{@{}>{\raggedright\arraybackslash}p{0.74\linewidth}@{\hspace{0.55em}}>{\raggedright\arraybackslash}p{0.16\linewidth}@{\hspace{0.25em}}}
\mbox{PGmc Final Z Deletion} & \emph{*fárna} \\
\end{tabular}
\end{minipage}
&
&
\begin{minipage}[t]{\linewidth}
\centering\textbf{Old English changes}\par
\vspace{0.35em}
\raggedright
\begin{tabular}{@{}>{\raggedright\arraybackslash}p{0.70\linewidth}@{\hspace{0.55em}}>{\raggedright\arraybackslash}p{0.16\linewidth}@{\hspace{0.25em}}}
PWGmc Final Bare A Loss & \emph{*fárn} \\
Anglo Frisian Brightening & \emph{*færn} \\
OE Breaking & \emph{*fearn} \\
\end{tabular}
\end{minipage}
\\
\end{tabularx}
\end{minipage}%
} \endgroup

Kroonen cites the noun as masculine \emph{*farna-} and gives Old English
\emph{fearn, fern}, while Orel gives the same lexeme as neuter
\emph{*farnan} with Old English \emph{fearn}
(\citeproc{ref-Kroonen2013}{Kroonen 2013, 169};
\citeproc{ref-Orel2003}{Orel 2003, 133}). Those are comparative headword
conventions rather than competing Old English outcomes; the modeled
input here is the nominative-style \emph{*fárnaz}. Clark Hall gives
\emph{fearn} as an Old English noun, and Bosworth-Toller records
\emph{fearn} with inflected forms such as \emph{fearnes}, \emph{fearna},
and \emph{fearne} (\citeproc{ref-ClarkHall1960}{Clark Hall 1960, 114};
\citeproc{ref-BosworthToller1898}{Bosworth and Toller 1898, 219}). From
\emph{*fárnaz}, loss of final \emph{-z} and final \emph{-a} gives
\emph{*fárn}; Anglo-Frisian brightening then yields \emph{*færn}, and
breaking before \emph{r} plus consonant gives \emph{fearn}
(\citeproc{ref-Campbell1959}{Campbell 1959};
\citeproc{ref-RingeTaylor2014}{Ringe and Taylor 2014}).

\subsubsection{field --- OE feld}\label{field-oe-feld}

Derivation: \emph{*félθuz} \textgreater{} \emph{feld} (regular).

\begingroup
\setlength{\fboxsep}{6pt}

\noindent\fbox{%
\begin{minipage}{0.97\linewidth}
\small
\begin{tabularx}{\linewidth}{@{}>{\raggedright\arraybackslash}p{0.460\linewidth}>{\centering\arraybackslash}X>{\raggedright\arraybackslash}p{0.460\linewidth}@{}}
\begin{minipage}[t]{\linewidth}
\centering\textbf{Earlier Germanic changes}\par
\vspace{0.35em}
\raggedright
\centering\textbf{West Germanic}\par
\raggedright
\vspace{0.2em}
\begin{tabular}{@{}>{\raggedright\arraybackslash}p{0.74\linewidth}@{\hspace{0.55em}}>{\raggedright\arraybackslash}p{0.16\linewidth}@{\hspace{0.25em}}}
PWGmc L Th Voicing & \emph{*félduz} \\
\end{tabular}
\vspace{0.6em}
\centering\textbf{Northwest Germanic}\par
\raggedright
\vspace{0.2em}
\begin{tabular}{@{}>{\raggedright\arraybackslash}p{0.74\linewidth}@{\hspace{0.55em}}>{\raggedright\arraybackslash}p{0.16\linewidth}@{\hspace{0.25em}}}
\mbox{PGmc Final Z Deletion} & \emph{*féldu} \\
\end{tabular}
\end{minipage}
&
&
\begin{minipage}[t]{\linewidth}
\centering\textbf{Old English changes}\par
\vspace{0.35em}
\raggedright
\begin{tabular}{@{}>{\raggedright\arraybackslash}p{0.74\linewidth}@{\hspace{0.55em}}>{\raggedright\arraybackslash}p{0.16\linewidth}@{\hspace{0.25em}}}
\mbox{OE High Vowel Apocope} & \emph{*féld} \\
\end{tabular}
\end{minipage}
\\
\end{tabularx}
\end{minipage}%
} \endgroup

Ringe and Taylor treat Old English \emph{feld} as one of the cases where
earlier \emph{*felþu-} \textasciitilde{} \emph{*feldu-} may reflect
either inherited \emph{*þ} \textasciitilde{} \emph{*d} alternation or
the regular West Germanic development \emph{*lþ} \textgreater{} ld
(\citeproc{ref-RingeTaylor2014}{Ringe and Taylor 2014, 170}). Clark Hall
records \emph{feld} with oblique forms such as \emph{felda} and
\emph{felde}, and Campbell notes early place-name spellings in
\emph{-felth} beside the later standard form
(\citeproc{ref-ClarkHall1960}{Clark Hall 1960, 114};
\citeproc{ref-Campbell1959}{Campbell 1959, 169}). In the modeled
pathway, medial \emph{*lþ} becomes \emph{ld}, final \emph{-z} is lost,
and high-vowel apocope then yields \emph{feld}.

\subsubsection{fly --- OE flēogan}\label{fly-oe-flux113ogan}

Derivation: \emph{*fléuganą} \textgreater{} \emph{flēogan} (regular).

\begingroup
\setlength{\fboxsep}{6pt}

\noindent\fbox{%
\begin{minipage}{0.97\linewidth}
\small
\begin{tabularx}{\linewidth}{@{}>{\raggedright\arraybackslash}p{0.300\linewidth}>{\centering\arraybackslash}X>{\raggedright\arraybackslash}p{0.540\linewidth}@{}}
\begin{minipage}[t]{\linewidth}
\centering\textbf{Earlier Germanic changes}\par
\vspace{0.35em}
\raggedright
\centering\textbf{West Germanic}\par
\raggedright
\vspace{0.2em}
\raggedright [no change]\par
\vspace{0.6em}
\centering\textbf{Northwest Germanic}\par
\raggedright
\vspace{0.2em}
\raggedright [no change]\par
\end{minipage}
&
&
\begin{minipage}[t]{\linewidth}
\centering\textbf{Old English changes}\par
\vspace{0.35em}
\raggedright
\begin{tabular}{@{}>{\raggedright\arraybackslash}p{0.68\linewidth}@{\hspace{0.55em}}>{\raggedright\arraybackslash}p{0.22\linewidth}@{\hspace{0.25em}}}
OE Diphthong Leveling & \emph{*flēoganą} \\
OE Heavy Syllable Nasal Apocope & \emph{*flēogan} \\
OE Secondary Nasalization & \emph{*flēogąn} \\
OE Weak Tail Reduction & \emph{*flēogan} \\
\end{tabular}
\end{minipage}
\\
\end{tabularx}
\end{minipage}%
} \endgroup

Ringe and Taylor derive the verb as \emph{*fleugana} \textgreater{} OE
\emph{fléogan} and elsewhere contrast West Saxon \emph{fléogan} with
Anglian \emph{flégan}, alongside related forms \emph{fléoge} /
\emph{flége} (\citeproc{ref-RingeTaylor2014}{Ringe and Taylor 2014, 189,
324}). The selected input \emph{*fléuganą} represents that inherited
strong verb in the notation used here. Bosworth-Toller records
\emph{flēogan} as the ordinary Old English strong verb
(\citeproc{ref-BosworthToller1898}{Bosworth and Toller 1898}). Clark
Hall likewise records \emph{flēogan} (\citeproc{ref-ClarkHall1960}{Clark
Hall 1960, 106}). Bright gives the familiar paradigm \emph{flēag},
flugon, flogen with present \emph{fleogeð}
(\citeproc{ref-BrightCassidyRingler1971}{Bright 1971, 363}). From
\emph{*fléuganą}, Old English diphthong leveling gives \emph{*flēoganą};
heavy-syllable nasal apocope and weak-tail reduction then yield
\emph{flēogan} (\citeproc{ref-RingeTaylor2014}{Ringe and Taylor 2014,
189}).

Form note. Ringe and Taylor also list related \emph{fléoge} /
\emph{flége} and Anglian \emph{flégan}, which belong to the same family
but do not replace the infinitive \emph{flēogan} treated here
(\citeproc{ref-RingeTaylor2014}{Ringe and Taylor 2014, 324}).

\subsubsection{forlorn --- OE lēosan}\label{forlorn-oe-lux113osan}

Derivation: \emph{*léusaną} \textgreater{} \emph{lēosan} (regular).

\begingroup
\setlength{\fboxsep}{6pt}

\noindent\fbox{%
\begin{minipage}{0.97\linewidth}
\small
\begin{tabularx}{\linewidth}{@{}>{\raggedright\arraybackslash}p{0.300\linewidth}>{\centering\arraybackslash}X>{\raggedright\arraybackslash}p{0.540\linewidth}@{}}
\begin{minipage}[t]{\linewidth}
\centering\textbf{Earlier Germanic changes}\par
\vspace{0.35em}
\raggedright
\centering\textbf{West Germanic}\par
\raggedright
\vspace{0.2em}
\raggedright [no change]\par
\vspace{0.6em}
\centering\textbf{Northwest Germanic}\par
\raggedright
\vspace{0.2em}
\raggedright [no change]\par
\end{minipage}
&
&
\begin{minipage}[t]{\linewidth}
\centering\textbf{Old English changes}\par
\vspace{0.35em}
\raggedright
\begin{tabular}{@{}>{\raggedright\arraybackslash}p{0.70\linewidth}@{\hspace{0.55em}}>{\raggedright\arraybackslash}p{0.20\linewidth}@{\hspace{0.25em}}}
OE Diphthong Leveling & \emph{*lēosaną} \\
OE Heavy Syllable Nasal Apocope & \emph{*lēosan} \\
OE Secondary Nasalization & \emph{*lēosąn} \\
OE Weak Tail Reduction & \emph{*lēosan} \\
\end{tabular}
\end{minipage}
\\
\end{tabularx}
\end{minipage}%
} \endgroup

Kroonen reconstructs the verb under \emph{*leusan-} and cites prefixed
daughters such as Gothic \emph{fra-liusan} and Old English
\emph{for-lēosan}; Orel likewise gives Old English \emph{for-leósan}
(\citeproc{ref-Kroonen2013}{Kroonen 2013}; \citeproc{ref-Orel2003}{Orel
2003}). The inherited verbal base is therefore clear, though the
daughter set often appears with the prefix. The direct Old English
evidence behind English \emph{forlorn} lies in the prefixed verb
\emph{forlēosan} and especially in the participle \emph{forloren},
recorded by Ringe and Taylor and in the dictionaries
(\citeproc{ref-RingeTaylor2014}{Ringe and Taylor 2014, 357};
\citeproc{ref-ClarkHall1960}{Clark Hall 1960};
\citeproc{ref-BosworthToller1898}{Bosworth and Toller 1898, 248}). From
\emph{*léusaną}, Old English diphthong leveling gives \emph{*lēosaną},
and later nasal apocope and weak-tail reduction yield \emph{lēosan}
(\citeproc{ref-RingeTaylor2014}{Ringe and Taylor 2014}).

Form note. As a base-form comparison, the simplex infinitive is
\emph{lēosan}, while the English adjective continues the prefixed Old
English family \emph{forlēosan} / forloren
(\citeproc{ref-RingeTaylor2014}{Ringe and Taylor 2014, 357}).

\subsubsection{gang --- OE gang}\label{gang-oe-gang}

Derivation: \emph{*gángaz} \textgreater{} \emph{gang} (regular).

\begingroup
\setlength{\fboxsep}{6pt}

\noindent\fbox{%
\begin{minipage}{0.97\linewidth}
\small
\begin{tabularx}{\linewidth}{@{}>{\raggedright\arraybackslash}p{0.540\linewidth}>{\centering\arraybackslash}X>{\raggedright\arraybackslash}p{0.300\linewidth}@{}}
\begin{minipage}[t]{\linewidth}
\centering\textbf{Earlier Germanic changes}\par
\vspace{0.35em}
\raggedright
\centering\textbf{West Germanic}\par
\raggedright
\vspace{0.2em}
\raggedright [no change]\par
\vspace{0.6em}
\centering\textbf{Northwest Germanic}\par
\raggedright
\vspace{0.2em}
\begin{tabular}{@{}>{\raggedright\arraybackslash}p{0.74\linewidth}@{\hspace{0.55em}}>{\raggedright\arraybackslash}p{0.16\linewidth}@{\hspace{0.25em}}}
\mbox{PGmc Final Z Deletion} & \emph{*gánga} \\
\end{tabular}
\end{minipage}
&
&
\begin{minipage}[t]{\linewidth}
\centering\textbf{Old English changes}\par
\vspace{0.35em}
\raggedright
\begin{tabular}{@{}>{\raggedright\arraybackslash}p{0.70\linewidth}@{\hspace{0.55em}}>{\raggedright\arraybackslash}p{0.16\linewidth}@{\hspace{0.25em}}}
PWGmc Final Bare A Loss & \emph{*gáng} \\
\end{tabular}
\end{minipage}
\\
\end{tabularx}
\end{minipage}%
} \endgroup

Orel reconstructs the noun as \emph{*gangaz} and cites Old English
\emph{gang} beside Old Norse \emph{gangr}, Old Frisian \emph{gang /
gong}, Old Saxon \emph{gang}, and Old High German \emph{gang}
(\citeproc{ref-Orel2003}{Orel 2003}). The selected input \emph{*gángaz}
is the same lexeme in the accent notation used here. Bosworth-Toller
records \emph{gang} as the noun `going, journey, way'
(\citeproc{ref-BosworthToller1898}{Bosworth and Toller 1898, 159}), and
Bright's glossary gives \emph{gong (gang), m., path, course}
(\citeproc{ref-BrightCassidyRingler1971}{Bright 1971, 392}). From
\emph{*gángaz}, loss of final \emph{-z} gives \emph{*gánga}, and later
loss of final bare \emph{-a} yields \emph{gang}.

Form note. This entry concerns the noun \emph{gang}, not the separate
verb \emph{gangan}.

\subsubsection{give --- OE ġiefan}\label{give-oe-ux121iefan}

Derivation: \emph{*gébaną} \textgreater{} \emph{ġiefan} (regular).

\begingroup
\setlength{\fboxsep}{6pt}

\noindent\fbox{%
\begin{minipage}{0.97\linewidth}
\small
\begin{tabularx}{\linewidth}{@{}>{\raggedright\arraybackslash}p{0.280\linewidth}>{\centering\arraybackslash}X>{\raggedright\arraybackslash}p{0.560\linewidth}@{}}
\begin{minipage}[t]{\linewidth}
\centering\textbf{Earlier Germanic changes}\par
\vspace{0.35em}
\raggedright
\centering\textbf{West Germanic}\par
\raggedright
\vspace{0.2em}
\raggedright [no change]\par
\vspace{0.6em}
\centering\textbf{Northwest Germanic}\par
\raggedright
\vspace{0.2em}
\raggedright [no change]\par
\end{minipage}
&
&
\begin{minipage}[t]{\linewidth}
\centering\textbf{Old English changes}\par
\vspace{0.35em}
\raggedright
\begin{tabular}{@{}>{\raggedright\arraybackslash}p{0.74\linewidth}@{\hspace{0.55em}}>{\raggedright\arraybackslash}p{0.16\linewidth}@{\hspace{0.25em}}}
OE Heavy Syllable Nasal Apocope & \emph{*géban} \\
OE Secondary Nasalization & \emph{*gébąn} \\
PGmc B Allophony & \emph{*géβąn} \\
OE Velar Palatalization & \emph{*ʤéβąn} \\
OE Ws Palatal Diphthongization & \emph{*ʤíeβąn} \\
OE Weak Tail Reduction & \emph{*ʤíeβan} \\
\end{tabular}
\end{minipage}
\\
\end{tabularx}
\end{minipage}%
} \endgroup

Kroonen reconstructs the strong verb as \emph{*geban-} and cites Old
English \emph{giefan} among its reflexes
(\citeproc{ref-Kroonen2013}{Kroonen 2013, 212}). Ringe and Taylor
contrast West Saxon \emph{giefan} with Mercian \emph{for-geofan} and
Northumbrian \emph{geafa}, showing that the inherited verb takes
different later dialectal shapes (\citeproc{ref-RingeTaylor2014}{Ringe
and Taylor 2014, 230}). Campbell gives \emph{gefan (W-S giefan)} among
examples of initial palatalization, and Clark Hall records the verb
under plain \emph{giefan} with forms such as \emph{geaf} and
\emph{giefen} (\citeproc{ref-Campbell1959}{Campbell 1959, sect. 428};
\citeproc{ref-ClarkHall1960}{Clark Hall 1960}). From \emph{*gébaną},
initial \emph{g} palatalizes before \emph{e}; West Saxon palatal
diphthongization then yields \emph{ie}, and later tail reduction gives
\emph{giefan} (\citeproc{ref-Campbell1959}{Campbell 1959, sect. 185};
\citeproc{ref-RingeTaylor2014}{Ringe and Taylor 2014, 230}).

Dialect note. West Saxon \emph{ie} here reflects palatal
diphthongization after initial palatalization; non-West-Saxon forms such
as \emph{geafa} or \emph{for-geofan} continue the same verb without the
West Saxon vocalism (\citeproc{ref-RingeTaylor2014}{Ringe and Taylor
2014, 230}).

\subsubsection{gold --- OE gold}\label{gold-oe-gold}

Derivation: \emph{*gúlθą} \textgreater{} \emph{gold} (regular).

\begingroup
\setlength{\fboxsep}{6pt}

\noindent\fbox{%
\begin{minipage}{0.97\linewidth}
\small
\begin{tabularx}{\linewidth}{@{}>{\raggedright\arraybackslash}p{0.460\linewidth}>{\centering\arraybackslash}X>{\raggedright\arraybackslash}p{0.460\linewidth}@{}}
\begin{minipage}[t]{\linewidth}
\centering\textbf{Earlier Germanic changes}\par
\vspace{0.35em}
\raggedright
\centering\textbf{West Germanic}\par
\raggedright
\vspace{0.2em}
\begin{tabular}{@{}>{\raggedright\arraybackslash}p{0.74\linewidth}@{\hspace{0.55em}}>{\raggedright\arraybackslash}p{0.16\linewidth}@{\hspace{0.25em}}}
PWGmc L Th Voicing & \emph{*gúldą} \\
\end{tabular}
\vspace{0.6em}
\centering\textbf{Northwest Germanic}\par
\raggedright
\vspace{0.2em}
\begin{tabular}{@{}>{\raggedright\arraybackslash}p{0.74\linewidth}@{\hspace{0.55em}}>{\raggedright\arraybackslash}p{0.16\linewidth}@{\hspace{0.25em}}}
\mbox{NWGmc U Lowering} & \emph{*góldą} \\
\end{tabular}
\end{minipage}
&
&
\begin{minipage}[t]{\linewidth}
\centering\textbf{Old English changes}\par
\vspace{0.35em}
\raggedright
\begin{tabular}{@{}>{\raggedright\arraybackslash}p{0.74\linewidth}@{\hspace{0.55em}}>{\raggedright\arraybackslash}p{0.16\linewidth}@{\hspace{0.25em}}}
OE Heavy Syllable Nasal Apocope & \emph{*góld} \\
\end{tabular}
\end{minipage}
\\
\end{tabularx}
\end{minipage}%
} \endgroup

Ringe and Taylor cite the noun as \emph{*gulþa-} / \emph{*gulda-}, and
Kroonen gives the same pair (\citeproc{ref-RingeTaylor2014}{Ringe and
Taylor 2014, 42}; \citeproc{ref-Kroonen2013}{Kroonen 2013}). The
selected input \emph{*gúlθą} preserves the older consonantal form while
leaving open whether the medial stop reflects inherited alternation or
regular West Germanic development. Bosworth-Toller and Clark Hall both
record \emph{gold} as the ordinary Old English neuter noun
(\citeproc{ref-BosworthToller1898}{Bosworth and Toller 1898, 121};
\citeproc{ref-ClarkHall1960}{Clark Hall 1960, 152}). From \emph{*gúlθą},
the regular consonant development gives \emph{*gúldą}; Northwest
Germanic / Old English lowering then yields \emph{*góldą}, and apocope
gives \emph{gold} (\citeproc{ref-Campbell1959}{Campbell 1959, sect.
414}; \citeproc{ref-RingeTaylor2014}{Ringe and Taylor 2014, 42}). Ringe
and Taylor note that the medial stop can be understood either as
alternation \emph{*gulþa-} / \emph{*gulda-} or as the ordinary West
Germanic change \emph{*lþ} \textgreater{} ld; both routes lead to the
same Old English consonantism (\citeproc{ref-RingeTaylor2014}{Ringe and
Taylor 2014, 42}).

\subsubsection{grave --- OE grafan}\label{grave-oe-grafan}

Derivation: \emph{*grábaną} \textgreater{} \emph{grafan} (regular).

\begingroup
\setlength{\fboxsep}{6pt}

\noindent\fbox{%
\begin{minipage}{0.97\linewidth}
\small
\begin{tabularx}{\linewidth}{@{}>{\raggedright\arraybackslash}p{0.280\linewidth}>{\centering\arraybackslash}X>{\raggedright\arraybackslash}p{0.560\linewidth}@{}}
\begin{minipage}[t]{\linewidth}
\centering\textbf{Earlier Germanic changes}\par
\vspace{0.35em}
\raggedright
\centering\textbf{West Germanic}\par
\raggedright
\vspace{0.2em}
\raggedright [no change]\par
\vspace{0.6em}
\centering\textbf{Northwest Germanic}\par
\raggedright
\vspace{0.2em}
\raggedright [no change]\par
\end{minipage}
&
&
\begin{minipage}[t]{\linewidth}
\centering\textbf{Old English changes}\par
\vspace{0.35em}
\raggedright
\begin{tabular}{@{}>{\raggedright\arraybackslash}p{0.70\linewidth}@{\hspace{0.55em}}>{\raggedright\arraybackslash}p{0.20\linewidth}@{\hspace{0.25em}}}
Anglo Frisian Brightening & \emph{*græbaną} \\
OE A Restoration & \emph{*grabaną} \\
OE Heavy Syllable Nasal Apocope & \emph{*graban} \\
OE Secondary Nasalization & \emph{*grabąn} \\
PGmc B Allophony & \emph{*graβąn} \\
OE Weak Tail Reduction & \emph{*graβan} \\
\end{tabular}
\end{minipage}
\\
\end{tabularx}
\end{minipage}%
} \endgroup

Campbell gives \emph{grafan} among the standard examples of Old English
a-restoration before a single consonant and a following back vowel, and
Ringe and Taylor describe the same development for Class VI infinitives
(\citeproc{ref-Campbell1959}{Campbell 1959, 61};
\citeproc{ref-RingeTaylor2014}{Ringe and Taylor 2014}). Clark Hall
records \emph{grafan} as the verb `to dig, grave' and separately records
noun \emph{græf} `grave, trench' (\citeproc{ref-ClarkHall1960}{Clark
Hall 1960, 153}). From \emph{*grábaną}, Anglo-Frisian brightening first
gives a fronted stem vowel.

Form note. Noun \emph{græf} and verbal forms such as \emph{græfð} or
past participial \emph{græfen} belong to other lexical or paradigm
positions and do not replace the infinitive \emph{grafan} as the target
here (\citeproc{ref-ClarkHall1960}{Clark Hall 1960, 153}).

\subsubsection{guest --- OE ġiest}\label{guest-oe-ux121iest}

Derivation: \emph{*gástiz} \textgreater{} \emph{ġiest} (regular).

\begingroup
\setlength{\fboxsep}{6pt}

\noindent\fbox{%
\begin{minipage}{0.97\linewidth}
\small
\begin{tabularx}{\linewidth}{@{}>{\raggedright\arraybackslash}p{0.460\linewidth}>{\centering\arraybackslash}X>{\raggedright\arraybackslash}p{0.460\linewidth}@{}}
\begin{minipage}[t]{\linewidth}
\centering\textbf{Earlier Germanic changes}\par
\vspace{0.35em}
\raggedright
\centering\textbf{West Germanic}\par
\raggedright
\vspace{0.2em}
\raggedright [no change]\par
\vspace{0.6em}
\centering\textbf{Northwest Germanic}\par
\raggedright
\vspace{0.2em}
\begin{tabular}{@{}>{\raggedright\arraybackslash}p{0.74\linewidth}@{\hspace{0.55em}}>{\raggedright\arraybackslash}p{0.16\linewidth}@{\hspace{0.25em}}}
\mbox{PGmc Final Z Deletion} & \emph{*gásti} \\
\end{tabular}
\end{minipage}
&
&
\begin{minipage}[t]{\linewidth}
\centering\textbf{Old English changes}\par
\vspace{0.35em}
\raggedright
\begin{tabular}{@{}>{\raggedright\arraybackslash}p{0.74\linewidth}@{\hspace{0.55em}}>{\raggedright\arraybackslash}p{0.16\linewidth}@{\hspace{0.25em}}}
Anglo Frisian Brightening & \emph{*gæsti} \\
OE Velar Palatalization & \emph{*ʤæsti} \\
OE I Umlaut & \emph{*ʤesti} \\
OE Ws Palatal Diphthongization & \emph{*ʤiesti} \\
\mbox{OE High Vowel Apocope} & \emph{*ʤiest} \\
\end{tabular}
\end{minipage}
\\
\end{tabularx}
\end{minipage}%
} \endgroup

Campbell and Ringe and Taylor treat this as an ordinary i-stem whose
West Saxon development yields \emph{ġiest}, while non-West-Saxon
evidence preserves forms of the \emph{gest} type
(\citeproc{ref-Campbell1959}{Campbell 1959, 91};
\citeproc{ref-RingeTaylor2014}{Ringe and Taylor 2014, 296}).
Bosworth-Toller and Clark Hall record spellings such as \emph{gist},
\emph{gest}, \emph{giest}, and \emph{gyst}
(\citeproc{ref-BosworthToller1898}{Bosworth and Toller 1898, 387};
\citeproc{ref-ClarkHall1960}{Clark Hall 1960, 149}). The derivation is
regular: brightening and i-mutation reshape the stem, and West Saxon
palatal diphthongization gives \emph{ie}.

Dialect note. Anglian \emph{gest} is genuine Old English evidence
alongside West Saxon \emph{ġiest} (\citeproc{ref-RingeTaylor2014}{Ringe
and Taylor 2014, 296}; \citeproc{ref-BosworthToller1898}{Bosworth and
Toller 1898, 387}).

\subsubsection{hair --- OE hǣr}\label{hair-oe-hux1e3r}

Derivation: \emph{*xḗrą} \textgreater{} \emph{hǣr} (regular).

\begingroup
\setlength{\fboxsep}{6pt}

\noindent\fbox{%
\begin{minipage}{0.97\linewidth}
\small
\begin{tabularx}{\linewidth}{@{}>{\raggedright\arraybackslash}p{0.460\linewidth}>{\centering\arraybackslash}X>{\raggedright\arraybackslash}p{0.460\linewidth}@{}}
\begin{minipage}[t]{\linewidth}
\centering\textbf{Earlier Germanic changes}\par
\vspace{0.35em}
\raggedright
\centering\textbf{West Germanic}\par
\raggedright
\vspace{0.2em}
\raggedright [no change]\par
\vspace{0.6em}
\centering\textbf{Northwest Germanic}\par
\raggedright
\vspace{0.2em}
\begin{tabular}{@{}>{\raggedright\arraybackslash}p{0.74\linewidth}@{\hspace{0.55em}}>{\raggedright\arraybackslash}p{0.16\linewidth}@{\hspace{0.25em}}}
\mbox{NWGmc Long E Lowering} & \emph{*xǣrą} \\
\end{tabular}
\end{minipage}
&
&
\begin{minipage}[t]{\linewidth}
\centering\textbf{Old English changes}\par
\vspace{0.35em}
\raggedright
\begin{tabular}{@{}>{\raggedright\arraybackslash}p{0.74\linewidth}@{\hspace{0.55em}}>{\raggedright\arraybackslash}p{0.16\linewidth}@{\hspace{0.25em}}}
OE Velar Fricative Palatalization & \emph{*çǣrą} \\
OE Heavy Syllable Nasal Apocope & \emph{*çǣr} \\
\end{tabular}
\end{minipage}
\\
\end{tabularx}
\end{minipage}%
} \endgroup

Kroonen cites the ordinary Proto-Germanic hair word as \emph{*hēra-}
(\citeproc{ref-Kroonen2013}{Kroonen 2013}). The selected input
\emph{*xḗrą} represents that same long-ē stem in the present derivation.
Clark Hall and Bosworth-Toller record \emph{hær} / \emph{hǣr} as the
ordinary Old English noun `hair' (\citeproc{ref-ClarkHall1960}{Clark
Hall 1960, 158}; \citeproc{ref-BosworthToller1898}{Bosworth and Toller
1898, 510}). From \emph{*xḗrą}, Northwest Germanic lowering gives a long
front vowel, and later loss of the final nasal leaves the Old English
form \emph{hǣr}.

Form note. Older references to \emph{*xazwăz} belong to a different
lexeme, and the separate haddr / heordan / \emph{hād-} material does not
displace the ordinary simplex \emph{hǣr} treated here
(\citeproc{ref-Kroonen2013}{Kroonen 2013};
\citeproc{ref-ClarkHall1960}{Clark Hall 1960, 158}).

\subsubsection{harvest --- OE hierfest}\label{harvest-oe-hierfest}

Derivation: \emph{*xárbistuz} \textgreater{} \emph{hierfest} (regular).

\begingroup
\setlength{\fboxsep}{6pt}

\noindent\fbox{%
\begin{minipage}{0.97\linewidth}
\footnotesize
\begin{tabularx}{\linewidth}{@{}>{\raggedright\arraybackslash}p{0.460\linewidth}>{\centering\arraybackslash}X>{\raggedright\arraybackslash}p{0.460\linewidth}@{}}
\begin{minipage}[t]{\linewidth}
\centering\textbf{Earlier Germanic changes}\par
\vspace{0.35em}
\raggedright
\centering\textbf{West Germanic}\par
\raggedright
\vspace{0.2em}
\raggedright [no change]\par
\vspace{0.6em}
\centering\textbf{Northwest Germanic}\par
\raggedright
\vspace{0.2em}
\begin{tabular}{@{}>{\raggedright\arraybackslash}p{0.68\linewidth}@{\hspace{0.55em}}>{\raggedright\arraybackslash}p{0.22\linewidth}@{\hspace{0.25em}}}
PGmc Final Z Deletion & \emph{*xárbistu} \\
\end{tabular}
\end{minipage}
&
&
\begin{minipage}[t]{\linewidth}
\centering\textbf{Old English changes}\par
\vspace{0.35em}
\raggedright
\begin{tabular}{@{}>{\raggedright\arraybackslash}p{0.64\linewidth}@{\hspace{0.55em}}>{\raggedright\arraybackslash}p{0.26\linewidth}@{\hspace{0.25em}}}
Anglo Frisian Brightening & \emph{*xærbistu} \\
OE Breaking & \emph{*xearbistu} \\
OE Velar Fricative Palatalization & \emph{*çearbistu} \\
PGmc B Allophony & \emph{*çearβistu} \\
OE I Umlaut & \emph{*çierβistu} \\
OE High Vowel Apocope & \emph{*çierβist} \\
OE Med Unstressed I Lowering1 & \emph{*çierβest} \\
\end{tabular}
\end{minipage}
\\
\end{tabularx}
\end{minipage}%
} \endgroup

Bammesberger and Ringe-Taylor treat \emph{*harbist-} as the inherited
base and explain that the regular native West Saxon development would be
of the \emph{hierfest / hyrfest} type
(\citeproc{ref-Bammesberger1997}{Bammesberger 1997, 224};
\citeproc{ref-RingeTaylor2014}{Ringe and Taylor 2014, 141}).
Bosworth-Toller and Clark Hall record \emph{hærfest}, with
\emph{herfest} as a variant in the lexical tradition
(\citeproc{ref-BosworthToller1898}{Bosworth and Toller 1898};
\citeproc{ref-ClarkHall1960}{Clark Hall 1960, 158};
\citeproc{ref-ClarkHall1960}{1960, 171}). From \emph{*xárbistuz},
Anglo-Frisian brightening, breaking, and i-mutation produce a
\emph{hierbist-} stage, and later lowering of unstressed medial \emph{i}
to \emph{e} gives \emph{hierfest}.

Source note. The selected target \emph{hierfest} represents the regular
native West Saxon outcome discussed by Bammesberger and Ringe-Taylor.
The attested Old English lexical tradition, however, is chiefly
\emph{hærfest} / herfest, commonly treated as non-West-Saxon or Anglian
material in West Saxon transmission
(\citeproc{ref-Bammesberger1997}{Bammesberger 1997, 230};
\citeproc{ref-RingeTaylor2014}{Ringe and Taylor 2014, 267}).

\subsubsection{hedge --- OE heġġ}\label{hedge-oe-heux121ux121}

Derivation: \emph{*xágjaz} \textgreater{} \emph{heġġ} (regular).

\begingroup
\setlength{\fboxsep}{6pt}

\noindent\fbox{%
\begin{minipage}{0.97\linewidth}
\small
\begin{tabularx}{\linewidth}{@{}>{\raggedright\arraybackslash}p{0.320\linewidth}>{\centering\arraybackslash}X>{\raggedright\arraybackslash}p{0.520\linewidth}@{}}
\begin{minipage}[t]{\linewidth}
\centering\textbf{Earlier Germanic changes}\par
\vspace{0.35em}
\raggedright
\centering\textbf{West Germanic}\par
\raggedright
\vspace{0.2em}
\begin{tabular}{@{}>{\raggedright\arraybackslash}p{0.70\linewidth}@{\hspace{0.55em}}>{\raggedright\arraybackslash}p{0.20\linewidth}@{\hspace{0.25em}}}
PWGmc J Gemination & \emph{*xággjaz} \\
\end{tabular}
\vspace{0.6em}
\centering\textbf{Northwest Germanic}\par
\raggedright
\vspace{0.2em}
\begin{tabular}{@{}>{\raggedright\arraybackslash}p{0.70\linewidth}@{\hspace{0.55em}}>{\raggedright\arraybackslash}p{0.20\linewidth}@{\hspace{0.25em}}}
\mbox{PGmc Final Z Deletion} & \emph{*xággja} \\
\end{tabular}
\end{minipage}
&
&
\begin{minipage}[t]{\linewidth}
\centering\textbf{Old English changes}\par
\vspace{0.35em}
\raggedright
\begin{tabular}{@{}>{\raggedright\arraybackslash}p{0.74\linewidth}@{\hspace{0.55em}}>{\raggedright\arraybackslash}p{0.16\linewidth}@{\hspace{0.25em}}}
PWGmc Final Bare A Loss & \emph{*xággj} \\
Anglo Frisian Brightening & \emph{*xæggj} \\
OE Velar Fricative Palatalization & \emph{*çæggj} \\
OE Velar Palatalization & \emph{*çæʤʤj} \\
OE I Umlaut & \emph{*çeʤʤj} \\
OE J Loss After Heavy & \emph{*çeʤʤ} \\
\end{tabular}
\end{minipage}
\\
\end{tabularx}
\end{minipage}%
} \endgroup

The selected input models a palatal \emph{*-gj-} noun whose Old English
development includes gemination, palatalization, and i-mutation. The
current derivation therefore reaches a palatal-geminate outcome of the
\emph{heġġ} type (\citeproc{ref-Campbell1959}{Campbell 1959}).
Bosworth-Toller and Clark Hall record the noun under standard spellings
\emph{hecg} / \emph{heċġ} (\citeproc{ref-BosworthToller1898}{Bosworth
and Toller 1898}; \citeproc{ref-ClarkHall1960}{Clark Hall 1960}). From
\emph{*xágjaz}, West Germanic j-gemination first yields a geminate stop,
and later Old English palatalization and loss of final \emph{j} produce
\emph{heġġ}.

Form note. Standard dictionary spelling is \emph{heċġ} or \emph{hecg}.
Normalized \emph{heġġ} is the selected target here, while the ordinary
lexicographic forms remain the main Old English citation evidence
(\citeproc{ref-BosworthToller1898}{Bosworth and Toller 1898};
\citeproc{ref-ClarkHall1960}{Clark Hall 1960}).

\subsubsection{helm --- OE helm}\label{helm-oe-helm}

Derivation: \emph{*xélmaz} \textgreater{} \emph{helm} (regular).

\begingroup
\setlength{\fboxsep}{6pt}

\noindent\fbox{%
\begin{minipage}{0.97\linewidth}
\small
\begin{tabularx}{\linewidth}{@{}>{\raggedright\arraybackslash}p{0.300\linewidth}>{\centering\arraybackslash}X>{\raggedright\arraybackslash}p{0.540\linewidth}@{}}
\begin{minipage}[t]{\linewidth}
\centering\textbf{Earlier Germanic changes}\par
\vspace{0.35em}
\raggedright
\centering\textbf{West Germanic}\par
\raggedright
\vspace{0.2em}
\raggedright [no change]\par
\vspace{0.6em}
\centering\textbf{Northwest Germanic}\par
\raggedright
\vspace{0.2em}
\begin{tabular}{@{}>{\raggedright\arraybackslash}p{0.74\linewidth}@{\hspace{0.55em}}>{\raggedright\arraybackslash}p{0.16\linewidth}@{\hspace{0.25em}}}
\mbox{PGmc Final Z Deletion} & \emph{*xélma} \\
\end{tabular}
\end{minipage}
&
&
\begin{minipage}[t]{\linewidth}
\centering\textbf{Old English changes}\par
\vspace{0.35em}
\raggedright
\begin{tabular}{@{}>{\raggedright\arraybackslash}p{0.74\linewidth}@{\hspace{0.55em}}>{\raggedright\arraybackslash}p{0.16\linewidth}@{\hspace{0.25em}}}
PWGmc Final Bare A Loss & \emph{*xélm} \\
OE Velar Fricative Palatalization & \emph{*çélm} \\
\end{tabular}
\end{minipage}
\\
\end{tabularx}
\end{minipage}%
} \endgroup

Kroonen cites the helmet noun as \emph{*helma-} and separately
distinguishes a different lexeme \emph{*helman-} `rudder'
(\citeproc{ref-Kroonen2013}{Kroonen 2013, 259}). The selected input
\emph{*xélmaz} is the nominative-style form used for the helmet noun
itself. Clark Hall and Bosworth-Toller record \emph{helm} as the
ordinary Old English noun for `helmet', while \emph{helma} belongs to a
separate rudder lexeme (\citeproc{ref-ClarkHall1960}{Clark Hall 1960,
168}; \citeproc{ref-BosworthToller1898}{Bosworth and Toller 1898, 542}).
From \emph{*xélmaz}, loss of final \emph{z} and later loss of the short
final vowel yield \emph{helm}.

Form note. Comparative \emph{*helma-} is headword notation for the
helmet cognate set. It should not be confused with Old English
\emph{helma}, which is a different noun meaning `helm, rudder'
(\citeproc{ref-Kroonen2013}{Kroonen 2013, 259};
\citeproc{ref-ClarkHall1960}{Clark Hall 1960, 168}).

\subsubsection{help --- OE helpan}\label{help-oe-helpan}

Derivation: \emph{*xélpaną} \textgreater{} \emph{helpan} (regular).

\begingroup
\setlength{\fboxsep}{6pt}

\noindent\fbox{%
\begin{minipage}{0.97\linewidth}
\small
\begin{tabularx}{\linewidth}{@{}>{\raggedright\arraybackslash}p{0.300\linewidth}>{\centering\arraybackslash}X>{\raggedright\arraybackslash}p{0.540\linewidth}@{}}
\begin{minipage}[t]{\linewidth}
\centering\textbf{Earlier Germanic changes}\par
\vspace{0.35em}
\raggedright
\centering\textbf{West Germanic}\par
\raggedright
\vspace{0.2em}
\raggedright [no change]\par
\vspace{0.6em}
\centering\textbf{Northwest Germanic}\par
\raggedright
\vspace{0.2em}
\raggedright [no change]\par
\end{minipage}
&
&
\begin{minipage}[t]{\linewidth}
\centering\textbf{Old English changes}\par
\vspace{0.35em}
\raggedright
\begin{tabular}{@{}>{\raggedright\arraybackslash}p{0.70\linewidth}@{\hspace{0.55em}}>{\raggedright\arraybackslash}p{0.20\linewidth}@{\hspace{0.25em}}}
OE Velar Fricative Palatalization & \emph{*çélpaną} \\
OE Heavy Syllable Nasal Apocope & \emph{*çélpan} \\
OE Secondary Nasalization & \emph{*çélpąn} \\
OE Weak Tail Reduction & \emph{*çélpan} \\
\end{tabular}
\end{minipage}
\\
\end{tabularx}
\end{minipage}%
} \endgroup

The entry treats the strong verb itself rather than the separate noun
\emph{help}. Bright's principal parts \emph{helpan, healp, hulpon,
holpen} show the ordinary Old English strong-verb family continued by
this input (\citeproc{ref-BrightCassidyRingler1971}{Bright 1971, 57}).
Clark Hall and Bosworth-Toller record \emph{helpan} as the verbal
headword `to help' (\citeproc{ref-ClarkHall1960}{Clark Hall 1960, 168};
\citeproc{ref-BosworthToller1898}{Bosworth and Toller 1898, 542}). From
\emph{*xélpaną}, no special repair is needed beyond the ordinary
reduction of the infinitive ending.

Form note. Noun \emph{help} belongs to a separate lexical line and
should not replace verbal \emph{helpan} as the target here
(\citeproc{ref-ClarkHall1960}{Clark Hall 1960, 168};
\citeproc{ref-BosworthToller1898}{Bosworth and Toller 1898, 542}).

\subsubsection{hind --- OE hind}\label{hind-oe-hind}

Derivation: \emph{*xéndjō} \textgreater{} \emph{hind} (regular).

\begingroup
\setlength{\fboxsep}{6pt}

\noindent\fbox{%
\begin{minipage}{0.97\linewidth}
\small
\begin{tabularx}{\linewidth}{@{}>{\raggedright\arraybackslash}p{0.460\linewidth}>{\centering\arraybackslash}X>{\raggedright\arraybackslash}p{0.460\linewidth}@{}}
\begin{minipage}[t]{\linewidth}
\centering\textbf{Earlier Germanic changes}\par
\vspace{0.35em}
\raggedright
\centering\textbf{West Germanic}\par
\raggedright
\vspace{0.2em}
\raggedright [no change]\par
\vspace{0.6em}
\centering\textbf{Northwest Germanic}\par
\raggedright
\vspace{0.2em}
\begin{tabular}{@{}>{\raggedright\arraybackslash}p{0.74\linewidth}@{\hspace{0.55em}}>{\raggedright\arraybackslash}p{0.16\linewidth}@{\hspace{0.25em}}}
\mbox{NWGmc Final Long O Raising} & \emph{*xéndju} \\
\end{tabular}
\end{minipage}
&
&
\begin{minipage}[t]{\linewidth}
\centering\textbf{Old English changes}\par
\vspace{0.35em}
\raggedright
\begin{tabular}{@{}>{\raggedright\arraybackslash}p{0.74\linewidth}@{\hspace{0.55em}}>{\raggedright\arraybackslash}p{0.16\linewidth}@{\hspace{0.25em}}}
OE Velar Fricative Palatalization & \emph{*çéndju} \\
OE I Umlaut & \emph{*çindju} \\
\mbox{OE High Vowel Apocope} & \emph{*çindj} \\
OE J Loss After Heavy & \emph{*çind} \\
\end{tabular}
\end{minipage}
\\
\end{tabularx}
\end{minipage}%
} \endgroup

Kroonen cites the animal name as \emph{*hindō-} f.~`hind'
(\citeproc{ref-Kroonen2013}{Kroonen 2013, 266}). The selected input
\emph{*xéndjō} represents that same noun in the present derivation.
Clark Hall and Bosworth-Toller record \emph{hind} as the noun `hind,
female deer' (\citeproc{ref-ClarkHall1960}{Clark Hall 1960, 173};
\citeproc{ref-BosworthToller1898}{Bosworth and Toller 1898, 554}). From
\emph{*xéndjō}, i-mutation produces the front-vocalic Old English stem,
and later apocope plus loss of final \emph{j} yield \emph{hind}.

Form note. \emph{hindan} `from behind, behind' is a different Old
English lexeme and does not belong to the noun history of \emph{hind}
(\citeproc{ref-ClarkHall1960}{Clark Hall 1960, 173};
\citeproc{ref-BosworthToller1898}{Bosworth and Toller 1898, 554}).

\subsubsection{hold --- OE healdan}\label{hold-oe-healdan}

Derivation: \emph{*xáldaną} \textgreater{} \emph{healdan} (regular).

\begingroup
\setlength{\fboxsep}{6pt}

\noindent\fbox{%
\begin{minipage}{0.97\linewidth}
\footnotesize
\begin{tabularx}{\linewidth}{@{}>{\raggedright\arraybackslash}p{0.280\linewidth}>{\centering\arraybackslash}X>{\raggedright\arraybackslash}p{0.560\linewidth}@{}}
\begin{minipage}[t]{\linewidth}
\centering\textbf{Earlier Germanic changes}\par
\vspace{0.35em}
\raggedright
\centering\textbf{West Germanic}\par
\raggedright
\vspace{0.2em}
\raggedright [no change]\par
\vspace{0.6em}
\centering\textbf{Northwest Germanic}\par
\raggedright
\vspace{0.2em}
\raggedright [no change]\par
\end{minipage}
&
&
\begin{minipage}[t]{\linewidth}
\centering\textbf{Old English changes}\par
\vspace{0.35em}
\raggedright
\begin{tabular}{@{}>{\raggedright\arraybackslash}p{0.66\linewidth}@{\hspace{0.55em}}>{\raggedright\arraybackslash}p{0.24\linewidth}@{\hspace{0.25em}}}
Anglo Frisian Brightening & \emph{*xældaną} \\
OE Breaking & \emph{*xealdaną} \\
OE Velar Fricative Palatalization & \emph{*çealdaną} \\
OE Heavy Syllable Nasal Apocope & \emph{*çealdan} \\
OE Secondary Nasalization & \emph{*çealdąn} \\
OE Weak Tail Reduction & \emph{*çealdan} \\
\end{tabular}
\end{minipage}
\\
\end{tabularx}
\end{minipage}%
} \endgroup

Campbell and Ringe-Taylor treat the verb as a regular \emph{*a} + lC
breaking case, with West Saxon \emph{healdan} opposed to Anglian and
Mercian \emph{haldan} (\citeproc{ref-Campbell1959}{Campbell 1959, sect.
144}; \citeproc{ref-RingeTaylor2014}{Ringe and Taylor 2014, 199}).
Bright gives the ordinary strong-verb citation form and principal parts
\emph{healdan, heold, heoldon, healden}
(\citeproc{ref-BrightCassidyRingler1971}{Bright 1971, 62}). From
\emph{*xáldaną}, Anglo-Frisian brightening first yields a fronted vowel,
and West Saxon breaking then produces \emph{ea} before \emph{ld}.

Dialect note. West Saxon \emph{healdan} is the selected target here.
Anglian and Mercian \emph{haldan} are genuine non-West-Saxon doublets
rather than corrections to that choice
(\citeproc{ref-Campbell1959}{Campbell 1959, sect. 144};
\citeproc{ref-RingeTaylor2014}{Ringe and Taylor 2014, 199}).

\subsubsection{horn --- OE horn}\label{horn-oe-horn}

Derivation: \emph{*xúrną} \textgreater{} \emph{horn} (regular).

\begingroup
\setlength{\fboxsep}{6pt}

\noindent\fbox{%
\begin{minipage}{0.97\linewidth}
\small
\begin{tabularx}{\linewidth}{@{}>{\raggedright\arraybackslash}p{0.460\linewidth}>{\centering\arraybackslash}X>{\raggedright\arraybackslash}p{0.460\linewidth}@{}}
\begin{minipage}[t]{\linewidth}
\centering\textbf{Earlier Germanic changes}\par
\vspace{0.35em}
\raggedright
\centering\textbf{West Germanic}\par
\raggedright
\vspace{0.2em}
\raggedright [no change]\par
\vspace{0.6em}
\centering\textbf{Northwest Germanic}\par
\raggedright
\vspace{0.2em}
\begin{tabular}{@{}>{\raggedright\arraybackslash}p{0.74\linewidth}@{\hspace{0.55em}}>{\raggedright\arraybackslash}p{0.16\linewidth}@{\hspace{0.25em}}}
\mbox{NWGmc U Lowering} & \emph{*xórną} \\
\end{tabular}
\end{minipage}
&
&
\begin{minipage}[t]{\linewidth}
\centering\textbf{Old English changes}\par
\vspace{0.35em}
\raggedright
\begin{tabular}{@{}>{\raggedright\arraybackslash}p{0.74\linewidth}@{\hspace{0.55em}}>{\raggedright\arraybackslash}p{0.16\linewidth}@{\hspace{0.25em}}}
OE Heavy Syllable Nasal Apocope & \emph{*xórn} \\
\end{tabular}
\end{minipage}
\\
\end{tabularx}
\end{minipage}%
} \endgroup

Kroonen and Orel cite lemma-style Proto-Germanic headwords of the
\emph{*hurna-} / \emph{*xurnan} type for this noun
(\citeproc{ref-Kroonen2013}{Kroonen 2013, 299};
\citeproc{ref-Orel2003}{Orel 2003, 234}). The selected input
\emph{*xúrną} is the nominative-style form used in the derivation here.
Clark Hall, Bosworth-Toller, and Bright all record \emph{horn} as the
ordinary Old English noun (\citeproc{ref-ClarkHall1960}{Clark Hall 1960,
179}; \citeproc{ref-BosworthToller1898}{Bosworth and Toller 1898, 108};
\citeproc{ref-BrightCassidyRingler1971}{Bright 1971, 400}). From
\emph{*xúrną}, Northwest Germanic u-lowering gives \emph{*xórną}, and
later loss of the final nasal leaves \emph{horn}.

Form note. The note's oblique \emph{*xurnăn} belongs to comparative stem
background only. It does not replace the selected input \emph{*xúrną} as
the derivational form used here (\citeproc{ref-Kroonen2013}{Kroonen
2013}; \citeproc{ref-Orel2003}{Orel 2003, 234}).

\subsubsection{lead --- OE lǣdan}\label{lead-oe-lux1e3dan}

Derivation: \emph{*láidijaną} \textgreater{} \emph{lǣdan} (regular).

\begingroup
\setlength{\fboxsep}{6pt}

\noindent\fbox{%
\begin{minipage}{0.97\linewidth}
\small
\begin{tabularx}{\linewidth}{@{}>{\raggedright\arraybackslash}p{0.280\linewidth}>{\centering\arraybackslash}X>{\raggedright\arraybackslash}p{0.560\linewidth}@{}}
\begin{minipage}[t]{\linewidth}
\centering\textbf{Earlier Germanic changes}\par
\vspace{0.35em}
\raggedright
\centering\textbf{West Germanic}\par
\raggedright
\vspace{0.2em}
\begin{tabular}{@{}>{\raggedright\arraybackslash}p{0.66\linewidth}@{\hspace{0.55em}}>{\raggedright\arraybackslash}p{0.22\linewidth}@{\hspace{0.25em}}}
PWGmc Ai Monophthongization & \emph{*lādijaną} \\
\end{tabular}
\vspace{0.6em}
\centering\textbf{Northwest Germanic}\par
\raggedright
\vspace{0.2em}
\raggedright [no change]\par
\end{minipage}
&
&
\begin{minipage}[t]{\linewidth}
\centering\textbf{Old English changes}\par
\vspace{0.35em}
\raggedright
\begin{tabular}{@{}>{\raggedright\arraybackslash}p{0.70\linewidth}@{\hspace{0.55em}}>{\raggedright\arraybackslash}p{0.20\linewidth}@{\hspace{0.25em}}}
OE Heavy Syllable Nasal Apocope & \emph{*lādijan} \\
OE Secondary Nasalization & \emph{*lādijąn} \\
Sievers Law Syncope & \emph{*lādjąn} \\
OE I Umlaut & \emph{*lǣdjąn} \\
OE Weak Tail Reduction & \emph{*lǣdjan} \\
OE J Loss After Heavy & \emph{*lǣdan} \\
\end{tabular}
\end{minipage}
\\
\end{tabularx}
\end{minipage}%
} \endgroup

Ringe and Taylor derive Old English \emph{lǣdan} from Proto-Germanic
\emph{*laidijaną}, and Kroonen likewise cites a weak verb of the
\emph{*laidjan-} type for `lead' (\citeproc{ref-RingeTaylor2014}{Ringe
and Taylor 2014, 249}; \citeproc{ref-Kroonen2013}{Kroonen 2013, 363}).
Clark Hall and Bosworth-Toller both record \emph{lædan} / \emph{lǣdan}
as the ordinary Old English verb `to lead, guide, conduct'
(\citeproc{ref-ClarkHall1960}{Clark Hall 1960, 194};
\citeproc{ref-BosworthToller1898}{Bosworth and Toller 1898}). The
selected form is the ordinary infinitive citation form. From
\emph{*láidijaną}, monophthongization of \emph{*ai} first gives a
\emph{*lād-} stage, and later syncope, i-mutation, weak-tail reduction,
and loss of \emph{j} after a heavy stem yield \emph{lǣdan}
(\citeproc{ref-RingeTaylor2014}{Ringe and Taylor 2014, 249}).

\subsubsection{learn --- OE liornian}\label{learn-oe-liornian}

Derivation: \emph{*líznōjaną} \textgreater{} \emph{liornian} (regular).

\begingroup
\setlength{\fboxsep}{6pt}

\noindent\fbox{%
\begin{minipage}{0.97\linewidth}
\footnotesize
\begin{tabularx}{\linewidth}{@{}>{\raggedright\arraybackslash}p{0.280\linewidth}>{\centering\arraybackslash}X>{\raggedright\arraybackslash}p{0.560\linewidth}@{}}
\begin{minipage}[t]{\linewidth}
\centering\textbf{Earlier Germanic changes}\par
\vspace{0.35em}
\raggedright
\centering\textbf{West Germanic}\par
\raggedright
\vspace{0.2em}
\raggedright [no change]\par
\vspace{0.6em}
\centering\textbf{Northwest Germanic}\par
\raggedright
\vspace{0.2em}
\raggedright [no change]\par
\end{minipage}
&
&
\begin{minipage}[t]{\linewidth}
\centering\textbf{Old English changes}\par
\vspace{0.35em}
\raggedright
\begin{tabular}{@{}>{\raggedright\arraybackslash}p{0.62\linewidth}@{\hspace{0.55em}}>{\raggedright\arraybackslash}p{0.28\linewidth}@{\hspace{0.25em}}}
OE Breaking & \emph{*líornōjaną} \\
OE Heavy Syllable Nasal Apocope & \emph{*líornōjan} \\
OE Secondary Nasalization & \emph{*líornōjąn} \\
OE I Umlaut & \emph{*líornējąn} \\
OE Unstressed Long Vowel Shortening & \emph{*líornejąn} \\
OE Weak Tail Reduction & \emph{*líornejan} \\
OE Intervocalic J Vocalization & \emph{*líorneian} \\
OE Unstressed EI Contraction & \emph{*líornian} \\
\end{tabular}
\end{minipage}
\\
\end{tabularx}
\end{minipage}%
} \endgroup

Ringe and Taylor place the verb in a class-II weak family of the
\emph{*liznō-} type, and Kroonen keeps the same comparative base
(\citeproc{ref-RingeTaylor2014}{Ringe and Taylor 2014, 318};
\citeproc{ref-Kroonen2013}{Kroonen 2013, 328}). Campbell records
Northumbrian \emph{liornian} and treats the \emph{eo} of \emph{leornian}
as secondary (\citeproc{ref-Campbell1959}{Campbell 1959, 123};
\citeproc{ref-Campbell1959}{1959, 296}). The Northumbrian spelling
preserves \emph{io} where original \emph{eo} and \emph{io} remain
distinct. The selected target \emph{liornian} therefore represents the
Northumbrian regular form.

Dialect note. Clark Hall and Bright give \emph{leornian} as the usual
headword (\citeproc{ref-ClarkHall1960}{Clark Hall 1960, 193};
\citeproc{ref-BrightCassidyRingler1971}{Bright 1971, 58}). From
\emph{*líznōjaną}, the regular development yields \emph{liornian};
\emph{leornian} reflects the later \emph{eo} development.

\subsubsection{lid --- OE hlid}\label{lid-oe-hlid}

Derivation: \emph{*xlídą} \textgreater{} \emph{hlid} (regular).

\begingroup
\setlength{\fboxsep}{6pt}

\noindent\fbox{%
\begin{minipage}{0.97\linewidth}
\small
\begin{tabularx}{\linewidth}{@{}>{\raggedright\arraybackslash}p{0.300\linewidth}>{\centering\arraybackslash}X>{\raggedright\arraybackslash}p{0.540\linewidth}@{}}
\begin{minipage}[t]{\linewidth}
\centering\textbf{Earlier Germanic changes}\par
\vspace{0.35em}
\raggedright
\centering\textbf{West Germanic}\par
\raggedright
\vspace{0.2em}
\raggedright [no change]\par
\vspace{0.6em}
\centering\textbf{Northwest Germanic}\par
\raggedright
\vspace{0.2em}
\raggedright [no change]\par
\end{minipage}
&
&
\begin{minipage}[t]{\linewidth}
\centering\textbf{Old English changes}\par
\vspace{0.35em}
\raggedright
\begin{tabular}{@{}>{\raggedright\arraybackslash}p{0.74\linewidth}@{\hspace{0.55em}}>{\raggedright\arraybackslash}p{0.16\linewidth}@{\hspace{0.25em}}}
OE Heavy Syllable Nasal Apocope & \emph{*xlíd} \\
\end{tabular}
\end{minipage}
\\
\end{tabularx}
\end{minipage}%
} \endgroup

Orel cites a neuter lexeme of the \emph{*xliđ-} type with Old English
\emph{hlid}, and Lloyd includes OE \emph{hlid} beside ON \emph{hliþó}
and OHG \emph{(h)lit} among forms that retain \emph{i}
(\citeproc{ref-Orel2003}{Orel 2003, 216}; \citeproc{ref-Lloyd1966}{Lloyd
1966, 738}). Clark Hall and Bosworth-Toller record \emph{hlid} as the
noun `lid, cover, door, gate' (\citeproc{ref-ClarkHall1960}{Clark Hall
1960, 175}; \citeproc{ref-BosworthToller1898}{Bosworth and Toller 1898,
563}). The selected input already represents the later Germanic
\emph{hliđ-} stage used for the derivation here.

Form note. An earlier etymological stage \emph{*liþuz} belongs to
comparative background only. The form represented here is the later
\emph{*xlídą} \textgreater{} hlid line that matches the attested Old
English noun (\citeproc{ref-Orel2003}{Orel 2003, 216};
\citeproc{ref-Lloyd1966}{Lloyd 1966, 738}).

\subsubsection{light --- OE līehtan}\label{light-oe-lux12behtan}

Derivation: \emph{*léuxtijaną} \textgreater{} \emph{līehtan} (regular).

\begingroup
\setlength{\fboxsep}{6pt}

\noindent\fbox{%
\begin{minipage}{0.97\linewidth}
\footnotesize
\begin{tabularx}{\linewidth}{@{}>{\raggedright\arraybackslash}p{0.280\linewidth}>{\centering\arraybackslash}X>{\raggedright\arraybackslash}p{0.560\linewidth}@{}}
\begin{minipage}[t]{\linewidth}
\centering\textbf{Earlier Germanic changes}\par
\vspace{0.35em}
\raggedright
\centering\textbf{West Germanic}\par
\raggedright
\vspace{0.2em}
\raggedright [no change]\par
\vspace{0.6em}
\centering\textbf{Northwest Germanic}\par
\raggedright
\vspace{0.2em}
\raggedright [no change]\par
\end{minipage}
&
&
\begin{minipage}[t]{\linewidth}
\centering\textbf{Old English changes}\par
\vspace{0.35em}
\raggedright
\begin{tabular}{@{}>{\raggedright\arraybackslash}p{0.62\linewidth}@{\hspace{0.55em}}>{\raggedright\arraybackslash}p{0.28\linewidth}@{\hspace{0.25em}}}
OE Diphthong Leveling & \emph{*lēoxtijaną} \\
OE Heavy Syllable Nasal Apocope & \emph{*lēoxtijan} \\
OE Secondary Nasalization & \emph{*lēoxtijąn} \\
Sievers Law Syncope & \emph{*lēoxtjąn} \\
OE I Umlaut & \emph{*līextjąn} \\
OE Weak Tail Reduction & \emph{*līextjan} \\
OE J Loss After Heavy & \emph{*līextan} \\
\end{tabular}
\end{minipage}
\\
\end{tabularx}
\end{minipage}%
} \endgroup

Fulk gives Proto-Germanic \emph{*liuxtijanan} with Old English
\emph{līehtan} `illuminate', and Ringe and Taylor likewise derive West
Saxon \emph{liehtan} from the same weak-verb formation
(\citeproc{ref-Fulk2018}{Fulk 2018, 81};
\citeproc{ref-RingeTaylor2014}{Ringe and Taylor 2014, 264}). Clark Hall
and Bosworth-Toller preserve the verb family under spellings such as
\emph{liehtan}, \emph{lihtan}, and \emph{līhtan}, distinct from the
related noun \emph{lēoht} and adjective \emph{leoht/liht}
(\citeproc{ref-ClarkHall1960}{Clark Hall 1960, 204};
\citeproc{ref-BosworthToller1898}{Bosworth and Toller 1898}). From
\emph{*léuxtijaną}, the regular verbal line preserves \emph{*xt}, passes
through a West Saxon \emph{liehtan} stage, and is represented here by
normalized \emph{līehtan}.

Dialect note. Ringe and Taylor and Campbell distinguish West Saxon
\emph{liehtan} from Anglian \emph{lihtan}, while later West Saxon also
shows \emph{lyhtan} (\citeproc{ref-RingeTaylor2014}{Ringe and Taylor
2014, 264}; \citeproc{ref-Campbell1959}{Campbell 1959, sect. 39}).

\subsubsection{linden --- OE lind}\label{linden-oe-lind}

Derivation: \emph{*líndō} \textgreater{} \emph{lind} (regular).

\begingroup
\setlength{\fboxsep}{6pt}

\noindent\fbox{%
\begin{minipage}{0.97\linewidth}
\small
\begin{tabularx}{\linewidth}{@{}>{\raggedright\arraybackslash}p{0.460\linewidth}>{\centering\arraybackslash}X>{\raggedright\arraybackslash}p{0.460\linewidth}@{}}
\begin{minipage}[t]{\linewidth}
\centering\textbf{Earlier Germanic changes}\par
\vspace{0.35em}
\raggedright
\centering\textbf{West Germanic}\par
\raggedright
\vspace{0.2em}
\raggedright [no change]\par
\vspace{0.6em}
\centering\textbf{Northwest Germanic}\par
\raggedright
\vspace{0.2em}
\begin{tabular}{@{}>{\raggedright\arraybackslash}p{0.74\linewidth}@{\hspace{0.55em}}>{\raggedright\arraybackslash}p{0.16\linewidth}@{\hspace{0.25em}}}
\mbox{NWGmc Final Long O Raising} & \emph{*líndu} \\
\end{tabular}
\end{minipage}
&
&
\begin{minipage}[t]{\linewidth}
\centering\textbf{Old English changes}\par
\vspace{0.35em}
\raggedright
\begin{tabular}{@{}>{\raggedright\arraybackslash}p{0.74\linewidth}@{\hspace{0.55em}}>{\raggedright\arraybackslash}p{0.16\linewidth}@{\hspace{0.25em}}}
\mbox{OE High Vowel Apocope} & \emph{*línd} \\
\end{tabular}
\end{minipage}
\\
\end{tabularx}
\end{minipage}%
} \endgroup

Kroonen reconstructs \emph{*lindō-}, and Clark Hall and Bosworth-Toller
record Old English \emph{lind} `linden, lime-tree'
(\citeproc{ref-Kroonen2013}{Kroonen 2013, 337};
\citeproc{ref-ClarkHall1960}{Clark Hall 1960, 198};
\citeproc{ref-BosworthToller1898}{Bosworth and Toller 1898, 630}). The
selected input \emph{*líndō} is the citation-form noun used here.
Northwest Germanic final \emph{ō} raises to \emph{*líndu}, and
high-vowel apocope yields \emph{lind}. Clark Hall also gives adjectival
\emph{linden} `made of linden-wood'; that is a related adjective, not
the noun treated here (\citeproc{ref-ClarkHall1960}{Clark Hall 1960,
198}).

\subsubsection{milk --- OE meoloc}\label{milk-oe-meoloc}

Derivation: \emph{*mélukz} \textgreater{} \emph{meoloc} (regular).

\begingroup
\setlength{\fboxsep}{6pt}

\noindent\fbox{%
\begin{minipage}{0.97\linewidth}
\small
\begin{tabularx}{\linewidth}{@{}>{\raggedright\arraybackslash}p{0.300\linewidth}>{\centering\arraybackslash}X>{\raggedright\arraybackslash}p{0.540\linewidth}@{}}
\begin{minipage}[t]{\linewidth}
\centering\textbf{Earlier Germanic changes}\par
\vspace{0.35em}
\raggedright
\centering\textbf{West Germanic}\par
\raggedright
\vspace{0.2em}
\raggedright [no change]\par
\vspace{0.6em}
\centering\textbf{Northwest Germanic}\par
\raggedright
\vspace{0.2em}
\begin{tabular}{@{}>{\raggedright\arraybackslash}p{0.74\linewidth}@{\hspace{0.55em}}>{\raggedright\arraybackslash}p{0.16\linewidth}@{\hspace{0.25em}}}
\mbox{PGmc Final Z Deletion} & \emph{*méluk} \\
\end{tabular}
\end{minipage}
&
&
\begin{minipage}[t]{\linewidth}
\centering\textbf{Old English changes}\par
\vspace{0.35em}
\raggedright
\begin{tabular}{@{}>{\raggedright\arraybackslash}p{0.74\linewidth}@{\hspace{0.55em}}>{\raggedright\arraybackslash}p{0.16\linewidth}@{\hspace{0.25em}}}
OE Med Unstressed U Lowering & \emph{*mélok} \\
OE Back Mutation & \emph{*méolok} \\
\end{tabular}
\end{minipage}
\\
\end{tabularx}
\end{minipage}%
} \endgroup

Kroonen and Orel reconstruct the noun as \emph{*meluk-} /
\emph{*melukz}, and the nominative-style input used here is
\emph{*mélukz} (\citeproc{ref-Kroonen2013}{Kroonen 2013, 404};
\citeproc{ref-Orel2003}{Orel 2003, 306}). Old English preserves a mixed
dossier for this noun. The unsyncopated line from \emph{*mélukz} loses
final \emph{*z}, lowers unstressed \emph{u} to \emph{o}, and with back
mutation yields \emph{meoloc}.

Comparison note. Campbell explicitly records \emph{meoluc} with dative
singular \emph{meoloc / meoloce} and contrasts that fuller line with
Anglian \emph{milc} (\citeproc{ref-Campbell1959}{Campbell 1959, sect.
574.5}). Syncopated \emph{meolc} and Anglian \emph{milc} belong to the
competing leveled tradition associated with oblique forms, whereas
\emph{meoloc / meoluc} preserves the fuller nominal shape
(\citeproc{ref-RingeTaylor2014}{Ringe and Taylor 2014, 268}).

\subsubsection{mother --- OE mōder}\label{mother-oe-mux14dder}

Derivation: \emph{*mōdēr} \textgreater{} \emph{mōder} (regular).

\begingroup
\setlength{\fboxsep}{6pt}

\noindent\fbox{%
\begin{minipage}{0.97\linewidth}
\small
\begin{tabularx}{\linewidth}{@{}>{\raggedright\arraybackslash}p{0.460\linewidth}>{\centering\arraybackslash}X>{\raggedright\arraybackslash}p{0.460\linewidth}@{}}
\begin{minipage}[t]{\linewidth}
\centering\textbf{Earlier Germanic changes}\par
\vspace{0.35em}
\raggedright
\centering\textbf{West Germanic}\par
\raggedright
\vspace{0.2em}
\raggedright [no change]\par
\vspace{0.6em}
\centering\textbf{Northwest Germanic}\par
\raggedright
\vspace{0.2em}
\begin{tabular}{@{}>{\raggedright\arraybackslash}p{0.74\linewidth}@{\hspace{0.55em}}>{\raggedright\arraybackslash}p{0.16\linewidth}@{\hspace{0.25em}}}
\mbox{NWGmc Long E Lowering} & \emph{*mōdǣr} \\
\end{tabular}
\end{minipage}
&
&
\begin{minipage}[t]{\linewidth}
\centering\textbf{Old English changes}\par
\vspace{0.35em}
\raggedright
\begin{tabular}{@{}>{\raggedright\arraybackslash}p{0.74\linewidth}@{\hspace{0.55em}}>{\raggedright\arraybackslash}p{0.16\linewidth}@{\hspace{0.25em}}}
OE Unstressed Long Vowel Shortening & \emph{*mōdær} \\
OE Unstressed AE Merger & \emph{*mōder} \\
\end{tabular}
\end{minipage}
\\
\end{tabularx}
\end{minipage}%
} \endgroup

Kroonen and Orel cite the Proto-Germanic r-stem as \emph{*mōdēr-} or
\emph{*mōdēr} (\citeproc{ref-Kroonen2013}{Kroonen 2013, 368};
\citeproc{ref-Orel2003}{Orel 2003, 274}). Clark Hall, Campbell, and
Ringe and Taylor preserve the Old English headword tradition \emph{mōdor
/ modor} with oblique \emph{mēder} in the paradigm
(\citeproc{ref-ClarkHall1960}{Clark Hall 1960, 219};
\citeproc{ref-Campbell1959}{Campbell 1959, 598};
\citeproc{ref-RingeTaylor2014}{Ringe and Taylor 2014, 312}). The
selected form here is nominative \emph{mōder}. From \emph{*mōdēr},
regular suffixal development yields that form.

Comparison note. \emph{Mōdor} is the usual dictionary form and reflects
later levelling within the r-stem paradigm. \emph{Mēder} preserves the
oblique line.

\subsubsection{net --- OE nett}\label{net-oe-nett}

Derivation: \emph{*nátją} \textgreater{} \emph{nett} (regular).

\begingroup
\setlength{\fboxsep}{6pt}

\noindent\fbox{%
\begin{minipage}{0.97\linewidth}
\small
\begin{tabularx}{\linewidth}{@{}>{\raggedright\arraybackslash}p{0.300\linewidth}>{\centering\arraybackslash}X>{\raggedright\arraybackslash}p{0.540\linewidth}@{}}
\begin{minipage}[t]{\linewidth}
\centering\textbf{Earlier Germanic changes}\par
\vspace{0.35em}
\raggedright
\centering\textbf{West Germanic}\par
\raggedright
\vspace{0.2em}
\begin{tabular}{@{}>{\raggedright\arraybackslash}p{0.64\linewidth}@{\hspace{0.55em}}>{\raggedright\arraybackslash}p{0.16\linewidth}@{\hspace{0.25em}}}
PWGmc J Gemination & \emph{*náttją} \\
\end{tabular}
\vspace{0.6em}
\centering\textbf{Northwest Germanic}\par
\raggedright
\vspace{0.2em}
\raggedright [no change]\par
\end{minipage}
&
&
\begin{minipage}[t]{\linewidth}
\centering\textbf{Old English changes}\par
\vspace{0.35em}
\raggedright
\begin{tabular}{@{}>{\raggedright\arraybackslash}p{0.74\linewidth}@{\hspace{0.55em}}>{\raggedright\arraybackslash}p{0.16\linewidth}@{\hspace{0.25em}}}
Anglo Frisian Brightening & \emph{*nættją} \\
OE Heavy Syllable Nasal Apocope & \emph{*nættj} \\
OE I Umlaut & \emph{*nettj} \\
OE J Loss After Heavy & \emph{*nett} \\
\end{tabular}
\end{minipage}
\\
\end{tabularx}
\end{minipage}%
} \endgroup

Orel gives \emph{*natjan} with Old English \emph{nett}, and Fulk's
account of West Germanic gemination before \emph{j} explains the
geminate outcome after a short vowel (\citeproc{ref-Orel2003}{Orel 2003,
282}; \citeproc{ref-Fulk2018}{Fulk 2018, sect. 6.15}). Bosworth-Toller
records \emph{nett} as the noun
(\citeproc{ref-BosworthToller1898}{Bosworth and Toller 1898, 29}), and
Campbell notes that final geminates are often graphically simplified in
Old English spelling (\citeproc{ref-Campbell1959}{Campbell 1959, sect.
66}). From \emph{*nátją}, West Germanic j-gemination first gives
\emph{*náttją}.

Form note. Spellings in \emph{net} can therefore be graphic
simplifications, but the lexical target supported by the dictionary
evidence is \emph{nett} (\citeproc{ref-Campbell1959}{Campbell 1959,
sect. 66}; \citeproc{ref-Orel2003}{Orel 2003, 282}).

\subsubsection{nightmare --- OE mare}\label{nightmare-oe-mare}

Derivation: \emph{*márōn} \textgreater{} \emph{mare} (regular).

\begingroup
\setlength{\fboxsep}{6pt}

\noindent\fbox{%
\begin{minipage}{0.97\linewidth}
\small
\begin{tabularx}{\linewidth}{@{}>{\raggedright\arraybackslash}p{0.300\linewidth}>{\centering\arraybackslash}X>{\raggedright\arraybackslash}p{0.540\linewidth}@{}}
\begin{minipage}[t]{\linewidth}
\centering\textbf{Earlier Germanic changes}\par
\vspace{0.35em}
\raggedright
\centering\textbf{West Germanic}\par
\raggedright
\vspace{0.2em}
\raggedright [no change]\par
\vspace{0.6em}
\centering\textbf{Northwest Germanic}\par
\raggedright
\vspace{0.2em}
\begin{tabular}{@{}>{\raggedright\arraybackslash}p{0.64\linewidth}@{\hspace{0.55em}}>{\raggedright\arraybackslash}p{0.16\linewidth}@{\hspace{0.25em}}}
NWGmc N Stem N Loss & \emph{*márǭ} \\
\end{tabular}
\end{minipage}
&
&
\begin{minipage}[t]{\linewidth}
\centering\textbf{Old English changes}\par
\vspace{0.35em}
\raggedright
\begin{tabular}{@{}>{\raggedright\arraybackslash}p{0.74\linewidth}@{\hspace{0.55em}}>{\raggedright\arraybackslash}p{0.16\linewidth}@{\hspace{0.25em}}}
Anglo Frisian Brightening & \emph{*mærǭ} \\
OE A Restoration & \emph{*marǭ} \\
OE Unstressed Long Vowel Shortening & \emph{*maræ} \\
OE Unstressed AE Merger & \emph{*mare} \\
\end{tabular}
\end{minipage}
\\
\end{tabularx}
\end{minipage}%
} \endgroup

Ringe and Taylor treat the lexeme as Proto-Germanic /
Proto-Northwest-Germanic \emph{*marōn-}, with Old English \emph{mare,
maran}, and variant \emph{mere}; Orel preserves the same comparative
lemma though with a different Old English headword tradition
(\citeproc{ref-RingeTaylor2014}{Ringe and Taylor 2014};
\citeproc{ref-Orel2003}{Orel 2003, 262}). Clark Hall records \emph{mare}
`nightmare, monster' and also preserves related variant forms \emph{mera
/ mere} (\citeproc{ref-ClarkHall1960}{Clark Hall 1960, 213}). The
selected simplex input \emph{*márōn} regularly gives \emph{mare} after
brightening, A-restoration before the n-stem ending, and later reduction
of the final vowel.

Form note. The concept corresponds to an unattested compound
\emph{*nihtmare}, but the Old English lexical evidence is for simplex
\emph{mare}, with oblique \emph{maran} and variant \emph{mere / mera}
(\citeproc{ref-RingeTaylor2014}{Ringe and Taylor 2014};
\citeproc{ref-ClarkHall1960}{Clark Hall 1960, 213}).

\subsubsection{coat --- OE rocc}\label{coat-oe-rocc}

Derivation: \emph{*rúkkaz} \textgreater{} \emph{rocc} (regular).

\begingroup
\setlength{\fboxsep}{6pt}

\noindent\fbox{%
\begin{minipage}{0.97\linewidth}
\small
\begin{tabularx}{\linewidth}{@{}>{\raggedright\arraybackslash}p{0.540\linewidth}>{\centering\arraybackslash}X>{\raggedright\arraybackslash}p{0.300\linewidth}@{}}
\begin{minipage}[t]{\linewidth}
\centering\textbf{Earlier Germanic changes}\par
\vspace{0.35em}
\raggedright
\centering\textbf{West Germanic}\par
\raggedright
\vspace{0.2em}
\raggedright [no change]\par
\vspace{0.6em}
\centering\textbf{Northwest Germanic}\par
\raggedright
\vspace{0.2em}
\begin{tabular}{@{}>{\raggedright\arraybackslash}p{0.74\linewidth}@{\hspace{0.55em}}>{\raggedright\arraybackslash}p{0.16\linewidth}@{\hspace{0.25em}}}
\mbox{NWGmc U Lowering} & \emph{*rókkaz} \\
\mbox{PGmc Final Z Deletion} & \emph{*rókka} \\
\end{tabular}
\end{minipage}
&
&
\begin{minipage}[t]{\linewidth}
\centering\textbf{Old English changes}\par
\vspace{0.35em}
\raggedright
\begin{tabular}{@{}>{\raggedright\arraybackslash}p{0.70\linewidth}@{\hspace{0.55em}}>{\raggedright\arraybackslash}p{0.16\linewidth}@{\hspace{0.25em}}}
PWGmc Final Bare A Loss & \emph{*rókk} \\
\end{tabular}
\end{minipage}
\\
\end{tabularx}
\end{minipage}%
} \endgroup

Orel cites a masculine \emph{*rukkaz} for the garment word, while
Kroonen gives \emph{*hrukka-}. Both treat this as the garment lexeme and
not as the separate stone word (\citeproc{ref-Orel2003}{Orel 2003, 347};
\citeproc{ref-Kroonen2013}{Kroonen 2013, 290}). Clark Hall and
Bosworth-Toller record \emph{rocc} as an over-garment or tunic and
preserve compounds such as \emph{bisceoprocc} and \emph{breóstrocc}
(\citeproc{ref-ClarkHall1960}{Clark Hall 1960, 259};
\citeproc{ref-BosworthToller1898}{Bosworth and Toller 1898}). With
\emph{*rúkkaz} as the selected input, Northwest Germanic u-lowering and
later loss of final \emph{-a} yield \emph{rocc} as a regular outcome.

Source note. This entry concerns the garment noun only. The stone word
seen in \emph{stānrocc} belongs to a different lexical history
(\citeproc{ref-ClarkHall1960}{Clark Hall 1960, 259};
\citeproc{ref-BosworthToller1898}{Bosworth and Toller 1898}).

\subsubsection{sheep --- OE sċēap}\label{sheep-oe-sux10bux113ap}

Derivation: \emph{*skḗpą} \textgreater{} \emph{sċēap} (regular).

\begingroup
\setlength{\fboxsep}{6pt}

\noindent\fbox{%
\begin{minipage}{0.97\linewidth}
\small
\begin{tabularx}{\linewidth}{@{}>{\raggedright\arraybackslash}p{0.460\linewidth}>{\centering\arraybackslash}X>{\raggedright\arraybackslash}p{0.460\linewidth}@{}}
\begin{minipage}[t]{\linewidth}
\centering\textbf{Earlier Germanic changes}\par
\vspace{0.35em}
\raggedright
\centering\textbf{West Germanic}\par
\raggedright
\vspace{0.2em}
\raggedright [no change]\par
\vspace{0.6em}
\centering\textbf{Northwest Germanic}\par
\raggedright
\vspace{0.2em}
\begin{tabular}{@{}>{\raggedright\arraybackslash}p{0.74\linewidth}@{\hspace{0.55em}}>{\raggedright\arraybackslash}p{0.16\linewidth}@{\hspace{0.25em}}}
\mbox{NWGmc Long E Lowering} & \emph{*skǣpą} \\
\end{tabular}
\end{minipage}
&
&
\begin{minipage}[t]{\linewidth}
\centering\textbf{Old English changes}\par
\vspace{0.35em}
\raggedright
\begin{tabular}{@{}>{\raggedright\arraybackslash}p{0.74\linewidth}@{\hspace{0.55em}}>{\raggedright\arraybackslash}p{0.16\linewidth}@{\hspace{0.25em}}}
OE Heavy Syllable Nasal Apocope & \emph{*skǣp} \\
OE Sk Palatalization & \emph{*ʃǣp} \\
OE Ws Palatal Diphthongization & \emph{*ʃēap} \\
\end{tabular}
\end{minipage}
\\
\end{tabularx}
\end{minipage}%
} \endgroup

Ringe and Taylor cite a later West Germanic \emph{*skap} \textgreater{}
WS \emph{scéap}, while Orel preserves a Proto-Germanic noun of the
\emph{*skēp-} type for the same lexeme
(\citeproc{ref-RingeTaylor2014}{Ringe and Taylor 2014, 142};
\citeproc{ref-Orel2003}{Orel 2003, 340}). Clark Hall records
\emph{scēap} with spelling variation, and Campbell likewise lists West
Saxon \emph{scéap} among the palatal-diphthongized forms
(\citeproc{ref-ClarkHall1960}{Clark Hall 1960, 266};
\citeproc{ref-Campbell1959}{Campbell 1959}). From \emph{*skḗpą},
Northwest Germanic lowering gives \emph{*skǣpą}; after apocope and
palatalization the West Saxon branch diphthongizes to \emph{sċēap}.

Dialect note. Ringe and Taylor contrast West Saxon \emph{scéap} with
Mercian and Kentish \emph{scép}, and Campbell also notes Northumbrian
\emph{scip}. The form represented here is the West Saxon headword
(\citeproc{ref-RingeTaylor2014}{Ringe and Taylor 2014, 142};
\citeproc{ref-Campbell1959}{Campbell 1959, sect. 186}).

\subsubsection{shilling --- OE sċilling}\label{shilling-oe-sux10billing}

Derivation: \emph{*skíllingaz} \textgreater{} \emph{sċilling} (regular).

\begingroup
\setlength{\fboxsep}{6pt}

\noindent\fbox{%
\begin{minipage}{0.97\linewidth}
\small
\begin{tabularx}{\linewidth}{@{}>{\raggedright\arraybackslash}p{0.300\linewidth}>{\centering\arraybackslash}X>{\raggedright\arraybackslash}p{0.540\linewidth}@{}}
\begin{minipage}[t]{\linewidth}
\centering\textbf{Earlier Germanic changes}\par
\vspace{0.35em}
\raggedright
\centering\textbf{West Germanic}\par
\raggedright
\vspace{0.2em}
\raggedright [no change]\par
\vspace{0.6em}
\centering\textbf{Northwest Germanic}\par
\raggedright
\vspace{0.2em}
\begin{tabular}{@{}>{\raggedright\arraybackslash}p{0.66\linewidth}@{\hspace{0.55em}}>{\raggedright\arraybackslash}p{0.24\linewidth}@{\hspace{0.25em}}}
PGmc Final Z Deletion & \emph{*skíllinga} \\
\end{tabular}
\end{minipage}
&
&
\begin{minipage}[t]{\linewidth}
\centering\textbf{Old English changes}\par
\vspace{0.35em}
\raggedright
\begin{tabular}{@{}>{\raggedright\arraybackslash}p{0.68\linewidth}@{\hspace{0.55em}}>{\raggedright\arraybackslash}p{0.22\linewidth}@{\hspace{0.25em}}}
PWGmc Final Bare A Loss & \emph{*skílling} \\
OE Sk Palatalization & \emph{*ʃílling} \\
OE Med Unstressed I Lowering1 & \emph{*ʃílleng} \\
OE Med Unstressed I Lowering & \emph{*ʃílling} \\
\end{tabular}
\end{minipage}
\\
\end{tabularx}
\end{minipage}%
} \endgroup

Kroonen treats the cognate set under \emph{*skellinga-}
\textasciitilde{} \emph{*skillinga-} and connects it with
\emph{*skeld-linga-}, while Orel likewise gives the coin word with OE
\emph{scilling} among the reflexes (\citeproc{ref-Kroonen2013}{Kroonen
2013}; \citeproc{ref-Orel2003}{Orel 2003, 377}). The selected input
\emph{*skíllingaz} is the nominative-style form used here to represent
that inherited \emph{*-ing-} derivative. Clark Hall records
\emph{scilling}, and Campbell cites \emph{scilling} among nouns whose
derivational \emph{-ing} keeps \emph{i} in unstressed syllables
(\citeproc{ref-ClarkHall1960}{Clark Hall 1960, 269};
\citeproc{ref-Campbell1959}{Campbell 1959, sect. 574.6}). From
\emph{*skíllingaz}, loss of final \emph{-az} yields \emph{*skílling}.

Form note. Kroonen's \emph{*skellinga-} \textasciitilde{}
\emph{*skillinga-} and his internal analysis \emph{*skeld-linga-} belong
to the etymological background of the cognate set. The selected input
\emph{*skíllingaz} is the specific form used for the derivation
represented here (\citeproc{ref-Kroonen2013}{Kroonen 2013}).

\subsubsection{show --- OE sċēawian}\label{show-oe-sux10bux113awian}

Derivation: \emph{*skáwōjaną} \textgreater{} \emph{sċēawian} (regular).

\begingroup
\setlength{\fboxsep}{6pt}

\noindent\fbox{%
\begin{minipage}{0.97\linewidth}
\footnotesize
\begin{tabularx}{\linewidth}{@{}>{\raggedright\arraybackslash}p{0.280\linewidth}>{\centering\arraybackslash}X>{\raggedright\arraybackslash}p{0.560\linewidth}@{}}
\begin{minipage}[t]{\linewidth}
\centering\textbf{Earlier Germanic changes}\par
\vspace{0.35em}
\raggedright
\centering\textbf{West Germanic}\par
\raggedright
\vspace{0.2em}
\raggedright [no change]\par
\vspace{0.6em}
\centering\textbf{Northwest Germanic}\par
\raggedright
\vspace{0.2em}
\raggedright [no change]\par
\end{minipage}
&
&
\begin{minipage}[t]{\linewidth}
\centering\textbf{Old English changes}\par
\vspace{0.35em}
\raggedright
\begin{tabular}{@{}>{\raggedright\arraybackslash}p{0.62\linewidth}@{\hspace{0.55em}}>{\raggedright\arraybackslash}p{0.28\linewidth}@{\hspace{0.25em}}}
OE Aw Long Diphthong & \emph{*skḗawōjaną} \\
OE Heavy Syllable Nasal Apocope & \emph{*skḗawōjan} \\
OE Secondary Nasalization & \emph{*skḗawōjąn} \\
OE Sk Palatalization & \emph{*ʃḗawōjąn} \\
OE I Umlaut & \emph{*ʃḗawējąn} \\
OE Unstressed Long Vowel Shortening & \emph{*ʃḗawejąn} \\
OE Weak Tail Reduction & \emph{*ʃḗawejan} \\
OE Intervocalic J Vocalization & \emph{*ʃḗaweian} \\
OE Unstressed EI Contraction & \emph{*ʃḗawian} \\
\end{tabular}
\end{minipage}
\\
\end{tabularx}
\end{minipage}%
} \endgroup

Orel and Kroonen place the verb in the ordinary class-II show-family,
and Brunner records Old English \emph{scēawian}, \emph{scāwian}
(\citeproc{ref-Orel2003}{Orel 2003, 339};
\citeproc{ref-Kroonen2013}{Kroonen 2013, 482};
\citeproc{ref-SieversBrunner1965}{Sievers 1965, 148}). From
\emph{*skáwōjaną}, \emph{aw} before a following vowel yields Old English
\emph{ēaw}, while the class-II suffix preserves \emph{ō} between
\emph{w} and \emph{j}. The result is \emph{sċēawian}
(\citeproc{ref-Campbell1959}{Campbell 1959, 138};
\citeproc{ref-Orel2003}{Orel 2003, 339}).

Form note. The spelling \emph{sċēawian} normalizes initial in dictionary
\emph{scēawian} (\citeproc{ref-Campbell1959}{Campbell 1959, 19};
\citeproc{ref-Hogg1992}{Hogg 1992, 236}).

\subsubsection{sleep --- OE slǣpan}\label{sleep-oe-slux1e3pan}

Derivation: \emph{*slḗpaną} \textgreater{} \emph{slǣpan} (regular).

\begingroup
\setlength{\fboxsep}{6pt}

\noindent\fbox{%
\begin{minipage}{0.97\linewidth}
\small
\begin{tabularx}{\linewidth}{@{}>{\raggedright\arraybackslash}p{0.460\linewidth}>{\centering\arraybackslash}X>{\raggedright\arraybackslash}p{0.460\linewidth}@{}}
\begin{minipage}[t]{\linewidth}
\centering\textbf{Earlier Germanic changes}\par
\vspace{0.35em}
\raggedright
\centering\textbf{West Germanic}\par
\raggedright
\vspace{0.2em}
\raggedright [no change]\par
\vspace{0.6em}
\centering\textbf{Northwest Germanic}\par
\raggedright
\vspace{0.2em}
\begin{tabular}{@{}>{\raggedright\arraybackslash}p{0.70\linewidth}@{\hspace{0.55em}}>{\raggedright\arraybackslash}p{0.20\linewidth}@{\hspace{0.25em}}}
\mbox{NWGmc Long E Lowering} & \emph{*slǣpaną} \\
\end{tabular}
\end{minipage}
&
&
\begin{minipage}[t]{\linewidth}
\centering\textbf{Old English changes}\par
\vspace{0.35em}
\raggedright
\begin{tabular}{@{}>{\raggedright\arraybackslash}p{0.74\linewidth}@{\hspace{0.55em}}>{\raggedright\arraybackslash}p{0.16\linewidth}@{\hspace{0.25em}}}
OE Heavy Syllable Nasal Apocope & \emph{*slǣpan} \\
OE Secondary Nasalization & \emph{*slǣpąn} \\
OE Weak Tail Reduction & \emph{*slǣpan} \\
\end{tabular}
\end{minipage}
\\
\end{tabularx}
\end{minipage}%
} \endgroup

Kroonen preserves the comparative verb as \emph{*slēpan-}, and Fulk
cites the same family under root \emph{*slēb-}
(\citeproc{ref-Kroonen2013}{Kroonen 2013}; \citeproc{ref-Fulk2018}{Fulk
2018, 120}). The selected input \emph{*slḗpaną} is the infinitive-style
form used here for that inherited sleep-verb. Clark Hall gives
\emph{slæpan} with preterite \emph{slēp}, \emph{slēap}, and Bright
likewise lists \emph{slæpan} \emph{(slāpan)}, \emph{slēp} \emph{slēpon}
\emph{slēpen} (\citeproc{ref-ClarkHall1960}{Clark Hall 1960, 281};
\citeproc{ref-BrightCassidyRingler1971}{Bright 1971, 435}). From
\emph{*slḗpaną}, Northwest Germanic lowering gives \emph{*slǣpaną}.

Form note. The note concerns lemma type rather than a special
derivational problem: this row represents the verb \emph{slǣpan},
whereas \emph{slǣp} belongs to noun or lookup background and
\emph{slēp}/\emph{slēap} are preterite forms
(\citeproc{ref-ClarkHall1960}{Clark Hall 1960, 281};
\citeproc{ref-BrightCassidyRingler1971}{Bright 1971, 435}).

\subsubsection{smear --- OE smierwan}\label{smear-oe-smierwan}

Derivation: \emph{*smérwijaną} \textgreater{} \emph{smierwan} (regular).

\begingroup
\setlength{\fboxsep}{6pt}

\noindent\fbox{%
\begin{minipage}{0.97\linewidth}
\footnotesize
\begin{tabularx}{\linewidth}{@{}>{\raggedright\arraybackslash}p{0.280\linewidth}>{\centering\arraybackslash}X>{\raggedright\arraybackslash}p{0.560\linewidth}@{}}
\begin{minipage}[t]{\linewidth}
\centering\textbf{Earlier Germanic changes}\par
\vspace{0.35em}
\raggedright
\centering\textbf{West Germanic}\par
\raggedright
\vspace{0.2em}
\raggedright [no change]\par
\vspace{0.6em}
\centering\textbf{Northwest Germanic}\par
\raggedright
\vspace{0.2em}
\raggedright [no change]\par
\end{minipage}
&
&
\begin{minipage}[t]{\linewidth}
\centering\textbf{Old English changes}\par
\vspace{0.35em}
\raggedright
\begin{tabular}{@{}>{\raggedright\arraybackslash}p{0.62\linewidth}@{\hspace{0.55em}}>{\raggedright\arraybackslash}p{0.28\linewidth}@{\hspace{0.25em}}}
OE Breaking & \emph{*sméorwijaną} \\
OE Heavy Syllable Nasal Apocope & \emph{*sméorwijan} \\
OE Secondary Nasalization & \emph{*sméorwijąn} \\
Sievers Law Syncope & \emph{*sméorwjąn} \\
OE I Umlaut & \emph{*smíerwjąn} \\
OE Weak Tail Reduction & \emph{*smíerwjan} \\
OE J Loss After Heavy & \emph{*smíerwan} \\
\end{tabular}
\end{minipage}
\\
\end{tabularx}
\end{minipage}%
} \endgroup

Kroonen gives the comparative headword as \emph{*smerwjan-}
(\citeproc{ref-Kroonen2013}{Kroonen 2013}). Ringe and Taylor instead
cite a later-stage \emph{*smirwijana}, from which they derive West Saxon
\emph{smierwan}, Mercian \emph{smirwan}, and Northumbrian \emph{smiriga}
(\citeproc{ref-RingeTaylor2014}{Ringe and Taylor 2014, 198, 263}). The
selected form is the West Saxon citation form \emph{smierwan}. From
\emph{*smérwijaną}, breaking before \emph{r + consonant} yields
\emph{eo}, and later i-umlaut produces \emph{ie}.

Dialect note. The entry therefore represents the West Saxon member of a
broader OE family: \emph{smierwan} in West Saxon, \emph{smirwan} in
Anglian or Mercian, and related later class-II forms such as
\emph{smirian} or \emph{smyrian} in the same lexical field
(\citeproc{ref-RingeTaylor2014}{Ringe and Taylor 2014};
\citeproc{ref-ClarkHall1960}{Clark Hall 1960, 283}).

\subsubsection{span --- OE spannan}\label{span-oe-spannan}

Derivation: \emph{*spánnaną} \textgreater{} \emph{spannan} (regular).

\begingroup
\setlength{\fboxsep}{6pt}

\noindent\fbox{%
\begin{minipage}{0.97\linewidth}
\small
\begin{tabularx}{\linewidth}{@{}>{\raggedright\arraybackslash}p{0.300\linewidth}>{\centering\arraybackslash}X>{\raggedright\arraybackslash}p{0.540\linewidth}@{}}
\begin{minipage}[t]{\linewidth}
\centering\textbf{Earlier Germanic changes}\par
\vspace{0.35em}
\raggedright
\centering\textbf{West Germanic}\par
\raggedright
\vspace{0.2em}
\raggedright [no change]\par
\vspace{0.6em}
\centering\textbf{Northwest Germanic}\par
\raggedright
\vspace{0.2em}
\raggedright [no change]\par
\end{minipage}
&
&
\begin{minipage}[t]{\linewidth}
\centering\textbf{Old English changes}\par
\vspace{0.35em}
\raggedright
\begin{tabular}{@{}>{\raggedright\arraybackslash}p{0.70\linewidth}@{\hspace{0.55em}}>{\raggedright\arraybackslash}p{0.20\linewidth}@{\hspace{0.25em}}}
OE Heavy Syllable Nasal Apocope & \emph{*spánnan} \\
OE Secondary Nasalization & \emph{*spánnąn} \\
OE Weak Tail Reduction & \emph{*spánnan} \\
\end{tabular}
\end{minipage}
\\
\end{tabularx}
\end{minipage}%
} \endgroup

Kroonen cites the inherited verb as \emph{*spannan-}, with OE
\emph{spannan} among the reflexes (\citeproc{ref-Kroonen2013}{Kroonen
2013, 505}). The selected input \emph{*spánnaną} is the infinitive-style
form used here for that same verbal lexeme. Clark Hall keeps noun
\emph{spann} and verb \emph{spannan} separate, and Brunner likewise
records \emph{sponnan, spannan stv.} (\citeproc{ref-ClarkHall1960}{Clark
Hall 1960, 286}; \citeproc{ref-SieversBrunner1965}{Sievers 1965}). From
\emph{*spánnaną}, the final nasal ending is lost and the regular OE
weak-tail steps surface \emph{spannan}.

Form note. English \emph{span} can also reach noun \emph{spann} in local
lookup material. The entry represented here is the verb \emph{spannan},
with the noun treated elsewhere (\citeproc{ref-ClarkHall1960}{Clark Hall
1960, 286}).

\subsubsection{spar --- OE spearra}\label{spar-oe-spearra}

Derivation: \emph{*spárrô} \textgreater{} \emph{spearra} (regular).

\begingroup
\setlength{\fboxsep}{6pt}

\noindent\fbox{%
\begin{minipage}{0.97\linewidth}
\small
\begin{tabularx}{\linewidth}{@{}>{\raggedright\arraybackslash}p{0.300\linewidth}>{\centering\arraybackslash}X>{\raggedright\arraybackslash}p{0.540\linewidth}@{}}
\begin{minipage}[t]{\linewidth}
\centering\textbf{Earlier Germanic changes}\par
\vspace{0.35em}
\raggedright
\centering\textbf{West Germanic}\par
\raggedright
\vspace{0.2em}
\raggedright [no change]\par
\vspace{0.6em}
\centering\textbf{Northwest Germanic}\par
\raggedright
\vspace{0.2em}
\raggedright [no change]\par
\end{minipage}
&
&
\begin{minipage}[t]{\linewidth}
\centering\textbf{Old English changes}\par
\vspace{0.35em}
\raggedright
\begin{tabular}{@{}>{\raggedright\arraybackslash}p{0.70\linewidth}@{\hspace{0.55em}}>{\raggedright\arraybackslash}p{0.20\linewidth}@{\hspace{0.25em}}}
Anglo Frisian Brightening & \emph{*spærrô} \\
OE Breaking & \emph{*spearrô} \\
OE Unstressed Long Vowel Shortening & \emph{*spearra} \\
\end{tabular}
\end{minipage}
\\
\end{tabularx}
\end{minipage}%
} \endgroup

Kroonen and Orel place this noun in the beam or rafter set
\emph{*spar(r)an-}, with cognates such as Old Saxon and Old High German
\emph{sparro} (\citeproc{ref-Kroonen2013}{Kroonen 2013, 506};
\citeproc{ref-Orel2003}{Orel 2003}). The selected input \emph{*spárrô}
is the OE-facing nominal form used here for that same lexeme. The noun
represented here is \emph{spearra}. From \emph{*spárrô}, Anglo-Frisian
brightening gives \emph{*spærrô}, and OE breaking before geminate
\emph{rr} yields \emph{*spearrô}, later \emph{spearra}.

Form note. This entry concerns the noun \emph{spearra} only. It should
be kept separate from verb \emph{sperran}, even though the Modern
English glosses overlap (\citeproc{ref-Kroonen2013}{Kroonen 2013, 506};
\citeproc{ref-Orel2003}{Orel 2003}).

\subsubsection{still --- OE stillan}\label{still-oe-stillan}

Derivation: \emph{*stéllijaną} \textgreater{} \emph{stillan} (regular).

\begingroup
\setlength{\fboxsep}{6pt}

\noindent\fbox{%
\begin{minipage}{0.97\linewidth}
\footnotesize
\begin{tabularx}{\linewidth}{@{}>{\raggedright\arraybackslash}p{0.280\linewidth}>{\centering\arraybackslash}X>{\raggedright\arraybackslash}p{0.560\linewidth}@{}}
\begin{minipage}[t]{\linewidth}
\centering\textbf{Earlier Germanic changes}\par
\vspace{0.35em}
\raggedright
\centering\textbf{West Germanic}\par
\raggedright
\vspace{0.2em}
\raggedright [no change]\par
\vspace{0.6em}
\centering\textbf{Northwest Germanic}\par
\raggedright
\vspace{0.2em}
\raggedright [no change]\par
\end{minipage}
&
&
\begin{minipage}[t]{\linewidth}
\centering\textbf{Old English changes}\par
\vspace{0.35em}
\raggedright
\begin{tabular}{@{}>{\raggedright\arraybackslash}p{0.64\linewidth}@{\hspace{0.55em}}>{\raggedright\arraybackslash}p{0.26\linewidth}@{\hspace{0.25em}}}
OE Heavy Syllable Nasal Apocope & \emph{*stéllijan} \\
OE Secondary Nasalization & \emph{*stéllijąn} \\
Sievers Law Syncope & \emph{*stélljąn} \\
OE I Umlaut & \emph{*stilljąn} \\
OE Weak Tail Reduction & \emph{*stilljan} \\
OE J Loss After Heavy & \emph{*stillan} \\
\end{tabular}
\end{minipage}
\\
\end{tabularx}
\end{minipage}%
} \endgroup

The wider West Germanic family includes adjective \emph{still} and verb
\emph{stillen} (\citeproc{ref-KlugeSeebold2011}{Kluge 2011}). The
selected input \emph{*stéllijaną} represents the verbal j-formation used
for the OE row. Clark Hall gives \emph{stillan} as the verb and
separately \emph{stille} as the adjective
(\citeproc{ref-ClarkHall1960}{Clark Hall 1960, 292}). As a heavy-stem
Class I weak verb, \emph{*stéllijaną} undergoes the expected syncope and
i-umlaut, and later loss of \emph{j} after a heavy stem yields
\emph{stillan}.

Form note. The note concerns lexical framing rather than sound law:
\emph{stillan} is the verb represented here, while \emph{stille} belongs
to the related adjectival branch of the family
(\citeproc{ref-ClarkHall1960}{Clark Hall 1960, 292};
\citeproc{ref-KlugeSeebold2011}{Kluge 2011}).

\subsubsection{summer --- OE sumer}\label{summer-oe-sumer}

Derivation: \emph{*súmaraz} \textgreater{} \emph{sumer} (regular).

\begingroup
\setlength{\fboxsep}{6pt}

\noindent\fbox{%
\begin{minipage}{0.97\linewidth}
\small
\begin{tabularx}{\linewidth}{@{}>{\raggedright\arraybackslash}p{0.300\linewidth}>{\centering\arraybackslash}X>{\raggedright\arraybackslash}p{0.540\linewidth}@{}}
\begin{minipage}[t]{\linewidth}
\centering\textbf{Earlier Germanic changes}\par
\vspace{0.35em}
\raggedright
\centering\textbf{West Germanic}\par
\raggedright
\vspace{0.2em}
\raggedright [no change]\par
\vspace{0.6em}
\centering\textbf{Northwest Germanic}\par
\raggedright
\vspace{0.2em}
\begin{tabular}{@{}>{\raggedright\arraybackslash}p{0.74\linewidth}@{\hspace{0.55em}}>{\raggedright\arraybackslash}p{0.16\linewidth}@{\hspace{0.25em}}}
\mbox{PGmc Final Z Deletion} & \emph{*súmara} \\
\end{tabular}
\end{minipage}
&
&
\begin{minipage}[t]{\linewidth}
\centering\textbf{Old English changes}\par
\vspace{0.35em}
\raggedright
\begin{tabular}{@{}>{\raggedright\arraybackslash}p{0.70\linewidth}@{\hspace{0.55em}}>{\raggedright\arraybackslash}p{0.16\linewidth}@{\hspace{0.25em}}}
PWGmc Final Bare A Loss & \emph{*súmar} \\
Anglo Frisian Brightening & \emph{*súmær} \\
OE Unstressed AE Merger & \emph{*súmer} \\
\end{tabular}
\end{minipage}
\\
\end{tabularx}
\end{minipage}%
} \endgroup

Kroonen gives the lexeme as \emph{*sumara-}, and Ringe and Taylor
likewise use \emph{*sumaraz}, while Orel preserves an alternate
\emph{*sumeraz} (\citeproc{ref-Kroonen2013}{Kroonen 2013, 492};
\citeproc{ref-RingeTaylor2014}{Ringe and Taylor 2014, 46};
\citeproc{ref-Orel2003}{Orel 2003, 425}). The selected input
\emph{*súmaraz} follows the \emph{*a} vocalism that underlies the
regular development represented here. Clark Hall gives \emph{sumor m.,
gs. sumeres, ds. sumera, sumere} (\citeproc{ref-ClarkHall1960}{Clark
Hall 1960, 281}). Bright likewise lists \emph{sumor (sumer)} with
genitive \emph{sumeres} (\citeproc{ref-BrightCassidyRingler1971}{Bright
1971, 440}). From \emph{*súmaraz}, loss of final \emph{-az} is followed
by fronting and merger in the unstressed second syllable, yielding
\emph{sumer}.

Form note. Old English also has \emph{sumor} as a common headword
spelling. Clark Hall's \emph{sumor \ldots{} sumeres \ldots{} sumere} and
Bright's \emph{sumor (sumer)} show that the \emph{e}-vocalism was also
real in Old English (\citeproc{ref-ClarkHall1960}{Clark Hall 1960, 281};
\citeproc{ref-BrightCassidyRingler1971}{Bright 1971, 440}).

\subsubsection{sunder --- OE sundrian}\label{sunder-oe-sundrian}

Derivation: \emph{*súndrōjaną} \textgreater{} \emph{sundrian} (regular).

\begingroup
\setlength{\fboxsep}{6pt}

\noindent\fbox{%
\begin{minipage}{0.97\linewidth}
\footnotesize
\begin{tabularx}{\linewidth}{@{}>{\raggedright\arraybackslash}p{0.280\linewidth}>{\centering\arraybackslash}X>{\raggedright\arraybackslash}p{0.560\linewidth}@{}}
\begin{minipage}[t]{\linewidth}
\centering\textbf{Earlier Germanic changes}\par
\vspace{0.35em}
\raggedright
\centering\textbf{West Germanic}\par
\raggedright
\vspace{0.2em}
\raggedright [no change]\par
\vspace{0.6em}
\centering\textbf{Northwest Germanic}\par
\raggedright
\vspace{0.2em}
\raggedright [no change]\par
\end{minipage}
&
&
\begin{minipage}[t]{\linewidth}
\centering\textbf{Old English changes}\par
\vspace{0.35em}
\raggedright
\begin{tabular}{@{}>{\raggedright\arraybackslash}p{0.64\linewidth}@{\hspace{0.55em}}>{\raggedright\arraybackslash}p{0.26\linewidth}@{\hspace{0.25em}}}
OE Heavy Syllable Nasal Apocope & \emph{*súndrōjan} \\
OE Secondary Nasalization & \emph{*súndrōjąn} \\
OE I Umlaut & \emph{*súndrējąn} \\
OE Unstressed Long Vowel Shortening & \emph{*súndrejąn} \\
OE Weak Tail Reduction & \emph{*súndrejan} \\
OE Intervocalic J Vocalization & \emph{*súndreian} \\
OE Unstressed EI Contraction & \emph{*súndrian} \\
\end{tabular}
\end{minipage}
\\
\end{tabularx}
\end{minipage}%
} \endgroup

Orel distinguishes three related formations: adverbial \emph{*sunþraz}
\textgreater{} sundor, Class I verbal \emph{*sunþrjanan} \textgreater{}
syndrian, and Class II verbal \emph{*sunþrōjanan} \textgreater{}
sundrian (\citeproc{ref-Orel2003}{Orel 2003}). Kluge-Seebold aligns the
cognate set with German \emph{sondern} and OE \emph{gesundrian}, so this
entry belongs with the Class II verb, not the adverb
(\citeproc{ref-KlugeSeebold2011}{Kluge 2011}). Clark Hall and
Bosworth-Toller keep \emph{sundrian} and \emph{syndrian} separate from
adverbial \emph{sundor}, and both preserve the prefixed verbal family
\emph{ā-sundrian} (\citeproc{ref-ClarkHall1960}{Clark Hall 1960, 296};
\citeproc{ref-BosworthToller1898}{Bosworth and Toller 1898}). From
\emph{*súndrōjaną}, the Class II weak-verb suffix yields regular OE
\emph{-ian}, producing \emph{sundrian}.

Form note. The earlier confusion was lexical, not phonological:
\emph{sundor} is the separate adverb, and \emph{syndrian} is a related
but different verb. The verb treated here is the Class II verb
\emph{sundrian} (\citeproc{ref-Orel2003}{Orel 2003};
\citeproc{ref-ClarkHall1960}{Clark Hall 1960, 296}).

\subsubsection{swallow --- OE swealwe}\label{swallow-oe-swealwe}

Derivation: \emph{*swálwōn} \textgreater{} \emph{swealwe} (regular).

\begingroup
\setlength{\fboxsep}{6pt}

\noindent\fbox{%
\begin{minipage}{0.97\linewidth}
\small
\begin{tabularx}{\linewidth}{@{}>{\raggedright\arraybackslash}p{0.300\linewidth}>{\centering\arraybackslash}X>{\raggedright\arraybackslash}p{0.540\linewidth}@{}}
\begin{minipage}[t]{\linewidth}
\centering\textbf{Earlier Germanic changes}\par
\vspace{0.35em}
\raggedright
\centering\textbf{West Germanic}\par
\raggedright
\vspace{0.2em}
\raggedright [no change]\par
\vspace{0.6em}
\centering\textbf{Northwest Germanic}\par
\raggedright
\vspace{0.2em}
\begin{tabular}{@{}>{\raggedright\arraybackslash}p{0.64\linewidth}@{\hspace{0.55em}}>{\raggedright\arraybackslash}p{0.16\linewidth}@{\hspace{0.25em}}}
NWGmc N Stem N Loss & \emph{*swálwǭ} \\
\end{tabular}
\end{minipage}
&
&
\begin{minipage}[t]{\linewidth}
\centering\textbf{Old English changes}\par
\vspace{0.35em}
\raggedright
\begin{tabular}{@{}>{\raggedright\arraybackslash}p{0.70\linewidth}@{\hspace{0.55em}}>{\raggedright\arraybackslash}p{0.20\linewidth}@{\hspace{0.25em}}}
Anglo Frisian Brightening & \emph{*swælwǭ} \\
OE Breaking & \emph{*swealwǭ} \\
OE Unstressed Long Vowel Shortening & \emph{*swealwæ} \\
OE Unstressed AE Merger & \emph{*swealwe} \\
\end{tabular}
\end{minipage}
\\
\end{tabularx}
\end{minipage}%
} \endgroup

Kroonen gives the bird name as \emph{*swalwōn-}, and Ringe and Taylor
cite the later West Germanic stage \emph{*swalwa}, from which West Saxon
\emph{swealwe} and Mercian \emph{swalwe} develop
(\citeproc{ref-Kroonen2013}{Kroonen 2013, 535};
\citeproc{ref-RingeTaylor2014}{Ringe and Taylor 2014, 200}). The
selected etymological comparison belongs to the swallow-bird family, not
to the verb \emph{swelgan}. Clark Hall records \emph{swealwe (a, o)} as
the noun headword (\citeproc{ref-ClarkHall1960}{Clark Hall 1960, 299}).
From \emph{*swálwōn}, brightening yields \emph{*swælw-}, and breaking
before \emph{lw} gives \emph{*swealw-}.

Form note. The final prose keeps the citation form \emph{swealwe}
separate from two different kinds of background material: the unrelated
verb \emph{swelgan}, and later or oblique spellings such as
\emph{swaluwe} (\citeproc{ref-Campbell1959}{Campbell 1959, sect. 365})
or \emph{swalewan} in the same broader spelling tradition.

\subsubsection{swine --- OE swīn}\label{swine-oe-swux12bn}

Derivation: citation reconstruction \emph{*swī́ną}; selected input
\emph{*swī́ną} \textgreater{} \emph{swīn} (regular).

\begingroup
\setlength{\fboxsep}{6pt}

\noindent\fbox{%
\begin{minipage}{0.97\linewidth}
\small
\begin{tabularx}{\linewidth}{@{}>{\raggedright\arraybackslash}p{0.300\linewidth}>{\centering\arraybackslash}X>{\raggedright\arraybackslash}p{0.540\linewidth}@{}}
\begin{minipage}[t]{\linewidth}
\centering\textbf{Earlier Germanic changes}\par
\vspace{0.35em}
\raggedright
\centering\textbf{West Germanic}\par
\raggedright
\vspace{0.2em}
\raggedright [no change]\par
\vspace{0.6em}
\centering\textbf{Northwest Germanic}\par
\raggedright
\vspace{0.2em}
\raggedright [no change]\par
\end{minipage}
&
&
\begin{minipage}[t]{\linewidth}
\centering\textbf{Old English changes}\par
\vspace{0.35em}
\raggedright
\begin{tabular}{@{}>{\raggedright\arraybackslash}p{0.74\linewidth}@{\hspace{0.55em}}>{\raggedright\arraybackslash}p{0.16\linewidth}@{\hspace{0.25em}}}
OE Heavy Syllable Nasal Apocope & \emph{*swī́n} \\
\end{tabular}
\end{minipage}
\\
\end{tabularx}
\end{minipage}%
} \endgroup

Kroonen cites the noun as \emph{*swina-} `pig'
(\citeproc{ref-Kroonen2013}{Kroonen 2013, 542}). The selected
comparative form here is the bare neuter citation cell \emph{*swī́ną},
which matches the singular noun aimed at in Old English rather than an
oblique stem form. Clark Hall records \emph{swin (y)} as the ordinary
noun headword (\citeproc{ref-ClarkHall1960}{Clark Hall 1960, 301}). From
selected input \emph{*swī́ną}, loss of the final nasal vowel yields
\emph{swīn}.

Source note. The selected input writes stressed long \emph{ī} as
\emph{*ī́}, so comparative \emph{*swī́ną} and derivational \emph{*swī́ną}
represent the same lexical form.

\subsubsection{think --- OE þenċan}\label{think-oe-uxfeenux10ban}

Derivation: \emph{*θánkijaną} \textgreater{} \emph{þenċan} (regular).

\begingroup
\setlength{\fboxsep}{6pt}

\noindent\fbox{%
\begin{minipage}{0.97\linewidth}
\footnotesize
\begin{tabularx}{\linewidth}{@{}>{\raggedright\arraybackslash}p{0.280\linewidth}>{\centering\arraybackslash}X>{\raggedright\arraybackslash}p{0.560\linewidth}@{}}
\begin{minipage}[t]{\linewidth}
\centering\textbf{Earlier Germanic changes}\par
\vspace{0.35em}
\raggedright
\centering\textbf{West Germanic}\par
\raggedright
\vspace{0.2em}
\raggedright [no change]\par
\vspace{0.6em}
\centering\textbf{Northwest Germanic}\par
\raggedright
\vspace{0.2em}
\raggedright [no change]\par
\end{minipage}
&
&
\begin{minipage}[t]{\linewidth}
\centering\textbf{Old English changes}\par
\vspace{0.35em}
\raggedright
\begin{tabular}{@{}>{\raggedright\arraybackslash}p{0.66\linewidth}@{\hspace{0.55em}}>{\raggedright\arraybackslash}p{0.24\linewidth}@{\hspace{0.25em}}}
OE Heavy Syllable Nasal Apocope & \emph{*θánkijan} \\
OE Secondary Nasalization & \emph{*θánkijąn} \\
Sievers Law Syncope & \emph{*θánkjąn} \\
OE Velar Palatalization & \emph{*θánʧjąn} \\
OE I Umlaut & \emph{*θenʧjąn} \\
OE Weak Tail Reduction & \emph{*θenʧjan} \\
OE J Loss After Heavy & \emph{*θenʧan} \\
\end{tabular}
\end{minipage}
\\
\end{tabularx}
\end{minipage}%
} \endgroup

Kroonen gives the verb as \emph{*þankjan-} `to think', and Ringe and
Taylor cite fully inflected \emph{*þankijaną} beside OE \emph{þenċan}
(\citeproc{ref-Kroonen2013}{Kroonen 2013};
\citeproc{ref-RingeTaylor2014}{Ringe and Taylor 2014}). The noun
\emph{*þankaz} belongs only to the wider derivational background.
Bosworth-Toller preserves the verb under \emph{þencan}/\emph{geþencan},
and the citation form here is the ordinary infinitive \emph{þenċan}
(\citeproc{ref-BosworthToller1898}{Bosworth and Toller 1898}). From
\emph{*θánkijaną}, palatalization before \emph{*j} and i-umlaut produce
\emph{þenċan}. Campbell's assibilation discussion uses the same verb
\emph{þencan}; the class-III relic \emph{hycgan} is a different lexeme
(\citeproc{ref-Campbell1959}{Campbell 1959, sect. 438};
\citeproc{ref-Hogg1992}{Hogg 1992, sect. 3.4.2.4}).

\subsubsection{thorn --- OE þorn}\label{thorn-oe-uxfeorn}

Derivation: \emph{*θúrnaz} \textgreater{} \emph{þorn} (regular).

\begingroup
\setlength{\fboxsep}{6pt}

\noindent\fbox{%
\begin{minipage}{0.97\linewidth}
\small
\begin{tabularx}{\linewidth}{@{}>{\raggedright\arraybackslash}p{0.540\linewidth}>{\centering\arraybackslash}X>{\raggedright\arraybackslash}p{0.300\linewidth}@{}}
\begin{minipage}[t]{\linewidth}
\centering\textbf{Earlier Germanic changes}\par
\vspace{0.35em}
\raggedright
\centering\textbf{West Germanic}\par
\raggedright
\vspace{0.2em}
\raggedright [no change]\par
\vspace{0.6em}
\centering\textbf{Northwest Germanic}\par
\raggedright
\vspace{0.2em}
\begin{tabular}{@{}>{\raggedright\arraybackslash}p{0.74\linewidth}@{\hspace{0.55em}}>{\raggedright\arraybackslash}p{0.16\linewidth}@{\hspace{0.25em}}}
\mbox{NWGmc U Lowering} & \emph{*θórnaz} \\
\mbox{PGmc Final Z Deletion} & \emph{*θórna} \\
\end{tabular}
\end{minipage}
&
&
\begin{minipage}[t]{\linewidth}
\centering\textbf{Old English changes}\par
\vspace{0.35em}
\raggedright
\begin{tabular}{@{}>{\raggedright\arraybackslash}p{0.70\linewidth}@{\hspace{0.55em}}>{\raggedright\arraybackslash}p{0.16\linewidth}@{\hspace{0.25em}}}
PWGmc Final Bare A Loss & \emph{*θórn} \\
\end{tabular}
\end{minipage}
\\
\end{tabularx}
\end{minipage}%
} \endgroup

Kroonen gives \emph{*þurna-} `thorn, briar', while Orel preserves the
masculine pair \emph{*þurnuz} \textasciitilde{} \emph{*þurnaz}
(\citeproc{ref-Kroonen2013}{Kroonen 2013, 33};
\citeproc{ref-Orel2003}{Orel 2003, 430}). The selected input
\emph{*θúrnaz} belongs to that same comparative family. Bright lists
\emph{þorn}, m., and Clark Hall likewise treats \emph{þorn} as the
ordinary noun headword (\citeproc{ref-BrightCassidyRingler1971}{Bright
1971, 450}; \citeproc{ref-ClarkHall1960}{Clark Hall 1960, 327}). The
inherited stem shows regular lowering of \emph{u} to \emph{o} before
\emph{r}, and final loss yields \emph{þorn}.

Source note. The comparative sources preserve more than one stem
formation, but the Old English target itself is simply the citation form
\emph{þorn}.

\subsubsection{tide --- OE tīd}\label{tide-oe-tux12bd}

Derivation: citation reconstruction \emph{*tī́diz}; selected input
\emph{*tī́diz} \textgreater{} \emph{tīd} (regular).

\begingroup
\setlength{\fboxsep}{6pt}

\noindent\fbox{%
\begin{minipage}{0.97\linewidth}
\small
\begin{tabularx}{\linewidth}{@{}>{\raggedright\arraybackslash}p{0.460\linewidth}>{\centering\arraybackslash}X>{\raggedright\arraybackslash}p{0.460\linewidth}@{}}
\begin{minipage}[t]{\linewidth}
\centering\textbf{Earlier Germanic changes}\par
\vspace{0.35em}
\raggedright
\centering\textbf{West Germanic}\par
\raggedright
\vspace{0.2em}
\raggedright [no change]\par
\vspace{0.6em}
\centering\textbf{Northwest Germanic}\par
\raggedright
\vspace{0.2em}
\begin{tabular}{@{}>{\raggedright\arraybackslash}p{0.74\linewidth}@{\hspace{0.55em}}>{\raggedright\arraybackslash}p{0.16\linewidth}@{\hspace{0.25em}}}
\mbox{PGmc Final Z Deletion} & \emph{*tī́di} \\
\end{tabular}
\end{minipage}
&
&
\begin{minipage}[t]{\linewidth}
\centering\textbf{Old English changes}\par
\vspace{0.35em}
\raggedright
\begin{tabular}{@{}>{\raggedright\arraybackslash}p{0.74\linewidth}@{\hspace{0.55em}}>{\raggedright\arraybackslash}p{0.16\linewidth}@{\hspace{0.25em}}}
\mbox{OE High Vowel Apocope} & \emph{*tī́d} \\
\end{tabular}
\end{minipage}
\\
\end{tabularx}
\end{minipage}%
} \endgroup

Kroonen gives \emph{*tīdi-} `time', and Orel's \emph{*tīđiz} points to
the same feminine noun (\citeproc{ref-Kroonen2013}{Kroonen 2013, 556};
\citeproc{ref-Orel2003}{Orel 2003, 446}). The related verb \emph{tīdan}
is separate. Bright records \emph{tīd} with singular \emph{tīde} and
plural \emph{tīda}, and Clark Hall treats \emph{tīd} as the ordinary
noun `time, period, season'
(\citeproc{ref-BrightCassidyRingler1971}{Bright 1971, 444};
\citeproc{ref-ClarkHall1960}{Clark Hall 1960, 309}). From \emph{*tī́diz},
final \emph{z} is lost and the high final vowel drops, leaving
\emph{tīd}. The note matters only because English \emph{tide} can pull
in the separate weak verb \emph{tīdan}; the noun targeted here is
\emph{tīd}.

\subsubsection{token --- OE tācn}\label{token-oe-tux101cn}

Derivation: \emph{*táikną} \textgreater{} \emph{tācn} (regular).

\begingroup
\setlength{\fboxsep}{6pt}

\noindent\fbox{%
\begin{minipage}{0.97\linewidth}
\small
\begin{tabularx}{\linewidth}{@{}>{\raggedright\arraybackslash}p{0.300\linewidth}>{\centering\arraybackslash}X>{\raggedright\arraybackslash}p{0.540\linewidth}@{}}
\begin{minipage}[t]{\linewidth}
\centering\textbf{Earlier Germanic changes}\par
\vspace{0.35em}
\raggedright
\centering\textbf{West Germanic}\par
\raggedright
\vspace{0.2em}
\begin{tabular}{@{}>{\raggedright\arraybackslash}p{0.70\linewidth}@{\hspace{0.55em}}>{\raggedright\arraybackslash}p{0.16\linewidth}@{\hspace{0.25em}}}
PWGmc Ai Monophthongization & \emph{*tākną} \\
\end{tabular}
\vspace{0.6em}
\centering\textbf{Northwest Germanic}\par
\raggedright
\vspace{0.2em}
\raggedright [no change]\par
\end{minipage}
&
&
\begin{minipage}[t]{\linewidth}
\centering\textbf{Old English changes}\par
\vspace{0.35em}
\raggedright
\begin{tabular}{@{}>{\raggedright\arraybackslash}p{0.74\linewidth}@{\hspace{0.55em}}>{\raggedright\arraybackslash}p{0.16\linewidth}@{\hspace{0.25em}}}
OE Heavy Syllable Nasal Apocope & \emph{*tākn} \\
\end{tabular}
\end{minipage}
\\
\end{tabularx}
\end{minipage}%
} \endgroup

Kroonen cites \emph{*taikna-} and Orel \emph{*taiknan} for the noun
`sign, token' (\citeproc{ref-Kroonen2013}{Kroonen 2013};
\citeproc{ref-Orel2003}{Orel 2003, 438}). The selected input
\emph{*táikną} is the simple citation-form noun used for the derivation.
Campbell preserves unbroken \emph{tācn} and oblique \emph{tācnes}
(\citeproc{ref-Campbell1959}{Campbell 1959, sect. 574}), and also
records the later broken form \emph{tācen}
(\citeproc{ref-Campbell1959}{Campbell 1959, sect. 365}). Sievers-Brunner
likewise preserves both unbroken and broken shapes
(\citeproc{ref-SieversBrunner1965}{Sievers 1965}). Monophthongization of
\emph{ai} yields \emph{ā}, and loss of the final nasal vowel leaves
\emph{tācn}.

Form note. \emph{tācn} is the attested unbroken citation form selected
here. Later West Saxon prose often prefers \emph{tācen}, but that does
not displace the older unbroken form
(\citeproc{ref-SieversBrunner1965}{Sievers 1965}).

\subsubsection{town --- OE tūn}\label{town-oe-tux16bn}

Derivation: \emph{*tūną} \textgreater{} \emph{tūn} (regular).

\begingroup
\setlength{\fboxsep}{6pt}

\noindent\fbox{%
\begin{minipage}{0.97\linewidth}
\small
\begin{tabularx}{\linewidth}{@{}>{\raggedright\arraybackslash}p{0.300\linewidth}>{\centering\arraybackslash}X>{\raggedright\arraybackslash}p{0.540\linewidth}@{}}
\begin{minipage}[t]{\linewidth}
\centering\textbf{Earlier Germanic changes}\par
\vspace{0.35em}
\raggedright
\centering\textbf{West Germanic}\par
\raggedright
\vspace{0.2em}
\raggedright [no change]\par
\vspace{0.6em}
\centering\textbf{Northwest Germanic}\par
\raggedright
\vspace{0.2em}
\raggedright [no change]\par
\end{minipage}
&
&
\begin{minipage}[t]{\linewidth}
\centering\textbf{Old English changes}\par
\vspace{0.35em}
\raggedright
\begin{tabular}{@{}>{\raggedright\arraybackslash}p{0.74\linewidth}@{\hspace{0.55em}}>{\raggedright\arraybackslash}p{0.16\linewidth}@{\hspace{0.25em}}}
OE Heavy Syllable Nasal Apocope & \emph{*tūn} \\
\end{tabular}
\end{minipage}
\\
\end{tabularx}
\end{minipage}%
} \endgroup

Kroonen cites \emph{*tūna-} `fenced area', while Orel gives
\emph{*tūnan} \textasciitilde{} \emph{*tūnaz}
(\citeproc{ref-Kroonen2013}{Kroonen 2013}; \citeproc{ref-Orel2003}{Orel
2003, 452}). The selected input \emph{*tūną} is the simple citation-form
noun used in the derivation. Clark Hall records \emph{tūn} as the
ordinary headword `enclosure, yard, village, town'
(\citeproc{ref-ClarkHall1960}{Clark Hall 1960}). The inherited long
\emph{ū} is preserved, and loss of the final nasal vowel yields
\emph{tūn} regularly (\citeproc{ref-SieversBrunner1965}{Sievers 1965}).

Source note. The comparative headwords vary, but the Old English target
here is the direct citation form \emph{tūn}, not an oblique
\emph{*tūnăn}.

\subsubsection{wade --- OE wadan}\label{wade-oe-wadan}

Derivation: \emph{*wádaną} \textgreater{} \emph{wadan} (regular).

\begingroup
\setlength{\fboxsep}{6pt}

\noindent\fbox{%
\begin{minipage}{0.97\linewidth}
\small
\begin{tabularx}{\linewidth}{@{}>{\raggedright\arraybackslash}p{0.300\linewidth}>{\centering\arraybackslash}X>{\raggedright\arraybackslash}p{0.540\linewidth}@{}}
\begin{minipage}[t]{\linewidth}
\centering\textbf{Earlier Germanic changes}\par
\vspace{0.35em}
\raggedright
\centering\textbf{West Germanic}\par
\raggedright
\vspace{0.2em}
\raggedright [no change]\par
\vspace{0.6em}
\centering\textbf{Northwest Germanic}\par
\raggedright
\vspace{0.2em}
\raggedright [no change]\par
\end{minipage}
&
&
\begin{minipage}[t]{\linewidth}
\centering\textbf{Old English changes}\par
\vspace{0.35em}
\raggedright
\begin{tabular}{@{}>{\raggedright\arraybackslash}p{0.74\linewidth}@{\hspace{0.55em}}>{\raggedright\arraybackslash}p{0.16\linewidth}@{\hspace{0.25em}}}
Anglo Frisian Brightening & \emph{*wædaną} \\
OE A Restoration & \emph{*wadaną} \\
OE Heavy Syllable Nasal Apocope & \emph{*wadan} \\
OE Secondary Nasalization & \emph{*wadąn} \\
OE Weak Tail Reduction & \emph{*wadan} \\
\end{tabular}
\end{minipage}
\\
\end{tabularx}
\end{minipage}%
} \endgroup

Campbell and Ringe and Taylor describe A-restoration before a following
back vowel, and Luick explicitly includes \emph{wadan} among the
standard open-syllable examples (\citeproc{ref-Campbell1959}{Campbell
1959, sect. 744}; \citeproc{ref-RingeTaylor2014}{Ringe and Taylor 2014};
\citeproc{ref-Luick1914}{Luick 1914-\/-1940, 239}). Clark Hall gives
\emph{wadan} as the verb `to go, move, stride, advance', and Bright
lists the same infinitive in the strong-verb paradigm
(\citeproc{ref-ClarkHall1960}{Clark Hall 1960, 353};
\citeproc{ref-BrightCassidyRingler1971}{Bright 1971, 456}). From
\emph{*wádaną}, Anglo-Frisian brightening first gives \emph{*wædaną}.
This infinitive belongs to the A-restoration class. The citation form is
therefore \emph{wadan}, not a fronted \emph{wæden}-type output.

\subsubsection{warp --- OE weorpan}\label{warp-oe-weorpan}

Derivation: \emph{*wérpaną} \textgreater{} \emph{weorpan} (regular).

\begingroup
\setlength{\fboxsep}{6pt}

\noindent\fbox{%
\begin{minipage}{0.97\linewidth}
\small
\begin{tabularx}{\linewidth}{@{}>{\raggedright\arraybackslash}p{0.300\linewidth}>{\centering\arraybackslash}X>{\raggedright\arraybackslash}p{0.540\linewidth}@{}}
\begin{minipage}[t]{\linewidth}
\centering\textbf{Earlier Germanic changes}\par
\vspace{0.35em}
\raggedright
\centering\textbf{West Germanic}\par
\raggedright
\vspace{0.2em}
\raggedright [no change]\par
\vspace{0.6em}
\centering\textbf{Northwest Germanic}\par
\raggedright
\vspace{0.2em}
\raggedright [no change]\par
\end{minipage}
&
&
\begin{minipage}[t]{\linewidth}
\centering\textbf{Old English changes}\par
\vspace{0.35em}
\raggedright
\begin{tabular}{@{}>{\raggedright\arraybackslash}p{0.68\linewidth}@{\hspace{0.55em}}>{\raggedright\arraybackslash}p{0.22\linewidth}@{\hspace{0.25em}}}
OE Breaking & \emph{*wéorpaną} \\
OE Heavy Syllable Nasal Apocope & \emph{*wéorpan} \\
OE Secondary Nasalization & \emph{*wéorpąn} \\
OE Weak Tail Reduction & \emph{*wéorpan} \\
\end{tabular}
\end{minipage}
\\
\end{tabularx}
\end{minipage}%
} \endgroup

Ringe and Taylor distinguish preterite \emph{*warp} from infinitive
\emph{*werpana}, and the selected input here is the verbal form
\emph{*wérpaną} (\citeproc{ref-RingeTaylor2014}{Ringe and Taylor 2014}).
Clark Hall records \emph{weorpan} as the strong verb headword and
separately lists \emph{wearp} as both noun and preterite. Breaking
before \emph{r + C} yields \emph{weor-}, and the infinitive develops
regularly to \emph{weorpan} (\citeproc{ref-Campbell1959}{Campbell 1959,
sect. 146}; \citeproc{ref-Hogg1992}{Hogg 1992, sect. 3.4.2.2}). English
\emph{warp} also points to related \emph{wearp} material. The selected
form here is the infinitive \emph{weorpan}.

\subsubsection{wash --- OE wascan}\label{wash-oe-wascan}

Derivation: \emph{*wáskaną} \textgreater{} \emph{wascan} (regular).

\begingroup
\setlength{\fboxsep}{6pt}

\noindent\fbox{%
\begin{minipage}{0.97\linewidth}
\small
\begin{tabularx}{\linewidth}{@{}>{\raggedright\arraybackslash}p{0.300\linewidth}>{\centering\arraybackslash}X>{\raggedright\arraybackslash}p{0.540\linewidth}@{}}
\begin{minipage}[t]{\linewidth}
\centering\textbf{Earlier Germanic changes}\par
\vspace{0.35em}
\raggedright
\centering\textbf{West Germanic}\par
\raggedright
\vspace{0.2em}
\raggedright [no change]\par
\vspace{0.6em}
\centering\textbf{Northwest Germanic}\par
\raggedright
\vspace{0.2em}
\raggedright [no change]\par
\end{minipage}
&
&
\begin{minipage}[t]{\linewidth}
\centering\textbf{Old English changes}\par
\vspace{0.35em}
\raggedright
\begin{tabular}{@{}>{\raggedright\arraybackslash}p{0.70\linewidth}@{\hspace{0.55em}}>{\raggedright\arraybackslash}p{0.20\linewidth}@{\hspace{0.25em}}}
Anglo Frisian Brightening & \emph{*wæskaną} \\
OE A Restoration & \emph{*waskaną} \\
OE Heavy Syllable Nasal Apocope & \emph{*waskan} \\
OE Secondary Nasalization & \emph{*waskąn} \\
OE Weak Tail Reduction & \emph{*waskan} \\
\end{tabular}
\end{minipage}
\\
\end{tabularx}
\end{minipage}%
} \endgroup

Kroonen cites \emph{*waskan-}, Orel \emph{*waskanan}, and Ringe and
Taylor likewise derive Old English \emph{wascan} from the same verb
family (\citeproc{ref-Kroonen2013}{Kroonen 2013, 575};
\citeproc{ref-Orel2003}{Orel 2003, 489};
\citeproc{ref-RingeTaylor2014}{Ringe and Taylor 2014, 142}). Clark Hall
heads the verb as \emph{wascan} (\citeproc{ref-ClarkHall1960}{Clark Hall
1960, 343}), while Sievers-Brunner also notes the variant \emph{wæscan}
(\citeproc{ref-SieversBrunner1965}{Sievers 1965}). From \emph{*wáskaną},
brightening gives \emph{*wæskaną}.

Form note. The conservative citation form \emph{wascan} is selected
here. Spellings such as \emph{wæscan} or \emph{wasċan} belong to variant
or normalized background rather than to the target of this entry.

\subsubsection{wax --- OE weaxan}\label{wax-oe-weaxan}

Derivation: \emph{*wáxsaną} \textgreater{} \emph{weaxan} (regular).

\begingroup
\setlength{\fboxsep}{6pt}

\noindent\fbox{%
\begin{minipage}{0.97\linewidth}
\footnotesize
\begin{tabularx}{\linewidth}{@{}>{\raggedright\arraybackslash}p{0.280\linewidth}>{\centering\arraybackslash}X>{\raggedright\arraybackslash}p{0.560\linewidth}@{}}
\begin{minipage}[t]{\linewidth}
\centering\textbf{Earlier Germanic changes}\par
\vspace{0.35em}
\raggedright
\centering\textbf{West Germanic}\par
\raggedright
\vspace{0.2em}
\raggedright [no change]\par
\vspace{0.6em}
\centering\textbf{Northwest Germanic}\par
\raggedright
\vspace{0.2em}
\raggedright [no change]\par
\end{minipage}
&
&
\begin{minipage}[t]{\linewidth}
\centering\textbf{Old English changes}\par
\vspace{0.35em}
\raggedright
\begin{tabular}{@{}>{\raggedright\arraybackslash}p{0.66\linewidth}@{\hspace{0.55em}}>{\raggedright\arraybackslash}p{0.24\linewidth}@{\hspace{0.25em}}}
Anglo Frisian Brightening & \emph{*wæxsaną} \\
OE Breaking & \emph{*weaxsaną} \\
OE Heavy Syllable Nasal Apocope & \emph{*weaxsan} \\
OE Secondary Nasalization & \emph{*weaxsąn} \\
OE Weak Tail Reduction & \emph{*weaxsan} \\
OE Xs Merge & \emph{*weaXSan} \\
\end{tabular}
\end{minipage}
\\
\end{tabularx}
\end{minipage}%
} \endgroup

Kroonen cites the verb as \emph{*wahs(j)an-}, Orel as \emph{*waxsanan},
and Ringe and Taylor discuss the prehistory of Old English \emph{weaxan}
within the same verbal family (\citeproc{ref-Kroonen2013}{Kroonen 2013,
606}; \citeproc{ref-Orel2003}{Orel 2003, 478};
\citeproc{ref-RingeTaylor2014}{Ringe and Taylor 2014, 191}). Clark Hall
gives \emph{weaxan} as the verb headword and separately records bare
\emph{wax} as a preterite form; Bright likewise treats \emph{weaxan} as
the infinitive (\citeproc{ref-ClarkHall1960}{Clark Hall 1960, 358};
\citeproc{ref-BrightCassidyRingler1971}{Bright 1971, 458}). From
\emph{*wáxsaną}, brightening and breaking yield \emph{weax-}, and the
infinitive develops regularly to \emph{weaxan}. The selected form is the
infinitive \emph{weaxan}. Noun \emph{weax} and preterite wax/\emph{wēox}
belong to different lexical or paradigm slots.

\subsubsection{way --- OE weġ}\label{way-oe-weux121}

Derivation: \emph{*wégaz} \textgreater{} \emph{weġ} (regular).

\begingroup
\setlength{\fboxsep}{6pt}

\noindent\fbox{%
\begin{minipage}{0.97\linewidth}
\small
\begin{tabularx}{\linewidth}{@{}>{\raggedright\arraybackslash}p{0.440\linewidth}>{\centering\arraybackslash}X>{\raggedright\arraybackslash}p{0.440\linewidth}@{}}
\begin{minipage}[t]{\linewidth}
\centering\textbf{Earlier Germanic changes}\par
\vspace{0.35em}
\raggedright
\centering\textbf{West Germanic}\par
\raggedright
\vspace{0.2em}
\raggedright [no change]\par
\vspace{0.6em}
\centering\textbf{Northwest Germanic}\par
\raggedright
\vspace{0.2em}
\begin{tabular}{@{}>{\raggedright\arraybackslash}p{0.74\linewidth}@{\hspace{0.55em}}>{\raggedright\arraybackslash}p{0.16\linewidth}@{\hspace{0.25em}}}
\mbox{PGmc Final Z Deletion} & \emph{*wéga} \\
\end{tabular}
\end{minipage}
&
&
\begin{minipage}[t]{\linewidth}
\centering\textbf{Old English changes}\par
\vspace{0.35em}
\raggedright
\begin{tabular}{@{}>{\raggedright\arraybackslash}p{0.70\linewidth}@{\hspace{0.55em}}>{\raggedright\arraybackslash}p{0.16\linewidth}@{\hspace{0.25em}}}
PWGmc Final Bare A Loss & \emph{*wég} \\
OE Velar Palatalization & \emph{*wéʤ} \\
\end{tabular}
\end{minipage}
\\
\end{tabularx}
\end{minipage}%
} \endgroup

Kroonen cites the noun as \emph{*wega-} `way, road', while the selected
derivational form here is nominative-singular \emph{*wégaz}
(\citeproc{ref-Kroonen2013}{Kroonen 2013}). Campbell, Hogg, and Ringe
and Taylor use the same word as the standard contrast between singular
palatal \emph{weġ} and inflected \emph{wegas/wegum}
(\citeproc{ref-Campbell1959}{Campbell 1959};
\citeproc{ref-Hogg1992}{Hogg 1992}; \citeproc{ref-RingeTaylor2014}{Ringe
and Taylor 2014, 341}). The Old English singular is the ordinary noun
\emph{weg}, here normalized as \emph{weġ} to show the palatal final.
From \emph{*wégaz}, final \emph{*z} is lost and the weak tail
apocopates, leaving word-final \emph{*g} after a front vowel.

Form note. Normalized \emph{weġ} and dictionary \emph{weg} represent the
same noun. \emph{wē} is not supported in the checked Old English
evidence for `way'.

\subsubsection{weapon --- OE wǣpn}\label{weapon-oe-wux1e3pn}

Derivation: \emph{*wḗpną} \textgreater{} \emph{wǣpn} (regular).

\begingroup
\setlength{\fboxsep}{6pt}

\noindent\fbox{%
\begin{minipage}{0.97\linewidth}
\small
\begin{tabularx}{\linewidth}{@{}>{\raggedright\arraybackslash}p{0.460\linewidth}>{\centering\arraybackslash}X>{\raggedright\arraybackslash}p{0.460\linewidth}@{}}
\begin{minipage}[t]{\linewidth}
\centering\textbf{Earlier Germanic changes}\par
\vspace{0.35em}
\raggedright
\centering\textbf{West Germanic}\par
\raggedright
\vspace{0.2em}
\raggedright [no change]\par
\vspace{0.6em}
\centering\textbf{Northwest Germanic}\par
\raggedright
\vspace{0.2em}
\begin{tabular}{@{}>{\raggedright\arraybackslash}p{0.74\linewidth}@{\hspace{0.55em}}>{\raggedright\arraybackslash}p{0.16\linewidth}@{\hspace{0.25em}}}
\mbox{NWGmc Long E Lowering} & \emph{*wǣpną} \\
\end{tabular}
\end{minipage}
&
&
\begin{minipage}[t]{\linewidth}
\centering\textbf{Old English changes}\par
\vspace{0.35em}
\raggedright
\begin{tabular}{@{}>{\raggedright\arraybackslash}p{0.74\linewidth}@{\hspace{0.55em}}>{\raggedright\arraybackslash}p{0.16\linewidth}@{\hspace{0.25em}}}
OE Heavy Syllable Nasal Apocope & \emph{*wǣpn} \\
\end{tabular}
\end{minipage}
\\
\end{tabularx}
\end{minipage}%
} \endgroup

Kroonen reconstructs a double-stem noun \emph{*wēbna- \textasciitilde{}
wēpna-} and cites Old English \emph{wæpn} among the reflexes
(\citeproc{ref-Kroonen2013}{Kroonen 2013, 617};
\citeproc{ref-Campbell1959}{Campbell 1959, 150};
\citeproc{ref-Campbell1959}{1959, 226}). Campbell preserves unbroken
\emph{wépn} beside broken \emph{wépen}-type forms. Bright contrasts
broken \emph{wǣpen/wapen} with unbroken \emph{wǣpnes}, and Clark Hall
lemmatizes the noun under \emph{wapen} while preserving unbroken forms
elsewhere (\citeproc{ref-BrightCassidyRingler1971}{Bright 1971, 29};
\citeproc{ref-ClarkHall1960}{Clark Hall 1960, 355}). The selected form
\emph{wǣpn} represents the unbroken line from \emph{*wḗpną}.

Form note. \emph{Wǣpen} is the usual late West Saxon simplex headword,
while \emph{wǣpnes} and related forms preserve the unbroken cluster
(\citeproc{ref-Campbell1959}{Campbell 1959, 150};
\citeproc{ref-Campbell1959}{1959, 226};
\citeproc{ref-BrightCassidyRingler1971}{Bright 1971, 29};
\citeproc{ref-ClarkHall1960}{Clark Hall 1960, 355}).

\subsubsection{will --- OE willa}\label{will-oe-willa}

Derivation: \emph{*wéljô} \textgreater{} \emph{willa} (regular).

\begingroup
\setlength{\fboxsep}{6pt}

\noindent\fbox{%
\begin{minipage}{0.97\linewidth}
\small
\begin{tabularx}{\linewidth}{@{}>{\raggedright\arraybackslash}p{0.300\linewidth}>{\centering\arraybackslash}X>{\raggedright\arraybackslash}p{0.540\linewidth}@{}}
\begin{minipage}[t]{\linewidth}
\centering\textbf{Earlier Germanic changes}\par
\vspace{0.35em}
\raggedright
\centering\textbf{West Germanic}\par
\raggedright
\vspace{0.2em}
\begin{tabular}{@{}>{\raggedright\arraybackslash}p{0.64\linewidth}@{\hspace{0.55em}}>{\raggedright\arraybackslash}p{0.16\linewidth}@{\hspace{0.25em}}}
PWGmc J Gemination & \emph{*wélljô} \\
\end{tabular}
\vspace{0.6em}
\centering\textbf{Northwest Germanic}\par
\raggedright
\vspace{0.2em}
\raggedright [no change]\par
\end{minipage}
&
&
\begin{minipage}[t]{\linewidth}
\centering\textbf{Old English changes}\par
\vspace{0.35em}
\raggedright
\begin{tabular}{@{}>{\raggedright\arraybackslash}p{0.74\linewidth}@{\hspace{0.55em}}>{\raggedright\arraybackslash}p{0.16\linewidth}@{\hspace{0.25em}}}
OE I Umlaut & \emph{*willjô} \\
OE Unstressed Long Vowel Shortening & \emph{*willja} \\
OE J Loss After Heavy & \emph{*willa} \\
\end{tabular}
\end{minipage}
\\
\end{tabularx}
\end{minipage}%
} \endgroup

Kroonen separates noun \emph{*weljan-} 2 `will, wish' from verb
\emph{*weljan-} 1 `to want', while Orel and Kluge represent the noun as
\emph{*weljōn} (\citeproc{ref-Kroonen2013}{Kroonen 2013};
\citeproc{ref-Orel2003}{Orel 2003};
\citeproc{ref-KlugeSeebold2011}{Kluge 2011}). The selected derivational
form \emph{*wéljô} is the noun-side input used for this row. Clark Hall
lemmatizes noun \emph{willa m.} separately from verb \emph{willan}
(\citeproc{ref-ClarkHall1960}{Clark Hall 1960, 368}). From
\emph{*wéljô}, j-gemination yields a heavy stem, i-umlaut gives
\emph{will-}, and later shortening plus j-loss produce \emph{willa}. The
selected form is the noun \emph{willa} `will, wish'. Related verb
\emph{willan} belongs to a separate lexeme and should not be substituted
for the noun row (\citeproc{ref-Kroonen2013}{Kroonen 2013};
\citeproc{ref-ClarkHall1960}{Clark Hall 1960, 368}).

\subsubsection{wind --- OE windan}\label{wind-oe-windan}

Derivation: \emph{*wíndaną} \textgreater{} \emph{windan} (regular).

\begingroup
\setlength{\fboxsep}{6pt}

\noindent\fbox{%
\begin{minipage}{0.97\linewidth}
\small
\begin{tabularx}{\linewidth}{@{}>{\raggedright\arraybackslash}p{0.300\linewidth}>{\centering\arraybackslash}X>{\raggedright\arraybackslash}p{0.540\linewidth}@{}}
\begin{minipage}[t]{\linewidth}
\centering\textbf{Earlier Germanic changes}\par
\vspace{0.35em}
\raggedright
\centering\textbf{West Germanic}\par
\raggedright
\vspace{0.2em}
\raggedright [no change]\par
\vspace{0.6em}
\centering\textbf{Northwest Germanic}\par
\raggedright
\vspace{0.2em}
\raggedright [no change]\par
\end{minipage}
&
&
\begin{minipage}[t]{\linewidth}
\centering\textbf{Old English changes}\par
\vspace{0.35em}
\raggedright
\begin{tabular}{@{}>{\raggedright\arraybackslash}p{0.74\linewidth}@{\hspace{0.55em}}>{\raggedright\arraybackslash}p{0.16\linewidth}@{\hspace{0.25em}}}
OE Heavy Syllable Nasal Apocope & \emph{*wíndan} \\
OE Secondary Nasalization & \emph{*wíndąn} \\
OE Weak Tail Reduction & \emph{*wíndan} \\
\end{tabular}
\end{minipage}
\\
\end{tabularx}
\end{minipage}%
} \endgroup

Kroonen distinguishes noun \emph{*winda-} from verb \emph{*windan-}, and
the present row belongs to the verb (\citeproc{ref-Kroonen2013}{Kroonen
2013}). Later handbook discussion keeps the dental original from PIE
\emph{*wendh-}, not a Verner alternant (\citeproc{ref-Fulk2018}{Fulk
2018}; \citeproc{ref-RingeTaylor2014}{Ringe and Taylor 2014}). Clark
Hall and Bosworth-Toller record \emph{windan} as the verb headword
(\citeproc{ref-ClarkHall1960}{Clark Hall 1960};
\citeproc{ref-BosworthToller1898}{Bosworth and Toller 1898, 101}). The
selected input \emph{*wíndaną} yields the regular infinitive
\emph{windan} by ordinary heavy-syllable apocope and weak-tail
reduction. English \emph{wind} also names the noun. This row targets the
class-III verb, not the noun (\citeproc{ref-Kroonen2013}{Kroonen 2013};
\citeproc{ref-ClarkHall1960}{Clark Hall 1960}).

\subsubsection{wold --- OE weald}\label{wold-oe-weald}

Derivation: \emph{*wálθuz} \textgreater{} \emph{weald} (regular).

\begingroup
\setlength{\fboxsep}{6pt}

\noindent\fbox{%
\begin{minipage}{0.97\linewidth}
\small
\begin{tabularx}{\linewidth}{@{}>{\raggedright\arraybackslash}p{0.460\linewidth}>{\centering\arraybackslash}X>{\raggedright\arraybackslash}p{0.460\linewidth}@{}}
\begin{minipage}[t]{\linewidth}
\centering\textbf{Earlier Germanic changes}\par
\vspace{0.35em}
\raggedright
\centering\textbf{West Germanic}\par
\raggedright
\vspace{0.2em}
\begin{tabular}{@{}>{\raggedright\arraybackslash}p{0.74\linewidth}@{\hspace{0.55em}}>{\raggedright\arraybackslash}p{0.16\linewidth}@{\hspace{0.25em}}}
PWGmc L Th Voicing & \emph{*wálduz} \\
\end{tabular}
\vspace{0.6em}
\centering\textbf{Northwest Germanic}\par
\raggedright
\vspace{0.2em}
\begin{tabular}{@{}>{\raggedright\arraybackslash}p{0.74\linewidth}@{\hspace{0.55em}}>{\raggedright\arraybackslash}p{0.16\linewidth}@{\hspace{0.25em}}}
\mbox{PGmc Final Z Deletion} & \emph{*wáldu} \\
\end{tabular}
\end{minipage}
&
&
\begin{minipage}[t]{\linewidth}
\centering\textbf{Old English changes}\par
\vspace{0.35em}
\raggedright
\begin{tabular}{@{}>{\raggedright\arraybackslash}p{0.74\linewidth}@{\hspace{0.55em}}>{\raggedright\arraybackslash}p{0.16\linewidth}@{\hspace{0.25em}}}
Anglo Frisian Brightening & \emph{*wældu} \\
OE Breaking & \emph{*wealdu} \\
\mbox{OE High Vowel Apocope} & \emph{*weald} \\
\end{tabular}
\end{minipage}
\\
\end{tabularx}
\end{minipage}%
} \endgroup

Kroonen reconstructs the noun as \emph{*walþu-} and gives OE
\emph{weald} beside other West Germanic \emph{wald} forms
(\citeproc{ref-Kroonen2013}{Kroonen 2013}). The selected input
\emph{*wálθuz} is the nominative singular used for the derivation. Clark
Hall makes \emph{weald} the main noun headword and cross-refers
\emph{wald} and \emph{wold} to it (\citeproc{ref-ClarkHall1960}{Clark
Hall 1960}). \emph{*lþ} voices to \emph{ld}, Anglo-Frisian brightening
yields \emph{wæld-}, and breaking before the cluster gives
\emph{weald-}; apocope then yields \emph{weald}.

Dialect note. \emph{wald} survives as an Anglian-type variant in the
same family. The selected target is normalized \emph{weald}, not the
variant form (\citeproc{ref-ClarkHall1960}{Clark Hall 1960};
\citeproc{ref-RingeTaylor2014}{Ringe and Taylor 2014}).

\subsubsection{yarn --- OE ġearn}\label{yarn-oe-ux121earn}

Derivation: \emph{*gárną} \textgreater{} \emph{ġearn} (regular).

\begingroup
\setlength{\fboxsep}{6pt}

\noindent\fbox{%
\begin{minipage}{0.97\linewidth}
\small
\begin{tabularx}{\linewidth}{@{}>{\raggedright\arraybackslash}p{0.300\linewidth}>{\centering\arraybackslash}X>{\raggedright\arraybackslash}p{0.540\linewidth}@{}}
\begin{minipage}[t]{\linewidth}
\centering\textbf{Earlier Germanic changes}\par
\vspace{0.35em}
\raggedright
\centering\textbf{West Germanic}\par
\raggedright
\vspace{0.2em}
\raggedright [no change]\par
\vspace{0.6em}
\centering\textbf{Northwest Germanic}\par
\raggedright
\vspace{0.2em}
\raggedright [no change]\par
\end{minipage}
&
&
\begin{minipage}[t]{\linewidth}
\centering\textbf{Old English changes}\par
\vspace{0.35em}
\raggedright
\begin{tabular}{@{}>{\raggedright\arraybackslash}p{0.74\linewidth}@{\hspace{0.55em}}>{\raggedright\arraybackslash}p{0.16\linewidth}@{\hspace{0.25em}}}
Anglo Frisian Brightening & \emph{*gærną} \\
OE Breaking & \emph{*gearną} \\
OE Heavy Syllable Nasal Apocope & \emph{*gearn} \\
OE Velar Palatalization & \emph{*ʤearn} \\
\end{tabular}
\end{minipage}
\\
\end{tabularx}
\end{minipage}%
} \endgroup

Kroonen cites the noun as \emph{*garna-}, and Ringe and Taylor give the
early chain \emph{*garna} \textgreater{} \emph{*geern} \textgreater{}
\emph{*gearn} \textgreater{} OE \emph{gearn}
(\citeproc{ref-Kroonen2013}{Kroonen 2013, 209};
\citeproc{ref-RingeTaylor2014}{Ringe and Taylor 2014, 220}). Clark Hall
records \emph{gearn (e) n.} `yarn, spun wool', and Bosworth-Toller
glosses \emph{gearn} as \emph{filatum}
(\citeproc{ref-ClarkHall1960}{Clark Hall 1960, 145};
\citeproc{ref-BosworthToller1898}{Bosworth and Toller 1898}). From
\emph{*gárną}, brightening and breaking before \emph{rn} yield
\emph{gearn}; palatalization of initial \emph{g} before the resulting
front-vocalic sequence gives normalized \emph{ġearn}.

Form note. Dictionary \emph{gearn} and normalized \emph{ġearn} refer to
the same noun. The comparative stem \emph{*garna-} and oblique
\emph{*garnăn} do not replace the selected input \emph{*gárną}.

\clearpage

\subsection{Part II. Attested variants and selected comparison
forms}\label{part-ii.-attested-variants-and-selected-comparison-forms}

These entries treat the selected Old English target as one member of an
attested or historically documented variant set. The target is therefore
anchored in the record, but the lexical comparison must account for
variation rather than for a single unproblematic citation form.

\subsubsection{cud --- OE cwedu}\label{cud-oe-cwedu}

Derivation: citation reconstruction \emph{*kwíθuz}; selected input
\emph{*kwéðuz} \textgreater{} \emph{cwedu} (attested variant).

\paragraph{Derivation trace}\label{derivation-trace}

Proto input: \emph{*kwéðuz}

\begingroup
\setlength{\fboxsep}{6pt}

\noindent\fbox{%
\begin{minipage}{0.97\linewidth}
\small
\begin{tabularx}{\linewidth}{@{}>{\raggedright\arraybackslash}p{0.540\linewidth}>{\centering\arraybackslash}X>{\raggedright\arraybackslash}p{0.300\linewidth}@{}}
\begin{minipage}[t]{\linewidth}
\centering\textbf{Earlier Germanic changes}\par
\vspace{0.35em}
\raggedright
\centering\textbf{West Germanic}\par
\raggedright
\vspace{0.2em}
\begin{tabular}{@{}>{\raggedright\arraybackslash}p{0.74\linewidth}@{\hspace{0.55em}}>{\raggedright\arraybackslash}p{0.16\linewidth}@{\hspace{0.25em}}}
PWGmc Dental Hardening & \emph{*kwéduz} \\
\end{tabular}
\vspace{0.6em}
\centering\textbf{Northwest Germanic}\par
\raggedright
\vspace{0.2em}
\begin{tabular}{@{}>{\raggedright\arraybackslash}p{0.74\linewidth}@{\hspace{0.55em}}>{\raggedright\arraybackslash}p{0.16\linewidth}@{\hspace{0.25em}}}
\mbox{PGmc Final Z Deletion} & \emph{*kwédu} \\
\end{tabular}
\end{minipage}
&
&
\begin{minipage}[t]{\linewidth}
\centering\textbf{Old English changes}\par
\vspace{0.35em}
\raggedright
\raggedright [no change]\par
\end{minipage}
\\
\end{tabularx}
\end{minipage}%
} \endgroup

Outcome: \emph{cwedu}

\paragraph{Reconstruction and comparative
evidence}\label{reconstruction-and-comparative-evidence}

Kroonen reconstructs the resin word as \emph{*kwedu-2} and gives Old
English variants \emph{cwidu}, \emph{cweodu}, and \emph{c(w)udu}
(\citeproc{ref-Kroonen2013}{Kroonen 2013, 355}). Orel likewise lists
\emph{cwidu} under the cognate set (\citeproc{ref-Orel2003}{Orel 2003,
266}). The selected input \emph{*kwéðuz} therefore represents the older
e-grade, voiced-dental form behind the chosen variant \emph{cwedu}.

\paragraph{Old English evidence}\label{old-english-evidence}

The Old English word survives in a wider variant set than one dictionary
headword suggests. Ringe and Taylor discuss \emph{cwidu} \textgreater{}
\emph{cwudu} \textgreater{} \emph{cudu} and also note late West Saxon
\emph{cweodu}; Clark Hall gives \emph{cwudu}, \emph{cweodu}, and
\emph{cudu} (\citeproc{ref-RingeTaylor2014}{Ringe and Taylor 2014, 338};
\citeproc{ref-ClarkHall1960}{Clark Hall 1960, 84}). Attested
\emph{cwedu} is treated here as the conservative variant within that
set.

\paragraph{Development to Old English}\label{development-to-old-english}

From \emph{*kwéðuz}, the West Germanic voiced dental hardens in the
expected way and the regular Old English development yields
\emph{cwedu}. The other Old English spellings belong to the same lexical
family, but reflect later leveling, back-umlaut, or further reduction
rather than a need to replace the selected input.

\paragraph{Variant comparison}\label{variant-comparison}

{\def\LTcaptype{none} % do not increment counter
\begin{longtable}[]{@{}
  >{\raggedright\arraybackslash}p{(\linewidth - 4\tabcolsep) * \real{0.3333}}
  >{\raggedright\arraybackslash}p{(\linewidth - 4\tabcolsep) * \real{0.3333}}
  >{\raggedright\arraybackslash}p{(\linewidth - 4\tabcolsep) * \real{0.3333}}@{}}
\toprule\noalign{}
\begin{minipage}[b]{\linewidth}\raggedright
Variant type
\end{minipage} & \begin{minipage}[b]{\linewidth}\raggedright
Old English form
\end{minipage} & \begin{minipage}[b]{\linewidth}\raggedright
Comment
\end{minipage} \\
\midrule\noalign{}
\endhead
\bottomrule\noalign{}
\endlastfoot
conservative target & \emph{cwedu} & selected attested variant
represented here \\
leveled i-grade form & \emph{cwidu} & common lexical variant in the same
family \\
back-umlauted forms & \emph{cweodu}, \emph{cwudu} & later developments
within the same OE tradition \\
reduced form & \emph{cudu} & further reduced member of the same variant
set \\
\end{longtable}
}

\subsubsection{ten --- OE tēon}\label{ten-oe-tux113on}

Derivation: \emph{*téxun} \textgreater{} \emph{tēon} (attested variant).

\paragraph{Derivation trace}\label{derivation-trace-1}

Proto input: \emph{*téxun}

\begingroup
\setlength{\fboxsep}{6pt}

\noindent\fbox{%
\begin{minipage}{0.97\linewidth}
\small
\begin{tabularx}{\linewidth}{@{}>{\raggedright\arraybackslash}p{0.300\linewidth}>{\centering\arraybackslash}X>{\raggedright\arraybackslash}p{0.540\linewidth}@{}}
\begin{minipage}[t]{\linewidth}
\centering\textbf{Earlier Germanic changes}\par
\vspace{0.35em}
\raggedright
\centering\textbf{West Germanic}\par
\raggedright
\vspace{0.2em}
\raggedright [no change]\par
\vspace{0.6em}
\centering\textbf{Northwest Germanic}\par
\raggedright
\vspace{0.2em}
\raggedright [no change]\par
\end{minipage}
&
&
\begin{minipage}[t]{\linewidth}
\centering\textbf{Old English changes}\par
\vspace{0.35em}
\raggedright
\begin{tabular}{@{}>{\raggedright\arraybackslash}p{0.74\linewidth}@{\hspace{0.55em}}>{\raggedright\arraybackslash}p{0.16\linewidth}@{\hspace{0.25em}}}
OE Med Unstressed U Lowering & \emph{*téxon} \\
OE Breaking & \emph{*téoxon} \\
OE H Loss & \emph{*téoon} \\
OE Contraction & \emph{*tḗon} \\
\end{tabular}
\end{minipage}
\\
\end{tabularx}
\end{minipage}%
} \endgroup

Outcome: \emph{tēon}

\paragraph{Reconstruction and comparative
evidence}\label{reconstruction-and-comparative-evidence-1}

Fulk states that Old English \emph{tien} shows umlaut from the inflected
forms, whereas the uninflected form without umlaut is reflected in
\emph{hund-tēon-tig} (\citeproc{ref-Fulk2018}{Fulk 2018, sect. 10.2}).
Brunner gives the same contrast more broadly: \emph{tēon} develops from
\emph{*tëhun}, while West Saxon tien, \emph{tȳn} belong to a different,
umlauted branch of the numeral history
(\citeproc{ref-SieversBrunner1965}{Sievers 1965, sects 129.2, 129} Anm.
6, 234).

The selected comparison form \emph{tēon} therefore represents the bare
cardinal's un-umlauted line, not the later umlauted simplex tradition.

\paragraph{Old English evidence}\label{old-english-evidence-1}

The attested simplex forms are varied. Campbell gives \emph{tien},
north-western West Saxon \emph{tēn}, and late Northumbrian \emph{tēo},
\emph{tēa} (\citeproc{ref-Campbell1959}{Campbell 1959, sect. 682}).
Brunner likewise lists West Saxon tien, \emph{tȳn} beside \emph{tēn},
\emph{tēo}, \emph{tēa} in other dialects
(\citeproc{ref-SieversBrunner1965}{Sievers 1965, sect. 325}).

Exact simplex \emph{tēon} is weaker as a directly cited headword than
those spellings. The un-umlauted stem is, however, explicit in
\emph{tēoða} and \emph{-tēontig}
(\citeproc{ref-SieversBrunner1965}{Sievers 1965, sect. 129.2};
\citeproc{ref-Fulk2018}{Fulk 2018, sect. 10.2}). The comparison form
\emph{tēon} is therefore a normalized spelling of that un-umlauted base.

\paragraph{Development to Old
English}\label{development-to-old-english-1}

From \emph{*téxun}, lowering of medial unstressed \emph{u} gives
\emph{*téxon}, breaking gives \emph{*téoxon}, loss of intervocalic
\emph{h}/\emph{x} yields \emph{*téoon}, and contraction produces
\emph{*tḗon}, written \emph{tēon}. This is the regular bare-cardinal
path.

The umlauted forms \emph{tien} / \emph{tīen} belong to a different
branch, created when the numeral was levelled from inflected forms with
a front-vocalic trigger (\citeproc{ref-Fulk2018}{Fulk 2018, sect. 10.2};
\citeproc{ref-SieversBrunner1965}{Sievers 1965, sect. 129} Anm. 6).

\paragraph{Variant comparison}\label{variant-comparison-1}

The comparison below is manual. It distinguishes the normalized
un-umlauted base from the attested simplex variants.

{\def\LTcaptype{none} % do not increment counter
\begin{longtable}[]{@{}
  >{\raggedright\arraybackslash}p{(\linewidth - 4\tabcolsep) * \real{0.3333}}
  >{\raggedright\arraybackslash}p{(\linewidth - 4\tabcolsep) * \real{0.3333}}
  >{\raggedright\arraybackslash}p{(\linewidth - 4\tabcolsep) * \real{0.3333}}@{}}
\toprule\noalign{}
\begin{minipage}[b]{\linewidth}\raggedright
Form or branch
\end{minipage} & \begin{minipage}[b]{\linewidth}\raggedright
Status
\end{minipage} & \begin{minipage}[b]{\linewidth}\raggedright
Relevance to this entry
\end{minipage} \\
\midrule\noalign{}
\endhead
\bottomrule\noalign{}
\endlastfoot
\emph{tēon} & normalized un-umlauted comparison form; trace-supported &
selected target \\
\emph{tien} / \emph{tīen} & attested West Saxon umlauted simplex forms &
genuine OE variants, but not the bare-cardinal line modeled here \\
\emph{tēn} / \emph{tēo} / \emph{tēa} & attested un-umlauted simplex
variants in other dialects & support the same branch as the selected
comparison form \\
\end{longtable}
}

\subsubsection{three --- OE þrīe}\label{three-oe-uxferux12be}

Derivation: \emph{*θréjez} \textgreater{} \emph{þrīe} (attested
variant).

\paragraph{Derivation trace}\label{derivation-trace-2}

Proto input: \emph{*θréjez}

\begingroup
\setlength{\fboxsep}{6pt}

\noindent\fbox{%
\begin{minipage}{0.97\linewidth}
\small
\begin{tabularx}{\linewidth}{@{}>{\raggedright\arraybackslash}p{0.300\linewidth}>{\centering\arraybackslash}X>{\raggedright\arraybackslash}p{0.540\linewidth}@{}}
\begin{minipage}[t]{\linewidth}
\centering\textbf{Earlier Germanic changes}\par
\vspace{0.35em}
\raggedright
\centering\textbf{West Germanic}\par
\raggedright
\vspace{0.2em}
\raggedright [no change]\par
\vspace{0.6em}
\centering\textbf{Northwest Germanic}\par
\raggedright
\vspace{0.2em}
\begin{tabular}{@{}>{\raggedright\arraybackslash}p{0.74\linewidth}@{\hspace{0.55em}}>{\raggedright\arraybackslash}p{0.16\linewidth}@{\hspace{0.25em}}}
\mbox{PGmc Final Z Deletion} & \emph{*θréje} \\
\end{tabular}
\end{minipage}
&
&
\begin{minipage}[t]{\linewidth}
\centering\textbf{Old English changes}\par
\vspace{0.35em}
\raggedright
\begin{tabular}{@{}>{\raggedright\arraybackslash}p{0.74\linewidth}@{\hspace{0.55em}}>{\raggedright\arraybackslash}p{0.16\linewidth}@{\hspace{0.25em}}}
OE I Umlaut & \emph{*θrije} \\
OE Intervocalic J Vocalization & \emph{*θriie} \\
OE Contraction & \emph{*θrīe} \\
\end{tabular}
\end{minipage}
\\
\end{tabularx}
\end{minipage}%
} \endgroup

Outcome: \emph{þrīe}

\paragraph{Reconstruction and comparative
evidence}\label{reconstruction-and-comparative-evidence-2}

Kroonen cites the numeral under a broader stem-style reconstruction
rather than under one Old English-ready paradigm cell
(\citeproc{ref-Kroonen2013}{Kroonen 2013}). The input \emph{*θréjez} is
therefore best understood as the inherited masculine
nominative-accusative singular.

That distinction matters because the Old English numeral does not have
one uniform citation form across the paradigm. The masculine singular
line must be kept apart from feminine-neuter \emph{þrēo} and from later
reduced spellings of the masculine form.

\paragraph{Old English evidence}\label{old-english-evidence-2}

Campbell gives masculine nominative-accusative \emph{þrīe}, feminine and
neuter nominative-accusative \emph{þrēo}, genitive \emph{þrēora}, and
dative \emph{þrim}, adding that late West Saxon has \emph{þry},
\emph{þri} for \emph{þrīe} (\citeproc{ref-Campbell1959}{Campbell 1959,
sect. 683}). Fulk presents the same masculine \emph{þrīe} beside the
wider numeral paradigm (\citeproc{ref-Fulk2018}{Fulk 2018, sect. 10.1}).

The target is therefore an attested Old English paradigm cell.
\emph{þrī} belongs to later reduction or headword-style citation,
whereas \emph{þrīe} is the conservative masculine nominative-accusative
form.

\paragraph{Development to Old
English}\label{development-to-old-english-2}

From \emph{*θréjez}, loss of final \emph{-z} leaves a form of the
\emph{*θréje} type. The following \emph{j} fronts the stem vowel, then
vocalizes between vowels, and contraction yields \emph{þrīe}. The
compact trace records the same sequence as \emph{*θrije} \textgreater{}
\emph{*θriie} \textgreater{} \emph{þrīe}.

\paragraph{Variant comparison}\label{variant-comparison-2}

The comparison below is manual. It separates the selected masculine cell
from the later reduced form and from the rest of the numeral paradigm.

{\def\LTcaptype{none} % do not increment counter
\begin{longtable}[]{@{}
  >{\raggedright\arraybackslash}p{(\linewidth - 4\tabcolsep) * \real{0.3333}}
  >{\raggedright\arraybackslash}p{(\linewidth - 4\tabcolsep) * \real{0.3333}}
  >{\raggedright\arraybackslash}p{(\linewidth - 4\tabcolsep) * \real{0.3333}}@{}}
\toprule\noalign{}
\begin{minipage}[b]{\linewidth}\raggedright
Form
\end{minipage} & \begin{minipage}[b]{\linewidth}\raggedright
Status
\end{minipage} & \begin{minipage}[b]{\linewidth}\raggedright
Relevance to this entry
\end{minipage} \\
\midrule\noalign{}
\endhead
\bottomrule\noalign{}
\endlastfoot
\emph{þrīe} & attested masculine nom./acc.; trace output & selected
target \\
\emph{þrī} / \emph{þry} & later reduced masculine variant & genuine OE
variant, but not the conservative comparison form \\
\emph{þrēo} & attested feminine-neuter nom./acc. & same numeral,
different paradigm cell \\
\emph{þrēora}, \emph{þrim} & attested genitive and dative forms &
confirm the wider paradigm, not the selected cell \\
\end{longtable}
}

\subsubsection{wasp --- OE wæfs}\label{wasp-oe-wuxe6fs}

Derivation: \emph{*wábsaz} \textgreater{} \emph{wæfs} (attested
variant).

\paragraph{Derivation trace}\label{derivation-trace-3}

Proto input: \emph{*wábsaz}

\begingroup
\setlength{\fboxsep}{6pt}

\noindent\fbox{%
\begin{minipage}{0.97\linewidth}
\small
\begin{tabularx}{\linewidth}{@{}>{\raggedright\arraybackslash}p{0.300\linewidth}>{\centering\arraybackslash}X>{\raggedright\arraybackslash}p{0.540\linewidth}@{}}
\begin{minipage}[t]{\linewidth}
\centering\textbf{Earlier Germanic changes}\par
\vspace{0.35em}
\raggedright
\centering\textbf{West Germanic}\par
\raggedright
\vspace{0.2em}
\raggedright [no change]\par
\vspace{0.6em}
\centering\textbf{Northwest Germanic}\par
\raggedright
\vspace{0.2em}
\begin{tabular}{@{}>{\raggedright\arraybackslash}p{0.74\linewidth}@{\hspace{0.55em}}>{\raggedright\arraybackslash}p{0.16\linewidth}@{\hspace{0.25em}}}
\mbox{PGmc Final Z Deletion} & \emph{*wábsa} \\
\end{tabular}
\end{minipage}
&
&
\begin{minipage}[t]{\linewidth}
\centering\textbf{Old English changes}\par
\vspace{0.35em}
\raggedright
\begin{tabular}{@{}>{\raggedright\arraybackslash}p{0.70\linewidth}@{\hspace{0.55em}}>{\raggedright\arraybackslash}p{0.16\linewidth}@{\hspace{0.25em}}}
PWGmc Final Bare A Loss & \emph{*wábs} \\
Anglo Frisian Brightening & \emph{*wæbs} \\
PGmc B Allophony & \emph{*wæβs} \\
\end{tabular}
\end{minipage}
\\
\end{tabularx}
\end{minipage}%
} \endgroup

Outcome: \emph{wæfs}

\paragraph{Reconstruction and comparative
evidence}\label{reconstruction-and-comparative-evidence-3}

The Proto-Germanic form \emph{*wábsaz} reaches Old English without any
special change of stem or paradigm cell. The question in this entry is
instead which attested Old English member of the variant set should
serve as the comparison form.

Fulk presents the Old English forms together as \emph{wæfs} with
variants \emph{wæsp} and \emph{wæps} (\citeproc{ref-Fulk2018}{Fulk 2018,
sect. 6.5}). Bülbring and Brunner then make the chronology more explicit
by deriving later \emph{wæps} and late West Saxon \emph{wasp} from
earlier \emph{waefs} / \emph{wæfs} through restricted metatheses
(\citeproc{ref-Bulbring1902}{Bülbring 1902, sect. 484} Anm. 3;
\citeproc{ref-SieversBrunner1965}{Sievers 1965, sects 193, 204}).

\paragraph{Old English evidence}\label{old-english-evidence-3}

The earliest directly cited Old English form is \emph{wæfs}, written
\emph{waefs} in the Épinal-Corpus material discussed by Bülbring and
Brunner (\citeproc{ref-Bulbring1902}{Bülbring 1902, sect. 484} Anm. 3;
\citeproc{ref-SieversBrunner1965}{Sievers 1965, sect. 193}). Later Old
English also shows \emph{wæps} and \emph{wæsp} / \emph{wasp}, and
dictionary practice often favors \emph{wæps} or later spellings as
headwords (\citeproc{ref-ClarkHall1960}{Clark Hall 1960, 341}).

This entry therefore distinguishes chronological priority from headword
habit. \emph{wæfs} is not a convenient reconstruction: it is an attested
Old English form and also the one that matches the regular development
most closely.

\paragraph{Development to Old
English}\label{development-to-old-english-3}

From \emph{*wábsaz}, the regular Old English path passes through loss of
final \emph{z}, Anglo-Frisian fronting, and the allophonic development
of \emph{b} to a fricative before \emph{s}, yielding \emph{wæfs}.

The later forms \emph{wæps} and \emph{wæsp} / \emph{wasp} belong to
subsequent, lexically restricted metatheses. They are genuine Old
English forms, but they are later within the variant history.

\paragraph{Variant comparison}\label{variant-comparison-3}

The comparison below is manual. It separates the earliest attested and
regular form from the later metathesized doublets.

{\def\LTcaptype{none} % do not increment counter
\begin{longtable}[]{@{}
  >{\raggedright\arraybackslash}p{(\linewidth - 4\tabcolsep) * \real{0.3333}}
  >{\raggedright\arraybackslash}p{(\linewidth - 4\tabcolsep) * \real{0.3333}}
  >{\raggedright\arraybackslash}p{(\linewidth - 4\tabcolsep) * \real{0.3333}}@{}}
\toprule\noalign{}
\begin{minipage}[b]{\linewidth}\raggedright
Form
\end{minipage} & \begin{minipage}[b]{\linewidth}\raggedright
Status
\end{minipage} & \begin{minipage}[b]{\linewidth}\raggedright
Relevance to this entry
\end{minipage} \\
\midrule\noalign{}
\endhead
\bottomrule\noalign{}
\endlastfoot
\emph{wæfs} & earliest attested OE form; regular trace output & selected
target \\
\emph{wæps} & later attested metathesized variant & genuine OE doublet,
but secondary \\
\emph{wæsp} / \emph{wasp} & later West Saxon metathesized variant &
genuine OE doublet, but not the selected form \\
\end{longtable}
}

\clearpage

\subsection{Part III. Early analogy and pre-Old-English input
selection}\label{part-iii.-early-analogy-and-pre-old-english-input-selection}

These entries involve a distinction between the lexeme-level citation
reconstruction and the form selected as input to the Old English
derivation. The issue is upstream of Old English: the selected input
represents the pre-Old-English form that gives the attested target under
the current cascade.

\subsubsection{bottom --- OE botm}\label{bottom-oe-botm}

Derivation: citation reconstruction \emph{*búdmaz}; selected input
\emph{*búttmaz} \textgreater{} \emph{botm} (early analogy).

\paragraph{Derivation trace}\label{derivation-trace-4}

Proto input: \emph{*búttmaz}

\begingroup
\setlength{\fboxsep}{6pt}

\noindent\fbox{%
\begin{minipage}{0.97\linewidth}
\small
\begin{tabularx}{\linewidth}{@{}>{\raggedright\arraybackslash}p{0.440\linewidth}>{\centering\arraybackslash}X>{\raggedright\arraybackslash}p{0.440\linewidth}@{}}
\begin{minipage}[t]{\linewidth}
\centering\textbf{Earlier Germanic changes}\par
\vspace{0.35em}
\raggedright
\centering\textbf{West Germanic}\par
\raggedright
\vspace{0.2em}
\raggedright [no change]\par
\vspace{0.6em}
\centering\textbf{Northwest Germanic}\par
\raggedright
\vspace{0.2em}
\begin{tabular}{@{}>{\raggedright\arraybackslash}p{0.70\linewidth}@{\hspace{0.55em}}>{\raggedright\arraybackslash}p{0.20\linewidth}@{\hspace{0.25em}}}
\mbox{NWGmc U Lowering} & \emph{*bóttmaz} \\
\mbox{PGmc Final Z Deletion} & \emph{*bóttma} \\
\end{tabular}
\end{minipage}
&
&
\begin{minipage}[t]{\linewidth}
\centering\textbf{Old English changes}\par
\vspace{0.35em}
\raggedright
\begin{tabular}{@{}>{\raggedright\arraybackslash}p{0.74\linewidth}@{\hspace{0.55em}}>{\raggedright\arraybackslash}p{0.16\linewidth}@{\hspace{0.25em}}}
PWGmc Final Bare A Loss & \emph{*bóttm} \\
OE Preconsonantal Degemination & \emph{*bótm} \\
\end{tabular}
\end{minipage}
\\
\end{tabularx}
\end{minipage}%
} \endgroup

Outcome: \emph{botm}

\paragraph{Reconstruction and comparative
evidence}\label{reconstruction-and-comparative-evidence-4}

Kroonen reconstructs the word as a stem complex \emph{*budmō}, gen.
\emph{*buttaz}, summarized as \emph{*budman-} \textasciitilde{}
\emph{*buttman-}, and gives Old English \emph{botm} as the reflex
(\citeproc{ref-Kroonen2013}{Kroonen 2013, 120}). The comparative label
\emph{*búdmaz} names the lexeme-level stem complex, while the selected
input \emph{*búttmaz} represents the pre-Old-English form with oblique
\emph{*butt-} generalized into the nominative formation.

Orel likewise preserves both sides of the comparison under
\emph{*budmaz} \emph{*butmaz} (\citeproc{ref-Orel2003}{Orel 2003, 100}).
The selected input is thus a historical stem choice, not an arbitrary
respelling.

\paragraph{Old English evidence}\label{old-english-evidence-4}

The Old English noun itself is secure. Clark Hall gives \emph{botm}
(\citeproc{ref-ClarkHall1960}{Clark Hall 1960, 63}). Bosworth-Toller
cross-references \emph{bodan} to \emph{botm}, showing the wider reflex
family without weakening the attested lemma
(\citeproc{ref-BosworthToller1898}{Bosworth and Toller 1898, 112}).

\paragraph{Development to Old
English}\label{development-to-old-english-4}

Once the oblique \emph{*butt-} stem has been generalized, the selected
input \emph{*búttmaz} develops regularly to \emph{botm}. The analogical
step is therefore early: it belongs to pre-Old-English stem formation
rather than to a later choice among Old English paradigm cells.

\subsubsection{brand --- OE brandes}\label{brand-oe-brandes}

Derivation: citation reconstruction \emph{*brándaz}; selected input
\emph{*brándas} \textgreater{} \emph{brandes} (early analogy).

\paragraph{Derivation trace}\label{derivation-trace-5}

Proto input: \emph{*brándas}

\begingroup
\setlength{\fboxsep}{6pt}

\noindent\fbox{%
\begin{minipage}{0.97\linewidth}
\small
\begin{tabularx}{\linewidth}{@{}>{\raggedright\arraybackslash}p{0.300\linewidth}>{\centering\arraybackslash}X>{\raggedright\arraybackslash}p{0.540\linewidth}@{}}
\begin{minipage}[t]{\linewidth}
\centering\textbf{Earlier Germanic changes}\par
\vspace{0.35em}
\raggedright
\centering\textbf{West Germanic}\par
\raggedright
\vspace{0.2em}
\raggedright [no change]\par
\vspace{0.6em}
\centering\textbf{Northwest Germanic}\par
\raggedright
\vspace{0.2em}
\raggedright [no change]\par
\end{minipage}
&
&
\begin{minipage}[t]{\linewidth}
\centering\textbf{Old English changes}\par
\vspace{0.35em}
\raggedright
\begin{tabular}{@{}>{\raggedright\arraybackslash}p{0.70\linewidth}@{\hspace{0.55em}}>{\raggedright\arraybackslash}p{0.20\linewidth}@{\hspace{0.25em}}}
Anglo Frisian Brightening & \emph{*brándæs} \\
OE Unstressed AE Merger & \emph{*brándes} \\
\end{tabular}
\end{minipage}
\\
\end{tabularx}
\end{minipage}%
} \endgroup

Outcome: \emph{brandes}

\paragraph{Reconstruction and comparative
evidence}\label{reconstruction-and-comparative-evidence-5}

The inherited noun is the masculine a-stem \emph{*brándaz}, continued by
Old English \emph{brand} and its continental cognates
(\citeproc{ref-Orel2003}{Orel 2003}; \citeproc{ref-Kroonen2013}{Kroonen
2013}). The selected input \emph{*brándas} is not a different lexeme but
the genitive singular of that same a-stem noun.

What matters here is therefore not a stem-class disagreement but the
difference between the citation form and a specific inherited
inflectional cell. The selected input preserves the same root and
declension as the headword while making the oblique ending explicit.

\paragraph{Old English evidence}\label{old-english-evidence-5}

Old English dictionaries lemmatize the noun as \emph{brand}
(\citeproc{ref-ClarkHall1960}{Clark Hall 1960};
\citeproc{ref-BosworthToller1898}{Bosworth and Toller 1898, 116}).
Bosworth-Toller also records inflectional forms such as \emph{brandas},
\emph{branda}, and \emph{brandum} under the same entry
(\citeproc{ref-BosworthToller1898}{Bosworth and Toller 1898, 116}).

The specific comparison form in this entry, \emph{brandes}, is the
expected genitive singular of that a-stem noun. It is therefore an
inferred Old English paradigm form rather than the ordinary dictionary
headword.

\paragraph{Development to Old
English}\label{development-to-old-english-5}

From \emph{*brándas}, the regular Old English development passes through
the usual unstressed-vowel weakening of the inflectional ending,
yielding \emph{brandes}. Nothing in the stem itself requires a special
repair. The root consonants and the stressed vowel are the same as in
the citation lemma \emph{brand}.

The analytical weight of the entry lies in the ending. By choosing the
oblique singular rather than the nominative citation form, the entry
presents the same lexeme in a different inherited cell.

\paragraph{Form comparison}\label{form-comparison}

The comparison below is manual. It separates the citation lemma from the
selected oblique singular.

{\def\LTcaptype{none} % do not increment counter
\begin{longtable}[]{@{}
  >{\raggedright\arraybackslash}p{(\linewidth - 8\tabcolsep) * \real{0.2000}}
  >{\raggedright\arraybackslash}p{(\linewidth - 8\tabcolsep) * \real{0.2000}}
  >{\raggedright\arraybackslash}p{(\linewidth - 8\tabcolsep) * \real{0.2000}}
  >{\raggedright\arraybackslash}p{(\linewidth - 8\tabcolsep) * \real{0.2000}}
  >{\raggedright\arraybackslash}p{(\linewidth - 8\tabcolsep) * \real{0.2000}}@{}}
\toprule\noalign{}
\begin{minipage}[b]{\linewidth}\raggedright
Form / interpretation
\end{minipage} & \begin{minipage}[b]{\linewidth}\raggedright
Candidate input
\end{minipage} & \begin{minipage}[b]{\linewidth}\raggedright
Expected or documented OE outcome
\end{minipage} & \begin{minipage}[b]{\linewidth}\raggedright
OE comparison form
\end{minipage} & \begin{minipage}[b]{\linewidth}\raggedright
Result
\end{minipage} \\
\midrule\noalign{}
\endhead
\bottomrule\noalign{}
\endlastfoot
citation nominative singular & *brándaz & expected OE lemma \emph{brand}
& brand & regular headword-level outcome \\
selected genitive singular & *brándas & compact-trace output:
\emph{brandes} & brandes & exact match for the chosen oblique cell \\
\end{longtable}
}

The noun itself is straightforwardly inherited. The main point of the
entry is that \emph{brandes} belongs to the same regular a-stem paradigm
as \emph{brand}, even though the citation lemma remains the nominative
singular.

\subsubsection{breast --- OE brēost}\label{breast-oe-brux113ost}

Derivation: citation reconstruction \emph{*brústz}; selected input
\emph{*bréustą} \textgreater{} \emph{brēost} (early analogy).

\paragraph{Derivation trace}\label{derivation-trace-6}

Proto input: \emph{*bréustą}

\begingroup
\setlength{\fboxsep}{6pt}

\noindent\fbox{%
\begin{minipage}{0.97\linewidth}
\small
\begin{tabularx}{\linewidth}{@{}>{\raggedright\arraybackslash}p{0.300\linewidth}>{\centering\arraybackslash}X>{\raggedright\arraybackslash}p{0.540\linewidth}@{}}
\begin{minipage}[t]{\linewidth}
\centering\textbf{Earlier Germanic changes}\par
\vspace{0.35em}
\raggedright
\centering\textbf{West Germanic}\par
\raggedright
\vspace{0.2em}
\raggedright [no change]\par
\vspace{0.6em}
\centering\textbf{Northwest Germanic}\par
\raggedright
\vspace{0.2em}
\raggedright [no change]\par
\end{minipage}
&
&
\begin{minipage}[t]{\linewidth}
\centering\textbf{Old English changes}\par
\vspace{0.35em}
\raggedright
\begin{tabular}{@{}>{\raggedright\arraybackslash}p{0.70\linewidth}@{\hspace{0.55em}}>{\raggedright\arraybackslash}p{0.20\linewidth}@{\hspace{0.25em}}}
OE Diphthong Leveling & \emph{*brēostą} \\
OE Heavy Syllable Nasal Apocope & \emph{*brēost} \\
\end{tabular}
\end{minipage}
\\
\end{tabularx}
\end{minipage}%
} \endgroup

Outcome: \emph{brēost}

\paragraph{Reconstruction and comparative
evidence}\label{reconstruction-and-comparative-evidence-6}

The word family shows two related but distinct Proto-Germanic
formations. The root noun \emph{*brust-} lies behind forms such as
Gothic \emph{brusts}, whereas Old English \emph{brēost} belongs to a
thematic formation \emph{*breusta-}, alongside Old Norse \emph{brjóst}
and Old Saxon \emph{briost} (\citeproc{ref-Kroonen2013}{Kroonen 2013};
\citeproc{ref-Orel2003}{Orel 2003, 95};
\citeproc{ref-RingeTaylor2014}{Ringe and Taylor 2014, 43}).

The selected input \emph{*bréustą} therefore differs from the citation
label \emph{*brústz} because Old English reflects the thematic branch
rather than the root noun. The morphological choice comes before the Old
English sound changes themselves.

\paragraph{Old English evidence}\label{old-english-evidence-6}

Old English dictionaries record the noun as \emph{brēost} /
\emph{breóst} (\citeproc{ref-BosworthToller1898}{Bosworth and Toller
1898}; \citeproc{ref-ClarkHall1960}{Clark Hall 1960, 65}). The form is
an established Old English lexeme, not a reconstructed target assembled
from comparative evidence alone.

What requires explanation is not the Old English attestation but the
relation between that attested noun and the broader Germanic word
family. The relevant comparison form is therefore the thematic Old
English noun \emph{brēost}.

\paragraph{Development to Old
English}\label{development-to-old-english-6}

From \emph{*bréustą}, the regular Old English development gives
\emph{brēost}, with the expected \emph{eu} \textgreater{} \emph{ēo}
vowel history (\citeproc{ref-Campbell1959}{Campbell 1959}). No special
repair is needed once the correct thematic formation is chosen.

The earlier mismatch arose only if the word was forced into the
root-noun line. The Old English noun itself continues the thematic
branch cleanly and directly.

\paragraph{Formation comparison}\label{formation-comparison}

The comparison below is manual. It separates the broader root-noun
family label from the thematic formation actually continued in Old
English.

{\def\LTcaptype{none} % do not increment counter
\begin{longtable}[]{@{}
  >{\raggedright\arraybackslash}p{(\linewidth - 8\tabcolsep) * \real{0.2000}}
  >{\raggedright\arraybackslash}p{(\linewidth - 8\tabcolsep) * \real{0.2000}}
  >{\raggedright\arraybackslash}p{(\linewidth - 8\tabcolsep) * \real{0.2000}}
  >{\raggedright\arraybackslash}p{(\linewidth - 8\tabcolsep) * \real{0.2000}}
  >{\raggedright\arraybackslash}p{(\linewidth - 8\tabcolsep) * \real{0.2000}}@{}}
\toprule\noalign{}
\begin{minipage}[b]{\linewidth}\raggedright
Formation
\end{minipage} & \begin{minipage}[b]{\linewidth}\raggedright
Candidate input
\end{minipage} & \begin{minipage}[b]{\linewidth}\raggedright
Expected or documented OE outcome
\end{minipage} & \begin{minipage}[b]{\linewidth}\raggedright
OE comparison form
\end{minipage} & \begin{minipage}[b]{\linewidth}\raggedright
Result
\end{minipage} \\
\midrule\noalign{}
\endhead
\bottomrule\noalign{}
\endlastfoot
broader root-noun family & *brústz & root-noun type outcomes outside OE
& non-OE comparanda & useful family label, but not the direct source of
\emph{brēost} \\
selected thematic formation & *bréustą & compact-trace output:
\emph{brēost} & brēost & exact match between formation and attested OE
noun \\
\end{longtable}
}

The relevant point is the formation split. \emph{brēost} is the regular
Old English outcome of the thematic \emph{*breusta-} branch, not of the
root noun \emph{*brust-}.

\subsubsection{craft --- OE cræft}\label{craft-oe-cruxe6ft}

Derivation: citation reconstruction \emph{*kráftiz}; selected input
\emph{*kráftaz} \textgreater{} \emph{cræft} (early analogy).

\paragraph{Derivation trace}\label{derivation-trace-7}

Proto input: \emph{*kráftaz}

\begingroup
\setlength{\fboxsep}{6pt}

\noindent\fbox{%
\begin{minipage}{0.97\linewidth}
\small
\begin{tabularx}{\linewidth}{@{}>{\raggedright\arraybackslash}p{0.300\linewidth}>{\centering\arraybackslash}X>{\raggedright\arraybackslash}p{0.540\linewidth}@{}}
\begin{minipage}[t]{\linewidth}
\centering\textbf{Earlier Germanic changes}\par
\vspace{0.35em}
\raggedright
\centering\textbf{West Germanic}\par
\raggedright
\vspace{0.2em}
\raggedright [no change]\par
\vspace{0.6em}
\centering\textbf{Northwest Germanic}\par
\raggedright
\vspace{0.2em}
\begin{tabular}{@{}>{\raggedright\arraybackslash}p{0.74\linewidth}@{\hspace{0.55em}}>{\raggedright\arraybackslash}p{0.16\linewidth}@{\hspace{0.25em}}}
\mbox{PGmc Final Z Deletion} & \emph{*kráfta} \\
\end{tabular}
\end{minipage}
&
&
\begin{minipage}[t]{\linewidth}
\centering\textbf{Old English changes}\par
\vspace{0.35em}
\raggedright
\begin{tabular}{@{}>{\raggedright\arraybackslash}p{0.70\linewidth}@{\hspace{0.55em}}>{\raggedright\arraybackslash}p{0.16\linewidth}@{\hspace{0.25em}}}
PWGmc Final Bare A Loss & \emph{*kráft} \\
Anglo Frisian Brightening & \emph{*kræft} \\
\end{tabular}
\end{minipage}
\\
\end{tabularx}
\end{minipage}%
} \endgroup

Outcome: \emph{cræft}

\paragraph{Reconstruction and comparative
evidence}\label{reconstruction-and-comparative-evidence-7}

Comparative sources disagree about the older stem class. Kroonen gives a
u-stem \emph{*kraftu-}, while Orel prints \emph{*kraftiz}
\textasciitilde{} \emph{*kraftuz} (\citeproc{ref-Kroonen2013}{Kroonen
2013, 340}; \citeproc{ref-Orel2003}{Orel 2003, 259}). The comparative
label \emph{*kráftiz} remains in view as a lexeme-level shorthand, while
\emph{*kráftaz} is the pre-Old-English form used for the Old English
derivation.

\paragraph{Old English evidence}\label{old-english-evidence-7}

The Old English noun itself is secure. Clark Hall and Bosworth-Toller
both give \emph{cræft} as the headword
(\citeproc{ref-ClarkHall1960}{Clark Hall 1960, 19};
\citeproc{ref-BosworthToller1898}{Bosworth and Toller 1898, 145}).

\paragraph{Development to Old
English}\label{development-to-old-english-7}

The comparison is between possible pre-Old-English inputs. The i-stem
comparator \emph{*kráftiz} gives \emph{creft}, while the u-stem
comparator \emph{*kráftuz} gives \emph{craft}. The a-stem-shaped input
\emph{*kráftaz} yields \emph{cræft} and is therefore the form used for
the Old English derivation. This does not require treating the
comparative dictionaries as identical; it shows the narrower point that
the Old English derivation needs a pre-Old-English form without the
i-umlaut trigger of \emph{*-iz} and without the back-vowel outcome
associated with the u-stem comparator.

\paragraph{Form comparison}\label{form-comparison-1}

{\def\LTcaptype{none} % do not increment counter
\begin{longtable}[]{@{}lll@{}}
\toprule\noalign{}
Candidate input & OE output & Result \\
\midrule\noalign{}
\endhead
\bottomrule\noalign{}
\endlastfoot
*kráftiz & \emph{creft} & non-match; i-stem comparator \\
*kráftuz & \emph{craft} & non-match; u-stem comparator \\
*kráftaz & \emph{cræft} & exact match; selected pre-OE input \\
\end{longtable}
}

\subsubsection{dill --- OE dile}\label{dill-oe-dile}

Derivation: citation reconstruction \emph{*déljaz}; selected input
\emph{*déliz} \textgreater{} \emph{dile} (early analogy).

\paragraph{Derivation trace}\label{derivation-trace-8}

Proto input: \emph{*déliz}

\begingroup
\setlength{\fboxsep}{6pt}

\noindent\fbox{%
\begin{minipage}{0.97\linewidth}
\small
\begin{tabularx}{\linewidth}{@{}>{\raggedright\arraybackslash}p{0.300\linewidth}>{\centering\arraybackslash}X>{\raggedright\arraybackslash}p{0.540\linewidth}@{}}
\begin{minipage}[t]{\linewidth}
\centering\textbf{Earlier Germanic changes}\par
\vspace{0.35em}
\raggedright
\centering\textbf{West Germanic}\par
\raggedright
\vspace{0.2em}
\raggedright [no change]\par
\vspace{0.6em}
\centering\textbf{Northwest Germanic}\par
\raggedright
\vspace{0.2em}
\begin{tabular}{@{}>{\raggedright\arraybackslash}p{0.74\linewidth}@{\hspace{0.55em}}>{\raggedright\arraybackslash}p{0.16\linewidth}@{\hspace{0.25em}}}
\mbox{PGmc Final Z Deletion} & \emph{*déli} \\
\end{tabular}
\end{minipage}
&
&
\begin{minipage}[t]{\linewidth}
\centering\textbf{Old English changes}\par
\vspace{0.35em}
\raggedright
\begin{tabular}{@{}>{\raggedright\arraybackslash}p{0.74\linewidth}@{\hspace{0.55em}}>{\raggedright\arraybackslash}p{0.16\linewidth}@{\hspace{0.25em}}}
OE I Umlaut & \emph{*dili} \\
OE Med Unstressed I Lowering1 & \emph{*dile} \\
\end{tabular}
\end{minipage}
\\
\end{tabularx}
\end{minipage}%
} \endgroup

Outcome: \emph{dile}

\paragraph{Reconstruction and comparative
evidence}\label{reconstruction-and-comparative-evidence-8}

Kroonen treats the word as preserving evidence for both an i-stem and a
ja-stem formation, with Old English \emph{dile} on one side and
continental forms such as Old Saxon \emph{dilli} and Old High German
\emph{tilli} on the other (\citeproc{ref-Kroonen2013}{Kroonen 2013}).
The selected input \emph{*déliz} therefore represents the i-stem side of
the paradigm, whereas the citation label \emph{*déljaz} is a broader
comparative headword.

That stem-class distinction matters for the Old English consonant shape.
A ja-stem with \emph{*-lj-} would be expected to produce gemination, but
the Old English noun shows a single \emph{l}. Fulk's discussion of
ja-stems transferred to the i-stems provides the relevant morphological
background for the OE side (\citeproc{ref-Fulk2018}{Fulk 2018}).

\paragraph{Old English evidence}\label{old-english-evidence-8}

Old English dictionaries record the plant name as \emph{dile}, alongside
the variant \emph{dili} (\citeproc{ref-BosworthToller1898}{Bosworth and
Toller 1898, 164}; \citeproc{ref-ClarkHall1960}{Clark Hall 1960}). The
form discussed here is therefore an attested Old English noun with
single \emph{l}.

The Old English evidence is the relevant point. Whatever broader
comparative headword is chosen for the family, the inherited form
reflected in OE is the i-stem type \emph{dile}, not a geminated
\emph{dill} outcome.

\paragraph{Development to Old
English}\label{development-to-old-english-8}

From \emph{*déliz}, regular loss of final \emph{z} and the later
lowering of unstressed \emph{i} yield \emph{dile}. The stem itself
remains ungeminated throughout that path.

The important contrast is negative rather than phonological. If the word
were forced through a ja-stem \emph{*-lj-} pathway, the expected result
would show \emph{ll}. The attested Old English noun instead matches the
i-stem development.

\paragraph{Formation comparison}\label{formation-comparison-1}

The comparison below is manual. It separates the broader comparative
headword from the stem class actually reflected in Old English.

{\def\LTcaptype{none} % do not increment counter
\begin{longtable}[]{@{}
  >{\raggedright\arraybackslash}p{(\linewidth - 8\tabcolsep) * \real{0.2000}}
  >{\raggedright\arraybackslash}p{(\linewidth - 8\tabcolsep) * \real{0.2000}}
  >{\raggedright\arraybackslash}p{(\linewidth - 8\tabcolsep) * \real{0.2000}}
  >{\raggedright\arraybackslash}p{(\linewidth - 8\tabcolsep) * \real{0.2000}}
  >{\raggedright\arraybackslash}p{(\linewidth - 8\tabcolsep) * \real{0.2000}}@{}}
\toprule\noalign{}
\begin{minipage}[b]{\linewidth}\raggedright
Formation
\end{minipage} & \begin{minipage}[b]{\linewidth}\raggedright
Candidate input
\end{minipage} & \begin{minipage}[b]{\linewidth}\raggedright
Expected or documented OE outcome
\end{minipage} & \begin{minipage}[b]{\linewidth}\raggedright
OE comparison form
\end{minipage} & \begin{minipage}[b]{\linewidth}\raggedright
Result
\end{minipage} \\
\midrule\noalign{}
\endhead
\bottomrule\noalign{}
\endlastfoot
comparative ja-stem label & *déljaz & ja-stem type outcome with
gemination & dill-type comparison & useful comparative label, but not
the OE form \\
selected i-stem formation & *déliz & compact-trace output: \emph{dile} &
dile & exact match between formation and attested OE noun \\
\end{longtable}
}

The single \emph{l} is the decisive diagnostic. It identifies
\emph{dile} with the i-stem formation rather than with the continental
ja-stem branch.

\subsubsection{fast --- OE festan}\label{fast-oe-festan}

Derivation: citation reconstruction \emph{*fastēną}; selected input
\emph{*fástijaną} \textgreater{} \emph{festan} (early analogy).

\paragraph{Derivation trace}\label{derivation-trace-9}

Proto input: \emph{*fástijaną}

\begingroup
\setlength{\fboxsep}{6pt}

\noindent\fbox{%
\begin{minipage}{0.97\linewidth}
\footnotesize
\begin{tabularx}{\linewidth}{@{}>{\raggedright\arraybackslash}p{0.280\linewidth}>{\centering\arraybackslash}X>{\raggedright\arraybackslash}p{0.560\linewidth}@{}}
\begin{minipage}[t]{\linewidth}
\centering\textbf{Earlier Germanic changes}\par
\vspace{0.35em}
\raggedright
\centering\textbf{West Germanic}\par
\raggedright
\vspace{0.2em}
\raggedright [no change]\par
\vspace{0.6em}
\centering\textbf{Northwest Germanic}\par
\raggedright
\vspace{0.2em}
\raggedright [no change]\par
\end{minipage}
&
&
\begin{minipage}[t]{\linewidth}
\centering\textbf{Old English changes}\par
\vspace{0.35em}
\raggedright
\begin{tabular}{@{}>{\raggedright\arraybackslash}p{0.64\linewidth}@{\hspace{0.55em}}>{\raggedright\arraybackslash}p{0.26\linewidth}@{\hspace{0.25em}}}
Anglo Frisian Brightening & \emph{*fæstijaną} \\
OE Heavy Syllable Nasal Apocope & \emph{*fæstijan} \\
OE Secondary Nasalization & \emph{*fæstijąn} \\
Sievers Law Syncope & \emph{*fæstjąn} \\
OE I Umlaut & \emph{*festjąn} \\
OE Weak Tail Reduction & \emph{*festjan} \\
OE J Loss After Heavy & \emph{*festan} \\
\end{tabular}
\end{minipage}
\\
\end{tabularx}
\end{minipage}%
} \endgroup

Outcome: \emph{festan}

\paragraph{Reconstruction and comparative
evidence}\label{reconstruction-and-comparative-evidence-9}

Kroonen places the verb under a comparative headword \emph{*fastēną},
the wider Germanic family behind meanings such as `make firm' and, in
Old English, `fast' (\citeproc{ref-Kroonen2013}{Kroonen 2013}). Ringe
and Taylor, however, distinguish the Old English verb more closely: they
treat OE `to fast' as originally a class-I weak verb that later acquired
the stative meaning through lexical association with that wider family
(\citeproc{ref-RingeTaylor2014}{Ringe and Taylor 2014}).

The selected input \emph{*fástijaną} therefore represents the inherited
class-I formation reflected in Old English, whereas the citation label
\emph{*fastēną} belongs to the broader comparative presentation of the
lexeme.

\paragraph{Old English evidence}\label{old-english-evidence-9}

Old English dictionaries record forms such as \emph{festan}, alongside
related \emph{fæstan} / \emph{fǣstan} spellings and meanings
(\citeproc{ref-BosworthToller1898}{Bosworth and Toller 1898, 213};
\citeproc{ref-ClarkHall1960}{Clark Hall 1960}). The form selected here
is \emph{festan}, which fits the regular class-I phonological
development.

The \emph{æ}-forms remain relevant, but they do not control the entry.
As Ringe and Taylor argue, their vowel belongs to later analogical
leveling under the adjective \emph{fæst}, whereas \emph{festan} reflects
the regular inherited class-I verb (\citeproc{ref-RingeTaylor2014}{Ringe
and Taylor 2014}).

\paragraph{Development to Old
English}\label{development-to-old-english-9}

From \emph{*fástijaną}, Anglo-Frisian brightening and subsequent
i-umlaut produce the fronted vowel seen in \emph{festan}. The later
weak-tail reductions and loss of \emph{j} after a heavy syllable
complete the regular Old English outcome.

What makes the entry non-regular is not the phonology of \emph{festan}
itself, but the choice of formation. Old English continues the class-I
verb, even though the comparative headword is often given under the
parallel \emph{*fastēn-} family.

\paragraph{Class comparison}\label{class-comparison}

The comparison below is manual. It distinguishes the comparative
class-III headword from the class-I formation actually reflected in Old
English.

{\def\LTcaptype{none} % do not increment counter
\begin{longtable}[]{@{}
  >{\raggedright\arraybackslash}p{(\linewidth - 8\tabcolsep) * \real{0.2000}}
  >{\raggedright\arraybackslash}p{(\linewidth - 8\tabcolsep) * \real{0.2000}}
  >{\raggedright\arraybackslash}p{(\linewidth - 8\tabcolsep) * \real{0.2000}}
  >{\raggedright\arraybackslash}p{(\linewidth - 8\tabcolsep) * \real{0.2000}}
  >{\raggedright\arraybackslash}p{(\linewidth - 8\tabcolsep) * \real{0.2000}}@{}}
\toprule\noalign{}
\begin{minipage}[b]{\linewidth}\raggedright
Formation / class
\end{minipage} & \begin{minipage}[b]{\linewidth}\raggedright
Candidate input
\end{minipage} & \begin{minipage}[b]{\linewidth}\raggedright
Expected or documented OE outcome
\end{minipage} & \begin{minipage}[b]{\linewidth}\raggedright
OE comparison form
\end{minipage} & \begin{minipage}[b]{\linewidth}\raggedright
Result
\end{minipage} \\
\midrule\noalign{}
\endhead
\bottomrule\noalign{}
\endlastfoot
comparative class-III headword & *fastēną & class-III type outcome, not
\emph{festan} & wider family context & useful family label, but not the
direct source of the target \\
selected class-I weak verb & *fástijaną & compact-trace output:
\emph{festan} & festan & exact match between formation and attested OE
verb \\
later analogical reshaping & adjective-driven \emph{fæst} influence &
\emph{fæstan} / \emph{fǣstan} type spellings & fæstan-type evidence &
genuine later OE reshaping, but secondary to the selected target \\
\end{longtable}
}

The relevant point is the class split. \emph{festan} is the regular Old
English outcome of the class-I formation, while the better-known
\emph{æ}-forms belong to a later analogical layer.

\subsubsection{flask --- OE flasce}\label{flask-oe-flasce}

Derivation: citation reconstruction \emph{*flaskō}; selected input
\emph{*fláskōn} \textgreater{} \emph{flasce} (early analogy).

\paragraph{Derivation trace}\label{derivation-trace-10}

Proto input: \emph{*fláskōn}

\begingroup
\setlength{\fboxsep}{6pt}

\noindent\fbox{%
\begin{minipage}{0.97\linewidth}
\small
\begin{tabularx}{\linewidth}{@{}>{\raggedright\arraybackslash}p{0.300\linewidth}>{\centering\arraybackslash}X>{\raggedright\arraybackslash}p{0.540\linewidth}@{}}
\begin{minipage}[t]{\linewidth}
\centering\textbf{Earlier Germanic changes}\par
\vspace{0.35em}
\raggedright
\centering\textbf{West Germanic}\par
\raggedright
\vspace{0.2em}
\raggedright [no change]\par
\vspace{0.6em}
\centering\textbf{Northwest Germanic}\par
\raggedright
\vspace{0.2em}
\begin{tabular}{@{}>{\raggedright\arraybackslash}p{0.64\linewidth}@{\hspace{0.55em}}>{\raggedright\arraybackslash}p{0.16\linewidth}@{\hspace{0.25em}}}
NWGmc N Stem N Loss & \emph{*fláskǭ} \\
\end{tabular}
\end{minipage}
&
&
\begin{minipage}[t]{\linewidth}
\centering\textbf{Old English changes}\par
\vspace{0.35em}
\raggedright
\begin{tabular}{@{}>{\raggedright\arraybackslash}p{0.74\linewidth}@{\hspace{0.55em}}>{\raggedright\arraybackslash}p{0.16\linewidth}@{\hspace{0.25em}}}
Anglo Frisian Brightening & \emph{*flæskǭ} \\
OE A Restoration & \emph{*flaskǭ} \\
OE Unstressed Long Vowel Shortening & \emph{*flaskæ} \\
OE Unstressed AE Merger & \emph{*flaske} \\
\end{tabular}
\end{minipage}
\\
\end{tabularx}
\end{minipage}%
} \endgroup

Outcome: \emph{flasce}

\paragraph{Reconstruction and comparative
evidence}\label{reconstruction-and-comparative-evidence-10}

The wider Germanic family is often cited under a form such as
\emph{*flaskō}, but the evidence relevant for Old English points instead
to a weak feminine formation \emph{*fláskōn} / \emph{*flaskǭ}
(\citeproc{ref-Orel2003}{Orel 2003}; \citeproc{ref-Kroonen2013}{Kroonen
2013}). That distinction is crucial for the suffixal history of the
noun.

The selected input therefore differs from the citation label in stem
class. Old English \emph{flasce} belongs with the weak feminine line,
and the plural or oblique forms \emph{flascan} support that analysis
(\citeproc{ref-RingeTaylor2014}{Ringe and Taylor 2014, 192}).

\paragraph{Old English evidence}\label{old-english-evidence-10}

Old English dictionaries record the noun as \emph{flasce}, with
inflectional support from forms such as \emph{flascan}; a later West
Saxon \emph{flaxe} is also noted as a secondary variant
(\citeproc{ref-BosworthToller1898}{Bosworth and Toller 1898, 235};
\citeproc{ref-ClarkHall1960}{Clark Hall 1960, 121}).

The relevant comparison form is therefore the weak feminine noun
\emph{flasce}. The plural and oblique evidence matters because it helps
explain why the vowel and ending are preserved as they are in the
singular.

\paragraph{Development to Old
English}\label{development-to-old-english-10}

From \emph{*fláskōn}, the weak feminine passes through the expected loss
of \emph{n} and the later Old English development of the unstressed
ending, reaching \emph{flasce}. Campbell cites restored \emph{a} in
exactly this environment, including \emph{flasce} after inflected
\emph{flascan} (\citeproc{ref-Campbell1959}{Campbell 1959, sect. 158}).
Once the weak feminine formation is chosen, the noun follows a regular
path to its Old English shape.

The decisive issue is morphological rather than phonological. A simple
strong feminine citation form does not capture the OE weak noun as
cleanly as the selected \emph{*fláskōn} does.

\paragraph{Formation comparison}\label{formation-comparison-2}

The comparison below is manual. It separates the broader comparative
headword from the weak feminine formation actually reflected in Old
English.

{\def\LTcaptype{none} % do not increment counter
\begin{longtable}[]{@{}
  >{\raggedright\arraybackslash}p{(\linewidth - 8\tabcolsep) * \real{0.2000}}
  >{\raggedright\arraybackslash}p{(\linewidth - 8\tabcolsep) * \real{0.2000}}
  >{\raggedright\arraybackslash}p{(\linewidth - 8\tabcolsep) * \real{0.2000}}
  >{\raggedright\arraybackslash}p{(\linewidth - 8\tabcolsep) * \real{0.2000}}
  >{\raggedright\arraybackslash}p{(\linewidth - 8\tabcolsep) * \real{0.2000}}@{}}
\toprule\noalign{}
\begin{minipage}[b]{\linewidth}\raggedright
Formation
\end{minipage} & \begin{minipage}[b]{\linewidth}\raggedright
Candidate input
\end{minipage} & \begin{minipage}[b]{\linewidth}\raggedright
Expected or documented OE outcome
\end{minipage} & \begin{minipage}[b]{\linewidth}\raggedright
OE comparison form
\end{minipage} & \begin{minipage}[b]{\linewidth}\raggedright
Result
\end{minipage} \\
\midrule\noalign{}
\endhead
\bottomrule\noalign{}
\endlastfoot
broader comparative headword & *flaskō & broader family label & wider
family context & useful lexeme label, but not the cleanest OE-facing
derivation \\
selected weak feminine formation & *fláskōn & compact-trace output:
\emph{flasce} & flasce & exact match between formation and attested OE
noun \\
\end{longtable}
}

The weak feminine suffix is the relevant point. It aligns the inherited
form with attested \emph{flasce} and its supporting paradigm forms.

\subsubsection{follow --- OE fylġan}\label{follow-oe-fylux121an}

Derivation: citation reconstruction \emph{*fulgēną}; selected input
\emph{*fúlgijaną} \textgreater{} \emph{fylġan} (early analogy).

\paragraph{Derivation trace}\label{derivation-trace-11}

Proto input: \emph{*fúlgijaną}

\begingroup
\setlength{\fboxsep}{6pt}

\noindent\fbox{%
\begin{minipage}{0.97\linewidth}
\footnotesize
\begin{tabularx}{\linewidth}{@{}>{\raggedright\arraybackslash}p{0.280\linewidth}>{\centering\arraybackslash}X>{\raggedright\arraybackslash}p{0.560\linewidth}@{}}
\begin{minipage}[t]{\linewidth}
\centering\textbf{Earlier Germanic changes}\par
\vspace{0.35em}
\raggedright
\centering\textbf{West Germanic}\par
\raggedright
\vspace{0.2em}
\raggedright [no change]\par
\vspace{0.6em}
\centering\textbf{Northwest Germanic}\par
\raggedright
\vspace{0.2em}
\raggedright [no change]\par
\end{minipage}
&
&
\begin{minipage}[t]{\linewidth}
\centering\textbf{Old English changes}\par
\vspace{0.35em}
\raggedright
\begin{tabular}{@{}>{\raggedright\arraybackslash}p{0.66\linewidth}@{\hspace{0.55em}}>{\raggedright\arraybackslash}p{0.24\linewidth}@{\hspace{0.25em}}}
OE Heavy Syllable Nasal Apocope & \emph{*fúlgijan} \\
OE Secondary Nasalization & \emph{*fúlgijąn} \\
Sievers Law Syncope & \emph{*fúlgjąn} \\
OE Velar Palatalization & \emph{*fúlʤjąn} \\
OE I Umlaut & \emph{*fylʤjąn} \\
OE Weak Tail Reduction & \emph{*fylʤjan} \\
OE J Loss After Heavy & \emph{*fylʤan} \\
\end{tabular}
\end{minipage}
\\
\end{tabularx}
\end{minipage}%
} \endgroup

Outcome: \emph{fylġan}

\paragraph{Reconstruction and comparative
evidence}\label{reconstruction-and-comparative-evidence-11}

Kroonen keeps the verb under \emph{*fulgen-} and gives Old English
\emph{fylgan, folgian}, adding that Old Norse \emph{fylgja} and Old
English \emph{fylg(e)an} continue a formation \emph{*fulgjan-}
(\citeproc{ref-Kroonen2013}{Kroonen 2013}). The comparative headword and
the class-I formation are therefore related but not identical.

Ringe and Taylor make the split explicit as PNWGmc \emph{*fulgija-}
\textasciitilde{} \emph{*fulgai-} \textgreater{} OE \emph{fylgan}
\textasciitilde{} \emph{folgian} and describe it as a dual formation
that probably reflects an older alternation between j-present and
e-stative (\citeproc{ref-RingeTaylor2014}{Ringe and Taylor 2014,
293--94}). This is a stem-class choice, not a spelling choice. The
selected input \emph{*fúlgijaną} belongs to the class-I \emph{*fulgija-}
/ \emph{*fulgjan-} branch; the citation form \emph{*fulgēną} belongs to
the parallel class-II history behind \emph{folgian}.

\paragraph{Old English evidence}\label{old-english-evidence-11}

The Old English evidence preserves both formations. Clark Hall gives
\emph{folgian} and cross-refers to \emph{fylgan}, while also listing
\emph{fylgan} with variant spellings \emph{fylgian} and \emph{fyligan}
(\citeproc{ref-ClarkHall1960}{Clark Hall 1960}). Bosworth-Toller
likewise has separate entries for \emph{folgian} and \emph{fylgean}
(\citeproc{ref-BosworthToller1898}{Bosworth and Toller 1898}).

Bright notes traces of the older conjugation in \emph{fylg(e)an}
(\citeproc{ref-BrightCassidyRingler1971}{Bright 1971, 77}) and lists
\emph{folgian (fylgean)} in the glossary
(\citeproc{ref-BrightCassidyRingler1971}{Bright 1971, 364}). The
relevant comparison form in this entry is therefore the class-I verb
\emph{fylgan} / \emph{fylgean}, here normalized as \emph{fylġan}. The
spelling with ġ represents the palatalized velar before a front-vocalic
environment.

\paragraph{Development to Old
English}\label{development-to-old-english-11}

\emph{*fúlgijaną} is a class-I weak-verb formation. In the class-I
branch the \emph{*j} blocks NWGmc lowering of \emph{u} to \emph{o},
since Ringe and Taylor formulate that lowering for environments in which
no \emph{*j} intervened (\citeproc{ref-RingeTaylor2014}{Ringe and Taylor
2014, 96}). The same front-vocalic environment then triggers i-umlaut,
so \emph{u} becomes \emph{y} (\citeproc{ref-RingeTaylor2014}{Ringe and
Taylor 2014, sect. 6.6.2}).

The subsequent Old English developments are palatalization of the velar,
weak-tail reduction, and loss of \emph{j} after a heavy syllable,
yielding \emph{fylġan}. This is the regular outcome of the class-I
formation. The class-II form \emph{folgian} belongs to the parallel
\emph{*-ē-} / \emph{*-ai-} branch and is not the form modeled here.

\paragraph{Class comparison}\label{class-comparison-1}

A class comparison identifies which inherited formation corresponds to
the established Old English form under discussion. The comparison below
is manual; no full automatic class probe is presented here.

{\def\LTcaptype{none} % do not increment counter
\begin{longtable}[]{@{}
  >{\raggedright\arraybackslash}p{(\linewidth - 8\tabcolsep) * \real{0.2000}}
  >{\raggedright\arraybackslash}p{(\linewidth - 8\tabcolsep) * \real{0.2000}}
  >{\raggedright\arraybackslash}p{(\linewidth - 8\tabcolsep) * \real{0.2000}}
  >{\raggedright\arraybackslash}p{(\linewidth - 8\tabcolsep) * \real{0.2000}}
  >{\raggedright\arraybackslash}p{(\linewidth - 8\tabcolsep) * \real{0.2000}}@{}}
\toprule\noalign{}
\begin{minipage}[b]{\linewidth}\raggedright
Formation / class
\end{minipage} & \begin{minipage}[b]{\linewidth}\raggedright
Candidate input
\end{minipage} & \begin{minipage}[b]{\linewidth}\raggedright
Expected or documented OE outcome
\end{minipage} & \begin{minipage}[b]{\linewidth}\raggedright
OE comparison form
\end{minipage} & \begin{minipage}[b]{\linewidth}\raggedright
Result
\end{minipage} \\
\midrule\noalign{}
\endhead
\bottomrule\noalign{}
\endlastfoot
citation class-II formation & *fulgēną & probe output: \emph{folgon} &
folgian & mismatch: the regular output is not the remodeled infinitive
\emph{folgian} \\
parallel class-II branch & PNWGmc *fulgai- & Ringe-Taylor: OE
\emph{folgian} & folgian & documents the separate class-II branch, but
not the target of this entry \\
selected class-I formation & *fúlgijaną & compact-trace output:
\emph{fylġan} & fylġan / fylgan & exact match between input, output, and
class \\
\end{longtable}
}

The relevant point is the class split. \emph{fylġan} is the regular Old
English outcome of the class-I \emph{*fulgija-} / \emph{*fulgjan-}
formation, whereas \emph{folgian} belongs to the parallel class-II
branch.

\subsubsection{gall --- OE ġealla}\label{gall-oe-ux121ealla}

Derivation: citation reconstruction \emph{*gállą}; selected input
\emph{*gállô} \textgreater{} \emph{ġealla} (early analogy).

\paragraph{Derivation trace}\label{derivation-trace-12}

Proto input: \emph{*gállô}

\begingroup
\setlength{\fboxsep}{6pt}

\noindent\fbox{%
\begin{minipage}{0.97\linewidth}
\small
\begin{tabularx}{\linewidth}{@{}>{\raggedright\arraybackslash}p{0.300\linewidth}>{\centering\arraybackslash}X>{\raggedright\arraybackslash}p{0.540\linewidth}@{}}
\begin{minipage}[t]{\linewidth}
\centering\textbf{Earlier Germanic changes}\par
\vspace{0.35em}
\raggedright
\centering\textbf{West Germanic}\par
\raggedright
\vspace{0.2em}
\raggedright [no change]\par
\vspace{0.6em}
\centering\textbf{Northwest Germanic}\par
\raggedright
\vspace{0.2em}
\raggedright [no change]\par
\end{minipage}
&
&
\begin{minipage}[t]{\linewidth}
\centering\textbf{Old English changes}\par
\vspace{0.35em}
\raggedright
\begin{tabular}{@{}>{\raggedright\arraybackslash}p{0.74\linewidth}@{\hspace{0.55em}}>{\raggedright\arraybackslash}p{0.16\linewidth}@{\hspace{0.25em}}}
Anglo Frisian Brightening & \emph{*gællô} \\
OE Breaking & \emph{*geallô} \\
OE Velar Palatalization & \emph{*ʤeallô} \\
OE Unstressed Long Vowel Shortening & \emph{*ʤealla} \\
\end{tabular}
\end{minipage}
\\
\end{tabularx}
\end{minipage}%
} \endgroup

Outcome: \emph{ġealla}

\paragraph{Reconstruction and comparative
evidence}\label{reconstruction-and-comparative-evidence-12}

The wider cognate family can be presented under a form such as
\emph{*gállą}, but the Old English noun itself belongs with a weak noun
\emph{*gallōn-}, cited here as \emph{*gállô}
(\citeproc{ref-Kroonen2013}{Kroonen 2013}). The selected input therefore
differs from the broader comparative headword in stem class.

That stem-class distinction matters directly for the Old English shape.
The weak masculine pathway preserves the ending needed for
\emph{ġealla}, whereas a simple strong-noun headword does not align as
closely with the attested OE noun.

\paragraph{Old English evidence}\label{old-english-evidence-12}

Old English dictionaries record the noun as \emph{gealla}, and Bright
also gives the dative \emph{geallan}, confirming a weak-noun paradigm
(\citeproc{ref-BosworthToller1898}{Bosworth and Toller 1898, 297};
\citeproc{ref-ClarkHall1960}{Clark Hall 1960, 145};
\citeproc{ref-BrightCassidyRingler1971}{Bright 1971, 372}). The form
used here, \emph{ġealla}, is a normalized spelling with macrons omitted
and palatal ġ made explicit.

Campbell also notes dialectal variation, contrasting West Saxon or
Kentish \emph{gealla} with Anglian \emph{galla}
(\citeproc{ref-Campbell1959}{Campbell 1959}). The target of this entry
is the West Saxon type \emph{ġealla}.

\paragraph{Development to Old
English}\label{development-to-old-english-12}

From \emph{*gállô}, the weak noun develops through the expected Old
English history of the suffix and the regular breaking environment
before \emph{ll}, yielding \emph{ġealla}
(\citeproc{ref-Campbell1959}{Campbell 1959}). Once the weak masculine
input is chosen, the noun follows a regular path to its attested Old
English form.

The decisive issue is therefore morphological. Old English reflects the
weak noun, while the broader family label belongs to a different way of
presenting the cognate set.

\paragraph{Stem comparison}\label{stem-comparison}

The comparison below is manual. It separates the broader comparative
headword from the weak noun formation actually reflected in Old English.

{\def\LTcaptype{none} % do not increment counter
\begin{longtable}[]{@{}
  >{\raggedright\arraybackslash}p{(\linewidth - 8\tabcolsep) * \real{0.2000}}
  >{\raggedright\arraybackslash}p{(\linewidth - 8\tabcolsep) * \real{0.2000}}
  >{\raggedright\arraybackslash}p{(\linewidth - 8\tabcolsep) * \real{0.2000}}
  >{\raggedright\arraybackslash}p{(\linewidth - 8\tabcolsep) * \real{0.2000}}
  >{\raggedright\arraybackslash}p{(\linewidth - 8\tabcolsep) * \real{0.2000}}@{}}
\toprule\noalign{}
\begin{minipage}[b]{\linewidth}\raggedright
Formation
\end{minipage} & \begin{minipage}[b]{\linewidth}\raggedright
Candidate input
\end{minipage} & \begin{minipage}[b]{\linewidth}\raggedright
Expected or documented OE outcome
\end{minipage} & \begin{minipage}[b]{\linewidth}\raggedright
OE comparison form
\end{minipage} & \begin{minipage}[b]{\linewidth}\raggedright
Result
\end{minipage} \\
\midrule\noalign{}
\endhead
\bottomrule\noalign{}
\endlastfoot
broader family label & *gállą & broader cognate-set headword & wider
family context & useful lexeme label, but not the direct source of
\emph{ġealla} \\
selected weak noun & *gállô & compact-trace output: \emph{ġealla} &
ġealla & exact match between formation and attested OE noun \\
dialectal Anglian continuation & weak noun branch & Anglian \emph{galla}
type & galla & genuine OE variant, but not the selected West Saxon
target \\
\end{longtable}
}

The weak-noun stem class is the relevant point. It gives a direct route
to attested \emph{ġealla}, while the broader comparative label serves
only as a family heading.

\subsubsection{knight --- OE cniht}\label{knight-oe-cniht}

Derivation: citation reconstruction \emph{*kníxtaz}; selected input
\emph{*knéxtaz} \textgreater{} \emph{cniht} (early analogy).

\paragraph{Derivation trace}\label{derivation-trace-13}

Proto input: \emph{*knéxtaz}

\begingroup
\setlength{\fboxsep}{6pt}

\noindent\fbox{%
\begin{minipage}{0.97\linewidth}
\small
\begin{tabularx}{\linewidth}{@{}>{\raggedright\arraybackslash}p{0.320\linewidth}>{\centering\arraybackslash}X>{\raggedright\arraybackslash}p{0.520\linewidth}@{}}
\begin{minipage}[t]{\linewidth}
\centering\textbf{Earlier Germanic changes}\par
\vspace{0.35em}
\raggedright
\centering\textbf{West Germanic}\par
\raggedright
\vspace{0.2em}
\raggedright [no change]\par
\vspace{0.6em}
\centering\textbf{Northwest Germanic}\par
\raggedright
\vspace{0.2em}
\begin{tabular}{@{}>{\raggedright\arraybackslash}p{0.74\linewidth}@{\hspace{0.55em}}>{\raggedright\arraybackslash}p{0.16\linewidth}@{\hspace{0.25em}}}
\mbox{PGmc Final Z Deletion} & \emph{*knéxta} \\
\end{tabular}
\end{minipage}
&
&
\begin{minipage}[t]{\linewidth}
\centering\textbf{Old English changes}\par
\vspace{0.35em}
\raggedright
\begin{tabular}{@{}>{\raggedright\arraybackslash}p{0.70\linewidth}@{\hspace{0.55em}}>{\raggedright\arraybackslash}p{0.16\linewidth}@{\hspace{0.25em}}}
PWGmc Final Bare A Loss & \emph{*knéxt} \\
OE Breaking & \emph{*knéoxt} \\
OE Ws Palatal Umlaut & \emph{*knixt} \\
\end{tabular}
\end{minipage}
\\
\end{tabularx}
\end{minipage}%
} \endgroup

Outcome: \emph{cniht}

\paragraph{Reconstruction and comparative
evidence}\label{reconstruction-and-comparative-evidence-13}

The comparative sources align on an \emph{e}-grade reconstruction for
this noun. Ringe and Taylor cite \emph{*kneht}, and Orel gives
\emph{*knextaz} (\citeproc{ref-RingeTaylor2014}{Ringe and Taylor 2014,
142}; \citeproc{ref-Orel2003}{Orel 2003, 256}). Kluge-Seebold likewise
points to \emph{*knehta-} (\citeproc{ref-KlugeSeebold2011}{Kluge 2011}).
The selected input \emph{*knéxtaz} follows that comparative evidence.

A competing citation reconstruction \emph{*kníxtaz} remains possible as
a label for the word family, but it is not the reconstruction followed
here. The Old English development discussed below is based on
\emph{*knéxtaz}.

\paragraph{Old English evidence}\label{old-english-evidence-13}

Old English dictionaries record the noun as \emph{cniht}
(\citeproc{ref-ClarkHall1960}{Clark Hall 1960};
\citeproc{ref-BosworthToller1898}{Bosworth and Toller 1898, 71}).
Campbell cites plural \emph{cneohtas} among the broken forms, showing
the same vowel environment from another point in the paradigm
(\citeproc{ref-Campbell1959}{Campbell 1959, sect. 146}).

The target is therefore an ordinary attested Old English noun. No
reconstructed OE comparator is needed here.

\paragraph{Development to Old
English}\label{development-to-old-english-13}

From \emph{*knéxtaz}, the relevant Old English changes include breaking
before the velar cluster and then the later reduction that yields
\emph{cniht}. Campbell later notes the early West-Saxon alternation
\emph{cniht} beside plural \emph{cneohtas}
(\citeproc{ref-Campbell1959}{Campbell 1959, sect. 305}). Sievers-Brunner
gives the same contrast as \emph{cniht \ldots{} cneohtas}
(\citeproc{ref-SieversBrunner1965}{Sievers 1965, sect. 122}). With that
corrected input, the derivation is straightforward.

\paragraph{Stem comparison}\label{stem-comparison-1}

The comparison below is manual. It separates the handbook-supported
\emph{e}-grade input from a competing citation reconstruction.

{\def\LTcaptype{none} % do not increment counter
\begin{longtable}[]{@{}
  >{\raggedright\arraybackslash}p{(\linewidth - 8\tabcolsep) * \real{0.2000}}
  >{\raggedright\arraybackslash}p{(\linewidth - 8\tabcolsep) * \real{0.2000}}
  >{\raggedright\arraybackslash}p{(\linewidth - 8\tabcolsep) * \real{0.2000}}
  >{\raggedright\arraybackslash}p{(\linewidth - 8\tabcolsep) * \real{0.2000}}
  >{\raggedright\arraybackslash}p{(\linewidth - 8\tabcolsep) * \real{0.2000}}@{}}
\toprule\noalign{}
\begin{minipage}[b]{\linewidth}\raggedright
Formation / label
\end{minipage} & \begin{minipage}[b]{\linewidth}\raggedright
Candidate input
\end{minipage} & \begin{minipage}[b]{\linewidth}\raggedright
Expected or documented OE outcome
\end{minipage} & \begin{minipage}[b]{\linewidth}\raggedright
OE comparison form
\end{minipage} & \begin{minipage}[b]{\linewidth}\raggedright
Result
\end{minipage} \\
\midrule\noalign{}
\endhead
\bottomrule\noalign{}
\endlastfoot
competing citation reconstruction & *kníxtaz & not the reconstruction
followed here & broader citation tradition & useful as a competing
label, but not the source-based choice used for the OE derivation \\
handbook-supported reconstruction & *knéxtaz & compact-trace output:
\emph{cniht} & cniht & exact match between comparative reconstruction
and attested OE noun \\
related plural evidence & same stem family & plural \emph{cneohtas} type
background & cneohtas & supports the vowel environment, but not the
selected target cell \\
\end{longtable}
}

\subsubsection{lade --- OE hladan}\label{lade-oe-hladan}

Derivation: citation reconstruction \emph{*laθōjaną}; selected input
\emph{*xláðaną} \textgreater{} \emph{hladan} (early analogy).

\paragraph{Derivation trace}\label{derivation-trace-14}

Proto input: \emph{*xláðaną}

\begingroup
\setlength{\fboxsep}{6pt}

\noindent\fbox{%
\begin{minipage}{0.97\linewidth}
\small
\begin{tabularx}{\linewidth}{@{}>{\raggedright\arraybackslash}p{0.300\linewidth}>{\centering\arraybackslash}X>{\raggedright\arraybackslash}p{0.540\linewidth}@{}}
\begin{minipage}[t]{\linewidth}
\centering\textbf{Earlier Germanic changes}\par
\vspace{0.35em}
\raggedright
\centering\textbf{West Germanic}\par
\raggedright
\vspace{0.2em}
\begin{tabular}{@{}>{\raggedright\arraybackslash}p{0.70\linewidth}@{\hspace{0.55em}}>{\raggedright\arraybackslash}p{0.20\linewidth}@{\hspace{0.25em}}}
PWGmc Dental Hardening & \emph{*xládaną} \\
\end{tabular}
\vspace{0.6em}
\centering\textbf{Northwest Germanic}\par
\raggedright
\vspace{0.2em}
\raggedright [no change]\par
\end{minipage}
&
&
\begin{minipage}[t]{\linewidth}
\centering\textbf{Old English changes}\par
\vspace{0.35em}
\raggedright
\begin{tabular}{@{}>{\raggedright\arraybackslash}p{0.70\linewidth}@{\hspace{0.55em}}>{\raggedright\arraybackslash}p{0.20\linewidth}@{\hspace{0.25em}}}
Anglo Frisian Brightening & \emph{*xlædaną} \\
OE A Restoration & \emph{*xladaną} \\
OE Heavy Syllable Nasal Apocope & \emph{*xladan} \\
OE Secondary Nasalization & \emph{*xladąn} \\
OE Weak Tail Reduction & \emph{*xladan} \\
\end{tabular}
\end{minipage}
\\
\end{tabularx}
\end{minipage}%
} \endgroup

Outcome: \emph{hladan}

\paragraph{Reconstruction and comparative
evidence}\label{reconstruction-and-comparative-evidence-14}

Ringe and Taylor cite the strong verb \emph{hladan} directly
(\citeproc{ref-RingeTaylor2014}{Ringe and Taylor 2014, 248}). Kroonen
treats the wider Germanic family under a weak-verb label such as
\emph{*laθōjaną} (\citeproc{ref-Kroonen2013}{Kroonen 2013}). The two
forms are related as members of one word family, but they do not play
the same role in the OE derivation.

The selected input therefore marks an early stem choice. The entry
follows the strong Verner-grade form that reaches Old English
\emph{hladan} directly.

\paragraph{Old English evidence}\label{old-english-evidence-14}

Bosworth-Toller records the verb as \emph{hladan} and preserves the
expected strong-verb paradigm material around it
(\citeproc{ref-BosworthToller1898}{Bosworth and Toller 1898, 559}).
Clark Hall likewise records \emph{hladan}
(\citeproc{ref-ClarkHall1960}{Clark Hall 1960, 159}). The target is an
attested infinitive rather than a reconstructed paradigm cell.

For this entry the relevant comparison form is the infinitive
\emph{hladan} itself. The question is how that attested strong verb
relates to the broader comparative family.

\paragraph{Development to Old
English}\label{development-to-old-english-14}

From \emph{*xláðaną}, the verb passes through the expected early voiced
stop stage, Anglo-Frisian brightening, and the later A-restoration that
returns the root vowel to \emph{a} before the full infinitival ending.
Campbell treats verbs such as \emph{hladan} as showing restored \emph{a}
in the present system (\citeproc{ref-Campbell1959}{Campbell 1959, sect.
744}). Ringe and Taylor likewise derive \emph{hladan} from the voiced
strong-verb line (\citeproc{ref-RingeTaylor2014}{Ringe and Taylor 2014,
248}). The resulting Old English infinitive is \emph{hladan}.

\paragraph{Class comparison}\label{class-comparison-2}

The comparison below is manual. It separates the wider weak-verb family
label from the strong verb actually reflected in Old English.

{\def\LTcaptype{none} % do not increment counter
\begin{longtable}[]{@{}
  >{\raggedright\arraybackslash}p{(\linewidth - 8\tabcolsep) * \real{0.2000}}
  >{\raggedright\arraybackslash}p{(\linewidth - 8\tabcolsep) * \real{0.2000}}
  >{\raggedright\arraybackslash}p{(\linewidth - 8\tabcolsep) * \real{0.2000}}
  >{\raggedright\arraybackslash}p{(\linewidth - 8\tabcolsep) * \real{0.2000}}
  >{\raggedright\arraybackslash}p{(\linewidth - 8\tabcolsep) * \real{0.2000}}@{}}
\toprule\noalign{}
\begin{minipage}[b]{\linewidth}\raggedright
Formation / class
\end{minipage} & \begin{minipage}[b]{\linewidth}\raggedright
Candidate input
\end{minipage} & \begin{minipage}[b]{\linewidth}\raggedright
Expected or documented OE outcome
\end{minipage} & \begin{minipage}[b]{\linewidth}\raggedright
OE comparison form
\end{minipage} & \begin{minipage}[b]{\linewidth}\raggedright
Result
\end{minipage} \\
\midrule\noalign{}
\endhead
\bottomrule\noalign{}
\endlastfoot
comparative weak-verb label & *laθōjaną & wider family background &
broader family context & useful family label, but not the direct source
of \emph{hladan} \\
selected strong Verner-grade input & *xláðaną & compact-trace output:
\emph{hladan} & hladan & exact match between formation and attested OE
infinitive \\
\end{longtable}
}

\subsubsection{lap --- OE lappa}\label{lap-oe-lappa}

Derivation: citation reconstruction \emph{*lábbaz}; selected input
\emph{*láppô} \textgreater{} \emph{lappa} (early analogy).

\paragraph{Derivation trace}\label{derivation-trace-15}

Proto input: \emph{*láppô}

\begingroup
\setlength{\fboxsep}{6pt}

\noindent\fbox{%
\begin{minipage}{0.97\linewidth}
\small
\begin{tabularx}{\linewidth}{@{}>{\raggedright\arraybackslash}p{0.300\linewidth}>{\centering\arraybackslash}X>{\raggedright\arraybackslash}p{0.540\linewidth}@{}}
\begin{minipage}[t]{\linewidth}
\centering\textbf{Earlier Germanic changes}\par
\vspace{0.35em}
\raggedright
\centering\textbf{West Germanic}\par
\raggedright
\vspace{0.2em}
\raggedright [no change]\par
\vspace{0.6em}
\centering\textbf{Northwest Germanic}\par
\raggedright
\vspace{0.2em}
\raggedright [no change]\par
\end{minipage}
&
&
\begin{minipage}[t]{\linewidth}
\centering\textbf{Old English changes}\par
\vspace{0.35em}
\raggedright
\begin{tabular}{@{}>{\raggedright\arraybackslash}p{0.74\linewidth}@{\hspace{0.55em}}>{\raggedright\arraybackslash}p{0.16\linewidth}@{\hspace{0.25em}}}
Anglo Frisian Brightening & \emph{*læppô} \\
OE A Restoration & \emph{*lappô} \\
OE Unstressed Long Vowel Shortening & \emph{*lappa} \\
\end{tabular}
\end{minipage}
\\
\end{tabularx}
\end{minipage}%
} \endgroup

Outcome: \emph{lappa}

\paragraph{Reconstruction and comparative
evidence}\label{reconstruction-and-comparative-evidence-15}

The comparative sources point to a weak noun with \emph{pp}: Orel gives
\emph{*lappōn}, Kroonen places the word in a weak \emph{*lappan-}
family, and Kluge-Seebold cites West Germanic \emph{lappa} forms
alongside Old English \emph{læppa} (\citeproc{ref-Orel2003}{Orel 2003};
\citeproc{ref-Kroonen2013}{Kroonen 2013};
\citeproc{ref-KlugeSeebold2011}{Kluge 2011}). The selected input
\emph{*láppô} follows that evidence.

A competing comparative label \emph{*lábbaz} has also circulated for the
word family, but the cited handbooks do not make it the direct source of
the Old English weak noun. The form relevant to the OE development is
the weak masculine input \emph{*láppô}.

\paragraph{Old English evidence}\label{old-english-evidence-15}

Campbell cites \emph{lappa} as a case of restored \emph{a}, while
Sievers-Brunner records \emph{lappa}, variant \emph{læppa}, and plural
or oblique \emph{leappan} (\citeproc{ref-Campbell1959}{Campbell 1959};
\citeproc{ref-SieversBrunner1965}{Sievers 1965}). The dictionary
tradition also preserves \emph{læppa}
(\citeproc{ref-ClarkHall1960}{Clark Hall 1960, 180};
\citeproc{ref-BosworthToller1898}{Bosworth and Toller 1898, 613}).

The target of this entry is the restored singular \emph{lappa}. The
variant \emph{læppa} and the oblique or plural \emph{leappan} remain
part of the Old English record and help frame the noun's vowel history.

\paragraph{Development to Old
English}\label{development-to-old-english-15}

From \emph{*láppô}, Anglo-Frisian brightening first yields \emph{æ}, and
later A-restoration returns the root vowel to \emph{a} before the
back-vocalic ending, after which shortening of the final vowel gives
\emph{lappa} (\citeproc{ref-Campbell1959}{Campbell 1959};
\citeproc{ref-SieversBrunner1965}{Sievers 1965}). With the weak
masculine input chosen, the OE development is regular.

\paragraph{Stem comparison}\label{stem-comparison-2}

The comparison below is manual. It separates the weak masculine
formation from a competing voiced comparative label.

{\def\LTcaptype{none} % do not increment counter
\begin{longtable}[]{@{}
  >{\raggedright\arraybackslash}p{(\linewidth - 8\tabcolsep) * \real{0.2000}}
  >{\raggedright\arraybackslash}p{(\linewidth - 8\tabcolsep) * \real{0.2000}}
  >{\raggedright\arraybackslash}p{(\linewidth - 8\tabcolsep) * \real{0.2000}}
  >{\raggedright\arraybackslash}p{(\linewidth - 8\tabcolsep) * \real{0.2000}}
  >{\raggedright\arraybackslash}p{(\linewidth - 8\tabcolsep) * \real{0.2000}}@{}}
\toprule\noalign{}
\begin{minipage}[b]{\linewidth}\raggedright
Formation / label
\end{minipage} & \begin{minipage}[b]{\linewidth}\raggedright
Candidate input
\end{minipage} & \begin{minipage}[b]{\linewidth}\raggedright
Expected or documented OE outcome
\end{minipage} & \begin{minipage}[b]{\linewidth}\raggedright
OE comparison form
\end{minipage} & \begin{minipage}[b]{\linewidth}\raggedright
Result
\end{minipage} \\
\midrule\noalign{}
\endhead
\bottomrule\noalign{}
\endlastfoot
competing voiced comparative label & *lábbaz & not the form followed for
the OE weak-noun derivation & broader comparative background & useful as
a competing label, but not the source-based choice used here \\
selected weak masculine noun & *láppô & compact-trace output:
\emph{lappa} & lappa & exact match between formation and attested OE
noun \\
attested OE variant line & same noun family & \emph{læppa},
\emph{leappan} & læppa / leappan & useful control forms within the same
OE tradition \\
\end{longtable}
}

\subsubsection{laugh --- OE hliehhan}\label{laugh-oe-hliehhan}

Derivation: citation reconstruction \emph{*lákaną}; selected input
\emph{*xláxjaną} \textgreater{} \emph{hliehhan} (early analogy).

\paragraph{Derivation trace}\label{derivation-trace-16}

Proto input: \emph{*xláxjaną}

\begingroup
\setlength{\fboxsep}{6pt}

\noindent\fbox{%
\begin{minipage}{0.97\linewidth}
\footnotesize
\begin{tabularx}{\linewidth}{@{}>{\raggedright\arraybackslash}p{0.470\linewidth}>{\centering\arraybackslash}X>{\raggedright\arraybackslash}p{0.470\linewidth}@{}}
\begin{minipage}[t]{\linewidth}
\centering\textbf{Earlier Germanic changes}\par
\vspace{0.35em}
\raggedright
\centering\textbf{West Germanic}\par
\raggedright
\vspace{0.2em}
\begin{tabular}{@{}>{\raggedright\arraybackslash}p{0.64\linewidth}@{\hspace{0.55em}}>{\raggedright\arraybackslash}p{0.24\linewidth}@{\hspace{0.25em}}}
PWGmc J Gemination & \emph{*xláxxjaną} \\
\end{tabular}
\vspace{0.6em}
\centering\textbf{Northwest Germanic}\par
\raggedright
\vspace{0.2em}
\raggedright [no change]\par
\end{minipage}
&
&
\begin{minipage}[t]{\linewidth}
\centering\textbf{Old English changes}\par
\vspace{0.35em}
\raggedright
\begin{tabular}{@{}>{\raggedright\arraybackslash}p{0.62\linewidth}@{\hspace{0.55em}}>{\raggedright\arraybackslash}p{0.28\linewidth}@{\hspace{0.25em}}}
Anglo Frisian Brightening & \emph{*xlæxxjaną} \\
OE Breaking & \emph{*xleaxxjaną} \\
OE Velar Fricative Palatalization & \emph{*xleaxçjaną} \\
OE Heavy Syllable Nasal Apocope & \emph{*xleaxçjan} \\
OE Secondary Nasalization & \emph{*xleaxçjąn} \\
OE I Umlaut & \emph{*xliexçjąn} \\
OE Weak Tail Reduction & \emph{*xliexçjan} \\
OE J Loss After Heavy & \emph{*xliexçan} \\
\end{tabular}
\end{minipage}
\\
\end{tabularx}
\end{minipage}%
} \endgroup

Outcome: \emph{hliehhan}

\paragraph{Reconstruction and comparative
evidence}\label{reconstruction-and-comparative-evidence-16}

The wider Germanic family includes a non-j branch represented here by
the citation label \emph{*lákaną}, but Kroonen and Ringe-Taylor both
distinguish a j-present formation behind Old English \emph{hliehhan}
(\citeproc{ref-Kroonen2013}{Kroonen 2013};
\citeproc{ref-RingeTaylor2014}{Ringe and Taylor 2014}). The selected
input \emph{*xláxjaną} reflects that j-present line.

This branch choice matters because it brings with it the geminate
fricative and the vowel development characteristic of the Old English
verb. The comparative family label and the OE-facing input are therefore
related but not identical.

\paragraph{Old English evidence}\label{old-english-evidence-16}

Old English dictionaries and readers record the verb as \emph{hliehhan},
while also preserving variants such as \emph{hlæhhan} and \emph{hlehhan}
(\citeproc{ref-BosworthToller1898}{Bosworth and Toller 1898};
\citeproc{ref-ClarkHall1960}{Clark Hall 1960};
\citeproc{ref-BrightCassidyRingler1971}{Bright 1971}). The target of
this entry is the West Saxon \emph{hliehhan}.

The variant set matters as background, but the argument of the entry
rests on the attested lemma \emph{hliehhan} itself.

\paragraph{Development to Old
English}\label{development-to-old-english-16}

From \emph{*xláxjaną}, regular j-gemination yields the doubled
fricative, and the subsequent Old English vowel developments lead to
\emph{hliehhan} (\citeproc{ref-Fulk2018}{Fulk 2018};
\citeproc{ref-Campbell1959}{Campbell 1959}). Ringe and Taylor discuss
the broken vowel of the Old English form as part of this same history,
with possible support from the related noun \emph{hleahtor}
(\citeproc{ref-RingeTaylor2014}{Ringe and Taylor 2014}).

\paragraph{Branch comparison}\label{branch-comparison}

The comparison below is manual. It separates the wider non-j family
label from the j-present branch actually reflected in Old English.

{\def\LTcaptype{none} % do not increment counter
\begin{longtable}[]{@{}
  >{\raggedright\arraybackslash}p{(\linewidth - 8\tabcolsep) * \real{0.2000}}
  >{\raggedright\arraybackslash}p{(\linewidth - 8\tabcolsep) * \real{0.2000}}
  >{\raggedright\arraybackslash}p{(\linewidth - 8\tabcolsep) * \real{0.2000}}
  >{\raggedright\arraybackslash}p{(\linewidth - 8\tabcolsep) * \real{0.2000}}
  >{\raggedright\arraybackslash}p{(\linewidth - 8\tabcolsep) * \real{0.2000}}@{}}
\toprule\noalign{}
\begin{minipage}[b]{\linewidth}\raggedright
Formation / branch
\end{minipage} & \begin{minipage}[b]{\linewidth}\raggedright
Candidate input
\end{minipage} & \begin{minipage}[b]{\linewidth}\raggedright
Expected or documented OE outcome
\end{minipage} & \begin{minipage}[b]{\linewidth}\raggedright
OE comparison form
\end{minipage} & \begin{minipage}[b]{\linewidth}\raggedright
Result
\end{minipage} \\
\midrule\noalign{}
\endhead
\bottomrule\noalign{}
\endlastfoot
wider non-j family & *lákaną & comparative background outside the
selected OE line & wider family context & useful family label, but not
the direct source of \emph{hliehhan} \\
selected j-present branch & *xláxjaną & compact-trace output:
\emph{hliehhan} & hliehhan & exact match between branch and attested OE
lemma \\
attested OE variants & same OE verb line & \emph{hlæhhan},
\emph{hlehhan} & hlæhhan / hlehhan & genuine variant evidence, but
secondary to the selected form \\
\end{longtable}
}

\subsubsection{loam --- OE lām}\label{loam-oe-lux101m}

Derivation: citation reconstruction \emph{*laimōn}; selected input
\emph{*láimą} \textgreater{} \emph{lām} (early analogy).

\paragraph{Derivation trace}\label{derivation-trace-17}

Proto input: \emph{*láimą}

\begingroup
\setlength{\fboxsep}{6pt}

\noindent\fbox{%
\begin{minipage}{0.97\linewidth}
\small
\begin{tabularx}{\linewidth}{@{}>{\raggedright\arraybackslash}p{0.300\linewidth}>{\centering\arraybackslash}X>{\raggedright\arraybackslash}p{0.540\linewidth}@{}}
\begin{minipage}[t]{\linewidth}
\centering\textbf{Earlier Germanic changes}\par
\vspace{0.35em}
\raggedright
\centering\textbf{West Germanic}\par
\raggedright
\vspace{0.2em}
\begin{tabular}{@{}>{\raggedright\arraybackslash}p{0.70\linewidth}@{\hspace{0.55em}}>{\raggedright\arraybackslash}p{0.16\linewidth}@{\hspace{0.25em}}}
PWGmc Ai Monophthongization & \emph{*lāmą} \\
\end{tabular}
\vspace{0.6em}
\centering\textbf{Northwest Germanic}\par
\raggedright
\vspace{0.2em}
\raggedright [no change]\par
\end{minipage}
&
&
\begin{minipage}[t]{\linewidth}
\centering\textbf{Old English changes}\par
\vspace{0.35em}
\raggedright
\begin{tabular}{@{}>{\raggedright\arraybackslash}p{0.74\linewidth}@{\hspace{0.55em}}>{\raggedright\arraybackslash}p{0.16\linewidth}@{\hspace{0.25em}}}
OE Heavy Syllable Nasal Apocope & \emph{*lām} \\
\end{tabular}
\end{minipage}
\\
\end{tabularx}
\end{minipage}%
} \endgroup

Outcome: \emph{lām}

\paragraph{Reconstruction and comparative
evidence}\label{reconstruction-and-comparative-evidence-17}

The inherited comparative noun is given as \emph{*laimōn} or
\emph{*laiman-}, and both Orel and Kroonen identify Old English
\emph{lām} as a neuter reflex of that family
(\citeproc{ref-Orel2003}{Orel 2003}; \citeproc{ref-Kroonen2013}{Kroonen
2013, 363}). The selected input \emph{*láimą} differs from the
comparative headword because it represents the stem class that matches
the Old English noun most directly.

This is therefore a class shift within the history of the English branch
rather than a dispute about the OE target itself.

\paragraph{Old English evidence}\label{old-english-evidence-17}

Old English dictionaries record the noun as \emph{lām}, a neuter word
for `loam, clay, mud' (\citeproc{ref-BosworthToller1898}{Bosworth and
Toller 1898}; \citeproc{ref-ClarkHall1960}{Clark Hall 1960, 196}). The
target is an attested citation form rather than a reconstructed
comparator.

The relevant question is not whether \emph{lām} is Old English, but
which inherited formation best accounts for that attested neuter noun.

\paragraph{Development to Old
English}\label{development-to-old-english-17}

From \emph{*láimą}, regular monophthongization of \emph{ai} and the
later loss of the final nasal syllable yield \emph{lām}. With that
OE-facing input, the phonological development is straightforward.

\paragraph{Class comparison}\label{class-comparison-3}

The comparison below is manual. It separates the inherited comparative
n-stem label from the OE-facing stem class used to derive the attested
noun.

{\def\LTcaptype{none} % do not increment counter
\begin{longtable}[]{@{}
  >{\raggedright\arraybackslash}p{(\linewidth - 8\tabcolsep) * \real{0.2000}}
  >{\raggedright\arraybackslash}p{(\linewidth - 8\tabcolsep) * \real{0.2000}}
  >{\raggedright\arraybackslash}p{(\linewidth - 8\tabcolsep) * \real{0.2000}}
  >{\raggedright\arraybackslash}p{(\linewidth - 8\tabcolsep) * \real{0.2000}}
  >{\raggedright\arraybackslash}p{(\linewidth - 8\tabcolsep) * \real{0.2000}}@{}}
\toprule\noalign{}
\begin{minipage}[b]{\linewidth}\raggedright
Formation / class
\end{minipage} & \begin{minipage}[b]{\linewidth}\raggedright
Candidate input
\end{minipage} & \begin{minipage}[b]{\linewidth}\raggedright
Expected or documented OE outcome
\end{minipage} & \begin{minipage}[b]{\linewidth}\raggedright
OE comparison form
\end{minipage} & \begin{minipage}[b]{\linewidth}\raggedright
Result
\end{minipage} \\
\midrule\noalign{}
\endhead
\bottomrule\noalign{}
\endlastfoot
inherited comparative noun & *laimōn & comparative family background &
wider family context & useful headword, but not the direct OE-facing
input \\
selected OE-facing stem class & *láimą & compact-trace output:
\emph{lām} & lām & exact match between input and attested OE noun \\
\end{longtable}
}

\subsubsection{lung --- OE lungen}\label{lung-oe-lungen}

Derivation: citation reconstruction \emph{*lungō}; selected input
\emph{*lúnganjō} \textgreater{} \emph{lungen} (early analogy).

\paragraph{Derivation trace}\label{derivation-trace-18}

Proto input: \emph{*lúnganjō}

\begingroup
\setlength{\fboxsep}{6pt}

\noindent\fbox{%
\begin{minipage}{0.97\linewidth}
\small
\begin{tabularx}{\linewidth}{@{}>{\raggedright\arraybackslash}p{0.470\linewidth}>{\centering\arraybackslash}X>{\raggedright\arraybackslash}p{0.470\linewidth}@{}}
\begin{minipage}[t]{\linewidth}
\centering\textbf{Earlier Germanic changes}\par
\vspace{0.35em}
\raggedright
\centering\textbf{West Germanic}\par
\raggedright
\vspace{0.2em}
\begin{tabular}{@{}>{\raggedright\arraybackslash}p{0.66\linewidth}@{\hspace{0.55em}}>{\raggedright\arraybackslash}p{0.24\linewidth}@{\hspace{0.25em}}}
PWGmc J Gemination & \emph{*lúngannjō} \\
\end{tabular}
\vspace{0.6em}
\centering\textbf{Northwest Germanic}\par
\raggedright
\vspace{0.2em}
\begin{tabular}{@{}>{\raggedright\arraybackslash}p{0.66\linewidth}@{\hspace{0.55em}}>{\raggedright\arraybackslash}p{0.24\linewidth}@{\hspace{0.25em}}}
NWGmc Final Long O Raising & \emph{*lúngannju} \\
\end{tabular}
\end{minipage}
&
&
\begin{minipage}[t]{\linewidth}
\centering\textbf{Old English changes}\par
\vspace{0.35em}
\raggedright
\begin{tabular}{@{}>{\raggedright\arraybackslash}p{0.64\linewidth}@{\hspace{0.55em}}>{\raggedright\arraybackslash}p{0.26\linewidth}@{\hspace{0.25em}}}
OE I Umlaut & \emph{*lúngennju} \\
OE High Vowel Apocope & \emph{*lúngennj} \\
OE J Loss After Heavy & \emph{*lúngenn} \\
OE Final Geminate Simplification & \emph{*lúngen} \\
\end{tabular}
\end{minipage}
\\
\end{tabularx}
\end{minipage}%
} \endgroup

Outcome: \emph{lungen}

\paragraph{Reconstruction and comparative
evidence}\label{reconstruction-and-comparative-evidence-18}

Kroonen treats the basic noun as \emph{*lungōn-} and also cites an
OE-facing derivative \emph{*lungunjō-}, continued by Old English
\emph{lungen} and close West Germanic cognates
(\citeproc{ref-Kroonen2013}{Kroonen 2013}). The selected input
\emph{*lúnganjō} models that derived feminine formation rather than the
base noun. The notation differs slightly from Kroonen's
\emph{*lungunjō-}, but both point to the same derived feminine line.

The difference between the citation label and the selected input is
therefore derivational. Old English \emph{lungen} is not a direct reflex
of the bare base noun \emph{*lungō}; it belongs to an expanded feminine
formation.

\paragraph{Old English evidence}\label{old-english-evidence-18}

Old English dictionaries record the noun as \emph{lungen}, with
inflected forms such as \emph{lungenne} and \emph{lungena}
(\citeproc{ref-BosworthToller1898}{Bosworth and Toller 1898, 634}).
Clark Hall also preserves a small family of compounds such as
\emph{lungenādl}, \emph{lungensealf}, and \emph{lungenwyrt}
(\citeproc{ref-ClarkHall1960}{Clark Hall 1960}).

The target is an attested Old English lexeme with its own paradigm, not
a rescued inflectional cell.

\paragraph{Development to Old
English}\label{development-to-old-english-18}

From the selected derived input, the expected derivational consonant and
vowel adjustments lead to \emph{lungen}. Once the expanded feminine
formation is chosen, the Old English outcome is regular.

\paragraph{Formation comparison}\label{formation-comparison-3}

The comparison below is manual. It separates the base noun from the
derived feminine formation reflected in Old English.

{\def\LTcaptype{none} % do not increment counter
\begin{longtable}[]{@{}
  >{\raggedright\arraybackslash}p{(\linewidth - 8\tabcolsep) * \real{0.2000}}
  >{\raggedright\arraybackslash}p{(\linewidth - 8\tabcolsep) * \real{0.2000}}
  >{\raggedright\arraybackslash}p{(\linewidth - 8\tabcolsep) * \real{0.2000}}
  >{\raggedright\arraybackslash}p{(\linewidth - 8\tabcolsep) * \real{0.2000}}
  >{\raggedright\arraybackslash}p{(\linewidth - 8\tabcolsep) * \real{0.2000}}@{}}
\toprule\noalign{}
\begin{minipage}[b]{\linewidth}\raggedright
Formation
\end{minipage} & \begin{minipage}[b]{\linewidth}\raggedright
Candidate input
\end{minipage} & \begin{minipage}[b]{\linewidth}\raggedright
Expected or documented OE outcome
\end{minipage} & \begin{minipage}[b]{\linewidth}\raggedright
OE comparison form
\end{minipage} & \begin{minipage}[b]{\linewidth}\raggedright
Result
\end{minipage} \\
\midrule\noalign{}
\endhead
\bottomrule\noalign{}
\endlastfoot
base noun & *lungō & base-noun outcome without the OE derivative suffix
& broader family context & useful headword, but not the direct source of
\emph{lungen} \\
derived OE-facing formation & *lúnganjō & compact-trace output:
\emph{lungen} & lungen & exact match between selected formation and
attested OE noun \\
Kroonen's cited derivative & \emph{*lungunjō-} & comparative support for
the same OE-facing formation & lungen and cognate set & supports the
derived feminine formation, with notation differing from the normalized
input form used here \\
\end{longtable}
}

\subsubsection{navel --- OE nafola}\label{navel-oe-nafola}

Derivation: citation reconstruction \emph{*nablô}; selected input
\emph{*nábulô} \textgreater{} \emph{nafola} (early analogy).

\paragraph{Derivation trace}\label{derivation-trace-19}

Proto input: \emph{*nábulô}

\begingroup
\setlength{\fboxsep}{6pt}

\noindent\fbox{%
\begin{minipage}{0.97\linewidth}
\small
\begin{tabularx}{\linewidth}{@{}>{\raggedright\arraybackslash}p{0.300\linewidth}>{\centering\arraybackslash}X>{\raggedright\arraybackslash}p{0.540\linewidth}@{}}
\begin{minipage}[t]{\linewidth}
\centering\textbf{Earlier Germanic changes}\par
\vspace{0.35em}
\raggedright
\centering\textbf{West Germanic}\par
\raggedright
\vspace{0.2em}
\raggedright [no change]\par
\vspace{0.6em}
\centering\textbf{Northwest Germanic}\par
\raggedright
\vspace{0.2em}
\raggedright [no change]\par
\end{minipage}
&
&
\begin{minipage}[t]{\linewidth}
\centering\textbf{Old English changes}\par
\vspace{0.35em}
\raggedright
\begin{tabular}{@{}>{\raggedright\arraybackslash}p{0.74\linewidth}@{\hspace{0.55em}}>{\raggedright\arraybackslash}p{0.16\linewidth}@{\hspace{0.25em}}}
OE Med Unstressed U Lowering & \emph{*nábolô} \\
Anglo Frisian Brightening & \emph{*næbolô} \\
OE A Restoration & \emph{*nabolô} \\
PGmc B Allophony & \emph{*naβolô} \\
OE Unstressed Long Vowel Shortening & \emph{*naβola} \\
\end{tabular}
\end{minipage}
\\
\end{tabularx}
\end{minipage}%
} \endgroup

Outcome: \emph{nafola}

\paragraph{Reconstruction and comparative
evidence}\label{reconstruction-and-comparative-evidence-19}

Kroonen lemmatizes the word with a syncopated comparative headword
\emph{*nablô}, while Ringe and Taylor give the derivational pathway
\emph{*nabulō} \textgreater{} \emph{*næbula} \textgreater{}
\emph{nafola} (\citeproc{ref-Kroonen2013}{Kroonen 2013};
\citeproc{ref-RingeTaylor2014}{Ringe and Taylor 2014, 270}). The
difference is one of stage and notation rather than of lexeme identity:
the selected input \emph{*nábulô} is the pre-syncope form needed for the
Old English development.

The literature also differs on the older history of the medial \emph{u},
whether it is inherited or secondary
(\citeproc{ref-Streitberg1896}{Streitberg 1896};
\citeproc{ref-Ringe2006}{Ringe 2006};
\citeproc{ref-Mayrhofer1992}{Mayrhofer 1992-\/-2001};
\citeproc{ref-SieversBrunner1965}{Sievers 1965}). For the Old English
comparison, however, both lines place a medial vowel in the pre-OE form.

\paragraph{Old English evidence}\label{old-english-evidence-19}

Ringe and Taylor note the early West Saxon shift \emph{nafola}
\textgreater{} \emph{nafela} (\citeproc{ref-RingeTaylor2014}{Ringe and
Taylor 2014, 336}). Campbell likewise records \emph{nafela} beside
Corpus \emph{nabula} (\citeproc{ref-Campbell1959}{Campbell 1959, sect.
159}). The target of this entry is the nominative singular
\emph{nafola}, the form that matches the selected derivational pathway
most directly.

\emph{nafela} is the better-known later West Saxon spelling, while
\emph{nabula} preserves a less reduced medial vowel. These forms belong
to the same lexical history, but this entry is centered on
\emph{nafola}.

\paragraph{Development to Old
English}\label{development-to-old-english-19}

Ringe and Taylor give the pre-OE line \emph{*nabulō} \textgreater{}
\emph{*næbula} \textgreater{} OE \emph{nafola}
(\citeproc{ref-RingeTaylor2014}{Ringe and Taylor 2014, 270}). The trace
represents that same development with stress-marked notation and
explicit intermediate weakening. Intervocalic \emph{b} then surfaces as
\emph{f}, and final weak-tail shortening gives \emph{nafola}.

The medial vowel is still present when A-restoration applies. That is
why the selected pre-syncope input differs from the syncopated
comparative headword.

\paragraph{Stage comparison}\label{stage-comparison}

The comparison below is manual. It separates the comparative citation
form from the pre-syncope input and from the later OE spellings.

{\def\LTcaptype{none} % do not increment counter
\begin{longtable}[]{@{}
  >{\raggedright\arraybackslash}p{(\linewidth - 8\tabcolsep) * \real{0.2000}}
  >{\raggedright\arraybackslash}p{(\linewidth - 8\tabcolsep) * \real{0.2000}}
  >{\raggedright\arraybackslash}p{(\linewidth - 8\tabcolsep) * \real{0.2000}}
  >{\raggedright\arraybackslash}p{(\linewidth - 8\tabcolsep) * \real{0.2000}}
  >{\raggedright\arraybackslash}p{(\linewidth - 8\tabcolsep) * \real{0.2000}}@{}}
\toprule\noalign{}
\begin{minipage}[b]{\linewidth}\raggedright
Formation / stage
\end{minipage} & \begin{minipage}[b]{\linewidth}\raggedright
Candidate input
\end{minipage} & \begin{minipage}[b]{\linewidth}\raggedright
Expected or documented OE outcome
\end{minipage} & \begin{minipage}[b]{\linewidth}\raggedright
OE comparison form
\end{minipage} & \begin{minipage}[b]{\linewidth}\raggedright
Result
\end{minipage} \\
\midrule\noalign{}
\endhead
\bottomrule\noalign{}
\endlastfoot
syncopated comparative headword & *nablô & reduced \emph{næfla}-type
outcome rather than \emph{nafola} & not the selected target & useful
citation form, but too reduced for the pathway modeled here \\
selected pre-syncope input & *nábulô & compact-trace output:
\emph{nafola} & nafola & exact match between selected input and
target \\
later OE reduction stages & same lexical history & attested
\emph{nafela}; Corpus \emph{nabula} & nafela / nabula & related OE
spellings, but not the chosen comparator \\
\end{longtable}
}

\subsubsection{neck --- OE hnecca}\label{neck-oe-hnecca}

Derivation: citation reconstruction \emph{*xnákkaz}; selected input
\emph{*xnékkô} \textgreater{} \emph{hnecca} (early analogy).

\paragraph{Derivation trace}\label{derivation-trace-20}

Proto input: \emph{*xnékkô}

\begingroup
\setlength{\fboxsep}{6pt}

\noindent\fbox{%
\begin{minipage}{0.97\linewidth}
\small
\begin{tabularx}{\linewidth}{@{}>{\raggedright\arraybackslash}p{0.300\linewidth}>{\centering\arraybackslash}X>{\raggedright\arraybackslash}p{0.540\linewidth}@{}}
\begin{minipage}[t]{\linewidth}
\centering\textbf{Earlier Germanic changes}\par
\vspace{0.35em}
\raggedright
\centering\textbf{West Germanic}\par
\raggedright
\vspace{0.2em}
\raggedright [no change]\par
\vspace{0.6em}
\centering\textbf{Northwest Germanic}\par
\raggedright
\vspace{0.2em}
\raggedright [no change]\par
\end{minipage}
&
&
\begin{minipage}[t]{\linewidth}
\centering\textbf{Old English changes}\par
\vspace{0.35em}
\raggedright
\begin{tabular}{@{}>{\raggedright\arraybackslash}p{0.74\linewidth}@{\hspace{0.55em}}>{\raggedright\arraybackslash}p{0.16\linewidth}@{\hspace{0.25em}}}
OE Unstressed Long Vowel Shortening & \emph{*xnékka} \\
\end{tabular}
\end{minipage}
\\
\end{tabularx}
\end{minipage}%
} \endgroup

Outcome: \emph{hnecca}

\paragraph{Reconstruction and comparative
evidence}\label{reconstruction-and-comparative-evidence-20}

The noun belongs to an ablauting n-stem family. Kroonen reconstructs a
paradigm with nominative \emph{*hnekkō}, genitive \emph{*hnukkaz}, and
accusative plural \emph{*hnakkuns}, and he places Old English
\emph{hnecca} among the e-grade descendants
(\citeproc{ref-Kroonen2011}{Kroonen 2011}). Kluge-Seebold likewise
identifies \emph{ae. hnecca} as an ablaut partner of the a-grade
\emph{Nacken} family (\citeproc{ref-KlugeSeebold2011}{Kluge 2011}).

A competing comparative label \emph{*xnákkaz} remains useful for the
wider family, and Orel also gives an a-grade headword line
(\citeproc{ref-Orel2003}{Orel 2003}). The selected input \emph{*xnékkô},
however, is the form that matches the Old English branch.

\paragraph{Old English evidence}\label{old-english-evidence-20}

Clark Hall records the weak masculine noun \emph{hnecca}
(\citeproc{ref-ClarkHall1960}{Clark Hall 1960, 162}). Bosworth-Toller
likewise records \emph{hnecca}
(\citeproc{ref-BosworthToller1898}{Bosworth and Toller 1898, 567}). The
target is therefore an attested citation form, not an oblique cell or a
reconstructed lemma.

The phonological question is upstream of the Old English evidence. The
attested noun already shows that the branch continued an e-grade form
rather than the a-grade seen in much of the continental material.

\paragraph{Development to Old
English}\label{development-to-old-english-20}

From \emph{*xnékkô}, the derivation is straightforward. The trace
shortens the final long vowel to \emph{*xnékka}, and Old English
orthography gives \emph{hnecca}.

The derivation depends on the earlier selection of the e-grade weak-noun
form continued by Old English.

\paragraph{Stem comparison}\label{stem-comparison-3}

The comparison below is manual. It separates the wider a-grade family
from the selected e-grade Old English branch.

{\def\LTcaptype{none} % do not increment counter
\begin{longtable}[]{@{}
  >{\raggedright\arraybackslash}p{(\linewidth - 8\tabcolsep) * \real{0.2000}}
  >{\raggedright\arraybackslash}p{(\linewidth - 8\tabcolsep) * \real{0.2000}}
  >{\raggedright\arraybackslash}p{(\linewidth - 8\tabcolsep) * \real{0.2000}}
  >{\raggedright\arraybackslash}p{(\linewidth - 8\tabcolsep) * \real{0.2000}}
  >{\raggedright\arraybackslash}p{(\linewidth - 8\tabcolsep) * \real{0.2000}}@{}}
\toprule\noalign{}
\begin{minipage}[b]{\linewidth}\raggedright
Formation / label
\end{minipage} & \begin{minipage}[b]{\linewidth}\raggedright
Candidate input
\end{minipage} & \begin{minipage}[b]{\linewidth}\raggedright
Expected or documented OE outcome
\end{minipage} & \begin{minipage}[b]{\linewidth}\raggedright
OE comparison form
\end{minipage} & \begin{minipage}[b]{\linewidth}\raggedright
Result
\end{minipage} \\
\midrule\noalign{}
\endhead
\bottomrule\noalign{}
\endlastfoot
competing comparative label & *xnákkaz & broader a-grade family rather
than the selected OE source & continental \emph{Nacken} line & useful
family label, but not the input followed for the Old English
derivation \\
weak noun with a-grade & *xnakkô & expected \emph{hnacca} type outcome &
hnacca & fixes the class, but not the vowel grade \\
selected e-grade nominative & *xnékkô & compact-trace output:
\emph{hnecca} & hnecca & exact match between selected input and attested
OE noun \\
oblique paradigm background & \emph{*hnukkaz}, \emph{*hnakkuns} &
ON/OHG/German a-grade continuation & hnakki / Nacken & shows the wider
ablaut family, but not the chosen OE branch \\
\end{longtable}
}

\subsubsection{needle --- OE nǣdl}\label{needle-oe-nux1e3dl}

Derivation: citation reconstruction \emph{*nḗθlō}; selected input
\emph{*nḗðlō} \textgreater{} \emph{nǣdl} (early analogy).

\paragraph{Derivation trace}\label{derivation-trace-21}

Proto input: \emph{*nḗðlō}

\begingroup
\setlength{\fboxsep}{6pt}

\noindent\fbox{%
\begin{minipage}{0.97\linewidth}
\small
\begin{tabularx}{\linewidth}{@{}>{\raggedright\arraybackslash}p{0.460\linewidth}>{\centering\arraybackslash}X>{\raggedright\arraybackslash}p{0.460\linewidth}@{}}
\begin{minipage}[t]{\linewidth}
\centering\textbf{Earlier Germanic changes}\par
\vspace{0.35em}
\raggedright
\centering\textbf{West Germanic}\par
\raggedright
\vspace{0.2em}
\begin{tabular}{@{}>{\raggedright\arraybackslash}p{0.74\linewidth}@{\hspace{0.55em}}>{\raggedright\arraybackslash}p{0.16\linewidth}@{\hspace{0.25em}}}
PWGmc Dental Hardening & \emph{*nḗdlō} \\
\end{tabular}
\vspace{0.6em}
\centering\textbf{Northwest Germanic}\par
\raggedright
\vspace{0.2em}
\begin{tabular}{@{}>{\raggedright\arraybackslash}p{0.74\linewidth}@{\hspace{0.55em}}>{\raggedright\arraybackslash}p{0.16\linewidth}@{\hspace{0.25em}}}
\mbox{NWGmc Final Long O Raising} & \emph{*nḗdlu} \\
\mbox{NWGmc Long E Lowering} & \emph{*nǣdlu} \\
\end{tabular}
\end{minipage}
&
&
\begin{minipage}[t]{\linewidth}
\centering\textbf{Old English changes}\par
\vspace{0.35em}
\raggedright
\begin{tabular}{@{}>{\raggedright\arraybackslash}p{0.74\linewidth}@{\hspace{0.55em}}>{\raggedright\arraybackslash}p{0.16\linewidth}@{\hspace{0.25em}}}
\mbox{OE High Vowel Apocope} & \emph{*nǣdl} \\
\end{tabular}
\end{minipage}
\\
\end{tabularx}
\end{minipage}%
} \endgroup

Outcome: \emph{nǣdl}

\paragraph{Reconstruction and comparative
evidence}\label{reconstruction-and-comparative-evidence-21}

Ringe and Taylor treat the word as a voiced/voiceless alternant, citing
PGmc \emph{*nēþlō}, \emph{*nēdlō-} `needle' \ldots{} \textgreater{} OE
\emph{nédl} (\citeproc{ref-RingeTaylor2014}{Ringe and Taylor 2014,
329}). The selected input \emph{*nḗðlō} is the voiced Verner-grade form
used for the Old English comparison, while the citation form
\emph{*nḗθlō} remains the broader lexeme label.

Kroonen and Orel preserve different comparative headwords for the family
(\citeproc{ref-Kroonen2013}{Kroonen 2013}; \citeproc{ref-Orel2003}{Orel
2003}). The development discussed here follows the Ringe-Taylor
alternant framework.

\paragraph{Old English evidence}\label{old-english-evidence-21}

Clark Hall records the attested citation form \emph{nǣdl}
(\citeproc{ref-ClarkHall1960}{Clark Hall 1960, 210}). Campbell lists
\emph{nédl} among the expected unbroken forms after \emph{t} and
\emph{d} (\citeproc{ref-Campbell1959}{Campbell 1959, sect. 367}). Hogg
also includes \emph{nidi} / \emph{nǣdl} in the same broader cluster
history (\citeproc{ref-Hogg1992}{Hogg 1992}).

The target is therefore an attested citation form. No oblique-cell
substitution is involved in this entry.

\paragraph{Development to Old
English}\label{development-to-old-english-21}

Ringe and Taylor give the historical line \emph{*nēþlō}, \emph{*nēdlō-}
\ldots{} \textgreater{} OE \emph{nédl}
(\citeproc{ref-RingeTaylor2014}{Ringe and Taylor 2014, 329}). Campbell
likewise lists \emph{nédl} among the expected unbroken forms after
\emph{t} and \emph{d} (\citeproc{ref-Campbell1959}{Campbell 1959, sect.
367}). The trace expresses the same pathway with the voiced alternant
selected for the Old English comparison.

The essential choice lies in the PGmc alternant selected for the
derivation. Once the voiced form is chosen, the rest of the pathway is
regular.

\paragraph{Alternant comparison}\label{alternant-comparison}

The comparison below is manual. It separates the broader citation
headword from the voiced alternant used for Old English.

{\def\LTcaptype{none} % do not increment counter
\begin{longtable}[]{@{}
  >{\raggedright\arraybackslash}p{(\linewidth - 8\tabcolsep) * \real{0.2000}}
  >{\raggedright\arraybackslash}p{(\linewidth - 8\tabcolsep) * \real{0.2000}}
  >{\raggedright\arraybackslash}p{(\linewidth - 8\tabcolsep) * \real{0.2000}}
  >{\raggedright\arraybackslash}p{(\linewidth - 8\tabcolsep) * \real{0.2000}}
  >{\raggedright\arraybackslash}p{(\linewidth - 8\tabcolsep) * \real{0.2000}}@{}}
\toprule\noalign{}
\begin{minipage}[b]{\linewidth}\raggedright
Formation / stage
\end{minipage} & \begin{minipage}[b]{\linewidth}\raggedright
Candidate input
\end{minipage} & \begin{minipage}[b]{\linewidth}\raggedright
Expected or documented OE outcome
\end{minipage} & \begin{minipage}[b]{\linewidth}\raggedright
OE comparison form
\end{minipage} & \begin{minipage}[b]{\linewidth}\raggedright
Result
\end{minipage} \\
\midrule\noalign{}
\endhead
\bottomrule\noalign{}
\endlastfoot
comparative voiceless headword & *nḗθlō & broader word-family label
rather than the OE-facing alternant & \emph{*nēþlō} line & useful
citation form, but not the selected input for the Old English
derivation \\
selected voiced Verner alternant & *nḗðlō & compact-trace output:
\emph{nǣdl} & nǣdl & exact match between selected input and attested OE
noun \\
later hardening stage & *nḗdlō & intermediate pre-OE stage in the same
derivation & nǣdl & genuine stage in the pathway, but not the selected
PGmc input \\
\end{longtable}
}

\subsubsection{nose --- OE nosu}\label{nose-oe-nosu}

Derivation: citation reconstruction \emph{*nasō}; selected input
\emph{*núsō} \textgreater{} \emph{nosu} (early analogy).

\paragraph{Derivation trace}\label{derivation-trace-22}

Proto input: \emph{*núsō}

\begingroup
\setlength{\fboxsep}{6pt}

\noindent\fbox{%
\begin{minipage}{0.97\linewidth}
\small
\begin{tabularx}{\linewidth}{@{}>{\raggedright\arraybackslash}p{0.540\linewidth}>{\centering\arraybackslash}X>{\raggedright\arraybackslash}p{0.300\linewidth}@{}}
\begin{minipage}[t]{\linewidth}
\centering\textbf{Earlier Germanic changes}\par
\vspace{0.35em}
\raggedright
\centering\textbf{West Germanic}\par
\raggedright
\vspace{0.2em}
\raggedright [no change]\par
\vspace{0.6em}
\centering\textbf{Northwest Germanic}\par
\raggedright
\vspace{0.2em}
\begin{tabular}{@{}>{\raggedright\arraybackslash}p{0.74\linewidth}@{\hspace{0.55em}}>{\raggedright\arraybackslash}p{0.16\linewidth}@{\hspace{0.25em}}}
\mbox{NWGmc U Lowering} & \emph{*nósō} \\
\mbox{NWGmc Final Long O Raising} & \emph{*nósu} \\
\end{tabular}
\end{minipage}
&
&
\begin{minipage}[t]{\linewidth}
\centering\textbf{Old English changes}\par
\vspace{0.35em}
\raggedright
\raggedright [no change]\par
\end{minipage}
\\
\end{tabularx}
\end{minipage}%
} \endgroup

Outcome: \emph{nosu}

\paragraph{Reconstruction and comparative
evidence}\label{reconstruction-and-comparative-evidence-22}

Kroonen reconstructs a Germanic ablaut pair \emph{*nasō-}
\textasciitilde{} \emph{*nusō-} and adds that the root \emph{*nus-} is
likely to have arisen as a secondary zero grade after a remodeling of
the older paradigm (\citeproc{ref-Kroonen2013}{Kroonen 2013}). Campbell
is more specific for Old English, citing \emph{nosu} \textless{}
\emph{*nusō} (\citeproc{ref-Campbell1959}{Campbell 1959, 44}).

The citation reconstruction \emph{*nasō} is therefore best treated as
the full-grade comparative headword, while the selected input
\emph{*núsō} represents the remodeled zero-grade line continued by the
Old English form discussed here. Orel's \emph{*nasō} \ldots{} OE
\emph{nasu} preserves the competing full-grade notation and shows that
the two lines should not be collapsed without comment
(\citeproc{ref-Orel2003}{Orel 2003}).

\paragraph{Old English evidence}\label{old-english-evidence-22}

\emph{Nosu} is an attested Old English noun. Ringe and Taylor list it
among the few surviving early Old English feminine u-stems
(\citeproc{ref-RingeTaylor2014}{Ringe and Taylor 2014, 385}). Clark Hall
likewise gives \emph{nosu f.}, with genitive-dative singular
\emph{nosa}, and cross-refers \emph{nasu} to \emph{nosu}
(\citeproc{ref-ClarkHall1960}{Clark Hall 1960, 810}).

The selected OE target is therefore an attested \emph{nosu}, not a
reconstructed placeholder. At the same time, the lexicographical record
keeps \emph{nasu} visible as a parallel notation belonging to the
full-grade side of the tradition.

\paragraph{Development to Old
English}\label{development-to-old-english-22}

From \emph{*núsō}, the regular path is the one documented by the current
trace: \emph{*núsō} \textgreater{} \emph{*nósō} \textgreater{}
\emph{*nósu} \textgreater{} \emph{nosu}. The early special step lies in
the choice of the zero-grade input, not in any late Old English repair.

With that input chosen, the OE development is straightforward. The
full-grade line behind \emph{*nasō} instead points toward \emph{nasu},
not to the form treated here.

\paragraph{Stem comparison}\label{stem-comparison-4}

The comparison below is manual. It separates the full-grade comparative
line from the remodeled zero-grade input that yields the Old English
form.

{\def\LTcaptype{none} % do not increment counter
\begin{longtable}[]{@{}
  >{\raggedright\arraybackslash}p{(\linewidth - 8\tabcolsep) * \real{0.2000}}
  >{\raggedright\arraybackslash}p{(\linewidth - 8\tabcolsep) * \real{0.2000}}
  >{\raggedright\arraybackslash}p{(\linewidth - 8\tabcolsep) * \real{0.2000}}
  >{\raggedright\arraybackslash}p{(\linewidth - 8\tabcolsep) * \real{0.2000}}
  >{\raggedright\arraybackslash}p{(\linewidth - 8\tabcolsep) * \real{0.2000}}@{}}
\toprule\noalign{}
\begin{minipage}[b]{\linewidth}\raggedright
Formation / label
\end{minipage} & \begin{minipage}[b]{\linewidth}\raggedright
Candidate input
\end{minipage} & \begin{minipage}[b]{\linewidth}\raggedright
Expected or documented OE outcome
\end{minipage} & \begin{minipage}[b]{\linewidth}\raggedright
OE comparison form
\end{minipage} & \begin{minipage}[b]{\linewidth}\raggedright
Result
\end{minipage} \\
\midrule\noalign{}
\endhead
\bottomrule\noalign{}
\endlastfoot
full-grade comparative line & *nasō & expected full-grade continuation
\emph{nasu} & nasu & useful comparative background, but not the selected
OE-facing input \\
remodeled zero-grade line & *núsō & compact-trace output: \emph{nosu} &
nosu & exact match between selected input and attested OE noun \\
\end{longtable}
}

\subsubsection{sap --- OE sæp}\label{sap-oe-suxe6p}

Derivation: citation reconstruction \emph{*sapōn}; selected input
\emph{*sápą} \textgreater{} \emph{sæp} (early analogy).

\paragraph{Derivation trace}\label{derivation-trace-23}

Proto input: \emph{*sápą}

\begingroup
\setlength{\fboxsep}{6pt}

\noindent\fbox{%
\begin{minipage}{0.97\linewidth}
\small
\begin{tabularx}{\linewidth}{@{}>{\raggedright\arraybackslash}p{0.300\linewidth}>{\centering\arraybackslash}X>{\raggedright\arraybackslash}p{0.540\linewidth}@{}}
\begin{minipage}[t]{\linewidth}
\centering\textbf{Earlier Germanic changes}\par
\vspace{0.35em}
\raggedright
\centering\textbf{West Germanic}\par
\raggedright
\vspace{0.2em}
\raggedright [no change]\par
\vspace{0.6em}
\centering\textbf{Northwest Germanic}\par
\raggedright
\vspace{0.2em}
\raggedright [no change]\par
\end{minipage}
&
&
\begin{minipage}[t]{\linewidth}
\centering\textbf{Old English changes}\par
\vspace{0.35em}
\raggedright
\begin{tabular}{@{}>{\raggedright\arraybackslash}p{0.74\linewidth}@{\hspace{0.55em}}>{\raggedright\arraybackslash}p{0.16\linewidth}@{\hspace{0.25em}}}
Anglo Frisian Brightening & \emph{*sæpą} \\
OE Heavy Syllable Nasal Apocope & \emph{*sæp} \\
\end{tabular}
\end{minipage}
\\
\end{tabularx}
\end{minipage}%
} \endgroup

Outcome: \emph{sæp}

\paragraph{Reconstruction and comparative
evidence}\label{reconstruction-and-comparative-evidence-23}

The comparative sources do not give one uniform inherited stem. Kroonen
derives the word family from material pointing to dialectal dissolution
of a primary n-stem \emph{*safō}, gen. \emph{*sappaz}
(\citeproc{ref-Kroonen2013}{Kroonen 2013}). Orel preserves the
comparative notation \emph{*sapōn} \textasciitilde{} \emph{*sapan}
(\citeproc{ref-Orel2003}{Orel 2003, 319}), while Kluge-Seebold instead
gives West Germanic \emph{*sapi-} and still cites Old English \emph{sæp}
n. (\citeproc{ref-KlugeSeebold2011}{Kluge 2011}).

The selected input \emph{*sápą} therefore does not replace those
comparative labels. It identifies the OE-facing stem shape that yields
the attested noun treated here.

\paragraph{Old English evidence}\label{old-english-evidence-23}

Clark Hall records \emph{sæp} (e) n. (\citeproc{ref-ClarkHall1960}{Clark
Hall 1960, 247}), and Kluge-Seebold likewise cites ae. \emph{sæp} n.
(\citeproc{ref-KlugeSeebold2011}{Kluge 2011}). The target is therefore
an attested neuter Old English noun. Orel's plain \emph{sap} notation
belongs to comparative normalization, not to the spelling adopted here
for the Old English form (\citeproc{ref-Orel2003}{Orel 2003, 319}).

\paragraph{Development to Old
English}\label{development-to-old-english-23}

From \emph{*sápą}, Anglo-Frisian brightening yields \emph{sæ}, and
heavy-syllable nasal apocope then produces \emph{sæp}. That is the
regular path documented by the current trace.

The competing comparative lines do not give the same result. The
inherited n-stem notation \emph{*sapōn} yields \emph{sape}, while an
i-stem continuation from the \emph{*sapi-} line leads to \emph{sep} /
\emph{sepe} rather than to \emph{sæp}. The special step in this entry is
therefore the early stem choice, not a late OE paradigm-cell selection.

\paragraph{Stem comparison}\label{stem-comparison-5}

The comparison below is manual. It separates the competing comparative
stem lines from the selected OE-facing input.

{\def\LTcaptype{none} % do not increment counter
\begin{longtable}[]{@{}
  >{\raggedright\arraybackslash}p{(\linewidth - 8\tabcolsep) * \real{0.2000}}
  >{\raggedright\arraybackslash}p{(\linewidth - 8\tabcolsep) * \real{0.2000}}
  >{\raggedright\arraybackslash}p{(\linewidth - 8\tabcolsep) * \real{0.2000}}
  >{\raggedright\arraybackslash}p{(\linewidth - 8\tabcolsep) * \real{0.2000}}
  >{\raggedright\arraybackslash}p{(\linewidth - 8\tabcolsep) * \real{0.2000}}@{}}
\toprule\noalign{}
\begin{minipage}[b]{\linewidth}\raggedright
Formation / label
\end{minipage} & \begin{minipage}[b]{\linewidth}\raggedright
Candidate input
\end{minipage} & \begin{minipage}[b]{\linewidth}\raggedright
Expected or documented OE outcome
\end{minipage} & \begin{minipage}[b]{\linewidth}\raggedright
OE comparison form
\end{minipage} & \begin{minipage}[b]{\linewidth}\raggedright
Result
\end{minipage} \\
\midrule\noalign{}
\endhead
\bottomrule\noalign{}
\endlastfoot
comparative n-stem line & *sapōn & local comparator output: \emph{sape}
& sape & useful comparative background, but not the source of attested
\emph{sæp} \\
inferred i-stem comparator from \emph{*sapi-} & *sapiz & local
comparator output: \emph{sepe} & sepe & confirms that an i-triggering
stem does not reach the target \\
selected a-stem input & *sápą & compact-trace output: \emph{sæp} & sæp &
exact match between selected input and attested OE noun \\
\end{longtable}
}

\subsubsection{sea --- OE sǣ}\label{sea-oe-sux1e3}

Derivation: citation reconstruction \emph{*sái}; selected input
\emph{*sáiwiz} \textgreater{} \emph{sǣ} (early analogy).

\paragraph{Derivation trace}\label{derivation-trace-24}

Proto input: \emph{*sáiwiz}

\begingroup
\setlength{\fboxsep}{6pt}

\noindent\fbox{%
\begin{minipage}{0.97\linewidth}
\small
\begin{tabularx}{\linewidth}{@{}>{\raggedright\arraybackslash}p{0.460\linewidth}>{\centering\arraybackslash}X>{\raggedright\arraybackslash}p{0.460\linewidth}@{}}
\begin{minipage}[t]{\linewidth}
\centering\textbf{Earlier Germanic changes}\par
\vspace{0.35em}
\raggedright
\centering\textbf{West Germanic}\par
\raggedright
\vspace{0.2em}
\begin{tabular}{@{}>{\raggedright\arraybackslash}p{0.74\linewidth}@{\hspace{0.55em}}>{\raggedright\arraybackslash}p{0.16\linewidth}@{\hspace{0.25em}}}
PWGmc Ai Monophthongization & \emph{*sāwiz} \\
\end{tabular}
\vspace{0.6em}
\centering\textbf{Northwest Germanic}\par
\raggedright
\vspace{0.2em}
\begin{tabular}{@{}>{\raggedright\arraybackslash}p{0.74\linewidth}@{\hspace{0.55em}}>{\raggedright\arraybackslash}p{0.16\linewidth}@{\hspace{0.25em}}}
\mbox{PGmc Final Z Deletion} & \emph{*sāwi} \\
\end{tabular}
\end{minipage}
&
&
\begin{minipage}[t]{\linewidth}
\centering\textbf{Old English changes}\par
\vspace{0.35em}
\raggedright
\begin{tabular}{@{}>{\raggedright\arraybackslash}p{0.74\linewidth}@{\hspace{0.55em}}>{\raggedright\arraybackslash}p{0.16\linewidth}@{\hspace{0.25em}}}
OE W Loss Before I & \emph{*sāi} \\
OE I Umlaut & \emph{*sǣi} \\
\mbox{OE High Vowel Apocope} & \emph{*sǣ} \\
\end{tabular}
\end{minipage}
\\
\end{tabularx}
\end{minipage}%
} \endgroup

Outcome: \emph{sǣ}

\paragraph{Reconstruction and comparative
evidence}\label{reconstruction-and-comparative-evidence-24}

Kroonen gives the noun in stem notation as \emph{*saiwi-}, an i-stem
whose English reflex is cited as OE \emph{sæ}
(\citeproc{ref-Kroonen2013}{Kroonen 2013}). Ringe and Taylor write the
fuller form \emph{*saiwiz} and derive it through \emph{*sawi}
\textgreater{} \emph{*sei} \textgreater{} OE \emph{sǣ}
(\citeproc{ref-RingeTaylor2014}{Ringe and Taylor 2014, sect. 6.7.1}).
The comparative headword is therefore shorter than the form required for
the English history: \emph{*sái} names the lexeme, but \emph{*sáiwiz}
preserves the medial \emph{*w} and the final high vowel that control the
later development.

\paragraph{Old English evidence}\label{old-english-evidence-24}

The Old English noun is the ordinary word for `sea'. Kroonen cites it as
\emph{sæ}; the normalized form here is \emph{sǣ}
(\citeproc{ref-Kroonen2013}{Kroonen 2013}). Campbell likewise treats
\emph{sea} as continuing the same \emph{*saiui-} \textgreater{}
\emph{*sǣi} history, with loss of \emph{u}/\emph{w} before \emph{i}
(\citeproc{ref-Campbell1959}{Campbell 1959, sect. 406}).

\paragraph{Development to Old
English}\label{development-to-old-english-24}

Once the fuller i-stem input is chosen, the development is regular.
After Proto-West-Germanic monophthongization \emph{*sáiwiz}
\textgreater{} \emph{*sāwiz} and final \emph{*-z} loss \emph{*sāwiz}
\textgreater{} \emph{*sāwi}, the non-initial \emph{*w} disappears before
unstressed \emph{*i}, and the following high vowel fronts the root vowel
before final apocope. The documented chain is \emph{*sáiwiz}
\textgreater{} \emph{*sāwiz} \textgreater{} \emph{*sāwi} \textgreater{}
\emph{*sāi} \textgreater{} \emph{*sǣi} \textgreater{} \emph{sǣ}
(\citeproc{ref-RingeTaylor2014}{Ringe and Taylor 2014, sect. 6.7.1};
\citeproc{ref-Campbell1959}{Campbell 1959, sect. 406}).

\paragraph{Stem and stage comparison}\label{stem-and-stage-comparison}

The comparison below is manual. It separates the abbreviated comparative
headword from the fuller i-stem input that yields the Old English form.

{\def\LTcaptype{none} % do not increment counter
\begin{longtable}[]{@{}
  >{\raggedright\arraybackslash}p{(\linewidth - 8\tabcolsep) * \real{0.2000}}
  >{\raggedright\arraybackslash}p{(\linewidth - 8\tabcolsep) * \real{0.2000}}
  >{\raggedright\arraybackslash}p{(\linewidth - 8\tabcolsep) * \real{0.2000}}
  >{\raggedright\arraybackslash}p{(\linewidth - 8\tabcolsep) * \real{0.2000}}
  >{\raggedright\arraybackslash}p{(\linewidth - 8\tabcolsep) * \real{0.2000}}@{}}
\toprule\noalign{}
\begin{minipage}[b]{\linewidth}\raggedright
Formation / label
\end{minipage} & \begin{minipage}[b]{\linewidth}\raggedright
Candidate input
\end{minipage} & \begin{minipage}[b]{\linewidth}\raggedright
OE output or comparison
\end{minipage} & \begin{minipage}[b]{\linewidth}\raggedright
OE comparison form
\end{minipage} & \begin{minipage}[b]{\linewidth}\raggedright
Result
\end{minipage} \\
\midrule\noalign{}
\endhead
\bottomrule\noalign{}
\endlastfoot
abbreviated comparative headword & *sái & too short to preserve the
\emph{*w} \ldots{} \emph{*i} environment needed for the documented
chronology & sǣ & useful comparative label, but not the selected
OE-facing input \\
selected i-stem input & *sáiwiz & documented trace output: \emph{sǣ} &
sǣ & exact match between selected input and Old English target \\
\end{longtable}
}

\subsubsection{sieve --- OE sife}\label{sieve-oe-sife}

Derivation: citation reconstruction \emph{*síbaz}; selected input
\emph{*síbi} \textgreater{} \emph{sife} (early analogy).

\paragraph{Derivation trace}\label{derivation-trace-25}

Proto input: \emph{*síbi}

\begingroup
\setlength{\fboxsep}{6pt}

\noindent\fbox{%
\begin{minipage}{0.97\linewidth}
\small
\begin{tabularx}{\linewidth}{@{}>{\raggedright\arraybackslash}p{0.300\linewidth}>{\centering\arraybackslash}X>{\raggedright\arraybackslash}p{0.540\linewidth}@{}}
\begin{minipage}[t]{\linewidth}
\centering\textbf{Earlier Germanic changes}\par
\vspace{0.35em}
\raggedright
\centering\textbf{West Germanic}\par
\raggedright
\vspace{0.2em}
\raggedright [no change]\par
\vspace{0.6em}
\centering\textbf{Northwest Germanic}\par
\raggedright
\vspace{0.2em}
\raggedright [no change]\par
\end{minipage}
&
&
\begin{minipage}[t]{\linewidth}
\centering\textbf{Old English changes}\par
\vspace{0.35em}
\raggedright
\begin{tabular}{@{}>{\raggedright\arraybackslash}p{0.74\linewidth}@{\hspace{0.55em}}>{\raggedright\arraybackslash}p{0.16\linewidth}@{\hspace{0.25em}}}
PGmc B Allophony & \emph{*síβi} \\
OE Med Unstressed I Lowering1 & \emph{*síβe} \\
\end{tabular}
\end{minipage}
\\
\end{tabularx}
\end{minipage}%
} \endgroup

Outcome: \emph{sife}

\paragraph{Reconstruction and comparative
evidence}\label{reconstruction-and-comparative-evidence-25}

Kluge-Seebold gives wg. \emph{*sibi-} n.~\ldots{} ae. sife, and Campbell
groups \emph{sife} with short neuter i-stems such as \emph{spere}
(\citeproc{ref-KlugeSeebold2011}{Kluge 2011};
\citeproc{ref-Campbell1959}{Campbell 1959, sect. 609}). The older
morphological background is the s-stem \emph{*sib-iz}, but the selected
input is the normalized i-stem form \emph{*síbi}.

Kroonen's relevant nearby entry is not `sieve' but \emph{*sebjō-}
`kinship', the source of Old English \emph{sibb}
(\citeproc{ref-Kroonen2013}{Kroonen 2013}). Orel's \emph{*sibaz}
\ldots{} OE \emph{sife} preserves a broader handbook notation, but that
a-stem shape does not fit the Old English form treated here
(\citeproc{ref-Orel2003}{Orel 2003}).

\paragraph{Old English evidence}\label{old-english-evidence-25}

Clark Hall gives \emph{sibi (GL) \ldots{} = sife} and also \emph{sife
n.~`sieve'} (\citeproc{ref-ClarkHall1960}{Clark Hall 1960}). Campbell
likewise cites Corpus Glossary \emph{sibi} and treats \emph{sife} as a
short neuter i-stem (\citeproc{ref-Campbell1959}{Campbell 1959, sects
444, 609}). The normalized Old English target is therefore \emph{sife},
while \emph{sibi} is an attested earlier spelling rather than a separate
lexeme.

\paragraph{Development to Old
English}\label{development-to-old-english-25}

From \emph{*síbi}, the documented trace gives \emph{*síβi}
\textgreater{} \emph{*síβe} \textgreater{} \emph{sife}. Medial \emph{b}
is realized as a spirant and later written \emph{f}, while the final
unstressed \emph{i} lowers to \emph{e}. The older s-stem background
\emph{*sib-iz} explains the morphology, but the selected input
\emph{*síbi} is the immediate pre-Old-English form.

\paragraph{Stem comparison}\label{stem-comparison-6}

The comparison below is manual. It distinguishes the accepted i-stem
line from its rejected competitors.

{\def\LTcaptype{none} % do not increment counter
\begin{longtable}[]{@{}
  >{\raggedright\arraybackslash}p{(\linewidth - 8\tabcolsep) * \real{0.2000}}
  >{\raggedright\arraybackslash}p{(\linewidth - 8\tabcolsep) * \real{0.2000}}
  >{\raggedright\arraybackslash}p{(\linewidth - 8\tabcolsep) * \real{0.2000}}
  >{\raggedright\arraybackslash}p{(\linewidth - 8\tabcolsep) * \real{0.2000}}
  >{\raggedright\arraybackslash}p{(\linewidth - 8\tabcolsep) * \real{0.2000}}@{}}
\toprule\noalign{}
\begin{minipage}[b]{\linewidth}\raggedright
Formation / label
\end{minipage} & \begin{minipage}[b]{\linewidth}\raggedright
Candidate input
\end{minipage} & \begin{minipage}[b]{\linewidth}\raggedright
Expected or documented OE outcome
\end{minipage} & \begin{minipage}[b]{\linewidth}\raggedright
OE comparison form
\end{minipage} & \begin{minipage}[b]{\linewidth}\raggedright
Result
\end{minipage} \\
\midrule\noalign{}
\endhead
\bottomrule\noalign{}
\endlastfoot
ja-stem kinship line & Kroonen \emph{*sebjō-} / comparator \emph{*sibja}
& OE \emph{sibb} & sibb & separate lexeme, not the target treated
here \\
a-stem handbook line & *síbaz & expected \emph{sif} & sif & wrong ending
for the attested noun \\
selected i-stem line from older \emph{*sib-iz} & *síbi & documented
trace output: \emph{sife} & sife; early spelling \emph{sibi} & exact
match between selected input and Old English evidence \\
\end{longtable}
}

\subsubsection{spare --- OE sparian}\label{spare-oe-sparian}

Derivation: citation reconstruction \emph{*sparēną}; selected input
\emph{*spárōjaną} \textgreater{} \emph{sparian} (early analogy).

\paragraph{Derivation trace}\label{derivation-trace-26}

Proto input: \emph{*spárōjaną}

\begingroup
\setlength{\fboxsep}{6pt}

\noindent\fbox{%
\begin{minipage}{0.97\linewidth}
\footnotesize
\begin{tabularx}{\linewidth}{@{}>{\raggedright\arraybackslash}p{0.280\linewidth}>{\centering\arraybackslash}X>{\raggedright\arraybackslash}p{0.560\linewidth}@{}}
\begin{minipage}[t]{\linewidth}
\centering\textbf{Earlier Germanic changes}\par
\vspace{0.35em}
\raggedright
\centering\textbf{West Germanic}\par
\raggedright
\vspace{0.2em}
\raggedright [no change]\par
\vspace{0.6em}
\centering\textbf{Northwest Germanic}\par
\raggedright
\vspace{0.2em}
\raggedright [no change]\par
\end{minipage}
&
&
\begin{minipage}[t]{\linewidth}
\centering\textbf{Old English changes}\par
\vspace{0.35em}
\raggedright
\begin{tabular}{@{}>{\raggedright\arraybackslash}p{0.64\linewidth}@{\hspace{0.55em}}>{\raggedright\arraybackslash}p{0.26\linewidth}@{\hspace{0.25em}}}
Anglo Frisian Brightening & \emph{*spærōjaną} \\
OE A Restoration & \emph{*sparōjaną} \\
OE Heavy Syllable Nasal Apocope & \emph{*sparōjan} \\
OE Secondary Nasalization & \emph{*sparōjąn} \\
OE I Umlaut & \emph{*sparējąn} \\
OE Unstressed Long Vowel Shortening & \emph{*sparejąn} \\
OE Weak Tail Reduction & \emph{*sparejan} \\
OE Intervocalic J Vocalization & \emph{*spareian} \\
OE Unstressed EI Contraction & \emph{*sparian} \\
\end{tabular}
\end{minipage}
\\
\end{tabularx}
\end{minipage}%
} \endgroup

Outcome: \emph{sparian}

\paragraph{Reconstruction and comparative
evidence}\label{reconstruction-and-comparative-evidence-26}

Kroonen and Orel keep the inherited verb under class-III \emph{*sparēn-}
/ \emph{*sparēnan} (\citeproc{ref-Kroonen2013}{Kroonen 2013};
\citeproc{ref-Orel2003}{Orel 2003}). Ringe and Taylor, however,
reconstruct \emph{*sparai-} \textasciitilde{} \emph{*sparja-} for the
English branch and derive the citation verb from a class-II line
(\citeproc{ref-RingeTaylor2014}{Ringe and Taylor 2014, 162, 191}). The
selected input \emph{*spárōjaną} therefore represents the refashioned
class-II formation behind Old English \emph{sparian}, while the citation
reconstruction \emph{*sparēną} remains the inherited comparative
headword.

\paragraph{Old English evidence}\label{old-english-evidence-26}

Campbell says that \emph{sparian} does not show the ordinary class-III
characteristics, but the Ritual forms, normalized here as \emph{spæria},
\emph{spær}, and \emph{spærede}, together with Vespasian Psalter
\emph{spearad}, point to primitive Old English forms both with and
without back vowels (\citeproc{ref-Campbell1959}{Campbell 1959, sect.
764}). Brunner likewise records Northumbrian \emph{spæria},
\emph{spærede} beside common Old English \emph{sparian} and Vespasian
Psalter \emph{spearad} (\citeproc{ref-SieversBrunner1965}{Sievers 1965,
sect. 364} Anm. 11). The citation form treated here is \emph{sparian};
the Anglian forms are relics of the older formation, not alternative
headwords of equal status.

\paragraph{Development to Old
English}\label{development-to-old-english-26}

Once the class-II formation \emph{*spárōjaną} is chosen, the remaining
development is regular. The documented trace shows brightening,
restoration of \emph{a} before the back vocalism of the suffix, later
i-mutation within the weak ending, weak-tail reduction, and contraction
to \emph{sparian}. By contrast, Brunner's rule against further apocope
of final \emph{-e} explains why Ritual \emph{spær} cannot be the regular
continuation of inherited \emph{*spárē}
(\citeproc{ref-SieversBrunner1965}{Sievers 1965, sect. 150}).

\paragraph{Formation comparison}\label{formation-comparison-4}

The comparison below is manual. It contrasts the inherited class-III
formation with the refashioned class-II one that yields the citation
verb.

{\def\LTcaptype{none} % do not increment counter
\begin{longtable}[]{@{}
  >{\raggedright\arraybackslash}p{(\linewidth - 8\tabcolsep) * \real{0.2000}}
  >{\raggedright\arraybackslash}p{(\linewidth - 8\tabcolsep) * \real{0.2000}}
  >{\raggedright\arraybackslash}p{(\linewidth - 8\tabcolsep) * \real{0.2000}}
  >{\raggedright\arraybackslash}p{(\linewidth - 8\tabcolsep) * \real{0.2000}}
  >{\raggedright\arraybackslash}p{(\linewidth - 8\tabcolsep) * \real{0.2000}}@{}}
\toprule\noalign{}
\begin{minipage}[b]{\linewidth}\raggedright
Formation / comparison
\end{minipage} & \begin{minipage}[b]{\linewidth}\raggedright
Candidate input
\end{minipage} & \begin{minipage}[b]{\linewidth}\raggedright
Expected or documented OE outcome
\end{minipage} & \begin{minipage}[b]{\linewidth}\raggedright
OE comparison form
\end{minipage} & \begin{minipage}[b]{\linewidth}\raggedright
Result
\end{minipage} \\
\midrule\noalign{}
\endhead
\bottomrule\noalign{}
\endlastfoot
inherited class-III infinitive & *spárēną & manual comparison / probe
output: \emph{sparen} & sparian & wrong class and wrong ending for the
citation verb \\
inherited class-III imperative singular & *spárē & manual comparison /
probe output: \emph{spære} & Ritual \emph{spær} & loss of final
\emph{-e} is not regular, so the relic form cannot control the entry \\
inherited class-III finite present & *spárēθi & manual comparison /
probe output: \emph{spæreþ} & \emph{spearad} & attested form is mixed,
not a direct continuation of the inherited cell \\
selected class-II formation & *spárōjaną & documented trace output:
\emph{sparian} & sparian & exact match between selected input and Old
English citation form \\
\end{longtable}
}

\subsubsection{staff --- OE stæf}\label{staff-oe-stuxe6f}

Derivation: citation reconstruction \emph{*stábiz}; selected input
\emph{*stábaz} \textgreater{} \emph{stæf} (early analogy).

\paragraph{Derivation trace}\label{derivation-trace-27}

Proto input: \emph{*stábaz}

\begingroup
\setlength{\fboxsep}{6pt}

\noindent\fbox{%
\begin{minipage}{0.97\linewidth}
\small
\begin{tabularx}{\linewidth}{@{}>{\raggedright\arraybackslash}p{0.300\linewidth}>{\centering\arraybackslash}X>{\raggedright\arraybackslash}p{0.540\linewidth}@{}}
\begin{minipage}[t]{\linewidth}
\centering\textbf{Earlier Germanic changes}\par
\vspace{0.35em}
\raggedright
\centering\textbf{West Germanic}\par
\raggedright
\vspace{0.2em}
\raggedright [no change]\par
\vspace{0.6em}
\centering\textbf{Northwest Germanic}\par
\raggedright
\vspace{0.2em}
\begin{tabular}{@{}>{\raggedright\arraybackslash}p{0.74\linewidth}@{\hspace{0.55em}}>{\raggedright\arraybackslash}p{0.16\linewidth}@{\hspace{0.25em}}}
\mbox{PGmc Final Z Deletion} & \emph{*stába} \\
\end{tabular}
\end{minipage}
&
&
\begin{minipage}[t]{\linewidth}
\centering\textbf{Old English changes}\par
\vspace{0.35em}
\raggedright
\begin{tabular}{@{}>{\raggedright\arraybackslash}p{0.70\linewidth}@{\hspace{0.55em}}>{\raggedright\arraybackslash}p{0.16\linewidth}@{\hspace{0.25em}}}
PWGmc Final Bare A Loss & \emph{*stáb} \\
Anglo Frisian Brightening & \emph{*stæb} \\
PGmc B Allophony & \emph{*stæβ} \\
\end{tabular}
\end{minipage}
\\
\end{tabularx}
\end{minipage}%
} \endgroup

Outcome: \emph{stæf}

\paragraph{Reconstruction and comparative
evidence}\label{reconstruction-and-comparative-evidence-27}

The comparative dictionaries do not give one uniform stem class. Kroonen
reconstructs an a-stem \emph{*staba-}
(\citeproc{ref-Kroonen2013}{Kroonen 2013}), Orel writes \emph{*stabiz}
\textasciitilde{} \emph{*stabaz} (\citeproc{ref-Orel2003}{Orel 2003}),
and Kluge-Seebold explicitly marks g. \emph{*stabi-}/a-
(\citeproc{ref-KlugeSeebold2011}{Kluge 2011}). That disagreement matters
because a direct i-stem input in \emph{*-iz} would predict i-mutation in
Old English, whereas the attested noun keeps \emph{æ}.

\paragraph{Old English evidence}\label{old-english-evidence-27}

The Old English noun itself is the ordinary citation form \emph{stæf}.
Luick lists \emph{stæf} among closed monosyllables with \emph{æ}
(\citeproc{ref-Luick1914}{Luick 1914-\/-1940, 176}), and Ringe and
Taylor pair singular \emph{stæf} with plural \emph{stafas}
(\citeproc{ref-RingeTaylor2014}{Ringe and Taylor 2014}). The normalized
form here is therefore \emph{stæf}; later English \emph{staff} with
\emph{a} belongs to a later stage of the word's history.

\paragraph{Development to Old
English}\label{development-to-old-english-27}

With the selected a-stem input, the development is regular. Final
\emph{*-z} disappears, bare final \emph{-a} is lost, Anglo-Frisian
brightening gives \emph{æ} in the closed monosyllable, and medial
\emph{b} surfaces as a fricative written \emph{f}. The documented chain
is \emph{*stábaz} \textgreater{} \emph{*stába} \textgreater{}
\emph{*stáb} \textgreater{} \emph{*stæb} \textgreater{} \emph{stæf}. A
direct continuation of \emph{*stábiz}, by contrast, would produce
i-mutated \emph{stefe} rather than the attested singular.

\paragraph{Formation comparison}\label{formation-comparison-5}

The comparison below is manual. It separates the rejected i-stem line
from the selected a-stem input.

{\def\LTcaptype{none} % do not increment counter
\begin{longtable}[]{@{}
  >{\raggedright\arraybackslash}p{(\linewidth - 8\tabcolsep) * \real{0.2000}}
  >{\raggedright\arraybackslash}p{(\linewidth - 8\tabcolsep) * \real{0.2000}}
  >{\raggedright\arraybackslash}p{(\linewidth - 8\tabcolsep) * \real{0.2000}}
  >{\raggedright\arraybackslash}p{(\linewidth - 8\tabcolsep) * \real{0.2000}}
  >{\raggedright\arraybackslash}p{(\linewidth - 8\tabcolsep) * \real{0.2000}}@{}}
\toprule\noalign{}
\begin{minipage}[b]{\linewidth}\raggedright
Formation / label
\end{minipage} & \begin{minipage}[b]{\linewidth}\raggedright
Candidate input
\end{minipage} & \begin{minipage}[b]{\linewidth}\raggedright
OE output or comparison
\end{minipage} & \begin{minipage}[b]{\linewidth}\raggedright
OE comparison form
\end{minipage} & \begin{minipage}[b]{\linewidth}\raggedright
Result
\end{minipage} \\
\midrule\noalign{}
\endhead
\bottomrule\noalign{}
\endlastfoot
comparative i-stem line & *stábiz & expected \emph{stefe} after
i-mutation & stæf & wrong vowel for the attested singular \\
mixed comparative notation & Orel \emph{*stabiz} \textasciitilde{}
\emph{*stabaz}; Kluge \emph{*stabi-}/a- & source-level stem-class
uncertainty & stæf & useful comparative background, but not a single
OE-facing input \\
selected a-stem input & *stábaz & documented trace output: \emph{stæf} &
stæf & exact match between selected input and Old English target \\
\end{longtable}
}

\subsubsection{stem --- OE stefn}\label{stem-oe-stefn}

Derivation: citation reconstruction \emph{*stámnaz}; selected input
\emph{*stébnō} \textgreater{} \emph{stefn} (early analogy).

\paragraph{Derivation trace}\label{derivation-trace-28}

Proto input: \emph{*stébnō}

\begingroup
\setlength{\fboxsep}{6pt}

\noindent\fbox{%
\begin{minipage}{0.97\linewidth}
\small
\begin{tabularx}{\linewidth}{@{}>{\raggedright\arraybackslash}p{0.460\linewidth}>{\centering\arraybackslash}X>{\raggedright\arraybackslash}p{0.460\linewidth}@{}}
\begin{minipage}[t]{\linewidth}
\centering\textbf{Earlier Germanic changes}\par
\vspace{0.35em}
\raggedright
\centering\textbf{West Germanic}\par
\raggedright
\vspace{0.2em}
\raggedright [no change]\par
\vspace{0.6em}
\centering\textbf{Northwest Germanic}\par
\raggedright
\vspace{0.2em}
\begin{tabular}{@{}>{\raggedright\arraybackslash}p{0.74\linewidth}@{\hspace{0.55em}}>{\raggedright\arraybackslash}p{0.16\linewidth}@{\hspace{0.25em}}}
\mbox{NWGmc Final Long O Raising} & \emph{*stébnu} \\
\end{tabular}
\end{minipage}
&
&
\begin{minipage}[t]{\linewidth}
\centering\textbf{Old English changes}\par
\vspace{0.35em}
\raggedright
\begin{tabular}{@{}>{\raggedright\arraybackslash}p{0.74\linewidth}@{\hspace{0.55em}}>{\raggedright\arraybackslash}p{0.16\linewidth}@{\hspace{0.25em}}}
PGmc B Allophony & \emph{*stéβnu} \\
\mbox{OE High Vowel Apocope} & \emph{*stéβn} \\
\end{tabular}
\end{minipage}
\\
\end{tabularx}
\end{minipage}%
} \endgroup

Outcome: \emph{stefn}

\paragraph{Reconstruction and comparative
evidence}\label{reconstruction-and-comparative-evidence-28}

The source tradition behind \emph{stefn} is not the same as the
comparative label \emph{*stámnaz}. Ringe and Taylor cite \emph{*stebnō}
for the noun continued by Gothic \emph{stibna} and Old English
\emph{stebn} \textgreater{} \emph{stefn} \textgreater{} \emph{stemn}
(\citeproc{ref-RingeTaylor2014}{Ringe and Taylor 2014, 330}). Orel
likewise gives \emph{*stebnō} \textasciitilde{} \emph{*stemnō}, whereas
Kroonen prefers \emph{*stimnō-}, and Fulk describes the etymology of
\emph{stefn, stemn} as insecure (\citeproc{ref-Orel2003}{Orel 2003,
374}; \citeproc{ref-Kroonen2013}{Kroonen 2013, 480};
\citeproc{ref-Fulk2018}{Fulk 2018, sect. 6.11} n.~6).

These forms belong to the Old English noun \emph{stefn} `voice, sound'.
The selected input \emph{*stébnō} is therefore best treated as the
OE-facing transponent supported by that source tradition. It does not
settle the deeper comparative reconstruction implied by the citation
label \emph{*stámnaz}.

\paragraph{Old English evidence}\label{old-english-evidence-28}

Clark Hall records \emph{stefn} as the noun `voice, sound' and
cross-refers \emph{stemn} to the same word
(\citeproc{ref-ClarkHall1960}{Clark Hall 1960, 276}). Ringe and Taylor
give the OE chronology directly as \emph{stebn} \textgreater{}
\emph{stefn} \textgreater{} \emph{stemn}
(\citeproc{ref-RingeTaylor2014}{Ringe and Taylor 2014, 330}).

Bülbring and Luick treat \emph{stemn} as a later West Saxon development
from older \emph{stefn}, produced by \emph{fn} \textgreater{} \emph{mn}
only after the earlier period of nasal influence on \emph{e}
(\citeproc{ref-Bulbring1902}{Bülbring 1902, sect. 62} Anm. 3, 445;
\citeproc{ref-Luick1914}{Luick 1914-\/-1940, sect. 75} Anm. 1). The
relevant comparison form is therefore the conservative \emph{stefn}, not
the later West Saxon doublet \emph{stemn}.

\paragraph{Development to Old
English}\label{development-to-old-english-28}

From \emph{*stébnō}, raising of final long \emph{ō} gives a
\emph{*stébnu} stage. Regular fricativization of \emph{b} before
\emph{n} then yields \emph{*stéβnu}, and loss of the final high vowel
leaves \emph{*stéβn}, written \emph{stefn} in Old English. The later
form \emph{stemn} belongs to a separate West Saxon assimilation after
this stage (\citeproc{ref-RingeTaylor2014}{Ringe and Taylor 2014, 330};
\citeproc{ref-Bulbring1902}{Bülbring 1902, sect. 445}).

\paragraph{Source comparison}\label{source-comparison}

The comparison below is manual. It keeps apart the broader comparative
label, the OE-facing transponent, and the later West Saxon variant
history.

{\def\LTcaptype{none} % do not increment counter
\begin{longtable}[]{@{}
  >{\raggedright\arraybackslash}p{(\linewidth - 6\tabcolsep) * \real{0.2500}}
  >{\raggedright\arraybackslash}p{(\linewidth - 6\tabcolsep) * \real{0.2500}}
  >{\raggedright\arraybackslash}p{(\linewidth - 6\tabcolsep) * \real{0.2500}}
  >{\raggedright\arraybackslash}p{(\linewidth - 6\tabcolsep) * \real{0.2500}}@{}}
\toprule\noalign{}
\begin{minipage}[b]{\linewidth}\raggedright
Form or label
\end{minipage} & \begin{minipage}[b]{\linewidth}\raggedright
Status
\end{minipage} & \begin{minipage}[b]{\linewidth}\raggedright
OE relation
\end{minipage} & \begin{minipage}[b]{\linewidth}\raggedright
Result
\end{minipage} \\
\midrule\noalign{}
\endhead
\bottomrule\noalign{}
\endlastfoot
\emph{*stámnaz} & comparative citation label for the broader stem/trunk
family & does not itself control the \emph{stefn} derivation discussed
here & broader lexical label only \\
\emph{*stébnō} & voice-noun transponent & trace output: \emph{stefn} &
selected OE-facing input \\
\emph{stemn} & later attested West Saxon doublet & secondary form from
\emph{stefn} by \emph{fn} \textgreater{} \emph{mn} & real OE variant,
but not the selected comparator \\
\end{longtable}
}

\subsubsection{swan --- OE swanes}\label{swan-oe-swanes}

Derivation: citation reconstruction \emph{*swánaz}; selected input
\emph{*swánas} \textgreater{} \emph{swanes} (early analogy).

\paragraph{Derivation trace}\label{derivation-trace-29}

Proto input: \emph{*swánas}

\begingroup
\setlength{\fboxsep}{6pt}

\noindent\fbox{%
\begin{minipage}{0.97\linewidth}
\small
\begin{tabularx}{\linewidth}{@{}>{\raggedright\arraybackslash}p{0.300\linewidth}>{\centering\arraybackslash}X>{\raggedright\arraybackslash}p{0.540\linewidth}@{}}
\begin{minipage}[t]{\linewidth}
\centering\textbf{Earlier Germanic changes}\par
\vspace{0.35em}
\raggedright
\centering\textbf{West Germanic}\par
\raggedright
\vspace{0.2em}
\raggedright [no change]\par
\vspace{0.6em}
\centering\textbf{Northwest Germanic}\par
\raggedright
\vspace{0.2em}
\raggedright [no change]\par
\end{minipage}
&
&
\begin{minipage}[t]{\linewidth}
\centering\textbf{Old English changes}\par
\vspace{0.35em}
\raggedright
\begin{tabular}{@{}>{\raggedright\arraybackslash}p{0.70\linewidth}@{\hspace{0.55em}}>{\raggedright\arraybackslash}p{0.16\linewidth}@{\hspace{0.25em}}}
Anglo Frisian Brightening & \emph{*swánæs} \\
OE Unstressed AE Merger & \emph{*swánes} \\
\end{tabular}
\end{minipage}
\\
\end{tabularx}
\end{minipage}%
} \endgroup

Outcome: \emph{swanes}

\paragraph{Reconstruction and comparative
evidence}\label{reconstruction-and-comparative-evidence-29}

The Germanic noun is ordinarily cited as the masculine a-stem
\emph{*swanaz} (\citeproc{ref-Orel2003}{Orel 2003, 367}). The selected
input \emph{*swánas} is not a competing lexeme reconstruction. It is the
genitive singular of the same paradigm.

The question here is therefore one of paradigm cell rather than stem
history. The citation form remains \emph{*swanaz} \textgreater{}
\emph{swan}; the selected comparison form is the genitive singular
\emph{*swánas} \textgreater{} \emph{swanes}.

\paragraph{Old English evidence}\label{old-english-evidence-29}

Old English dictionaries give the ordinary headword as \emph{swan}
(\citeproc{ref-ClarkHall1960}{Clark Hall 1960}). Bright's glossary,
however, also records the exact inflected form \emph{swanes}; it glosses
the noun as \emph{swan, m.} with genitive singular \emph{swanes} and
cites the phrase \emph{swanes feðre}
(\citeproc{ref-BrightCassidyRingler1971}{Bright 1971, 441}).

The target is therefore an attested Old English genitive singular, not a
reconstruction. It is also not the ordinary citation lemma. The entry
must keep those two facts distinct.

\paragraph{Development to Old
English}\label{development-to-old-english-29}

From \emph{*swánas}, Anglo-Frisian brightening gives \emph{*swánæs}, and
subsequent merger of unstressed \emph{æ} with \emph{e} yields
\emph{swanes}. The comparison is straightforward once the genitive
singular is chosen as the relevant cell.

\paragraph{Paradigm-cell comparison}\label{paradigm-cell-comparison}

The comparison below is manual. It separates the ordinary citation form
from the selected inflected cell.

{\def\LTcaptype{none} % do not increment counter
\begin{longtable}[]{@{}
  >{\raggedright\arraybackslash}p{(\linewidth - 8\tabcolsep) * \real{0.2000}}
  >{\raggedright\arraybackslash}p{(\linewidth - 8\tabcolsep) * \real{0.2000}}
  >{\raggedright\arraybackslash}p{(\linewidth - 8\tabcolsep) * \real{0.2000}}
  >{\raggedright\arraybackslash}p{(\linewidth - 8\tabcolsep) * \real{0.2000}}
  >{\raggedright\arraybackslash}p{(\linewidth - 8\tabcolsep) * \real{0.2000}}@{}}
\toprule\noalign{}
\begin{minipage}[b]{\linewidth}\raggedright
PGmc cell / interpretation
\end{minipage} & \begin{minipage}[b]{\linewidth}\raggedright
Candidate input
\end{minipage} & \begin{minipage}[b]{\linewidth}\raggedright
OE output or comparison
\end{minipage} & \begin{minipage}[b]{\linewidth}\raggedright
OE comparison form
\end{minipage} & \begin{minipage}[b]{\linewidth}\raggedright
Result
\end{minipage} \\
\midrule\noalign{}
\endhead
\bottomrule\noalign{}
\endlastfoot
citation nominative singular & *swánaz & OE headword \emph{swan} & swan
& ordinary lexeme line \\
selected genitive singular & *swánas & trace output: \emph{swanes} &
swanes & selected attested cell \\
\end{longtable}
}

\subsubsection{thousand --- OE þūsend}\label{thousand-oe-uxfeux16bsend}

Derivation: citation reconstruction \emph{*θūs-undī}; selected input
\emph{*θūs-èndi} \textgreater{} \emph{þūsend} (early analogy).

\paragraph{Derivation trace}\label{derivation-trace-30}

Proto input: \emph{*θūsèndi}

\begingroup
\setlength{\fboxsep}{6pt}

\noindent\fbox{%
\begin{minipage}{0.97\linewidth}
\small
\begin{tabularx}{\linewidth}{@{}>{\raggedright\arraybackslash}p{0.300\linewidth}>{\centering\arraybackslash}X>{\raggedright\arraybackslash}p{0.540\linewidth}@{}}
\begin{minipage}[t]{\linewidth}
\centering\textbf{Earlier Germanic changes}\par
\vspace{0.35em}
\raggedright
\centering\textbf{West Germanic}\par
\raggedright
\vspace{0.2em}
\begin{tabular}{@{}>{\raggedright\arraybackslash}p{0.64\linewidth}@{\hspace{0.55em}}>{\raggedright\arraybackslash}p{0.16\linewidth}@{\hspace{0.25em}}}
PWGmc Early I Apocope & \emph{*θūsènd} \\
\end{tabular}
\vspace{0.6em}
\centering\textbf{Northwest Germanic}\par
\raggedright
\vspace{0.2em}
\raggedright [no change]\par
\end{minipage}
&
&
\begin{minipage}[t]{\linewidth}
\centering\textbf{Old English changes}\par
\vspace{0.35em}
\raggedright
\begin{tabular}{@{}>{\raggedright\arraybackslash}p{0.70\linewidth}@{\hspace{0.55em}}>{\raggedright\arraybackslash}p{0.16\linewidth}@{\hspace{0.25em}}}
OE Strip Secondary Stress & \emph{*θūsend} \\
\end{tabular}
\end{minipage}
\\
\end{tabularx}
\end{minipage}%
} \endgroup

Outcome: \emph{þūsend}

\paragraph{Reconstruction and comparative
evidence}\label{reconstruction-and-comparative-evidence-30}

Kroonen reconstructs the Germanic numeral as \emph{*þūsundī-} and cites
Old English \emph{þūsend} among its continuations
(\citeproc{ref-Kroonen2013}{Kroonen 2013, 554}). The selected input
\emph{*θūs-èndi} is not the same claim. It is an OE-oriented transponent
with the second-member vowel already resolved to \emph{e} and the final
high vowel already shortened for apocope.

The important question is therefore chronological. Why does Old English
show \emph{þūsend}, while related languages such as Old Saxon and Old
High German keep \emph{u} in the second syllable?
(\citeproc{ref-Kroonen2013}{Kroonen 2013, 554}).

\paragraph{Old English evidence}\label{old-english-evidence-30}

Old English \emph{þūsend} is an ordinary citation form, not a selected
oblique or paradigm cell. Campbell treats it as a neuter noun with
normal case forms (\citeproc{ref-Campbell1959}{Campbell 1959, sect.
689}). The problem lies in the internal history of the word, not in its
lexical status.

\paragraph{Development to Old
English}\label{development-to-old-english-30}

If the old final \emph{-ī} had remained long enough to trigger ordinary
double umlaut, Campbell's rule would point toward a form of
\emph{*þȳsend} type rather than attested \emph{þūsend}
(\citeproc{ref-Campbell1959}{Campbell 1959, sect. 203}). Preserved root
\emph{ū} therefore argues that the umlaut-triggering vowel was lost or
neutralized before the ordinary OE umlaut outcome could develop.

That early loss, however, does not by itself explain the medial
\emph{e}. Luick compares the word with \emph{ærende} and later groups
\emph{thousand} with forms reshaped on that pattern
(\citeproc{ref-Luick1914}{Luick 1914-\/-1940, sects 198, 492}). Viredaz
is more cautious, arguing that Old English \emph{e} in this weak
position may simply write schwa and so need not prove a unique
\emph{ærende}-type analogy
(\citeproc{ref-GermanicSlavicBaltic2025}{Viredaz 2025, sect. 2.1.4}).

The selected transponent \emph{*θūs-èndi} captures the OE-side state
from which the documented trace reaches \emph{þūsend}.

\paragraph{Stage comparison}\label{stage-comparison-1}

The comparison below is manual. It separates the secure chronology from
the more interpretive account of the second-syllable vowel.

{\def\LTcaptype{none} % do not increment counter
\begin{longtable}[]{@{}
  >{\raggedright\arraybackslash}p{(\linewidth - 6\tabcolsep) * \real{0.2500}}
  >{\raggedright\arraybackslash}p{(\linewidth - 6\tabcolsep) * \real{0.2500}}
  >{\raggedright\arraybackslash}p{(\linewidth - 6\tabcolsep) * \real{0.2500}}
  >{\raggedright\arraybackslash}p{(\linewidth - 6\tabcolsep) * \real{0.2500}}@{}}
\toprule\noalign{}
\begin{minipage}[b]{\linewidth}\raggedright
Stage / interpretation
\end{minipage} & \begin{minipage}[b]{\linewidth}\raggedright
Candidate form
\end{minipage} & \begin{minipage}[b]{\linewidth}\raggedright
OE relation
\end{minipage} & \begin{minipage}[b]{\linewidth}\raggedright
Result
\end{minipage} \\
\midrule\noalign{}
\endhead
\bottomrule\noalign{}
\endlastfoot
surviving \emph{-ī} with ordinary double umlaut & \emph{*þūsundī-}
treated as still umlaut-active in OE & would point toward \emph{*þȳsend}
& excluded by preserved \emph{ū} \\
early loss of the trigger without further reshaping & \emph{*þūsund-}
type & explains \emph{ū}, but not why OE alone has medial \emph{e} &
incomplete account \\
selected OE-oriented transponent & \emph{*θūs-èndi} & trace output:
\emph{þūsend} & selected modeling input \\
\end{longtable}
}

\subsubsection{timber --- OE timber}\label{timber-oe-timber}

Derivation: citation reconstruction \emph{*tímrą}; selected input
\emph{*tímbrą} \textgreater{} \emph{timber} (early analogy).

\paragraph{Derivation trace}\label{derivation-trace-31}

Proto input: \emph{*tímbrą}

\begingroup
\setlength{\fboxsep}{6pt}

\noindent\fbox{%
\begin{minipage}{0.97\linewidth}
\small
\begin{tabularx}{\linewidth}{@{}>{\raggedright\arraybackslash}p{0.300\linewidth}>{\centering\arraybackslash}X>{\raggedright\arraybackslash}p{0.540\linewidth}@{}}
\begin{minipage}[t]{\linewidth}
\centering\textbf{Earlier Germanic changes}\par
\vspace{0.35em}
\raggedright
\centering\textbf{West Germanic}\par
\raggedright
\vspace{0.2em}
\raggedright [no change]\par
\vspace{0.6em}
\centering\textbf{Northwest Germanic}\par
\raggedright
\vspace{0.2em}
\raggedright [no change]\par
\end{minipage}
&
&
\begin{minipage}[t]{\linewidth}
\centering\textbf{Old English changes}\par
\vspace{0.35em}
\raggedright
\begin{tabular}{@{}>{\raggedright\arraybackslash}p{0.74\linewidth}@{\hspace{0.55em}}>{\raggedright\arraybackslash}p{0.16\linewidth}@{\hspace{0.25em}}}
OE Heavy Syllable Nasal Apocope & \emph{*tímbr} \\
OE Epenthetic Vowel & \emph{*tímber} \\
\end{tabular}
\end{minipage}
\\
\end{tabularx}
\end{minipage}%
} \endgroup

Outcome: \emph{timber}

\paragraph{Reconstruction and comparative
evidence}\label{reconstruction-and-comparative-evidence-31}

Kroonen reconstructs the noun as \emph{*timbra-} and cites Old English
\emph{timber} among its continuations
(\citeproc{ref-Kroonen2013}{Kroonen 2013}). Ringe and Taylor instead
state the history from PGmc \emph{*timra} through West Germanic
\emph{*timbr} to Old English \emph{timber}
(\citeproc{ref-RingeTaylor2014}{Ringe and Taylor 2014, 327}).

The difference is therefore not over the Old English noun itself. It
concerns whether medial \emph{b} belongs in the comparative citation
form or appears in an early pre-Old-English stage of the cluster.

\paragraph{Old English evidence}\label{old-english-evidence-31}

Clark Hall lemmatizes the noun as \emph{timber} and also records
\emph{timbor} as a variant spelling (\citeproc{ref-ClarkHall1960}{Clark
Hall 1960}). The Old English form is thus an ordinary citation noun, not
a selected oblique cell or a reconstructed convenience form.

\paragraph{Development to Old
English}\label{development-to-old-english-31}

With the consonantal frame \emph{timbr-} in place, the rest of the
development is straightforward. Loss of final \emph{-ą} leaves
\emph{*tímbr}, and epenthetic \emph{e} in the final cluster yields
\emph{timber}. Ringe and Taylor's \emph{*timra} \textgreater{}
\emph{*timbr} \textgreater{} OE \emph{timber} and the handbook treatment
of this epenthetic vowel point to the same Old English result
(\citeproc{ref-RingeTaylor2014}{Ringe and Taylor 2014, 327};
\citeproc{ref-Campbell1959}{Campbell 1959, sects 463--464}).

\paragraph{Formation comparison}\label{formation-comparison-6}

The comparison below is manual. It separates the comparative headword
from the OE-facing consonantal input.

{\def\LTcaptype{none} % do not increment counter
\begin{longtable}[]{@{}
  >{\raggedright\arraybackslash}p{(\linewidth - 6\tabcolsep) * \real{0.2500}}
  >{\raggedright\arraybackslash}p{(\linewidth - 6\tabcolsep) * \real{0.2500}}
  >{\raggedright\arraybackslash}p{(\linewidth - 6\tabcolsep) * \real{0.2500}}
  >{\raggedright\arraybackslash}p{(\linewidth - 6\tabcolsep) * \real{0.2500}}@{}}
\toprule\noalign{}
\begin{minipage}[b]{\linewidth}\raggedright
Formation or notation
\end{minipage} & \begin{minipage}[b]{\linewidth}\raggedright
Candidate form
\end{minipage} & \begin{minipage}[b]{\linewidth}\raggedright
OE relation
\end{minipage} & \begin{minipage}[b]{\linewidth}\raggedright
Result
\end{minipage} \\
\midrule\noalign{}
\endhead
\bottomrule\noalign{}
\endlastfoot
Kroonen's comparative citation & \emph{*timbra-} & already matches the
consonantal frame of OE \emph{timber} & closest comparative support for
the selected input \\
Ringe-Taylor citation line & \emph{*timra} \textgreater{} \emph{*timbr}
& reaches the same OE noun through early cluster expansion & compatible
comparative background \\
modeled input & \emph{*tímbrą} & trace output: \emph{timber} & selected
OE-facing input \\
\end{longtable}
}

\subsubsection{wake --- OE wacan}\label{wake-oe-wacan}

Derivation: citation reconstruction \emph{*wakēną}; selected input
\emph{*wákaną} \textgreater{} \emph{wacan} (early analogy).

\paragraph{Derivation trace}\label{derivation-trace-32}

Proto input: \emph{*wákaną}

\begingroup
\setlength{\fboxsep}{6pt}

\noindent\fbox{%
\begin{minipage}{0.97\linewidth}
\small
\begin{tabularx}{\linewidth}{@{}>{\raggedright\arraybackslash}p{0.300\linewidth}>{\centering\arraybackslash}X>{\raggedright\arraybackslash}p{0.540\linewidth}@{}}
\begin{minipage}[t]{\linewidth}
\centering\textbf{Earlier Germanic changes}\par
\vspace{0.35em}
\raggedright
\centering\textbf{West Germanic}\par
\raggedright
\vspace{0.2em}
\raggedright [no change]\par
\vspace{0.6em}
\centering\textbf{Northwest Germanic}\par
\raggedright
\vspace{0.2em}
\raggedright [no change]\par
\end{minipage}
&
&
\begin{minipage}[t]{\linewidth}
\centering\textbf{Old English changes}\par
\vspace{0.35em}
\raggedright
\begin{tabular}{@{}>{\raggedright\arraybackslash}p{0.74\linewidth}@{\hspace{0.55em}}>{\raggedright\arraybackslash}p{0.16\linewidth}@{\hspace{0.25em}}}
Anglo Frisian Brightening & \emph{*wækaną} \\
OE A Restoration & \emph{*wakaną} \\
OE Heavy Syllable Nasal Apocope & \emph{*wakan} \\
OE Secondary Nasalization & \emph{*wakąn} \\
OE Weak Tail Reduction & \emph{*wakan} \\
\end{tabular}
\end{minipage}
\\
\end{tabularx}
\end{minipage}%
} \endgroup

Outcome: \emph{wacan}

\paragraph{Reconstruction and comparative
evidence}\label{reconstruction-and-comparative-evidence-32}

Kroonen gives the strong verb as \emph{*wakan-} with Old English
\emph{wacan} (\citeproc{ref-Kroonen2013}{Kroonen 2013}). Ringe and
Taylor separately derive Old English \emph{wacian} from weak
\emph{*wakai-} \textasciitilde{} \emph{*wakja-}
(\citeproc{ref-RingeTaylor2014}{Ringe and Taylor 2014, sect. 3.3.2}).

The difference is therefore lexical and class-based, not graphic. Strong
\emph{wacan} `wake up, arise' and weak \emph{wacian} `be awake, watch'
belong to related but distinct histories.

\paragraph{Old English evidence}\label{old-english-evidence-32}

Clark Hall keeps \emph{wacan} and \emph{wacian} as separate headwords
(\citeproc{ref-ClarkHall1960}{Clark Hall 1960}). Bosworth-Toller adds an
important caution under \emph{wacan}: the simplex infinitive itself does
not occur, its place seeming to be taken by \emph{wæcnan}
(\citeproc{ref-BosworthToller1898}{Bosworth and Toller 1898, 226}).

The target \emph{wacan} is therefore best understood as a normalized
strong headword for the verb family, not as a directly quoted simplex
infinitive. It still remains the correct Old English comparison form for
the strong branch.

\paragraph{Development to Old
English}\label{development-to-old-english-32}

With strong \emph{*wákaną}, Anglo-Frisian brightening first gives a form
of the \emph{*wækaną} type. A-restoration then returns \emph{a}, and the
ordinary tail reductions yield \emph{wacan}. The weak verb \emph{wacian}
belongs to a different prehistory and is not the expected outcome of
this input.

\paragraph{Class comparison}\label{class-comparison-4}

The comparison below is manual. It separates the strong and weak verb
lines.

{\def\LTcaptype{none} % do not increment counter
\begin{longtable}[]{@{}
  >{\raggedright\arraybackslash}p{(\linewidth - 6\tabcolsep) * \real{0.2500}}
  >{\raggedright\arraybackslash}p{(\linewidth - 6\tabcolsep) * \real{0.2500}}
  >{\raggedright\arraybackslash}p{(\linewidth - 6\tabcolsep) * \real{0.2500}}
  >{\raggedright\arraybackslash}p{(\linewidth - 6\tabcolsep) * \real{0.2500}}@{}}
\toprule\noalign{}
\begin{minipage}[b]{\linewidth}\raggedright
Formation / class
\end{minipage} & \begin{minipage}[b]{\linewidth}\raggedright
Candidate input
\end{minipage} & \begin{minipage}[b]{\linewidth}\raggedright
OE outcome or comparison
\end{minipage} & \begin{minipage}[b]{\linewidth}\raggedright
Result
\end{minipage} \\
\midrule\noalign{}
\endhead
\bottomrule\noalign{}
\endlastfoot
weak class-III / class-II branch & \emph{*wakēną}, \emph{*wakai-}
\textasciitilde{} \emph{*wakja-} & OE \emph{wacian} and related weak
forms & related lexeme, but not the target of this entry \\
strong class-VI branch & \emph{*wákaną} & trace output: \emph{wacan} &
selected OE-facing input \\
strong normalized headword & \emph{wacan} & dictionary comparison form
beside attested strong-family forms & correct Old English comparator,
though not a directly quoted simplex infinitive \\
\end{longtable}
}

\subsubsection{water --- OE wæter}\label{water-oe-wuxe6ter}

Derivation: citation reconstruction \emph{*wátną}; selected input
\emph{*wátōr} \textgreater{} \emph{wæter} (early analogy).

\paragraph{Derivation trace}\label{derivation-trace-33}

Proto input: \emph{*wátōr}

\begingroup
\setlength{\fboxsep}{6pt}

\noindent\fbox{%
\begin{minipage}{0.97\linewidth}
\small
\begin{tabularx}{\linewidth}{@{}>{\raggedright\arraybackslash}p{0.300\linewidth}>{\centering\arraybackslash}X>{\raggedright\arraybackslash}p{0.540\linewidth}@{}}
\begin{minipage}[t]{\linewidth}
\centering\textbf{Earlier Germanic changes}\par
\vspace{0.35em}
\raggedright
\centering\textbf{West Germanic}\par
\raggedright
\vspace{0.2em}
\begin{tabular}{@{}>{\raggedright\arraybackslash}p{0.70\linewidth}@{\hspace{0.55em}}>{\raggedright\arraybackslash}p{0.16\linewidth}@{\hspace{0.25em}}}
PWGmc Final Or Lowering & \emph{*wátar} \\
\end{tabular}
\vspace{0.6em}
\centering\textbf{Northwest Germanic}\par
\raggedright
\vspace{0.2em}
\raggedright [no change]\par
\end{minipage}
&
&
\begin{minipage}[t]{\linewidth}
\centering\textbf{Old English changes}\par
\vspace{0.35em}
\raggedright
\begin{tabular}{@{}>{\raggedright\arraybackslash}p{0.70\linewidth}@{\hspace{0.55em}}>{\raggedright\arraybackslash}p{0.16\linewidth}@{\hspace{0.25em}}}
Anglo Frisian Brightening & \emph{*wætær} \\
OE Unstressed AE Merger & \emph{*wæter} \\
\end{tabular}
\end{minipage}
\\
\end{tabularx}
\end{minipage}%
} \endgroup

Outcome: \emph{wæter}

\paragraph{Reconstruction and comparative
evidence}\label{reconstruction-and-comparative-evidence-33}

Kroonen reconstructs a heteroclitic noun \emph{*watar-}
\textasciitilde{} \emph{*watan-} and states that the Proto-Germanic
material points to \emph{*watōr}, \emph{*watenaz}
(\citeproc{ref-Kroonen2013}{Kroonen 2013, 616}). Ringe and Taylor
likewise start from singular \emph{*wator} before the Old English branch
(\citeproc{ref-RingeTaylor2014}{Ringe and Taylor 2014, sect. 3.1.4}).

The generalized comparative label is therefore broader than the singular
form that actually corresponds to Old English \emph{wæter}. The relevant
comparator is the inherited nominative-accusative singular
\emph{*wátōr}.

\paragraph{Old English evidence}\label{old-english-evidence-33}

Bright gives the noun as \emph{wæter} with the regular paradigm
\emph{wæteres}, \emph{wætere}, \emph{wæter(u)}, \emph{wætera},
\emph{wæterum} (\citeproc{ref-BrightCassidyRingler1971}{Bright 1971,
29}). Ringe and Taylor add the dialectal contrast between West Saxon
\emph{weeter} and Mercian \emph{weter}
(\citeproc{ref-RingeTaylor2014}{Ringe and Taylor 2014, sect. 6.5.2}).

The target is therefore an attested Old English citation form within a
normal paradigm. The complication lies on the comparative side of the
lexeme, not in Old English attestation.

\paragraph{Development to Old
English}\label{development-to-old-english-33}

From \emph{*wátōr}, pre-final \emph{*ō} becomes \emph{a} before final
\emph{r}, yielding \emph{*watar} (\citeproc{ref-RingeTaylor2014}{Ringe
and Taylor 2014, sect. 3.1.4}). Anglo-Frisian brightening then gives
\emph{*wætær}, and merger of unstressed \emph{æ}/\emph{e} yields
\emph{wæter}.

\paragraph{Stage comparison}\label{stage-comparison-2}

The comparison below is manual. It separates the generalized lexeme
label from the singular input that matches the Old English citation
form.

{\def\LTcaptype{none} % do not increment counter
\begin{longtable}[]{@{}
  >{\raggedright\arraybackslash}p{(\linewidth - 6\tabcolsep) * \real{0.2500}}
  >{\raggedright\arraybackslash}p{(\linewidth - 6\tabcolsep) * \real{0.2500}}
  >{\raggedright\arraybackslash}p{(\linewidth - 6\tabcolsep) * \real{0.2500}}
  >{\raggedright\arraybackslash}p{(\linewidth - 6\tabcolsep) * \real{0.2500}}@{}}
\toprule\noalign{}
\begin{minipage}[b]{\linewidth}\raggedright
Stage or notation
\end{minipage} & \begin{minipage}[b]{\linewidth}\raggedright
Candidate form
\end{minipage} & \begin{minipage}[b]{\linewidth}\raggedright
OE relation
\end{minipage} & \begin{minipage}[b]{\linewidth}\raggedright
Result
\end{minipage} \\
\midrule\noalign{}
\endhead
\bottomrule\noalign{}
\endlastfoot
generalized comparative label & \emph{*wátną} & broader lexeme
shorthand, not the singular that corresponds directly to \emph{wæter} &
useful background only \\
heteroclitic stem notation & \emph{*watar-} \textasciitilde{}
\emph{*watan-} & source-faithful comparative reconstruction & explains
why a singular comparator is needed \\
inherited singular input & \emph{*wátōr} & trace output: \emph{wæter} &
selected OE-facing input \\
\end{longtable}
}

\subsubsection{whale --- OE hwæl}\label{whale-oe-hwuxe6l}

Derivation: citation reconstruction \emph{*wálaz}; selected input
\emph{*xwálaz} \textgreater{} \emph{hwæl} (early analogy).

\paragraph{Derivation trace}\label{derivation-trace-34}

Proto input: \emph{*xwálaz}

\begingroup
\setlength{\fboxsep}{6pt}

\noindent\fbox{%
\begin{minipage}{0.97\linewidth}
\small
\begin{tabularx}{\linewidth}{@{}>{\raggedright\arraybackslash}p{0.300\linewidth}>{\centering\arraybackslash}X>{\raggedright\arraybackslash}p{0.540\linewidth}@{}}
\begin{minipage}[t]{\linewidth}
\centering\textbf{Earlier Germanic changes}\par
\vspace{0.35em}
\raggedright
\centering\textbf{West Germanic}\par
\raggedright
\vspace{0.2em}
\raggedright [no change]\par
\vspace{0.6em}
\centering\textbf{Northwest Germanic}\par
\raggedright
\vspace{0.2em}
\begin{tabular}{@{}>{\raggedright\arraybackslash}p{0.74\linewidth}@{\hspace{0.55em}}>{\raggedright\arraybackslash}p{0.16\linewidth}@{\hspace{0.25em}}}
\mbox{PGmc Final Z Deletion} & \emph{*xwála} \\
\end{tabular}
\end{minipage}
&
&
\begin{minipage}[t]{\linewidth}
\centering\textbf{Old English changes}\par
\vspace{0.35em}
\raggedright
\begin{tabular}{@{}>{\raggedright\arraybackslash}p{0.70\linewidth}@{\hspace{0.55em}}>{\raggedright\arraybackslash}p{0.16\linewidth}@{\hspace{0.25em}}}
PWGmc Final Bare A Loss & \emph{*xwál} \\
Anglo Frisian Brightening & \emph{*xwæl} \\
\end{tabular}
\end{minipage}
\\
\end{tabularx}
\end{minipage}%
} \endgroup

Outcome: \emph{hwæl}

\paragraph{Reconstruction and comparative
evidence}\label{reconstruction-and-comparative-evidence-34}

The comparative sources are not uniform. Orel gives \emph{*xwalaz} and
notes some mixed \emph{*xwaliz} evidence (\citeproc{ref-Orel2003}{Orel
2003, 197}). Kroonen instead cites \emph{*hwali-}
(\citeproc{ref-Kroonen2013}{Kroonen 2013}).

Both notations agree on inherited initial \emph{hw-/xw-}, but they
differ in stem label. The a-stem-like input followed here is closer to
Orel's notation than to Kroonen's citation form.

\paragraph{Old English evidence}\label{old-english-evidence-34}

Clark Hall lemmatizes the noun as \emph{hwal}, and Bosworth-Toller
preserves the plural \emph{hwalas} (\citeproc{ref-ClarkHall1960}{Clark
Hall 1960, 170}; \citeproc{ref-BosworthToller1898}{Bosworth and Toller
1898, 326}). The comparison form is normalized here as \emph{hwæl} for
the singular citation form with Anglo- Frisian fronting.

The plural \emph{hwalas} remains important control evidence. It shows
the same lexeme with \emph{a} in an open syllable, beside singular
\emph{hwæl} in the closed monosyllable.

\paragraph{Development to Old
English}\label{development-to-old-english-34}

From \emph{*xwálaz}, final \emph{-z} disappears and bare final \emph{-a}
is lost. Anglo-Frisian fronting then yields \emph{æ} in the closed
monosyllable, and Old English orthography writes \emph{hwæl}.

\paragraph{Formation comparison}\label{formation-comparison-7}

The comparison below is manual. It separates the competing comparative
notations from the normalized Old English singular.

{\def\LTcaptype{none} % do not increment counter
\begin{longtable}[]{@{}
  >{\raggedright\arraybackslash}p{(\linewidth - 6\tabcolsep) * \real{0.2500}}
  >{\raggedright\arraybackslash}p{(\linewidth - 6\tabcolsep) * \real{0.2500}}
  >{\raggedright\arraybackslash}p{(\linewidth - 6\tabcolsep) * \real{0.2500}}
  >{\raggedright\arraybackslash}p{(\linewidth - 6\tabcolsep) * \real{0.2500}}@{}}
\toprule\noalign{}
\begin{minipage}[b]{\linewidth}\raggedright
Comparative line
\end{minipage} & \begin{minipage}[b]{\linewidth}\raggedright
Candidate form
\end{minipage} & \begin{minipage}[b]{\linewidth}\raggedright
OE relation
\end{minipage} & \begin{minipage}[b]{\linewidth}\raggedright
Result
\end{minipage} \\
\midrule\noalign{}
\endhead
\bottomrule\noalign{}
\endlastfoot
Orel's citation & \emph{*xwalaz} & same stem notation as the modeled
singular line & closest comparative support for the selected input \\
Kroonen's citation & \emph{*hwali-} & same initial cluster, different
stem label & important comparative rival, but not the notation followed
here \\
modeled input & \emph{*xwálaz} & trace output: \emph{hwæl} & selected
OE-facing input \\
plural control & \emph{hwalas} & attested open-syllable plural beside
singular \emph{hwæl} & confirms that the lexeme also preserves an
\emph{a}-vocalism branch \\
\end{longtable}
}

\subsubsection{whine --- OE hwīnan}\label{whine-oe-hwux12bnan}

Derivation: citation reconstruction \emph{*wainōjaną}; selected input
\emph{*xwī́naną} \textgreater{} \emph{hwīnan} (early analogy).

\paragraph{Derivation trace}\label{derivation-trace-35}

Proto input: \emph{*xwī́naną}

\begingroup
\setlength{\fboxsep}{6pt}

\noindent\fbox{%
\begin{minipage}{0.97\linewidth}
\small
\begin{tabularx}{\linewidth}{@{}>{\raggedright\arraybackslash}p{0.300\linewidth}>{\centering\arraybackslash}X>{\raggedright\arraybackslash}p{0.540\linewidth}@{}}
\begin{minipage}[t]{\linewidth}
\centering\textbf{Earlier Germanic changes}\par
\vspace{0.35em}
\raggedright
\centering\textbf{West Germanic}\par
\raggedright
\vspace{0.2em}
\raggedright [no change]\par
\vspace{0.6em}
\centering\textbf{Northwest Germanic}\par
\raggedright
\vspace{0.2em}
\raggedright [no change]\par
\end{minipage}
&
&
\begin{minipage}[t]{\linewidth}
\centering\textbf{Old English changes}\par
\vspace{0.35em}
\raggedright
\begin{tabular}{@{}>{\raggedright\arraybackslash}p{0.74\linewidth}@{\hspace{0.55em}}>{\raggedright\arraybackslash}p{0.16\linewidth}@{\hspace{0.25em}}}
OE Heavy Syllable Nasal Apocope & \emph{*xwī́nan} \\
OE Secondary Nasalization & \emph{*xwī́nąn} \\
OE Weak Tail Reduction & \emph{*xwī́nan} \\
\end{tabular}
\end{minipage}
\\
\end{tabularx}
\end{minipage}%
} \endgroup

Outcome: \emph{hwīnan}

\paragraph{Reconstruction and comparative
evidence}\label{reconstruction-and-comparative-evidence-35}

The citation reconstruction preserved in the header belongs to the
lament-family verb seen in German \emph{weinen} and Old English
\emph{wānian}. Kroonen instead separates Old English \emph{hwīnan} under
\emph{*hwinan-} (\citeproc{ref-Kroonen2013}{Kroonen 2013, 267}). Orel
likewise distinguishes strong \emph{*xwinanan} from weak
\emph{*wainōjanan} (\citeproc{ref-Orel2003}{Orel 2003, 201}). Ringe and
Taylor make the same split at the Northwest Germanic level, linking Old
Norse \emph{hvina} and Old English \emph{hwinan} to the same strong verb
(\citeproc{ref-RingeTaylor2014}{Ringe and Taylor 2014, 130}).

The difference affects both phonology and morphology. The lament family
has initial \emph{w-}, diphthongal \emph{ai}, and weak-II morphology,
whereas the verb behind Old English \emph{hwīnan} has initial
\emph{hw-/xw-}, long \emph{ī}, and strong-verb inflection
(\citeproc{ref-Kroonen2013}{Kroonen 2013}; \citeproc{ref-Orel2003}{Orel
2003}). The selected input \emph{*xwī́naną} therefore represents a
competing comparative identification rather than a hidden cell of
\emph{*wainōjaną}.

\paragraph{Old English evidence}\label{old-english-evidence-35}

Clark Hall records \emph{hwinan} with the gloss `to hiss, whizz,
whistle' (\citeproc{ref-ClarkHall1960}{Clark Hall 1960, 171}). Seebold
keeps the verb among the strong verbs and notes that only a
present-tense attestation is directly preserved, while Sievers-Brunner
likewise lists \emph{hwinan stv.} (\citeproc{ref-Seebold1970}{Seebold
1970}; \citeproc{ref-SieversBrunner1965}{Sievers 1965}).

The Old English form is normalized here as \emph{hwīnan}. That
normalization adds the usual vowel length marking to the dictionary
spelling \emph{hwinan}; it does not turn an unattested verb into a
reconstructed one.

\paragraph{Development to Old
English}\label{development-to-old-english-35}

Once the strong-verb input \emph{*xwī́naną} is selected, the path to Old
English is straightforward. The compact trace shows heavy-syllable nasal
apocope, secondary nasalization, and weak-tail reduction, after which
the form surfaces as \emph{hwīnan}.

No special paradigm maneuver is needed for this verb. The comparison is
between two different Germanic verb families: the Old English form
belongs with the strong verb \emph{*hwīnan-}, not with the weak lament
verb.

\paragraph{Verb-family comparison}\label{verb-family-comparison}

The comparison below is manual. It separates the competing comparative
labels that stand behind the inherited Old English forms.

{\def\LTcaptype{none} % do not increment counter
\begin{longtable}[]{@{}
  >{\raggedright\arraybackslash}p{(\linewidth - 8\tabcolsep) * \real{0.2000}}
  >{\raggedright\arraybackslash}p{(\linewidth - 8\tabcolsep) * \real{0.2000}}
  >{\raggedright\arraybackslash}p{(\linewidth - 8\tabcolsep) * \real{0.2000}}
  >{\raggedright\arraybackslash}p{(\linewidth - 8\tabcolsep) * \real{0.2000}}
  >{\raggedright\arraybackslash}p{(\linewidth - 8\tabcolsep) * \real{0.2000}}@{}}
\toprule\noalign{}
\begin{minipage}[b]{\linewidth}\raggedright
Verb family / interpretation
\end{minipage} & \begin{minipage}[b]{\linewidth}\raggedright
Candidate input
\end{minipage} & \begin{minipage}[b]{\linewidth}\raggedright
Old English outcome or comparison
\end{minipage} & \begin{minipage}[b]{\linewidth}\raggedright
OE comparison form
\end{minipage} & \begin{minipage}[b]{\linewidth}\raggedright
Result
\end{minipage} \\
\midrule\noalign{}
\endhead
\bottomrule\noalign{}
\endlastfoot
lament-family weak verb & \emph{*wainōjaną} & comparative continuation
in OE \emph{wānian} & wānian & competing citation reconstruction, but
not the source of \emph{hwīnan} \\
selected strong verb & \emph{*xwī́naną} & compact-trace output:
\emph{hwīnan} & hwīnan & exact match between selected input and OE
verb \\
comparative North Germanic cognate & Northwest Germanic strong verb
behind ON \emph{hvina} / OE \emph{hwinan} & ON \emph{hvina} / OE
\emph{hwinan} & hwīnan & supports the strong-verb identification \\
\end{longtable}
}

\subsubsection{withy --- OE wīþiġ}\label{withy-oe-wux12buxfeiux121}

Derivation: citation reconstruction \emph{*wáiθiz}; selected input
\emph{*wī́θagą} \textgreater{} \emph{wīþiġ} (early analogy).

\paragraph{Derivation trace}\label{derivation-trace-36}

Proto input: \emph{*wī́θagą}

\begingroup
\setlength{\fboxsep}{6pt}

\noindent\fbox{%
\begin{minipage}{0.97\linewidth}
\small
\begin{tabularx}{\linewidth}{@{}>{\raggedright\arraybackslash}p{0.300\linewidth}>{\centering\arraybackslash}X>{\raggedright\arraybackslash}p{0.540\linewidth}@{}}
\begin{minipage}[t]{\linewidth}
\centering\textbf{Earlier Germanic changes}\par
\vspace{0.35em}
\raggedright
\centering\textbf{West Germanic}\par
\raggedright
\vspace{0.2em}
\raggedright [no change]\par
\vspace{0.6em}
\centering\textbf{Northwest Germanic}\par
\raggedright
\vspace{0.2em}
\raggedright [no change]\par
\end{minipage}
&
&
\begin{minipage}[t]{\linewidth}
\centering\textbf{Old English changes}\par
\vspace{0.35em}
\raggedright
\begin{tabular}{@{}>{\raggedright\arraybackslash}p{0.74\linewidth}@{\hspace{0.55em}}>{\raggedright\arraybackslash}p{0.16\linewidth}@{\hspace{0.25em}}}
Anglo Frisian Brightening & \emph{*wī́θægą} \\
OE Heavy Syllable Nasal Apocope & \emph{*wī́θæg} \\
OE Velar Palatalization & \emph{*wī́θæʤ} \\
OE Unstressed AE Merger & \emph{*wī́θeʤ} \\
OE Late Unstressed Ag Suffix & \emph{*wī́θiʤ} \\
\end{tabular}
\end{minipage}
\\
\end{tabularx}
\end{minipage}%
} \endgroup

Outcome: \emph{wīþiġ}

\paragraph{Reconstruction and comparative
evidence}\label{reconstruction-and-comparative-evidence-36}

The comparative evidence groups the word with Germanic forms of the
\emph{*wīþja}/\emph{ō-} or \emph{*wiþ-} type
(\citeproc{ref-KlugeSeebold2011}{Kluge 2011};
\citeproc{ref-Orel2003}{Orel 2003}). That material is useful for the
cognate set, but it does not by itself explain the Old English suffix of
\emph{wīþiġ}.

For Old English, the relevant point is the suffix history. Campbell's
account of OE \emph{-ig}, including forms such as \emph{hunig}, supports
an analysis in which the \emph{-iġ} of \emph{wīþiġ} continues a
derivational \emph{*-ag-} sequence rather than a heavy ja-stem
\emph{*-ij-} (\citeproc{ref-Campbell1959}{Campbell 1959, sects 275,
376}). The earlier ja-stem pathway remains useful as a comparative
possibility, but the Campbell-Adamczyk line on heavy ja-stems points
instead toward \emph{-e} or zero type outcomes
(\citeproc{ref-Campbell1959}{Campbell 1959};
\citeproc{ref-Adamczyk2001}{Adamczyk 2001}). The selected input
\emph{*wī́θagą} is thus a formation choice rather than a mere respelling
of the comparative headword.

\paragraph{Old English evidence}\label{old-english-evidence-36}

Clark Hall records the noun as \emph{wiðig}, with related inflected
forms of the same lexical base (\citeproc{ref-ClarkHall1960}{Clark Hall
1960}). The form used here, \emph{wīþiġ}, is a normalized Old English
spelling with macrons and palatal ġ marked explicitly.

The relevant comparison form is therefore not a reconstructed dictionary
convenience but an established Old English noun. What requires
explanation is why the selected Proto-Germanic input is \emph{*wī́θagą}
rather than a comparative headword of the \emph{*wīþja-} type.

\paragraph{Development to Old
English}\label{development-to-old-english-36}

From \emph{*wī́θagą}, Anglo-Frisian brightening gives a fronted vowel in
the suffixal syllable, and, on the Campbell analysis adopted here, the
later Old English development of \emph{*-ag-} yields \emph{-iġ}
(\citeproc{ref-Campbell1959}{Campbell 1959, sects 275, 376}).
Palatalization supplies the final \emph{ġ}, and the full development
reaches \emph{wīþiġ}.

This derivation is regular for the selected formation. The central claim
of the entry is therefore morphological: Old English \emph{wīþiġ}
belongs with an \emph{*-ag-} derivative, whereas the comparative
\emph{*wīþja-} label belongs to a different way of presenting the
cognate family.

\paragraph{Formation comparison}\label{formation-comparison-8}

The comparison below is manual. It distinguishes the comparative
headword from the Old English-facing formation that actually yields the
attested noun.

{\def\LTcaptype{none} % do not increment counter
\begin{longtable}[]{@{}
  >{\raggedright\arraybackslash}p{(\linewidth - 8\tabcolsep) * \real{0.2000}}
  >{\raggedright\arraybackslash}p{(\linewidth - 8\tabcolsep) * \real{0.2000}}
  >{\raggedright\arraybackslash}p{(\linewidth - 8\tabcolsep) * \real{0.2000}}
  >{\raggedright\arraybackslash}p{(\linewidth - 8\tabcolsep) * \real{0.2000}}
  >{\raggedright\arraybackslash}p{(\linewidth - 8\tabcolsep) * \real{0.2000}}@{}}
\toprule\noalign{}
\begin{minipage}[b]{\linewidth}\raggedright
Formation
\end{minipage} & \begin{minipage}[b]{\linewidth}\raggedright
Candidate input
\end{minipage} & \begin{minipage}[b]{\linewidth}\raggedright
Expected or documented OE outcome
\end{minipage} & \begin{minipage}[b]{\linewidth}\raggedright
OE comparison form
\end{minipage} & \begin{minipage}[b]{\linewidth}\raggedright
Result
\end{minipage} \\
\midrule\noalign{}
\endhead
\bottomrule\noalign{}
\endlastfoot
comparative family label & *wáiθiz & broader cognate-set headword & OE
family context & useful lexeme label, but not the direct source of
\emph{wīþiġ} \\
heavy ja-stem analysis & \emph{*wīþja-} type & Campbell/Adamczyk-style
heavy ja-stem \emph{-e} / zero outcome & \emph{wīþiġ} & does not account
cleanly for the OE suffix \\
selected \emph{*-ag-} derivative & *wī́θagą & compact-trace output:
\emph{wīþiġ} & wīþiġ & exact match between formation and target \\
\end{longtable}
}

\subsubsection{world --- OE weorold}\label{world-oe-weorold}

Derivation: citation reconstruction \emph{*wíra-àldiz}; selected input
\emph{*wír-àldu} \textgreater{} \emph{weorold} (early analogy).

\paragraph{Derivation trace}\label{derivation-trace-37}

Proto input: \emph{*wíràldu}

\begingroup
\setlength{\fboxsep}{6pt}

\noindent\fbox{%
\begin{minipage}{0.97\linewidth}
\small
\begin{tabularx}{\linewidth}{@{}>{\raggedright\arraybackslash}p{0.300\linewidth}>{\centering\arraybackslash}X>{\raggedright\arraybackslash}p{0.540\linewidth}@{}}
\begin{minipage}[t]{\linewidth}
\centering\textbf{Earlier Germanic changes}\par
\vspace{0.35em}
\raggedright
\centering\textbf{West Germanic}\par
\raggedright
\vspace{0.2em}
\raggedright [no change]\par
\vspace{0.6em}
\centering\textbf{Northwest Germanic}\par
\raggedright
\vspace{0.2em}
\begin{tabular}{@{}>{\raggedright\arraybackslash}p{0.64\linewidth}@{\hspace{0.55em}}>{\raggedright\arraybackslash}p{0.20\linewidth}@{\hspace{0.25em}}}
NWGmc I Lowering & \emph{*wéràldu} \\
\end{tabular}
\end{minipage}
&
&
\begin{minipage}[t]{\linewidth}
\centering\textbf{Old English changes}\par
\vspace{0.35em}
\raggedright
\begin{tabular}{@{}>{\raggedright\arraybackslash}p{0.68\linewidth}@{\hspace{0.55em}}>{\raggedright\arraybackslash}p{0.22\linewidth}@{\hspace{0.25em}}}
OE Inter Stress Raising & \emph{*wéruldu} \\
OE Med Unstressed U Lowering & \emph{*wéroldu} \\
OE Back Mutation & \emph{*wéoroldu} \\
OE High Vowel Apocope & \emph{*wéorold} \\
\end{tabular}
\end{minipage}
\\
\end{tabularx}
\end{minipage}%
} \endgroup

Outcome: \emph{weorold}

\paragraph{Reconstruction and comparative
evidence}\label{reconstruction-and-comparative-evidence-37}

The word is the old compound `age of men'. Comparative sources preserve
two slightly different views of its first element. Orel and the
\emph{*wira-} tradition keep the older \emph{i}-vocalism, while Ringe
and Taylor discuss the lowered form \emph{*weraldiz} and its
pre-Old-English chain \emph{*weraldu} \textgreater{} \emph{*weruld}
(\citeproc{ref-Orel2003}{Orel 2003};
\citeproc{ref-RingeTaylor2014}{Ringe and Taylor 2014, 341}).
Kluge-Seebold preserves both tendencies by giving compound
\emph{*wira-aldō} beside simplex \emph{*wera-}
(\citeproc{ref-KlugeSeebold2011}{Kluge 2011}).

The selected input \emph{*wír-àldu} therefore differs from the citation
label in two ways. It keeps the older \emph{*wir-} vowel of the
comparative headword, but it also presupposes the early shift of the
compound into the ō-stems that Ringe and Taylor note for this lexeme
(\citeproc{ref-RingeTaylor2014}{Ringe and Taylor 2014, 341}). The early
analogical step lies in that stem-class reassignment; the later
phonological developments can then run regularly.

\paragraph{Old English evidence}\label{old-english-evidence-37}

Old English does not preserve a single isolated form. Ringe and Taylor
give West Saxon \emph{weorold} \textasciitilde{} \emph{worold}, Mercian
\emph{weoruld}, Northumbrian \emph{woruld}, and Kentish \emph{wiarald}
(\citeproc{ref-RingeTaylor2014}{Ringe and Taylor 2014, 341}).
Sievers-Brunner and Bright present the same wider set, including the
syncopated \emph{world} and later rounded \emph{wurold}
(\citeproc{ref-SieversBrunner1965}{Sievers 1965};
\citeproc{ref-BrightCassidyRingler1971}{Bright 1971, 465}).

The selected target here is the West Saxon form \emph{weorold}. It is an
attested Old English form within that broader variant cluster, not the
only form the lexeme ever shows.

\paragraph{Development to Old
English}\label{development-to-old-english-37}

From the selected input \emph{*wír-àldu}, Northwest Germanic
\emph{i}-lowering gives \emph{*wér-àldu}. Inter-stress raising then
changes the medial \emph{a} to \emph{u}, producing \emph{*wér-uldu}. In
the Old English branch that unstressed \emph{u} lowers to \emph{o}, and
back mutation yields \emph{*wéor-oldu}; final high-vowel apocope then
gives \emph{weorold}.

This sequence matches the comparative background in Ringe and Taylor's
\emph{*weraldiz} \textgreater{} \emph{*weraldu} \textgreater{}
\emph{*weruld} chain while preserving the \emph{*wir-} notation of the
selected comparative label (\citeproc{ref-RingeTaylor2014}{Ringe and
Taylor 2014, 341}). The modeled Old English form therefore stands at the
meeting point of an early stem-class reshaping and later regular sound
change.

\paragraph{Stage comparison}\label{stage-comparison-3}

The comparison below is manual. It separates the comparative headword
from the OE-facing stage chosen for the derivation.

{\def\LTcaptype{none} % do not increment counter
\begin{longtable}[]{@{}
  >{\raggedright\arraybackslash}p{(\linewidth - 6\tabcolsep) * \real{0.2500}}
  >{\raggedright\arraybackslash}p{(\linewidth - 6\tabcolsep) * \real{0.2500}}
  >{\raggedright\arraybackslash}p{(\linewidth - 6\tabcolsep) * \real{0.2500}}
  >{\raggedright\arraybackslash}p{(\linewidth - 6\tabcolsep) * \real{0.2500}}@{}}
\toprule\noalign{}
\begin{minipage}[b]{\linewidth}\raggedright
Stage / interpretation
\end{minipage} & \begin{minipage}[b]{\linewidth}\raggedright
Candidate form
\end{minipage} & \begin{minipage}[b]{\linewidth}\raggedright
Old English outcome or comparison
\end{minipage} & \begin{minipage}[b]{\linewidth}\raggedright
Relevance to this entry
\end{minipage} \\
\midrule\noalign{}
\endhead
\bottomrule\noalign{}
\endlastfoot
comparative compound with older first-element vowel & \emph{*wíra-àldiz}
& citation reconstruction / lexeme label & preserves the older
\emph{*wir-} tradition of the compound \\
literature-stage lowered compound after early stem-class shift &
\emph{*weraldiz} \textgreater{} \emph{*weraldu} \textgreater{}
\emph{*weruld} & Ringe-Taylor background chain to OE \emph{weorold}
\textasciitilde{} \emph{worold} & explains the older comparative
literature cited for the word \\
selected OE-facing input & \emph{*wír-àldu} & compact-trace output:
\emph{weorold} & exact match for the selected West Saxon target \\
broader OE variant cluster & --- & \emph{worold}, \emph{weoruld},
\emph{woruld}, \emph{wiarald}, \emph{world} & real attested comparanda
that remain outside the selected target line \\
\end{longtable}
}

\subsubsection{youth --- OE ġeoguþ}\label{youth-oe-ux121eoguuxfe}

Derivation: citation reconstruction \emph{*júgunθiz}; selected input
\emph{*júgunθ} \textgreater{} \emph{ġeoguþ} (early analogy).

\paragraph{Derivation trace}\label{derivation-trace-38}

Proto input: \emph{*júgunθ}

\begingroup
\setlength{\fboxsep}{6pt}

\noindent\fbox{%
\begin{minipage}{0.97\linewidth}
\small
\begin{tabularx}{\linewidth}{@{}>{\raggedright\arraybackslash}p{0.440\linewidth}>{\centering\arraybackslash}X>{\raggedright\arraybackslash}p{0.440\linewidth}@{}}
\begin{minipage}[t]{\linewidth}
\centering\textbf{Earlier Germanic changes}\par
\vspace{0.35em}
\raggedright
\centering\textbf{West Germanic}\par
\raggedright
\vspace{0.2em}
\raggedright [no change]\par
\vspace{0.6em}
\centering\textbf{Northwest Germanic}\par
\raggedright
\vspace{0.2em}
\begin{tabular}{@{}>{\raggedright\arraybackslash}p{0.70\linewidth}@{\hspace{0.55em}}>{\raggedright\arraybackslash}p{0.20\linewidth}@{\hspace{0.25em}}}
OE Ws Palatal Glide & \emph{*jéugunθ} \\
NWGmc Nasal Spirant Lengthening & \emph{*jéugūnθ} \\
NWGmc Nasal Spirant Loss & \emph{*jéugūθ} \\
\end{tabular}
\end{minipage}
&
&
\begin{minipage}[t]{\linewidth}
\centering\textbf{Old English changes}\par
\vspace{0.35em}
\raggedright
\begin{tabular}{@{}>{\raggedright\arraybackslash}p{0.74\linewidth}@{\hspace{0.55em}}>{\raggedright\arraybackslash}p{0.16\linewidth}@{\hspace{0.25em}}}
OE Diphthong Leveling & \emph{*jéogūθ} \\
OE Unstressed Long Vowel Shortening & \emph{*jéoguθ} \\
\end{tabular}
\end{minipage}
\\
\end{tabularx}
\end{minipage}%
} \endgroup

Outcome: \emph{ġeoguþ}

\paragraph{Reconstruction and comparative
evidence}\label{reconstruction-and-comparative-evidence-38}

The wider etymological tradition reconstructs an earlier form of the
word as \emph{*ju(w)unþi-} (\citeproc{ref-Kroonen2013}{Kroonen 2013};
\citeproc{ref-Fulk2018}{Fulk 2018}). The selected comparative label
\emph{*júgunθiz} already stands at a later Germanic stage with \emph{g},
and the chosen input \emph{*júgunθ} is later again: it represents the
form after final \emph{-i} has been lost.

That staging matters because Ringe and Taylor explicitly give the
sequence \emph{*jugunþi} \textgreater{} \emph{*juguþ} \textgreater{} OE
\emph{geoguþ} \textasciitilde{} \emph{iuguþ}
(\citeproc{ref-RingeTaylor2014}{Ringe and Taylor 2014, 141}). The
selected input therefore differs from the broader comparative headword
because the Old English development must begin after early loss of final
\emph{-i}.

\paragraph{Old English evidence}\label{old-english-evidence-38}

The Old English noun is attested with varying spellings. Ringe and
Taylor cite \emph{geoguþ} \textasciitilde{} \emph{iuguþ}
(\citeproc{ref-RingeTaylor2014}{Ringe and Taylor 2014, 141}). The form
is normalized here as \emph{ġeoguþ}: the initial palatal is written with
\emph{ġ}, and the attested spelling variation is treated as orthographic
rather than lexical.

Nothing in the source stack suggests that a different paradigm cell
should be chosen. The relevant Old English comparison form is the noun
\emph{ġeoguþ} itself.

\paragraph{Development to Old
English}\label{development-to-old-english-38}

The decisive early step is the loss of final \emph{-i} before the Old
English umlaut stage. If that high vowel remained, the word would
develop an over-umlauted \emph{y}-type vowel instead of the attested
form (\citeproc{ref-RingeTaylor2014}{Ringe and Taylor 2014, 141}).

From the selected input \emph{*júgunθ}, the later development is
regular: palatal fronting yields \emph{*jéugunθ}; nasal-spirant
lengthening and loss give \emph{*jéogūθ}; unstressed long-vowel
shortening then produces \emph{*jéoguθ}, which surfaces as
\emph{ġeoguþ}. Campbell preserves \emph{u} after accented \emph{u} in
forms such as \emph{duguþ} and \emph{munuc}
(\citeproc{ref-Campbell1959}{Campbell 1959, sect. 374}). Brunner
likewise cites \emph{iuzuð} \emph{Jugend} and \emph{munuc} \emph{Mönch}
in the same environment (\citeproc{ref-SieversBrunner1965}{Sievers 1965,
sect. 150.3}).

\paragraph{Stage comparison}\label{stage-comparison-4}

The comparison below is manual. It separates the broader comparative
headword from the later stages relevant to the Old English noun.

{\def\LTcaptype{none} % do not increment counter
\begin{longtable}[]{@{}
  >{\raggedright\arraybackslash}p{(\linewidth - 6\tabcolsep) * \real{0.2500}}
  >{\raggedright\arraybackslash}p{(\linewidth - 6\tabcolsep) * \real{0.2500}}
  >{\raggedright\arraybackslash}p{(\linewidth - 6\tabcolsep) * \real{0.2500}}
  >{\raggedright\arraybackslash}p{(\linewidth - 6\tabcolsep) * \real{0.2500}}@{}}
\toprule\noalign{}
\begin{minipage}[b]{\linewidth}\raggedright
Stage / interpretation
\end{minipage} & \begin{minipage}[b]{\linewidth}\raggedright
Candidate form
\end{minipage} & \begin{minipage}[b]{\linewidth}\raggedright
Old English outcome or comparison
\end{minipage} & \begin{minipage}[b]{\linewidth}\raggedright
Relevance to this entry
\end{minipage} \\
\midrule\noalign{}
\endhead
\bottomrule\noalign{}
\endlastfoot
earlier etymological headword & \emph{*ju(w)unþi-} & comparative family
background & older comparative reconstruction of the lexeme \\
later g-bearing comparative label & \emph{*júgunθiz} & citation
reconstruction / lexeme label & preserves the later Germanic stage
behind the selected entry \\
selected OE-facing input & \emph{*júgunθ} & compact-trace output:
\emph{ġeoguþ} & exact match for the chosen Old English form \\
full \emph{-i} stage retained too long & \emph{*jugunþi} & expected
over-umlauted \emph{y}-type result & negative control showing why early
\emph{-i} loss must precede the OE umlaut stage \\
\end{longtable}
}

\clearpage

\subsection{Part IV. Late analogy and paradigm-cell
selection}\label{part-iv.-late-analogy-and-paradigm-cell-selection}

These entries involve a later paradigm-cell or analogical comparison.
The citation reconstruction remains relevant to the lexeme, but the
selected target is best explained through a particular inflectional or
analogical form rather than through the citation form alone.

\subsubsection{ban --- OE bannes}\label{ban-oe-bannes}

Derivation: citation reconstruction \emph{*bánną}; selected input
\emph{*bánnas} \textgreater{} \emph{bannes} (late analogy).

\paragraph{Derivation trace}\label{derivation-trace-39}

Proto input: \emph{*bánnas}

\begingroup
\setlength{\fboxsep}{6pt}

\noindent\fbox{%
\begin{minipage}{0.97\linewidth}
\small
\begin{tabularx}{\linewidth}{@{}>{\raggedright\arraybackslash}p{0.300\linewidth}>{\centering\arraybackslash}X>{\raggedright\arraybackslash}p{0.540\linewidth}@{}}
\begin{minipage}[t]{\linewidth}
\centering\textbf{Earlier Germanic changes}\par
\vspace{0.35em}
\raggedright
\centering\textbf{West Germanic}\par
\raggedright
\vspace{0.2em}
\raggedright [no change]\par
\vspace{0.6em}
\centering\textbf{Northwest Germanic}\par
\raggedright
\vspace{0.2em}
\raggedright [no change]\par
\end{minipage}
&
&
\begin{minipage}[t]{\linewidth}
\centering\textbf{Old English changes}\par
\vspace{0.35em}
\raggedright
\begin{tabular}{@{}>{\raggedright\arraybackslash}p{0.70\linewidth}@{\hspace{0.55em}}>{\raggedright\arraybackslash}p{0.16\linewidth}@{\hspace{0.25em}}}
Anglo Frisian Brightening & \emph{*bánnæs} \\
OE Unstressed AE Merger & \emph{*bánnes} \\
\end{tabular}
\end{minipage}
\\
\end{tabularx}
\end{minipage}%
} \endgroup

Outcome: \emph{bannes}

\paragraph{Reconstruction and comparative
evidence}\label{reconstruction-and-comparative-evidence-39}

Orel cites a bann-noun under \emph{*bannan}, while Seebold distinguishes
bann-stems of both masculine and neuter type and gives Old English
\emph{gebann} as the noun reflex (\citeproc{ref-Orel2003}{Orel 2003};
\citeproc{ref-Seebold1970}{Seebold 1970}). The citation reconstruction
\emph{*bánną} therefore names the lexeme, but the selected input
\emph{*bánnas} is a specific genitive singular cell.

That distinction matters because the analysis depends on medial, not
final, gemination.

\paragraph{Old English evidence}\label{old-english-evidence-39}

Old English lexicographic evidence securely supports the noun itself.
Clark Hall gives \emph{+bann}, and Bosworth-Toller records
\emph{ge-bann} with oblique usage such as \emph{gebanne}
(\citeproc{ref-ClarkHall1960}{Clark Hall 1960};
\citeproc{ref-BosworthToller1898}{Bosworth and Toller 1898, 303}). The
exact unprefixed genitive \emph{bannes} is less directly cited in the
reviewed material, so it is best treated here as the selected regular
genitive comparison form rather than as a dictionary headword.

\paragraph{Development to Old
English}\label{development-to-old-english-39}

From \emph{*bánnas}, the geminate remains medial before the case ending
and the unstressed vowel develops regularly to give \emph{bannes}. By
contrast, citation \emph{*bánną} loses its final vowel and simplifies
the word-final geminate, so the regular nominative outcome is \emph{ban}
(\citeproc{ref-Campbell1959}{Campbell 1959}).

\paragraph{Paradigm comparison}\label{paradigm-comparison}

The comparison below is manual. It shows why the genitive singular is
the conservative cell used for the entry.

{\def\LTcaptype{none} % do not increment counter
\begin{longtable}[]{@{}
  >{\raggedright\arraybackslash}p{(\linewidth - 8\tabcolsep) * \real{0.2000}}
  >{\raggedright\arraybackslash}p{(\linewidth - 8\tabcolsep) * \real{0.2000}}
  >{\raggedright\arraybackslash}p{(\linewidth - 8\tabcolsep) * \real{0.2000}}
  >{\raggedright\arraybackslash}p{(\linewidth - 8\tabcolsep) * \real{0.2000}}
  >{\raggedright\arraybackslash}p{(\linewidth - 8\tabcolsep) * \real{0.2000}}@{}}
\toprule\noalign{}
\begin{minipage}[b]{\linewidth}\raggedright
PGmc cell / interpretation
\end{minipage} & \begin{minipage}[b]{\linewidth}\raggedright
Candidate input
\end{minipage} & \begin{minipage}[b]{\linewidth}\raggedright
OE output or comparison
\end{minipage} & \begin{minipage}[b]{\linewidth}\raggedright
OE comparison form
\end{minipage} & \begin{minipage}[b]{\linewidth}\raggedright
Result
\end{minipage} \\
\midrule\noalign{}
\endhead
\bottomrule\noalign{}
\endlastfoot
citation nominative singular & *bánną & compact-trace output: \emph{ban}
& ban & regular nominative outcome, but not the selected target \\
selected genitive singular & *bánnas & compact-trace output:
\emph{bannes} & bannes & exact match for the chosen conservative cell \\
\end{longtable}
}

\subsubsection{berry --- OE berġes}\label{berry-oe-berux121es}

Derivation: citation reconstruction \emph{*bázją}; selected input
\emph{*bázjas} \textgreater{} \emph{berġes} (late analogy).

\paragraph{Derivation trace}\label{derivation-trace-40}

Proto input: \emph{*bázjas}

\begingroup
\setlength{\fboxsep}{6pt}

\noindent\fbox{%
\begin{minipage}{0.97\linewidth}
\small
\begin{tabularx}{\linewidth}{@{}>{\raggedright\arraybackslash}p{0.300\linewidth}>{\centering\arraybackslash}X>{\raggedright\arraybackslash}p{0.540\linewidth}@{}}
\begin{minipage}[t]{\linewidth}
\centering\textbf{Earlier Germanic changes}\par
\vspace{0.35em}
\raggedright
\centering\textbf{West Germanic}\par
\raggedright
\vspace{0.2em}
\raggedright [no change]\par
\vspace{0.6em}
\centering\textbf{Northwest Germanic}\par
\raggedright
\vspace{0.2em}
\raggedright [no change]\par
\end{minipage}
&
&
\begin{minipage}[t]{\linewidth}
\centering\textbf{Old English changes}\par
\vspace{0.35em}
\raggedright
\begin{tabular}{@{}>{\raggedright\arraybackslash}p{0.70\linewidth}@{\hspace{0.55em}}>{\raggedright\arraybackslash}p{0.16\linewidth}@{\hspace{0.25em}}}
Anglo Frisian Brightening & \emph{*bærjæs} \\
OE I Umlaut & \emph{*berjæs} \\
OE Unstressed AE Merger & \emph{*berjes} \\
\end{tabular}
\end{minipage}
\\
\end{tabularx}
\end{minipage}%
} \endgroup

Outcome: \emph{berġes}

\paragraph{Reconstruction and comparative
evidence}\label{reconstruction-and-comparative-evidence-40}

Kroonen reconstructs the berry noun as \emph{*basja-} \textasciitilde{}
\emph{*bazja-} (\citeproc{ref-Kroonen2013}{Kroonen 2013, 54}). The
selected input \emph{*bázjas} is therefore not a rival lexeme headword,
but a specific genitive singular cell drawn from that paradigm.

The relevant point is that \emph{*rj} did not geminate in Proto-West
Germanic. Ringe and Taylor's \emph{here, herges} comparison shows the
same \emph{rj} environment in an Old English paradigm without any hidden
gemination repair (\citeproc{ref-RingeTaylor2014}{Ringe and Taylor 2014,
181}).

\paragraph{Old English evidence}\label{old-english-evidence-40}

Campbell cites feminine \emph{berige} `berry' and notes that \emph{-j-}
is retained after \emph{r} in this type
(\citeproc{ref-Campbell1959}{Campbell 1959, 250}). The reviewed evidence
therefore supports the citation form more directly than the exact
genitive \emph{berġes}, which is best read here as the selected regular
genitive comparison form rather than as a dictionary headword.

\paragraph{Development to Old
English}\label{development-to-old-english-40}

Citation \emph{*bázją} gives \emph{bere}, not the selected target. The
genitive singular \emph{*bázjas}, however, gives \emph{berġes}, with
medial \emph{-rġ-} preserved in the same way that Ringe and Taylor cite
\emph{herges} beside \emph{here} (\citeproc{ref-RingeTaylor2014}{Ringe
and Taylor 2014, 181}). This points to paradigm choice rather than to an
extra phonological rule.

\paragraph{Paradigm comparison}\label{paradigm-comparison-1}

The comparison below is manual. It shows the contrast between the
citation form and the selected genitive singular cell.

{\def\LTcaptype{none} % do not increment counter
\begin{longtable}[]{@{}
  >{\raggedright\arraybackslash}p{(\linewidth - 8\tabcolsep) * \real{0.2000}}
  >{\raggedright\arraybackslash}p{(\linewidth - 8\tabcolsep) * \real{0.2000}}
  >{\raggedright\arraybackslash}p{(\linewidth - 8\tabcolsep) * \real{0.2000}}
  >{\raggedright\arraybackslash}p{(\linewidth - 8\tabcolsep) * \real{0.2000}}
  >{\raggedright\arraybackslash}p{(\linewidth - 8\tabcolsep) * \real{0.2000}}@{}}
\toprule\noalign{}
\begin{minipage}[b]{\linewidth}\raggedright
PGmc cell / interpretation
\end{minipage} & \begin{minipage}[b]{\linewidth}\raggedright
Candidate input
\end{minipage} & \begin{minipage}[b]{\linewidth}\raggedright
OE output or comparison
\end{minipage} & \begin{minipage}[b]{\linewidth}\raggedright
OE comparison form
\end{minipage} & \begin{minipage}[b]{\linewidth}\raggedright
Result
\end{minipage} \\
\midrule\noalign{}
\endhead
\bottomrule\noalign{}
\endlastfoot
citation nominative singular & *bázją & compact-trace output:
\emph{bere} & berige / berġe & useful citation-form background, but not
the selected target \\
selected genitive singular & *bázjas & compact-trace output:
\emph{berġes} & berġes & exact match for the chosen conservative cell \\
\end{longtable}
}

\subsubsection{bow --- OE bēag}\label{bow-oe-bux113ag}

Derivation: citation reconstruction \emph{*béuganą}; selected input
\emph{*báug} \textgreater{} \emph{bēag} (late analogy).

\paragraph{Derivation trace}\label{derivation-trace-41}

Proto input: \emph{*báug}

\begingroup
\setlength{\fboxsep}{6pt}

\noindent\fbox{%
\begin{minipage}{0.97\linewidth}
\small
\begin{tabularx}{\linewidth}{@{}>{\raggedright\arraybackslash}p{0.320\linewidth}>{\centering\arraybackslash}X>{\raggedright\arraybackslash}p{0.520\linewidth}@{}}
\begin{minipage}[t]{\linewidth}
\centering\textbf{Earlier Germanic changes}\par
\vspace{0.35em}
\raggedright
\centering\textbf{West Germanic}\par
\raggedright
\vspace{0.2em}
\raggedright [no change]\par
\vspace{0.6em}
\centering\textbf{Northwest Germanic}\par
\raggedright
\vspace{0.2em}
\raggedright [no change]\par
\end{minipage}
&
&
\begin{minipage}[t]{\linewidth}
\centering\textbf{Old English changes}\par
\vspace{0.35em}
\raggedright
\begin{tabular}{@{}>{\raggedright\arraybackslash}p{0.68\linewidth}@{\hspace{0.55em}}>{\raggedright\arraybackslash}p{0.16\linewidth}@{\hspace{0.25em}}}
OE Au Fronting & \emph{*báeug} \\
OE Diphthong Leveling & \emph{*bēag} \\
\end{tabular}
\end{minipage}
\\
\end{tabularx}
\end{minipage}%
} \endgroup

Outcome: \emph{bēag}

\paragraph{Reconstruction and comparative
evidence}\label{reconstruction-and-comparative-evidence-41}

The inherited verb belongs to the class-II strong-verb family
\emph{*béuganą} (\citeproc{ref-RingeTaylor2014}{Ringe and Taylor 2014,
55}). Within that paradigm, however, the infinitive and the singular
preterite continue different ablaut grades. The selected input
\emph{*báug} is the singular preterite cell, whereas the citation form
\emph{*béuganą} is the infinitive.

Campbell's account of Old English class-II strong verbs treats the
singular preterite \emph{au} \textgreater{} \emph{ēa} development as
regular in this environment (\citeproc{ref-Campbell1959}{Campbell 1959,
53}). That is the phonological path relevant for \emph{bēag}, whereas
the analogical \emph{ū} of the present stem belongs to the separate
history behind \emph{būgan} (\citeproc{ref-RingeTaylor2014}{Ringe and
Taylor 2014, 55}).

\paragraph{Old English evidence}\label{old-english-evidence-41}

Bosworth-Toller and Clark Hall both record \emph{bēag} as a preterite
form of \emph{būgan} (\citeproc{ref-BosworthToller1898}{Bosworth and
Toller 1898, 122}; \citeproc{ref-ClarkHall1960}{Clark Hall 1960, 45}).
The form discussed here is therefore an attested Old English verbal
form, not a reconstructed substitute for the infinitive.

The ordinary dictionary headword remains \emph{būgan}, but the relevant
comparison form for this entry is the singular preterite \emph{bēag}.
That is the paradigm cell in which the inherited \emph{*au} grade is
preserved most directly.

\paragraph{Development to Old
English}\label{development-to-old-english-41}

From \emph{*báug}, Anglo-Frisian fronting and the later leveling of the
diphthong produce \emph{bēag} (\citeproc{ref-Campbell1959}{Campbell
1959, 53}). No special analogical repair is needed for that cell. The
form is the regular Old English outcome of the singular-preterite grade.

The analogical element in the wider lexeme belongs instead to the
present stem seen in \emph{būgan}. The selected input differs from the
citation form because the regular inherited pathway survives more
transparently in the preterite than in the infinitive.

\paragraph{Paradigm comparison}\label{paradigm-comparison-2}

The comparison below is manual. It distinguishes the regular singular
preterite from the more familiar infinitival citation form.

{\def\LTcaptype{none} % do not increment counter
\begin{longtable}[]{@{}
  >{\raggedright\arraybackslash}p{(\linewidth - 8\tabcolsep) * \real{0.2000}}
  >{\raggedright\arraybackslash}p{(\linewidth - 8\tabcolsep) * \real{0.2000}}
  >{\raggedright\arraybackslash}p{(\linewidth - 8\tabcolsep) * \real{0.2000}}
  >{\raggedright\arraybackslash}p{(\linewidth - 8\tabcolsep) * \real{0.2000}}
  >{\raggedright\arraybackslash}p{(\linewidth - 8\tabcolsep) * \real{0.2000}}@{}}
\toprule\noalign{}
\begin{minipage}[b]{\linewidth}\raggedright
PGmc cell / interpretation
\end{minipage} & \begin{minipage}[b]{\linewidth}\raggedright
Candidate input
\end{minipage} & \begin{minipage}[b]{\linewidth}\raggedright
Expected or documented OE outcome
\end{minipage} & \begin{minipage}[b]{\linewidth}\raggedright
OE comparison form
\end{minipage} & \begin{minipage}[b]{\linewidth}\raggedright
Result
\end{minipage} \\
\midrule\noalign{}
\endhead
\bottomrule\noalign{}
\endlastfoot
citation infinitive & *béuganą & inherited present-stem history behind
\emph{būgan} & būgan & establishes the lexeme, but not the selected
target \\
singular preterite & *báug & compact-trace output: \emph{bēag} & bēag &
exact match between input, output, and attested cell \\
past participial branch & participial \emph{*bugan-} type & later
participial outcomes & bogen-type evidence & relevant to the paradigm,
but not the clearest match for this entry \\
\end{longtable}
}

The singular preterite is the relevant comparison form. It gives a
direct lautgesetzlich path to attested \emph{bēag}, while the citation
form \emph{būgan} belongs to a paradigm whose present stem has already
undergone later leveling.

\subsubsection{cow --- OE cȳ}\label{cow-oe-cux233}

Derivation: citation reconstruction \emph{*kōz}; selected input
\emph{*kūi} \textgreater{} \emph{cȳ} (late analogy).

\paragraph{Derivation trace}\label{derivation-trace-42}

Proto input: \emph{*kūi}

\begingroup
\setlength{\fboxsep}{6pt}

\noindent\fbox{%
\begin{minipage}{0.97\linewidth}
\small
\begin{tabularx}{\linewidth}{@{}>{\raggedright\arraybackslash}p{0.300\linewidth}>{\centering\arraybackslash}X>{\raggedright\arraybackslash}p{0.540\linewidth}@{}}
\begin{minipage}[t]{\linewidth}
\centering\textbf{Earlier Germanic changes}\par
\vspace{0.35em}
\raggedright
\centering\textbf{West Germanic}\par
\raggedright
\vspace{0.2em}
\raggedright [no change]\par
\vspace{0.6em}
\centering\textbf{Northwest Germanic}\par
\raggedright
\vspace{0.2em}
\raggedright [no change]\par
\end{minipage}
&
&
\begin{minipage}[t]{\linewidth}
\centering\textbf{Old English changes}\par
\vspace{0.35em}
\raggedright
\begin{tabular}{@{}>{\raggedright\arraybackslash}p{0.74\linewidth}@{\hspace{0.55em}}>{\raggedright\arraybackslash}p{0.16\linewidth}@{\hspace{0.25em}}}
OE I Umlaut & \emph{*kȳi} \\
\mbox{OE High Vowel Apocope} & \emph{*kȳ} \\
\end{tabular}
\end{minipage}
\\
\end{tabularx}
\end{minipage}%
} \endgroup

Outcome: \emph{cȳ}

\paragraph{Reconstruction and comparative
evidence}\label{reconstruction-and-comparative-evidence-42}

Kroonen reconstructs a root noun with the inherited alternation
\emph{*kō-} \textasciitilde{} \emph{*ku-}, explicitly nom. \emph{*kōz},
obl. \emph{*kū-} (\citeproc{ref-Kroonen2013}{Kroonen 2013, 299}). The
citation form therefore belongs to the nominative singular, whereas the
selected input \emph{*kūi} belongs to the oblique stem.

Ringe and Taylor also posit a later PNWGmc nominative \emph{*kūaz}
\textgreater{} \emph{*kūz} behind Old English \emph{cū}
(\citeproc{ref-RingeTaylor2014}{Ringe and Taylor 2014, sect. 3.1.3}).
That nominative history matters for the headword, but the form \emph{cȳ}
depends on the oblique \emph{*kū-} stem and a following \emph{*i}.

\paragraph{Old English evidence}\label{old-english-evidence-42}

Clark Hall lemmatizes the noun under \emph{cū} and records gen.sg.
\emph{cū(e)}, \emph{cȳ}, \emph{cūs}, dat.sg. \emph{cȳ}, nom.-acc.pl.
\emph{cȳ}, and dat.pl. \emph{cūm} (\citeproc{ref-ClarkHall1960}{Clark
Hall 1960}). Ringe and Taylor likewise state that the root-noun
\emph{cū} exhibits dat.sg., nom.-acc.pl. \emph{cȳ} \textless{}
\emph{*cūi}, dat.pl. \emph{cūm} \textless{} \emph{*cūm}, and apparently
gen.sg. \emph{cā} \textless{} \emph{*cūiz}
(\citeproc{ref-RingeTaylor2014}{Ringe and Taylor 2014, sect. 6.6.1}).

The dictionary headword is therefore \emph{cū}, but \emph{cȳ} is an
established Old English paradigm form rather than a convenient
reconstruction. It serves as datative singular and also as
nominative-accusative plural. The genitive singular is less stable, with
\emph{cā}, \emph{cȳ}, \emph{cū(e)}, and \emph{cūs} all appearing in the
local source record.

\paragraph{Development to Old
English}\label{development-to-old-english-42}

\emph{*kūi} is the dative singular of the oblique \emph{*kū-} stem. The
following \emph{*i} triggers i-umlaut, so \emph{ū} becomes \emph{ȳ}, and
loss of the final high vowel leaves \emph{cȳ}. The development is
\emph{*kūi} \textgreater{} \emph{*kȳi} \textgreater{} \emph{*kȳ}
\textgreater{} \emph{cȳ}.

This is the regular oblique-cell path recognized by Ringe and Taylor's
paradigm statement (\citeproc{ref-RingeTaylor2014}{Ringe and Taylor
2014, sect. 6.6.1}). It also explains why \emph{cȳ} is the cleanest
comparison form for the present entry, even though the ordinary headword
is \emph{cū}.

\paragraph{Paradigm comparison}\label{paradigm-comparison-3}

A paradigm comparison identifies the Proto-Germanic inflectional cell
that corresponds to an established Old English paradigm form. The
comparison below is manual; no full automatic paradigm-generation run is
presented here.

{\def\LTcaptype{none} % do not increment counter
\begin{longtable}[]{@{}
  >{\raggedright\arraybackslash}p{(\linewidth - 8\tabcolsep) * \real{0.2000}}
  >{\raggedright\arraybackslash}p{(\linewidth - 8\tabcolsep) * \real{0.2000}}
  >{\raggedright\arraybackslash}p{(\linewidth - 8\tabcolsep) * \real{0.2000}}
  >{\raggedright\arraybackslash}p{(\linewidth - 8\tabcolsep) * \real{0.2000}}
  >{\raggedright\arraybackslash}p{(\linewidth - 8\tabcolsep) * \real{0.2000}}@{}}
\toprule\noalign{}
\begin{minipage}[b]{\linewidth}\raggedright
PGmc cell / interpretation
\end{minipage} & \begin{minipage}[b]{\linewidth}\raggedright
Candidate input
\end{minipage} & \begin{minipage}[b]{\linewidth}\raggedright
OE output or comparison
\end{minipage} & \begin{minipage}[b]{\linewidth}\raggedright
OE comparison form
\end{minipage} & \begin{minipage}[b]{\linewidth}\raggedright
Result
\end{minipage} \\
\midrule\noalign{}
\endhead
\bottomrule\noalign{}
\endlastfoot
citation nominative singular & *kōz & OE headword \emph{cū} belongs to
the nominative history of the lexeme & cū & useful background, but not
the chosen comparison for \emph{cȳ} \\
later generalized nominative & PNWGmc \emph{kūaz \textgreater{} }kūz &
inferred nominative \emph{cū} & cū & explains the leveled headword, not
the oblique target \\
dative singular oblique & *kūi & compact-trace output: \emph{cȳ} & cȳ &
exact match between input, output, and paradigm cell \\
genitive singular oblique & *kūiz & Ringe-Taylor: apparently \emph{cā};
Hall: \emph{cū(e)}, \emph{cȳ}, \emph{cūs} & gen.sg. variable & too
unstable to control the entry \\
\end{longtable}
}

The dative singular is the relevant comparison form. It gives a regular
path to attested \emph{cȳ}, while the broader Old English paradigm shows
how far the oblique \emph{*kū-} grade spread beyond that one cell.

\subsubsection{find --- OE fundene}\label{find-oe-fundene}

Derivation: citation reconstruction \emph{*fínθaną}; selected input
\emph{*fúnðanǭ} \textgreater{} \emph{fundene} (late analogy).

\paragraph{Derivation trace}\label{derivation-trace-43}

Proto input: \emph{*fúnðanǭ}

\begingroup
\setlength{\fboxsep}{6pt}

\noindent\fbox{%
\begin{minipage}{0.97\linewidth}
\small
\begin{tabularx}{\linewidth}{@{}>{\raggedright\arraybackslash}p{0.300\linewidth}>{\centering\arraybackslash}X>{\raggedright\arraybackslash}p{0.540\linewidth}@{}}
\begin{minipage}[t]{\linewidth}
\centering\textbf{Earlier Germanic changes}\par
\vspace{0.35em}
\raggedright
\centering\textbf{West Germanic}\par
\raggedright
\vspace{0.2em}
\begin{tabular}{@{}>{\raggedright\arraybackslash}p{0.70\linewidth}@{\hspace{0.55em}}>{\raggedright\arraybackslash}p{0.20\linewidth}@{\hspace{0.25em}}}
PWGmc Dental Hardening & \emph{*fúndanǭ} \\
\end{tabular}
\vspace{0.6em}
\centering\textbf{Northwest Germanic}\par
\raggedright
\vspace{0.2em}
\raggedright [no change]\par
\end{minipage}
&
&
\begin{minipage}[t]{\linewidth}
\centering\textbf{Old English changes}\par
\vspace{0.35em}
\raggedright
\begin{tabular}{@{}>{\raggedright\arraybackslash}p{0.70\linewidth}@{\hspace{0.55em}}>{\raggedright\arraybackslash}p{0.20\linewidth}@{\hspace{0.25em}}}
OE Unstressed Fronting Early & \emph{*fúndænǭ} \\
OE Unstressed Long Vowel Shortening & \emph{*fúndænæ} \\
OE Unstressed AE Merger & \emph{*fúndene} \\
\end{tabular}
\end{minipage}
\\
\end{tabularx}
\end{minipage}%
} \endgroup

Outcome: \emph{fundene}

\paragraph{Reconstruction and comparative
evidence}\label{reconstruction-and-comparative-evidence-43}

The inherited verb is the strong verb \emph{*fínθaną}, continued by Old
English \emph{findan} (\citeproc{ref-RingeTaylor2014}{Ringe and Taylor
2014}). The selected input \emph{*fúnðanǭ}, however, belongs to the
past-participial paradigm rather than to the infinitive. It represents
an oblique singular form of the participle.

That distinction matters because the familiar dictionary form
\emph{funden} is not the cleanest inherited comparison. Campbell, Luick,
and Brunner treat the better-known nominative participial forms as
analogically leveled from inflected cases of the paradigm
(\citeproc{ref-Campbell1959}{Campbell 1959};
\citeproc{ref-Luick1914}{Luick 1914-\/-1940};
\citeproc{ref-SieversBrunner1965}{Sievers 1965}). The selected input
therefore targets the cell in which the regular development is most
transparent.

\paragraph{Old English evidence}\label{old-english-evidence-43}

Bosworth-Toller records \emph{fundene} under \emph{findan} as an Old
English participial form (\citeproc{ref-BosworthToller1898}{Bosworth and
Toller 1898}). Clark Hall likewise preserves the participial stem in
forms such as \emph{funden} and \emph{tō-fundennes}
(\citeproc{ref-ClarkHall1960}{Clark Hall 1960}).

The ordinary dictionary headword for the participle is \emph{funden},
but the relevant comparison form for this entry is the attested oblique
\emph{fundene}. It is an Old English form in its own right, not a merely
convenient probe.

\paragraph{Development to Old
English}\label{development-to-old-english-43}

From \emph{*fúnðanǭ}, the participial oblique develops through regular
loss and weakening of the final ending, yielding \emph{fundene}. In that
cell both the consonantism and the medial vowel history remain regular.

The broader participial paradigm then matters for interpretation. The
more familiar nominative \emph{funden} represents a later analogical
leveling, whereas the selected oblique form preserves the inherited
pathway more directly.

\paragraph{Paradigm comparison}\label{paradigm-comparison-4}

The comparison below is manual. It distinguishes the attested oblique
participle from the more familiar but analogically leveled participial
forms.

{\def\LTcaptype{none} % do not increment counter
\begin{longtable}[]{@{}
  >{\raggedright\arraybackslash}p{(\linewidth - 8\tabcolsep) * \real{0.2000}}
  >{\raggedright\arraybackslash}p{(\linewidth - 8\tabcolsep) * \real{0.2000}}
  >{\raggedright\arraybackslash}p{(\linewidth - 8\tabcolsep) * \real{0.2000}}
  >{\raggedright\arraybackslash}p{(\linewidth - 8\tabcolsep) * \real{0.2000}}
  >{\raggedright\arraybackslash}p{(\linewidth - 8\tabcolsep) * \real{0.2000}}@{}}
\toprule\noalign{}
\begin{minipage}[b]{\linewidth}\raggedright
PGmc cell / interpretation
\end{minipage} & \begin{minipage}[b]{\linewidth}\raggedright
Candidate input
\end{minipage} & \begin{minipage}[b]{\linewidth}\raggedright
Expected or documented OE outcome
\end{minipage} & \begin{minipage}[b]{\linewidth}\raggedright
OE comparison form
\end{minipage} & \begin{minipage}[b]{\linewidth}\raggedright
Result
\end{minipage} \\
\midrule\noalign{}
\endhead
\bottomrule\noalign{}
\endlastfoot
citation infinitive & *fínθaną & inherited verb \emph{findan} & findan &
establishes the lexeme, but not the selected target \\
nominative participial line & *fúnðanaz & later leveled \emph{funden}
type & funden & important paradigm background, but not the cleanest
regular cell \\
selected oblique participle & *fúnðanǭ & compact-trace output:
\emph{fundene} & fundene & exact match between input, output, and
attested cell \\
\end{longtable}
}

The oblique participle is the relevant comparison form. It preserves the
inherited development most directly, while the nominative participial
headword belongs to a later analogical leveling within the paradigm.

\subsubsection{fright --- OE fyrhte}\label{fright-oe-fyrhte}

Derivation: citation reconstruction \emph{*furxtīn}; selected input
\emph{*fúrxtīnaz} \textgreater{} \emph{fyrhte} (late analogy).

\paragraph{Derivation trace}\label{derivation-trace-44}

Proto input: \emph{*fúrxtīnaz}

\begingroup
\setlength{\fboxsep}{6pt}

\noindent\fbox{%
\begin{minipage}{0.97\linewidth}
\small
\begin{tabularx}{\linewidth}{@{}>{\raggedright\arraybackslash}p{0.460\linewidth}>{\centering\arraybackslash}X>{\raggedright\arraybackslash}p{0.460\linewidth}@{}}
\begin{minipage}[t]{\linewidth}
\centering\textbf{Earlier Germanic changes}\par
\vspace{0.35em}
\raggedright
\centering\textbf{West Germanic}\par
\raggedright
\vspace{0.2em}
\raggedright [no change]\par
\vspace{0.6em}
\centering\textbf{Northwest Germanic}\par
\raggedright
\vspace{0.2em}
\begin{tabular}{@{}>{\raggedright\arraybackslash}p{0.68\linewidth}@{\hspace{0.55em}}>{\raggedright\arraybackslash}p{0.22\linewidth}@{\hspace{0.25em}}}
PGmc Final Z Deletion & \emph{*fúrxtīna} \\
\end{tabular}
\end{minipage}
&
&
\begin{minipage}[t]{\linewidth}
\centering\textbf{Old English changes}\par
\vspace{0.35em}
\raggedright
\begin{tabular}{@{}>{\raggedright\arraybackslash}p{0.70\linewidth}@{\hspace{0.55em}}>{\raggedright\arraybackslash}p{0.20\linewidth}@{\hspace{0.25em}}}
PWGmc Final Bare A Loss & \emph{*fúrxtīn} \\
OE I Umlaut & \emph{*fyrxtīn} \\
NWGmc In Stem N Loss & \emph{*fyrxtī} \\
OE Unstressed Long Vowel Shortening & \emph{*fyrxti} \\
OE Med Unstressed I Lowering1 & \emph{*fyrxte} \\
\end{tabular}
\end{minipage}
\\
\end{tabularx}
\end{minipage}%
} \endgroup

Outcome: \emph{fyrhte}

\paragraph{Reconstruction and comparative
evidence}\label{reconstruction-and-comparative-evidence-44}

The noun belongs to the inherited in-stem abstract \emph{*furxtīn}, the
same family as Gothic \emph{faurhtei} (\citeproc{ref-Orel2003}{Orel
2003}). The selected input \emph{*fúrxtīnaz} is not a different lexeme
but an oblique singular cell within that in-stem paradigm.

Ringe and Taylor treat the later nominative forms with \emph{-u} or
\emph{-o} as analogically remodeled, whereas the oblique in-stem forms
continue the older history more directly
(\citeproc{ref-RingeTaylor2014}{Ringe and Taylor 2014}). The selected
input therefore differs from the citation form because the oblique cell
preserves the inherited pathway more clearly than the better-known lemma
forms do.

\paragraph{Old English evidence}\label{old-english-evidence-44}

Bosworth-Toller records \emph{fyrhte} with multiple textual
attestations, and it also records nominative forms such as \emph{fyrhtu}
and \emph{fyrhto} (\citeproc{ref-BosworthToller1898}{Bosworth and Toller
1898}). The noun must be kept distinct from the adjective \emph{forht},
a distinction also reflected in Clark Hall
(\citeproc{ref-ClarkHall1960}{Clark Hall 1960}).

The relevant comparison form is therefore the attested oblique
\emph{fyrhte}. The nominative lemma forms remain part of the Old English
evidence, but they are not the cleanest inherited comparison for this
entry.

\paragraph{Development to Old
English}\label{development-to-old-english-44}

From \emph{*fúrxtīnaz}, the oblique in-stem develops through the loss
and weakening of the final ending and the regular OE history summarized
by Campbell as \emph{-e \textless{} -i \textless{} -in} in this class of
abstracts (\citeproc{ref-Campbell1959}{Campbell 1959}). That yields
\emph{fyrhte}.

The later nominative forms with \emph{-u} or \emph{-o} belong to a
subsequent analogical reshaping of the paradigm. The selected target is
earlier in that sense: it is the attested OE cell in which the inherited
in-stem development remains most transparent.

\paragraph{Paradigm comparison}\label{paradigm-comparison-5}

The comparison below is manual. It separates the attested oblique form
from the later remodeled nominative line.

{\def\LTcaptype{none} % do not increment counter
\begin{longtable}[]{@{}
  >{\raggedright\arraybackslash}p{(\linewidth - 8\tabcolsep) * \real{0.2000}}
  >{\raggedright\arraybackslash}p{(\linewidth - 8\tabcolsep) * \real{0.2000}}
  >{\raggedright\arraybackslash}p{(\linewidth - 8\tabcolsep) * \real{0.2000}}
  >{\raggedright\arraybackslash}p{(\linewidth - 8\tabcolsep) * \real{0.2000}}
  >{\raggedright\arraybackslash}p{(\linewidth - 8\tabcolsep) * \real{0.2000}}@{}}
\toprule\noalign{}
\begin{minipage}[b]{\linewidth}\raggedright
PGmc cell / interpretation
\end{minipage} & \begin{minipage}[b]{\linewidth}\raggedright
Candidate input
\end{minipage} & \begin{minipage}[b]{\linewidth}\raggedright
Expected or documented OE outcome
\end{minipage} & \begin{minipage}[b]{\linewidth}\raggedright
OE comparison form
\end{minipage} & \begin{minipage}[b]{\linewidth}\raggedright
Result
\end{minipage} \\
\midrule\noalign{}
\endhead
\bottomrule\noalign{}
\endlastfoot
citation in-stem headword & *furxtīn & broader noun-class label & wider
family context & useful lexeme label, but not the selected cell \\
remodeled nominative line & nominative in-stem forms & \emph{fyrhtu} /
\emph{fyrhto} type lemma forms & fyrhtu / fyrhto & genuine OE evidence,
but later remodeled \\
selected oblique singular & *fúrxtīnaz & compact-trace output:
\emph{fyrhte} & fyrhte & exact match between input, output, and attested
cell \\
\end{longtable}
}

The oblique in-stem form is the relevant comparison form. It yields
attested \emph{fyrhte} directly, while the more familiar nominative
forms belong to a later analogical layer.

\subsubsection{hammer --- OE hameres}\label{hammer-oe-hameres}

Derivation: citation reconstruction \emph{*xámaraz}; selected input
\emph{*xámaras} \textgreater{} \emph{hameres} (late analogy).

\paragraph{Derivation trace}\label{derivation-trace-45}

Proto input: \emph{*xámaras}

\begingroup
\setlength{\fboxsep}{6pt}

\noindent\fbox{%
\begin{minipage}{0.97\linewidth}
\small
\begin{tabularx}{\linewidth}{@{}>{\raggedright\arraybackslash}p{0.300\linewidth}>{\centering\arraybackslash}X>{\raggedright\arraybackslash}p{0.540\linewidth}@{}}
\begin{minipage}[t]{\linewidth}
\centering\textbf{Earlier Germanic changes}\par
\vspace{0.35em}
\raggedright
\centering\textbf{West Germanic}\par
\raggedright
\vspace{0.2em}
\raggedright [no change]\par
\vspace{0.6em}
\centering\textbf{Northwest Germanic}\par
\raggedright
\vspace{0.2em}
\raggedright [no change]\par
\end{minipage}
&
&
\begin{minipage}[t]{\linewidth}
\centering\textbf{Old English changes}\par
\vspace{0.35em}
\raggedright
\begin{tabular}{@{}>{\raggedright\arraybackslash}p{0.70\linewidth}@{\hspace{0.55em}}>{\raggedright\arraybackslash}p{0.20\linewidth}@{\hspace{0.25em}}}
Anglo Frisian Brightening & \emph{*xámæræs} \\
OE Unstressed AE Merger & \emph{*xámeres} \\
\end{tabular}
\end{minipage}
\\
\end{tabularx}
\end{minipage}%
} \endgroup

Outcome: \emph{hameres}

\paragraph{Reconstruction and comparative
evidence}\label{reconstruction-and-comparative-evidence-45}

The inherited noun is the masculine a-stem \emph{*xámaraz}, reflected in
Old English citation forms such as \emph{hamor} and \emph{hamer}
(\citeproc{ref-Kroonen2013}{Kroonen 2013, 206};
\citeproc{ref-Orel2003}{Orel 2003, 197};
\citeproc{ref-ClarkHall1960}{Clark Hall 1960, 160}). The selected input
\emph{*xámaras} is the genitive singular of that same noun rather than a
different lexeme.

The genitive matters because the citation tradition is already mixed in
its unstressed vowel, while the oblique singular gives a cleaner
comparison form. This is a cell choice within one paradigm, not a change
of stem class.

\paragraph{Old English evidence}\label{old-english-evidence-45}

Bosworth-Toller directly records \emph{hameres} in an Old English
genitival phrase (\citeproc{ref-BosworthToller1898}{Bosworth and Toller
1898}). The same dictionary tradition and Clark Hall also preserve the
simplex headword as \emph{hamor} or \emph{hamer}
(\citeproc{ref-BosworthToller1898}{Bosworth and Toller 1898};
\citeproc{ref-ClarkHall1960}{Clark Hall 1960, 160}).

Sievers-Brunner gives a paradigm line \emph{hamor --- hamores}, which
shows that the oblique tradition itself was not entirely uniform
(\citeproc{ref-SieversBrunner1965}{Sievers 1965, sect. 245}). The
relevant comparison form here is the attested genitive singular
\emph{hameres}.

\paragraph{Development to Old
English}\label{development-to-old-english-45}

From \emph{*xámaras}, Anglo-Frisian brightening and the subsequent
merger of unstressed \emph{æ} with \emph{e} yield \emph{hameres}. The
derivation of that oblique form is straightforward once the genitive
singular cell is selected.

The noun as a whole retains a mixed citation tradition in \emph{hamor}
and \emph{hamer}, and the selected oblique cell avoids making that
variation carry the argument of the entry.

\paragraph{Paradigm comparison}\label{paradigm-comparison-6}

The comparison below is manual. It separates the attested genitive
singular from the less stable citation tradition.

{\def\LTcaptype{none} % do not increment counter
\begin{longtable}[]{@{}
  >{\raggedright\arraybackslash}p{(\linewidth - 8\tabcolsep) * \real{0.2000}}
  >{\raggedright\arraybackslash}p{(\linewidth - 8\tabcolsep) * \real{0.2000}}
  >{\raggedright\arraybackslash}p{(\linewidth - 8\tabcolsep) * \real{0.2000}}
  >{\raggedright\arraybackslash}p{(\linewidth - 8\tabcolsep) * \real{0.2000}}
  >{\raggedright\arraybackslash}p{(\linewidth - 8\tabcolsep) * \real{0.2000}}@{}}
\toprule\noalign{}
\begin{minipage}[b]{\linewidth}\raggedright
PGmc cell / interpretation
\end{minipage} & \begin{minipage}[b]{\linewidth}\raggedright
Candidate input
\end{minipage} & \begin{minipage}[b]{\linewidth}\raggedright
Expected or documented OE outcome
\end{minipage} & \begin{minipage}[b]{\linewidth}\raggedright
OE comparison form
\end{minipage} & \begin{minipage}[b]{\linewidth}\raggedright
Result
\end{minipage} \\
\midrule\noalign{}
\endhead
\bottomrule\noalign{}
\endlastfoot
citation nominative singular & *xámaraz & regular citation form
\emph{hamer} / \emph{hamor} & hamor / hamer & good lexical background,
but not the selected target \\
selected genitive singular & *xámaras & compact-trace output:
\emph{hameres} & hameres & exact match between input, output, and
attested cell \\
later oblique tradition & oblique a-stem forms & \emph{hamores} type
evidence & hamores & attested background variant, but not the chosen
comparison form \\
\end{longtable}
}

\subsubsection{have --- OE hæfeþ}\label{have-oe-huxe6feuxfe}

Derivation: citation reconstruction \emph{*xabēną}; selected input
\emph{*xábēθi} \textgreater{} \emph{hæfeþ} (late analogy).

\paragraph{Derivation trace}\label{derivation-trace-46}

Proto input: \emph{*xábēθi}

\begingroup
\setlength{\fboxsep}{6pt}

\noindent\fbox{%
\begin{minipage}{0.97\linewidth}
\small
\begin{tabularx}{\linewidth}{@{}>{\raggedright\arraybackslash}p{0.460\linewidth}>{\centering\arraybackslash}X>{\raggedright\arraybackslash}p{0.460\linewidth}@{}}
\begin{minipage}[t]{\linewidth}
\centering\textbf{Earlier Germanic changes}\par
\vspace{0.35em}
\raggedright
\centering\textbf{West Germanic}\par
\raggedright
\vspace{0.2em}
\begin{tabular}{@{}>{\raggedright\arraybackslash}p{0.74\linewidth}@{\hspace{0.55em}}>{\raggedright\arraybackslash}p{0.16\linewidth}@{\hspace{0.25em}}}
PWGmc Early I Apocope & \emph{*xábēθ} \\
\end{tabular}
\vspace{0.6em}
\centering\textbf{Northwest Germanic}\par
\raggedright
\vspace{0.2em}
\begin{tabular}{@{}>{\raggedright\arraybackslash}p{0.74\linewidth}@{\hspace{0.55em}}>{\raggedright\arraybackslash}p{0.16\linewidth}@{\hspace{0.25em}}}
\mbox{NWGmc Long E Lowering} & \emph{*xábǣθ} \\
\end{tabular}
\end{minipage}
&
&
\begin{minipage}[t]{\linewidth}
\centering\textbf{Old English changes}\par
\vspace{0.35em}
\raggedright
\begin{tabular}{@{}>{\raggedright\arraybackslash}p{0.74\linewidth}@{\hspace{0.55em}}>{\raggedright\arraybackslash}p{0.16\linewidth}@{\hspace{0.25em}}}
Anglo Frisian Brightening & \emph{*xæbǣθ} \\
OE Velar Fricative Palatalization & \emph{*çæbǣθ} \\
PGmc B Allophony & \emph{*çæβǣθ} \\
OE Unstressed Long Vowel Shortening & \emph{*çæβæθ} \\
OE Unstressed AE Merger & \emph{*çæβeθ} \\
\end{tabular}
\end{minipage}
\\
\end{tabularx}
\end{minipage}%
} \endgroup

Outcome: \emph{hæfeþ}

\paragraph{Reconstruction and comparative
evidence}\label{reconstruction-and-comparative-evidence-46}

The verb belongs to the inherited class-III weak paradigm usually cited
under \emph{*xabēną} and Old English \emph{habban}
(\citeproc{ref-Kroonen2013}{Kroonen 2013};
\citeproc{ref-RingeTaylor2014}{Ringe and Taylor 2014, 93}). Within that
paradigm, however, the infinitive and the singular present indicative do
not continue the same stem. Ringe and Taylor distinguish a \emph{-ja-}
stem in the infinitive from a non-geminating -ai- / \emph{-ē-} stem in
the 2sg and 3sg present forms (\citeproc{ref-RingeTaylor2014}{Ringe and
Taylor 2014, 93}).

The selected input \emph{*xábēθi} is therefore the 3sg present cell
rather than a rephrasing of the infinitive. Fulk's discussion of
\emph{habban} treats the ordinary citation form as analogically leveled,
which is why the finite cell is the cleaner comparator here
(\citeproc{ref-Fulk2018}{Fulk 2018}).

\paragraph{Old English evidence}\label{old-english-evidence-46}

The ordinary Old English headword is \emph{habban}, while the present
paradigm also shows forms of the \emph{hæf-} type
(\citeproc{ref-BosworthToller1898}{Bosworth and Toller 1898};
\citeproc{ref-ClarkHall1960}{Clark Hall 1960, 157}). Campbell's Anglian
paradigm includes unsyncopated 3sg forms of the \emph{hæfed} type,
supporting the normalized target \emph{hæfeþ}
(\citeproc{ref-Campbell1959}{Campbell 1959, sect. 762}).

The target form is therefore a normalized finite cell rather than the
ordinary dictionary lemma. It represents the inherited non-geminating
present stem more directly than \emph{habban} does.

\paragraph{Development to Old
English}\label{development-to-old-english-46}

From \emph{*xábēθi}, the finite form yields \emph{hæfeþ} regularly.
Ringe and Taylor discuss this non-geminating present stem under
\emph{habban} (\citeproc{ref-RingeTaylor2014}{Ringe and Taylor 2014,
364}). Campbell's Anglian paradigms include unsyncopated 3sg forms of
the \emph{hæfeþ} / \emph{hæfed} type
(\citeproc{ref-Campbell1959}{Campbell 1959, sect. 762}).

The wider lexeme is less straightforward only because the infinitive
\emph{habban} shows later leveling. That difference in paradigm history
is what makes the 3sg present cell the more useful comparison form.

\paragraph{Paradigm comparison}\label{paradigm-comparison-7}

The comparison below is manual. It separates the analogically leveled
citation form from the regular 3sg present line.

{\def\LTcaptype{none} % do not increment counter
\begin{longtable}[]{@{}
  >{\raggedright\arraybackslash}p{(\linewidth - 8\tabcolsep) * \real{0.2000}}
  >{\raggedright\arraybackslash}p{(\linewidth - 8\tabcolsep) * \real{0.2000}}
  >{\raggedright\arraybackslash}p{(\linewidth - 8\tabcolsep) * \real{0.2000}}
  >{\raggedright\arraybackslash}p{(\linewidth - 8\tabcolsep) * \real{0.2000}}
  >{\raggedright\arraybackslash}p{(\linewidth - 8\tabcolsep) * \real{0.2000}}@{}}
\toprule\noalign{}
\begin{minipage}[b]{\linewidth}\raggedright
PGmc cell / interpretation
\end{minipage} & \begin{minipage}[b]{\linewidth}\raggedright
Candidate input
\end{minipage} & \begin{minipage}[b]{\linewidth}\raggedright
Expected or documented OE outcome
\end{minipage} & \begin{minipage}[b]{\linewidth}\raggedright
OE comparison form
\end{minipage} & \begin{minipage}[b]{\linewidth}\raggedright
Result
\end{minipage} \\
\midrule\noalign{}
\endhead
\bottomrule\noalign{}
\endlastfoot
citation infinitive & \emph{-ja-} stem of \emph{*xabēną} & citation form
\emph{habban} & habban & important headword, but shaped by later
leveling \\
selected 3sg present & *xábēθi & compact-trace output: \emph{hæfeþ} &
hæfeþ & exact match between input, output, and selected finite cell \\
syncopated finite tradition & same present stem & \emph{hæfþ} type
evidence & hæfþ & genuine later OE finite form, but not the normalized
target used here \\
\end{longtable}
}

\subsubsection{heaven --- OE heofon}\label{heaven-oe-heofon}

Derivation: citation reconstruction \emph{*xémenaz}; selected input
\emph{*xémonų} \textgreater{} \emph{heofon} (late analogy).

\paragraph{Derivation trace}\label{derivation-trace-47}

Proto input: \emph{*xémonų}

\begingroup
\setlength{\fboxsep}{6pt}

\noindent\fbox{%
\begin{minipage}{0.97\linewidth}
\small
\begin{tabularx}{\linewidth}{@{}>{\raggedright\arraybackslash}p{0.320\linewidth}>{\centering\arraybackslash}X>{\raggedright\arraybackslash}p{0.520\linewidth}@{}}
\begin{minipage}[t]{\linewidth}
\centering\textbf{Earlier Germanic changes}\par
\vspace{0.35em}
\raggedright
\centering\textbf{West Germanic}\par
\raggedright
\vspace{0.2em}
\raggedright [no change]\par
\vspace{0.6em}
\centering\textbf{Northwest Germanic}\par
\raggedright
\vspace{0.2em}
\begin{tabular}{@{}>{\raggedright\arraybackslash}p{0.70\linewidth}@{\hspace{0.55em}}>{\raggedright\arraybackslash}p{0.16\linewidth}@{\hspace{0.25em}}}
NWGmc Unstressed O Raising & \emph{*xémunų} \\
NWGmc Mn Dissimilation & \emph{*xéβunų} \\
\end{tabular}
\end{minipage}
&
&
\begin{minipage}[t]{\linewidth}
\centering\textbf{Old English changes}\par
\vspace{0.35em}
\raggedright
\begin{tabular}{@{}>{\raggedright\arraybackslash}p{0.70\linewidth}@{\hspace{0.55em}}>{\raggedright\arraybackslash}p{0.20\linewidth}@{\hspace{0.25em}}}
OE Med Unstressed U Lowering & \emph{*xéβonų} \\
OE Velar Fricative Palatalization & \emph{*çéβonų} \\
OE Back Mutation & \emph{*çéoβonų} \\
\mbox{OE High Vowel Apocope} & \emph{*çéoβon} \\
\end{tabular}
\end{minipage}
\\
\end{tabularx}
\end{minipage}%
} \endgroup

Outcome: \emph{heofon}

\paragraph{Reconstruction and comparative
evidence}\label{reconstruction-and-comparative-evidence-47}

The inherited noun belongs to the mn-stem family cited by Kroonen as
\emph{*hemina-} \textasciitilde{} \emph{*hemna-}
(\citeproc{ref-Kroonen2013}{Kroonen 2013, 220}). The selected input
\emph{*xémonų} is an oblique singular form within that paradigm rather
than the lexeme-level citation form \emph{*xémenaz}.

That difference matters for the West Saxon target. Ringe and Taylor give
northern WGmc \emph{*hebun} \textgreater{} West Saxon and Northumbrian
\emph{heofon}, Mercian \emph{heofen}
(\citeproc{ref-RingeTaylor2014}{Ringe and Taylor 2014, 324}). Campbell
likewise gives \emph{heofon} beside \emph{hefen} in the same West-Saxon
\emph{u}-umlaut environment (\citeproc{ref-Campbell1959}{Campbell 1959,
sect. 210.1}).

\paragraph{Old English evidence}\label{old-english-evidence-47}

Old English dictionaries record the standard West Saxon noun as
\emph{heofon}, alongside Anglian or Mercian \emph{hefen} material
(\citeproc{ref-ClarkHall1960}{Clark Hall 1960};
\citeproc{ref-BosworthToller1898}{Bosworth and Toller 1898, 43}).
Campbell also cites an earlier stage \emph{hefzen} in the history of the
word (\citeproc{ref-Campbell1959}{Campbell 1959, sect. 381}).

The target of this entry is the West Saxon citation form \emph{heofon}.
Its vowel history points toward the oblique stem rather than the
front-vocalic nominative line.

\paragraph{Development to Old
English}\label{development-to-old-english-47}

From \emph{*xémonų}, the relevant pathway includes early
\emph{o}-raising before \emph{u}, dissimilation in the m/n cluster, and
later lowering of unstressed \emph{u} (\citeproc{ref-Fulk2018}{Fulk
2018}; \citeproc{ref-Campbell1959}{Campbell 1959}). Back mutation then
yields \emph{heo-} before the labial plus back-vowel sequence, and loss
of the final high vowel gives \emph{heofon}.

The front-vocalic nominative line remains important as background
because it explains the dialectal \emph{hefen} type. West Saxon
\emph{heofon} reflects the oblique stem that was generalized into the
nominative position.

\paragraph{Paradigm comparison}\label{paradigm-comparison-8}

The comparison below is manual. It distinguishes the front-vocalic
nominative line from the oblique stem selected for West Saxon
\emph{heofon}.

{\def\LTcaptype{none} % do not increment counter
\begin{longtable}[]{@{}
  >{\raggedright\arraybackslash}p{(\linewidth - 8\tabcolsep) * \real{0.2000}}
  >{\raggedright\arraybackslash}p{(\linewidth - 8\tabcolsep) * \real{0.2000}}
  >{\raggedright\arraybackslash}p{(\linewidth - 8\tabcolsep) * \real{0.2000}}
  >{\raggedright\arraybackslash}p{(\linewidth - 8\tabcolsep) * \real{0.2000}}
  >{\raggedright\arraybackslash}p{(\linewidth - 8\tabcolsep) * \real{0.2000}}@{}}
\toprule\noalign{}
\begin{minipage}[b]{\linewidth}\raggedright
PGmc cell / interpretation
\end{minipage} & \begin{minipage}[b]{\linewidth}\raggedright
Candidate input
\end{minipage} & \begin{minipage}[b]{\linewidth}\raggedright
Expected or documented OE outcome
\end{minipage} & \begin{minipage}[b]{\linewidth}\raggedright
OE comparison form
\end{minipage} & \begin{minipage}[b]{\linewidth}\raggedright
Result
\end{minipage} \\
\midrule\noalign{}
\endhead
\bottomrule\noalign{}
\endlastfoot
citation nominative singular & *xémenaz & front-vocalic \emph{hefen}
type outcome & hefen / heofen & useful control, but not the selected
West Saxon target \\
selected oblique singular & *xémonų & compact-trace output:
\emph{heofon} & heofon & exact match between input, output, and
target \\
older pre-OE stage & inherited oblique line & earlier \emph{hefzen}
stage & hefzen & historical background for the same West Saxon
development \\
\end{longtable}
}

\subsubsection{live --- OE lifeþ}\label{live-oe-lifeuxfe}

Derivation: citation reconstruction \emph{*libēną}; selected input
\emph{*líbēθi} \textgreater{} \emph{lifeþ} (late analogy).

\paragraph{Derivation trace}\label{derivation-trace-48}

Proto input: \emph{*líbēθi}

\begingroup
\setlength{\fboxsep}{6pt}

\noindent\fbox{%
\begin{minipage}{0.97\linewidth}
\small
\begin{tabularx}{\linewidth}{@{}>{\raggedright\arraybackslash}p{0.460\linewidth}>{\centering\arraybackslash}X>{\raggedright\arraybackslash}p{0.460\linewidth}@{}}
\begin{minipage}[t]{\linewidth}
\centering\textbf{Earlier Germanic changes}\par
\vspace{0.35em}
\raggedright
\centering\textbf{West Germanic}\par
\raggedright
\vspace{0.2em}
\begin{tabular}{@{}>{\raggedright\arraybackslash}p{0.74\linewidth}@{\hspace{0.55em}}>{\raggedright\arraybackslash}p{0.16\linewidth}@{\hspace{0.25em}}}
PWGmc Early I Apocope & \emph{*líbēθ} \\
\end{tabular}
\vspace{0.6em}
\centering\textbf{Northwest Germanic}\par
\raggedright
\vspace{0.2em}
\begin{tabular}{@{}>{\raggedright\arraybackslash}p{0.74\linewidth}@{\hspace{0.55em}}>{\raggedright\arraybackslash}p{0.16\linewidth}@{\hspace{0.25em}}}
\mbox{NWGmc Long E Lowering} & \emph{*líbǣθ} \\
\end{tabular}
\end{minipage}
&
&
\begin{minipage}[t]{\linewidth}
\centering\textbf{Old English changes}\par
\vspace{0.35em}
\raggedright
\begin{tabular}{@{}>{\raggedright\arraybackslash}p{0.74\linewidth}@{\hspace{0.55em}}>{\raggedright\arraybackslash}p{0.16\linewidth}@{\hspace{0.25em}}}
PGmc B Allophony & \emph{*líβǣθ} \\
OE Unstressed Long Vowel Shortening & \emph{*líβæθ} \\
OE Unstressed AE Merger & \emph{*líβeθ} \\
\end{tabular}
\end{minipage}
\\
\end{tabularx}
\end{minipage}%
} \endgroup

Outcome: \emph{lifeþ}

\paragraph{Reconstruction and comparative
evidence}\label{reconstruction-and-comparative-evidence-48}

The inherited verb belongs to the class-III weak family cited by Kroonen
under \emph{*libēn-}, reflected in Old English \emph{libban}
(\citeproc{ref-Kroonen2013}{Kroonen 2013, 336}). Ringe and Taylor show
that the paradigm also contained a separate 3sg present stem, continued
in late Northumbrian \emph{lifed}, which they treat as an archaism
(\citeproc{ref-RingeTaylor2014}{Ringe and Taylor 2014, 364}).

The selected input \emph{*líbēθi} therefore represents a finite present
cell rather than the citation infinitive. That distinction matters
because the ordinary later lemma tradition also includes remodeled forms
such as \emph{lifian}.

\paragraph{Old English evidence}\label{old-english-evidence-48}

The ordinary dictionary headwords are \emph{libban} and, in later
remodeling, \emph{lifian} (\citeproc{ref-BosworthToller1898}{Bosworth
and Toller 1898}; \citeproc{ref-ClarkHall1960}{Clark Hall 1960}). For
this entry, however, the relevant comparison form is the archaic 3sg
present attested as \emph{lifed}, here normalized as \emph{lifeþ}
(\citeproc{ref-RingeTaylor2014}{Ringe and Taylor 2014, 364};
\citeproc{ref-Campbell1959}{Campbell 1959, sect. 762}).

The target is thus a normalized finite form, not the ordinary dictionary
lemma. Its value lies in preserving the older present-stem history more
clearly than the remodeled lemma tradition does.

\paragraph{Development to Old
English}\label{development-to-old-english-48}

From \emph{*líbēθi}, regular reduction of the final syllable and later
weakening of the unstressed vowel yield \emph{lifeþ}. The attested
spelling \emph{lifed} belongs to the same finite form in late
Northumbrian orthography (\citeproc{ref-Campbell1959}{Campbell 1959,
sect. 762}; \citeproc{ref-RingeTaylor2014}{Ringe and Taylor 2014, 364}).

\paragraph{Paradigm comparison}\label{paradigm-comparison-9}

The comparison below is manual. It separates the archaic finite cell
from the ordinary infinitival and later remodeled lemma lines.

{\def\LTcaptype{none} % do not increment counter
\begin{longtable}[]{@{}
  >{\raggedright\arraybackslash}p{(\linewidth - 8\tabcolsep) * \real{0.2000}}
  >{\raggedright\arraybackslash}p{(\linewidth - 8\tabcolsep) * \real{0.2000}}
  >{\raggedright\arraybackslash}p{(\linewidth - 8\tabcolsep) * \real{0.2000}}
  >{\raggedright\arraybackslash}p{(\linewidth - 8\tabcolsep) * \real{0.2000}}
  >{\raggedright\arraybackslash}p{(\linewidth - 8\tabcolsep) * \real{0.2000}}@{}}
\toprule\noalign{}
\begin{minipage}[b]{\linewidth}\raggedright
PGmc cell / interpretation
\end{minipage} & \begin{minipage}[b]{\linewidth}\raggedright
Candidate input
\end{minipage} & \begin{minipage}[b]{\linewidth}\raggedright
Expected or documented OE outcome
\end{minipage} & \begin{minipage}[b]{\linewidth}\raggedright
OE comparison form
\end{minipage} & \begin{minipage}[b]{\linewidth}\raggedright
Result
\end{minipage} \\
\midrule\noalign{}
\endhead
\bottomrule\noalign{}
\endlastfoot
citation infinitive line & *libēną & OE \emph{libban} headword tradition
& libban & establishes the lexeme, but not the selected target \\
selected 3sg present & *líbēθi & compact-trace output: \emph{lifeþ};
attested \emph{lifed} & lifed, normalized here as lifeþ & selected
archaic finite cell \\
later remodeled present tradition & later class-II-type forms &
\emph{lifian} and related finite remodeling & lifian & genuine OE
development, but secondary to the selected cell \\
\end{longtable}
}

\subsubsection{man --- OE mannes}\label{man-oe-mannes}

Derivation: citation reconstruction \emph{*mánnaz}; selected input
\emph{*mánnas} \textgreater{} \emph{mannes} (late analogy).

\paragraph{Derivation trace}\label{derivation-trace-49}

Proto input: \emph{*mánnas}

\begingroup
\setlength{\fboxsep}{6pt}

\noindent\fbox{%
\begin{minipage}{0.97\linewidth}
\small
\begin{tabularx}{\linewidth}{@{}>{\raggedright\arraybackslash}p{0.300\linewidth}>{\centering\arraybackslash}X>{\raggedright\arraybackslash}p{0.540\linewidth}@{}}
\begin{minipage}[t]{\linewidth}
\centering\textbf{Earlier Germanic changes}\par
\vspace{0.35em}
\raggedright
\centering\textbf{West Germanic}\par
\raggedright
\vspace{0.2em}
\raggedright [no change]\par
\vspace{0.6em}
\centering\textbf{Northwest Germanic}\par
\raggedright
\vspace{0.2em}
\raggedright [no change]\par
\end{minipage}
&
&
\begin{minipage}[t]{\linewidth}
\centering\textbf{Old English changes}\par
\vspace{0.35em}
\raggedright
\begin{tabular}{@{}>{\raggedright\arraybackslash}p{0.70\linewidth}@{\hspace{0.55em}}>{\raggedright\arraybackslash}p{0.16\linewidth}@{\hspace{0.25em}}}
Anglo Frisian Brightening & \emph{*mánnæs} \\
OE Unstressed AE Merger & \emph{*mánnes} \\
\end{tabular}
\end{minipage}
\\
\end{tabularx}
\end{minipage}%
} \endgroup

Outcome: \emph{mannes}

\paragraph{Reconstruction and comparative
evidence}\label{reconstruction-and-comparative-evidence-49}

The lexeme-level reconstruction is not uniform. Kroonen cites
\emph{*mannan-}, Orel has \emph{*mannz}, and Ringe and Taylor summarize
the inherited noun as \emph{*mann-} (\citeproc{ref-Kroonen2013}{Kroonen
2013, 354}; \citeproc{ref-Orel2003}{Orel 2003, 299};
\citeproc{ref-RingeTaylor2014}{Ringe and Taylor 2014}). The selected
input \emph{*mánnas} belongs to a different level: it is the
genitive-singular cell chosen for the Old English comparison.

That distinction matters because the target of this entry is not the
ordinary citation form. The selected cell is the one that keeps the
geminate medial before the ending.

\paragraph{Old English evidence}\label{old-english-evidence-49}

Campbell gives the paradigm mann, \emph{man} / \emph{mannes} /
\emph{menn} (\citeproc{ref-Campbell1959}{Campbell 1959}).
Sievers-Brunner likewise cites \emph{man mannes}
(\citeproc{ref-SieversBrunner1965}{Sievers 1965, sect. 226}). He also
explains that word-final simplification underlies forms such as
\emph{man} beside inflected \emph{monnes}
(\citeproc{ref-SieversBrunner1965}{Sievers 1965, sect. 231}). Clark Hall
keeps the dictionary headword under \emph{mann}
(\citeproc{ref-ClarkHall1960}{Clark Hall 1960, 197}).

The relevant comparison form is therefore the attested genitive singular
\emph{mannes}, not the citation lemma \emph{mann}.

\paragraph{Development to Old
English}\label{development-to-old-english-49}

From \emph{*mánnas}, Anglo-Frisian brightening yields \emph{*mánnæs},
and the later unstressed merger gives \emph{*mánnes}, hence
\emph{mannes} (\citeproc{ref-Campbell1959}{Campbell 1959}). In this cell
the geminate remains medial before \emph{-es}. The citation form behaves
differently because word-final gemination was simplified in Old English
(\citeproc{ref-SieversBrunner1965}{Sievers 1965, sect. 231}).

\paragraph{Paradigm comparison}\label{paradigm-comparison-10}

The comparison below is manual. It separates the citation-form line from
the selected genitive singular.

{\def\LTcaptype{none} % do not increment counter
\begin{longtable}[]{@{}
  >{\raggedright\arraybackslash}p{(\linewidth - 8\tabcolsep) * \real{0.2000}}
  >{\raggedright\arraybackslash}p{(\linewidth - 8\tabcolsep) * \real{0.2000}}
  >{\raggedright\arraybackslash}p{(\linewidth - 8\tabcolsep) * \real{0.2000}}
  >{\raggedright\arraybackslash}p{(\linewidth - 8\tabcolsep) * \real{0.2000}}
  >{\raggedright\arraybackslash}p{(\linewidth - 8\tabcolsep) * \real{0.2000}}@{}}
\toprule\noalign{}
\begin{minipage}[b]{\linewidth}\raggedright
PGmc cell / interpretation
\end{minipage} & \begin{minipage}[b]{\linewidth}\raggedright
Candidate input
\end{minipage} & \begin{minipage}[b]{\linewidth}\raggedright
Expected or documented OE outcome
\end{minipage} & \begin{minipage}[b]{\linewidth}\raggedright
OE comparison form
\end{minipage} & \begin{minipage}[b]{\linewidth}\raggedright
Result
\end{minipage} \\
\midrule\noalign{}
\endhead
\bottomrule\noalign{}
\endlastfoot
citation nominative singular & *mannăz & expected citation-form outcome
\emph{man} & mann / monn & establishes the lexeme, but not the selected
target \\
accusative singular & *manną & expected \emph{man} & man & same
word-final simplification as the nominative \\
dative singular & *mannăi & expected \emph{manne} & manne & preserves
medial \emph{nn}, but not the chosen cell \\
selected genitive singular & *mánnas & compact-trace output:
\emph{mannes} & mannes & exact match between input, output, and attested
comparator \\
\end{longtable}
}

\subsubsection{meed --- OE meorde}\label{meed-oe-meorde}

Derivation: citation reconstruction \emph{*mizdō}; selected input
\emph{*mízdai} \textgreater{} \emph{meorde} (late analogy).

\paragraph{Derivation trace}\label{derivation-trace-50}

Proto input: \emph{*mízdai}

\begingroup
\setlength{\fboxsep}{6pt}

\noindent\fbox{%
\begin{minipage}{0.97\linewidth}
\small
\begin{tabularx}{\linewidth}{@{}>{\raggedright\arraybackslash}p{0.440\linewidth}>{\centering\arraybackslash}X>{\raggedright\arraybackslash}p{0.440\linewidth}@{}}
\begin{minipage}[t]{\linewidth}
\centering\textbf{Earlier Germanic changes}\par
\vspace{0.35em}
\raggedright
\centering\textbf{West Germanic}\par
\raggedright
\vspace{0.2em}
\begin{tabular}{@{}>{\raggedright\arraybackslash}p{0.70\linewidth}@{\hspace{0.55em}}>{\raggedright\arraybackslash}p{0.16\linewidth}@{\hspace{0.25em}}}
PWGmc Ai Monophthongization & \emph{*mírdē} \\
\end{tabular}
\vspace{0.6em}
\centering\textbf{Northwest Germanic}\par
\raggedright
\vspace{0.2em}
\begin{tabular}{@{}>{\raggedright\arraybackslash}p{0.70\linewidth}@{\hspace{0.55em}}>{\raggedright\arraybackslash}p{0.16\linewidth}@{\hspace{0.25em}}}
NWGmc I Lowering & \emph{*mérdē} \\
\end{tabular}
\end{minipage}
&
&
\begin{minipage}[t]{\linewidth}
\centering\textbf{Old English changes}\par
\vspace{0.35em}
\raggedright
\begin{tabular}{@{}>{\raggedright\arraybackslash}p{0.74\linewidth}@{\hspace{0.55em}}>{\raggedright\arraybackslash}p{0.16\linewidth}@{\hspace{0.25em}}}
OE Breaking & \emph{*méordē} \\
OE Unstressed Long Vowel Shortening & \emph{*méorde} \\
\end{tabular}
\end{minipage}
\\
\end{tabularx}
\end{minipage}%
} \endgroup

Outcome: \emph{meorde}

\paragraph{Reconstruction and comparative
evidence}\label{reconstruction-and-comparative-evidence-50}

The lexeme-level reconstruction is \emph{*mizdō}, but the selected input
\emph{*mízdai} is a dative-singular cell rather than the citation form.
That distinction is important because the Old English evidence for the
\emph{meord} side is oblique.

The wider history of competing \emph{mēd} remains disputed. Crist,
Kroonen, Ringe and Taylor, and Fulk explain it through some form of
\emph{z}-loss and compensatory lengthening
(\citeproc{ref-Crist2002}{Crist 2002};
\citeproc{ref-Kroonen2013}{Kroonen 2013};
\citeproc{ref-RingeTaylor2014}{Ringe and Taylor 2014};
\citeproc{ref-Fulk2018}{Fulk 2018}), Orel keeps a doublet analysis
(\citeproc{ref-Orel2003}{Orel 2003}), and Kilday instead argues that
West Saxon \emph{mēd} is a Saxono-Frisian loan
(\citeproc{ref-Kilday2024}{Kilday 2024}). The comparison here concerns
the attested oblique line \emph{meorde}.

\paragraph{Old English evidence}\label{old-english-evidence-50}

The directly attested forms are obliques: \emph{meorde} as a dative
singular and \emph{meorda} as a genitive plural
(\citeproc{ref-BrightCassidyRingler1971}{Bright 1971, 328};
\citeproc{ref-BosworthToller1898}{Bosworth and Toller 1898}).
Lexicographers reconstruct a bare nominative \emph{meord} from those
obliques, while West Saxon prose more commonly shows the competing
doublet \emph{mēd} (\citeproc{ref-ClarkHall1960}{Clark Hall 1960};
\citeproc{ref-BosworthToller1898}{Bosworth and Toller 1898}).

The target of this entry is therefore the attested oblique
\emph{meorde}, not the reconstructed lemma \emph{meord} and not the
better-known West Saxon citation form \emph{mēd}.

\paragraph{Development to Old
English}\label{development-to-old-english-50}

From \emph{*mízdai}, rhotacism gives \emph{*mírdai}. Proto-West-Germanic
monophthongization then yields \emph{*mírdē}, Northwest-Germanic
lowering gives \emph{*mérdē}, Old English breaking yields
\emph{*méordē}, and unstressed shortening gives \emph{*méorde}, hence
\emph{meorde} (\citeproc{ref-RingeTaylor2014}{Ringe and Taylor 2014}).

This is the regular oblique-cell path modeled by the current trace. The
entry therefore depends on the attested oblique line rather than on a
full decision about the history of the competing \emph{mēd} tradition.

\paragraph{Paradigm comparison}\label{paradigm-comparison-11}

The comparison below is manual. It distinguishes the attested oblique
target from the broader lemma history.

{\def\LTcaptype{none} % do not increment counter
\begin{longtable}[]{@{}
  >{\raggedright\arraybackslash}p{(\linewidth - 8\tabcolsep) * \real{0.2000}}
  >{\raggedright\arraybackslash}p{(\linewidth - 8\tabcolsep) * \real{0.2000}}
  >{\raggedright\arraybackslash}p{(\linewidth - 8\tabcolsep) * \real{0.2000}}
  >{\raggedright\arraybackslash}p{(\linewidth - 8\tabcolsep) * \real{0.2000}}
  >{\raggedright\arraybackslash}p{(\linewidth - 8\tabcolsep) * \real{0.2000}}@{}}
\toprule\noalign{}
\begin{minipage}[b]{\linewidth}\raggedright
PGmc cell / interpretation
\end{minipage} & \begin{minipage}[b]{\linewidth}\raggedright
Candidate input
\end{minipage} & \begin{minipage}[b]{\linewidth}\raggedright
Expected or documented OE outcome
\end{minipage} & \begin{minipage}[b]{\linewidth}\raggedright
OE comparison form
\end{minipage} & \begin{minipage}[b]{\linewidth}\raggedright
Result
\end{minipage} \\
\midrule\noalign{}
\endhead
\bottomrule\noalign{}
\endlastfoot
citation nominative singular & *mizdō & inferred lemma outcome
\emph{meord} & meord & useful background, but the bare lemma is
reconstructed rather than directly attested \\
selected dative singular & *mízdai & compact-trace output: \emph{meorde}
& meorde & exact match between selected input and attested target \\
genitive singular & *mizdōz & compact-trace output: \emph{meorde} &
meorde & converges on the same attested string, but the dat.sg. has the
clearest direct support \\
genitive plural control & plural oblique line & attested \emph{meorda} &
meorda & confirms the broader oblique tradition, but not the chosen
singular target \\
\end{longtable}
}

\subsubsection{night --- OE niht}\label{night-oe-niht}

Derivation: citation reconstruction \emph{*náxtz}; selected input
\emph{*náxti} \textgreater{} \emph{niht} (late analogy).

\paragraph{Derivation trace}\label{derivation-trace-51}

Proto input: \emph{*náxti}

\begingroup
\setlength{\fboxsep}{6pt}

\noindent\fbox{%
\begin{minipage}{0.97\linewidth}
\small
\begin{tabularx}{\linewidth}{@{}>{\raggedright\arraybackslash}p{0.300\linewidth}>{\centering\arraybackslash}X>{\raggedright\arraybackslash}p{0.540\linewidth}@{}}
\begin{minipage}[t]{\linewidth}
\centering\textbf{Earlier Germanic changes}\par
\vspace{0.35em}
\raggedright
\centering\textbf{West Germanic}\par
\raggedright
\vspace{0.2em}
\raggedright [no change]\par
\vspace{0.6em}
\centering\textbf{Northwest Germanic}\par
\raggedright
\vspace{0.2em}
\raggedright [no change]\par
\end{minipage}
&
&
\begin{minipage}[t]{\linewidth}
\centering\textbf{Old English changes}\par
\vspace{0.35em}
\raggedright
\begin{tabular}{@{}>{\raggedright\arraybackslash}p{0.74\linewidth}@{\hspace{0.55em}}>{\raggedright\arraybackslash}p{0.16\linewidth}@{\hspace{0.25em}}}
Anglo Frisian Brightening & \emph{*næxti} \\
OE Breaking & \emph{*neaxti} \\
OE I Umlaut & \emph{*niexti} \\
OE Ws Palatal Umlaut & \emph{*nixti} \\
\mbox{OE High Vowel Apocope} & \emph{*nixt} \\
\end{tabular}
\end{minipage}
\\
\end{tabularx}
\end{minipage}%
} \endgroup

Outcome: \emph{niht}

\paragraph{Reconstruction and comparative
evidence}\label{reconstruction-and-comparative-evidence-51}

Ringe and Taylor cite gen.sg. \emph{*nahtiz}, dat.sg. \emph{*nahti}, and
nom.pl. \emph{*nahtiz} for the high-vowel side of the paradigm, and
derive West Saxon \emph{niht} from that side
(\citeproc{ref-RingeTaylor2014}{Ringe and Taylor 2014, 240}). The
citation reconstruction \emph{*náxtz} therefore belongs to the
nominative-like headword, while the selected input \emph{*náxti}
represents the dative-singular cell.

That distinction matters because the word later became the model for
endingless datives. Ringe and Taylor explicitly explain forms such as
\emph{dæg} by analogy with dat. sg. \emph{niht} \textless{}
\emph{*nahti} (\citeproc{ref-RingeTaylor2014}{Ringe and Taylor 2014,
380}).

\paragraph{Old English evidence}\label{old-english-evidence-51}

Clark Hall lemmatizes \emph{niht} and records the spelling range
\emph{æ}, e, ea, ie, y, while cross-referencing forms such as
\emph{neaht}, \emph{neht}, and \emph{nieht}
(\citeproc{ref-ClarkHall1960}{Clark Hall 1960}). Campbell likewise
preserves the fluctuation between \emph{neaht} and \emph{niht}, giving
genitive \emph{nihte, nihtes}, dative \emph{niht, nihte}, nominative
plural \emph{niht}, and the contrasting plural-side forms represented by
\emph{neahtas} (\citeproc{ref-Campbell1959}{Campbell 1959, sect.
628.3}).

The comparison form used here is therefore an attested Old English
\emph{niht}, not a reconstructed substitute. The broader lexical record
still preserves the non-umlauted side of the paradigm in
\emph{neaht}-type forms.

\paragraph{Development to Old
English}\label{development-to-old-english-51}

Ringe and Taylor derive West Saxon \emph{niht} from \emph{*nahti} via
\emph{*nehti} and \emph{*neahti} (\citeproc{ref-RingeTaylor2014}{Ringe
and Taylor 2014, 240}). Campbell and Brunner preserve the contrasting
non-umlauted \emph{neaht}-type forms elsewhere in the paradigm
(\citeproc{ref-Campbell1959}{Campbell 1959, sect. 628.3};
\citeproc{ref-SieversBrunner1965}{Sievers 1965, sect. 284}).

The modeled path is therefore \emph{*náxti} \textgreater{} \emph{*næxti}
\textgreater{} \emph{*neaxti} \textgreater{} \emph{*niexti}
\textgreater{} \emph{*nixti} \textgreater{} \emph{niht}.

\paragraph{Paradigm comparison}\label{paradigm-comparison-12}

The comparison below is manual. It identifies the inherited cell that
matches the attested Old English form.

{\def\LTcaptype{none} % do not increment counter
\begin{longtable}[]{@{}
  >{\raggedright\arraybackslash}p{(\linewidth - 8\tabcolsep) * \real{0.2000}}
  >{\raggedright\arraybackslash}p{(\linewidth - 8\tabcolsep) * \real{0.2000}}
  >{\raggedright\arraybackslash}p{(\linewidth - 8\tabcolsep) * \real{0.2000}}
  >{\raggedright\arraybackslash}p{(\linewidth - 8\tabcolsep) * \real{0.2000}}
  >{\raggedright\arraybackslash}p{(\linewidth - 8\tabcolsep) * \real{0.2000}}@{}}
\toprule\noalign{}
\begin{minipage}[b]{\linewidth}\raggedright
PGmc cell / interpretation
\end{minipage} & \begin{minipage}[b]{\linewidth}\raggedright
Candidate input
\end{minipage} & \begin{minipage}[b]{\linewidth}\raggedright
OE output or comparison
\end{minipage} & \begin{minipage}[b]{\linewidth}\raggedright
OE comparison form
\end{minipage} & \begin{minipage}[b]{\linewidth}\raggedright
Result
\end{minipage} \\
\midrule\noalign{}
\endhead
\bottomrule\noalign{}
\endlastfoot
citation nominative singular & *náxtz & expected non-umlauted outcome
\emph{neaht} & neaht & useful background, but not the selected
comparison for \emph{niht} \\
selected dative singular & *náxti & compact-trace output: \emph{niht} &
niht & exact match between input, output, and paradigm cell \\
\end{longtable}
}

\subsubsection{rest --- OE ræste}\label{rest-oe-ruxe6ste}

Derivation: citation reconstruction \emph{*rastō}; selected input
\emph{*rástōz} \textgreater{} \emph{ræste} (late analogy).

\paragraph{Derivation trace}\label{derivation-trace-52}

Proto input: \emph{*rástōz}

\begingroup
\setlength{\fboxsep}{6pt}

\noindent\fbox{%
\begin{minipage}{0.97\linewidth}
\small
\begin{tabularx}{\linewidth}{@{}>{\raggedright\arraybackslash}p{0.460\linewidth}>{\centering\arraybackslash}X>{\raggedright\arraybackslash}p{0.460\linewidth}@{}}
\begin{minipage}[t]{\linewidth}
\centering\textbf{Earlier Germanic changes}\par
\vspace{0.35em}
\raggedright
\centering\textbf{West Germanic}\par
\raggedright
\vspace{0.2em}
\raggedright [no change]\par
\vspace{0.6em}
\centering\textbf{Northwest Germanic}\par
\raggedright
\vspace{0.2em}
\begin{tabular}{@{}>{\raggedright\arraybackslash}p{0.74\linewidth}@{\hspace{0.55em}}>{\raggedright\arraybackslash}p{0.16\linewidth}@{\hspace{0.25em}}}
\mbox{PGmc Final Z Deletion} & \emph{*rástō} \\
\end{tabular}
\end{minipage}
&
&
\begin{minipage}[t]{\linewidth}
\centering\textbf{Old English changes}\par
\vspace{0.35em}
\raggedright
\begin{tabular}{@{}>{\raggedright\arraybackslash}p{0.74\linewidth}@{\hspace{0.55em}}>{\raggedright\arraybackslash}p{0.16\linewidth}@{\hspace{0.25em}}}
PWGmc Surviving Bimoric O Unrounding & \emph{*rástā} \\
Anglo Frisian Brightening & \emph{*ræstǣ} \\
OE Unstressed Long Vowel Shortening & \emph{*ræstæ} \\
OE Unstressed AE Merger & \emph{*ræste} \\
\end{tabular}
\end{minipage}
\\
\end{tabularx}
\end{minipage}%
} \endgroup

Outcome: \emph{ræste}

\paragraph{Reconstruction and comparative
evidence}\label{reconstruction-and-comparative-evidence-52}

Kroonen treats the noun as a feminine ō-stem \emph{*rastō-}, continued
by Old English \emph{ræst} (\citeproc{ref-Kroonen2013}{Kroonen 2013}).
The selected input \emph{*rástōz} therefore does not replace the
lexeme-level headword. It identifies one oblique singular cell on the
side of the paradigm that yields \emph{ræste}.

The source tradition used here labels that cell specifically as genitive
singular, but the broader local synthesis of the ō-stem paradigm shows
that the oblique singulars converge on the same front-vocalic
\emph{ræste} side, in contrast to a nominative singular that would
remain \emph{rast}.

\paragraph{Old English evidence}\label{old-english-evidence-52}

The ordinary Old English citation form is \emph{ræst}
(\citeproc{ref-Kroonen2013}{Kroonen 2013};
\citeproc{ref-ClarkHall1960}{Clark Hall 1960}). Bosworth-Toller also
preserves oblique uses of \emph{ræste}, including prepositional examples
such as on \emph{ræste} and \emph{tó} \emph{ræste}
(\citeproc{ref-BosworthToller1898}{Bosworth and Toller 1898, 121}).

The comparison form used here is therefore an attested oblique
\emph{ræste}, not a reconstructed surrogate. The dictionary headword
\emph{ræst} remains an equally real part of the Old English record.

\paragraph{Development to Old
English}\label{development-to-old-english-52}

Once final \emph{*z} is lost, the selected input moves through the
front-vocalic oblique side of the paradigm rather than the back-vocalic
nominative side. In the modeled derivation, the surviving final long
vowel is first exposed, unrounded, fronted, shortened, and then reduced
to the final \emph{-e} of \emph{ræste}.

The key point is the paradigm split. Nominative \emph{*rastō} yields a
regular \emph{rast}, whereas the selected oblique input yields
\emph{ræste}. The later citation form \emph{ræst} is best explained as
leveling from that oblique \emph{ræst-} stem.

\paragraph{Paradigm comparison}\label{paradigm-comparison-13}

The comparison below is manual. It distinguishes the nominative citation
form from the oblique singular chosen here.

{\def\LTcaptype{none} % do not increment counter
\begin{longtable}[]{@{}
  >{\raggedright\arraybackslash}p{(\linewidth - 8\tabcolsep) * \real{0.2000}}
  >{\raggedright\arraybackslash}p{(\linewidth - 8\tabcolsep) * \real{0.2000}}
  >{\raggedright\arraybackslash}p{(\linewidth - 8\tabcolsep) * \real{0.2000}}
  >{\raggedright\arraybackslash}p{(\linewidth - 8\tabcolsep) * \real{0.2000}}
  >{\raggedright\arraybackslash}p{(\linewidth - 8\tabcolsep) * \real{0.2000}}@{}}
\toprule\noalign{}
\begin{minipage}[b]{\linewidth}\raggedright
PGmc cell / interpretation
\end{minipage} & \begin{minipage}[b]{\linewidth}\raggedright
Candidate input
\end{minipage} & \begin{minipage}[b]{\linewidth}\raggedright
OE output or comparison
\end{minipage} & \begin{minipage}[b]{\linewidth}\raggedright
OE comparison form
\end{minipage} & \begin{minipage}[b]{\linewidth}\raggedright
Result
\end{minipage} \\
\midrule\noalign{}
\endhead
\bottomrule\noalign{}
\endlastfoot
citation nominative singular & *rastō & expected regular outcome
\emph{rast} & ræst & useful background, but not the cell that matches
attested oblique \emph{ræste} \\
selected oblique singular & *rástōz & compact-trace output: \emph{ræste}
& ræste & exact match between selected input and attested OE oblique
form \\
\end{longtable}
}

\subsubsection{shoulder --- OE sċuldrum}\label{shoulder-oe-sux10buldrum}

Derivation: citation reconstruction \emph{*skuldrō}; selected input
\emph{*skúldramiz} \textgreater{} \emph{sċuldrum} (late analogy).

\paragraph{Derivation trace}\label{derivation-trace-53}

Proto input: \emph{*skúldramiz}

\begingroup
\setlength{\fboxsep}{6pt}

\noindent\fbox{%
\begin{minipage}{0.97\linewidth}
\small
\begin{tabularx}{\linewidth}{@{}>{\raggedright\arraybackslash}p{0.560\linewidth}>{\centering\arraybackslash}X>{\raggedright\arraybackslash}p{0.280\linewidth}@{}}
\begin{minipage}[t]{\linewidth}
\centering\textbf{Earlier Germanic changes}\par
\vspace{0.35em}
\raggedright
\centering\textbf{West Germanic}\par
\raggedright
\vspace{0.2em}
\begin{tabular}{@{}>{\raggedright\arraybackslash}p{0.64\linewidth}@{\hspace{0.55em}}>{\raggedright\arraybackslash}p{0.26\linewidth}@{\hspace{0.25em}}}
NWGmc A To U Before M & \emph{*skúldrumiz} \\
PWGmc Early I Apocope & \emph{*skúldrumz} \\
\end{tabular}
\vspace{0.6em}
\centering\textbf{Northwest Germanic}\par
\raggedright
\vspace{0.2em}
\begin{tabular}{@{}>{\raggedright\arraybackslash}p{0.64\linewidth}@{\hspace{0.55em}}>{\raggedright\arraybackslash}p{0.26\linewidth}@{\hspace{0.25em}}}
PGmc Final Z Deletion & \emph{*skúldrum} \\
\end{tabular}
\end{minipage}
&
&
\begin{minipage}[t]{\linewidth}
\centering\textbf{Old English changes}\par
\vspace{0.35em}
\raggedright
\begin{tabular}{@{}>{\raggedright\arraybackslash}p{0.64\linewidth}@{\hspace{0.55em}}>{\raggedright\arraybackslash}p{0.20\linewidth}@{\hspace{0.25em}}}
OE Sk Palatalization & \emph{*ʃúldrum} \\
\end{tabular}
\end{minipage}
\\
\end{tabularx}
\end{minipage}%
} \endgroup

Outcome: \emph{sċuldrum}

\paragraph{Reconstruction and comparative
evidence}\label{reconstruction-and-comparative-evidence-53}

The handbooks do not agree on the reconstruction of the Germanic word.
Orel gives \emph{*skuldr(j)ō}, a feminine ō-/jō-stem, and explicitly
notes that Old English \emph{sculdor} is masculine beside OFrisian
\emph{skulder}, Middle Low German \emph{schulder}, and Old High German
\emph{scultra, scultirra} (\citeproc{ref-Orel2003}{Orel 2003, 345}).
Kroonen reconstructs \emph{*skuldra-}, a masculine a-stem, and derives
the Old High German feminine forms from \emph{*skuldrjōn-}
(\citeproc{ref-Kroonen2013}{Kroonen 2013, 478}). Ringe and Taylor cite
PWGmc \emph{*skuldru} for the Old English branch
(\citeproc{ref-RingeTaylor2014}{Ringe and Taylor 2014, 142}).

These forms imply different stem classes and different expectations for
the Old English inflection. The question is which inflectional cell best
aligns with the Old English evidence.

A dative/instrumental plural form \emph{*skúldramiz} aligns with the
inherited plural ending that later yields Old English \emph{-um}, and it
corresponds directly to the attested dative plural discussed below.

\paragraph{Old English evidence}\label{old-english-evidence-53}

The ordinary Old English headword is \emph{sculdor}. Bosworth-Toller and
Clark Hall both lemmatize \emph{sculdor} as the normal dictionary form,
and the Bosworth-Toller material also preserves plural and oblique forms
such as \emph{sculdru}, \emph{sculdra}, and \emph{sculdrum}
(\citeproc{ref-BosworthToller1898}{Bosworth and Toller 1898};
\citeproc{ref-ClarkHall1960}{Clark Hall 1960, 257}). The dative plural
\emph{sculdrum} is directly attested in the dictionary material
(\citeproc{ref-BosworthToller1898}{Bosworth and Toller 1898}).

Bosworth-Toller's Supplement records a weak-feminine \emph{sculdra, an},
so \emph{sculdra} belongs to the Old English record beside the stronger
masculine paradigm headed by \emph{sculdor}
(\citeproc{ref-BosworthToller1898}{Bosworth and Toller 1898}). Brunner
and Luick also record later spellings such as \emph{sceoldor} and the
i-mutated dative plural \emph{scyldrum}, which reflect secondary
phonological and analogical reshaping within Old English
(\citeproc{ref-SieversBrunner1965}{Sievers 1965, sect. 92.2.a};
\citeproc{ref-Luick1914}{Luick 1914-\/-1940, 230}).

The singular and plural evidence point to different parts of the
paradigm. The relevant comparison form here is the attested dative
plural \emph{sċuldrum}. The spelling with \emph{sċ-} is a normalized
representation of the same Old English initial cluster.

\paragraph{Development to Old
English}\label{development-to-old-english-53}

Proto-Germanic \emph{*skúldramiz} can be interpreted as a
dative/instrumental plural form. In this environment the post-tonic
\emph{a} before \emph{m} is raised to \emph{u}, giving a form of the
\emph{*skúldrumiz} type. Unstressed \emph{u} is regularly preserved
before \emph{m}, especially in the dative plural ending \emph{-um}:
Campbell states this explicitly, and Hogg formulates the same condition
for the dative plural inflexion (\citeproc{ref-Campbell1959}{Campbell
1959, sect. 373}; \citeproc{ref-Hogg1992}{Hogg 1992, sect. 3.3.1.3}).
Brunner points in the same direction by excluding \emph{m} from the
environments in which medial \emph{o} became general in West Saxon
(\citeproc{ref-SieversBrunner1965}{Sievers 1965, sect. 44} Anm. 7).

Subsequent reduction of the ending removes the final \emph{*i} and
\emph{*z}, so that the inflectional ending appears in Old English as
\emph{-um}. The initial cluster is written here as \emph{sċ-}, and the
development is \emph{*skúldramiz} \textgreater{} \emph{*skúldrumiz}
\textgreater{} \emph{*skúldrum} \textgreater{} \emph{sċuldrum}.

\paragraph{Paradigm comparison}\label{paradigm-comparison-14}

A paradigm comparison identifies the Proto-Germanic inflectional cell
that corresponds to an established Old English paradigm form. The
comparison below is manual; no full automatic paradigm-generation run is
presented here.

{\def\LTcaptype{none} % do not increment counter
\begin{longtable}[]{@{}
  >{\raggedright\arraybackslash}p{(\linewidth - 8\tabcolsep) * \real{0.2000}}
  >{\raggedright\arraybackslash}p{(\linewidth - 8\tabcolsep) * \real{0.2000}}
  >{\raggedright\arraybackslash}p{(\linewidth - 8\tabcolsep) * \real{0.2000}}
  >{\raggedright\arraybackslash}p{(\linewidth - 8\tabcolsep) * \real{0.2000}}
  >{\raggedright\arraybackslash}p{(\linewidth - 8\tabcolsep) * \real{0.2000}}@{}}
\toprule\noalign{}
\begin{minipage}[b]{\linewidth}\raggedright
PGmc cell / interpretation
\end{minipage} & \begin{minipage}[b]{\linewidth}\raggedright
Candidate input
\end{minipage} & \begin{minipage}[b]{\linewidth}\raggedright
OE output or comparison
\end{minipage} & \begin{minipage}[b]{\linewidth}\raggedright
OE comparison form
\end{minipage} & \begin{minipage}[b]{\linewidth}\raggedright
Result
\end{minipage} \\
\midrule\noalign{}
\endhead
\bottomrule\noalign{}
\endlastfoot
singular-oriented citation input & *skúldrō & probe output:
\emph{sċoldor} & sculdor & fails: the singular output has root \emph{o},
not the attested \emph{u} \\
serious plural-based singular alternative & *skúldru & probe output:
\emph{sċuldor} & sculdor & close formally, but it compares a
plural-stage input with a singular form \\
dat./inst.pl. input & *skúldramiz & compact-trace output:
\emph{sċuldrum} & sculdrum & matches both the output and the dative
plural comparison form \\
later weak-feminine singular & --- & OE \emph{sculdra} & sculdra &
secondary doublet, useful as a control rather than the inherited
target \\
\end{longtable}
}

The dative plural line is decisive because it matches both the output
and the paradigm cell of Old English \emph{sculdrum}. Singular-oriented
candidates either lower the root vowel or compare unlike cells.

\subsubsection{shove --- OE sċēaf}\label{shove-oe-sux10bux113af}

Derivation: citation reconstruction \emph{*skéubaną}; selected input
\emph{*skáub} \textgreater{} \emph{sċēaf} (late analogy).

\paragraph{Derivation trace}\label{derivation-trace-54}

Proto input: \emph{*skáub}

\begingroup
\setlength{\fboxsep}{6pt}

\noindent\fbox{%
\begin{minipage}{0.97\linewidth}
\small
\begin{tabularx}{\linewidth}{@{}>{\raggedright\arraybackslash}p{0.320\linewidth}>{\centering\arraybackslash}X>{\raggedright\arraybackslash}p{0.520\linewidth}@{}}
\begin{minipage}[t]{\linewidth}
\centering\textbf{Earlier Germanic changes}\par
\vspace{0.35em}
\raggedright
\centering\textbf{West Germanic}\par
\raggedright
\vspace{0.2em}
\raggedright [no change]\par
\vspace{0.6em}
\centering\textbf{Northwest Germanic}\par
\raggedright
\vspace{0.2em}
\raggedright [no change]\par
\end{minipage}
&
&
\begin{minipage}[t]{\linewidth}
\centering\textbf{Old English changes}\par
\vspace{0.35em}
\raggedright
\begin{tabular}{@{}>{\raggedright\arraybackslash}p{0.68\linewidth}@{\hspace{0.55em}}>{\raggedright\arraybackslash}p{0.16\linewidth}@{\hspace{0.25em}}}
OE Au Fronting & \emph{*skáeub} \\
OE Diphthong Leveling & \emph{*skēab} \\
PGmc B Allophony & \emph{*skēaβ} \\
OE Sk Palatalization & \emph{*ʃēaβ} \\
\end{tabular}
\end{minipage}
\\
\end{tabularx}
\end{minipage}%
} \endgroup

Outcome: \emph{sċēaf}

\paragraph{Reconstruction and comparative
evidence}\label{reconstruction-and-comparative-evidence-54}

Kroonen reconstructs the strong verb as \emph{*skeuban-}
\textasciitilde{} \emph{*skūban-} and cites Old English present forms
\emph{scēofan}, \emph{scūfan} (\citeproc{ref-Kroonen2013}{Kroonen
2013}). Ringe and Taylor also show that the English present system
belongs to a wider class-II split that is not identical with the
preterite grade (\citeproc{ref-RingeTaylor2014}{Ringe and Taylor 2014}).
The selected input \emph{*skáub} is therefore not a spelling variant of
the infinitive but the singular preterite cell.

\paragraph{Old English evidence}\label{old-english-evidence-54}

The ordinary dictionary verb is \emph{scūfan}/\emph{scēofan}, but the
preterite itself is well attested. Bright gives the principal parts
\emph{scufan, sceaf, scufon, scofen}
(\citeproc{ref-BrightCassidyRingler1971}{Bright 1971, 347}). Sweet gives
the same paradigm (\citeproc{ref-Sweet1953}{Sweet 1953}). The normalized
form here is \emph{sċēaf}, regularizing the attested spellings
\emph{sceaf} and prefixed \emph{āsceaf}.

\paragraph{Development to Old
English}\label{development-to-old-english-54}

From \emph{*skáub}, the documented trace is straightforward. \emph{*au}
fronts and levels to \emph{ēa}, final \emph{*b} becomes a fricative and
is written \emph{f}, and initial \emph{*sk-} undergoes the usual Old
English palatalized spelling in this environment. The trace therefore
gives \emph{*skáub} \textgreater{} \emph{*skáeub} \textgreater{}
\emph{*skēab} \textgreater{} \emph{*skēaβ} \textgreater{} \emph{sċēaf}.

\paragraph{Paradigm comparison}\label{paradigm-comparison-15}

A paradigm comparison is required here because the ordinary citation
verb and the selected Old English target belong to different cells of
the same strong paradigm. The comparison below is manual.

{\def\LTcaptype{none} % do not increment counter
\begin{longtable}[]{@{}
  >{\raggedright\arraybackslash}p{(\linewidth - 8\tabcolsep) * \real{0.2000}}
  >{\raggedright\arraybackslash}p{(\linewidth - 8\tabcolsep) * \real{0.2000}}
  >{\raggedright\arraybackslash}p{(\linewidth - 8\tabcolsep) * \real{0.2000}}
  >{\raggedright\arraybackslash}p{(\linewidth - 8\tabcolsep) * \real{0.2000}}
  >{\raggedright\arraybackslash}p{(\linewidth - 8\tabcolsep) * \real{0.2000}}@{}}
\toprule\noalign{}
\begin{minipage}[b]{\linewidth}\raggedright
PGmc cell / interpretation
\end{minipage} & \begin{minipage}[b]{\linewidth}\raggedright
Candidate input
\end{minipage} & \begin{minipage}[b]{\linewidth}\raggedright
OE output or comparison
\end{minipage} & \begin{minipage}[b]{\linewidth}\raggedright
OE comparison form
\end{minipage} & \begin{minipage}[b]{\linewidth}\raggedright
Result
\end{minipage} \\
\midrule\noalign{}
\endhead
\bottomrule\noalign{}
\endlastfoot
citation infinitive & *skéubaną & inherited infinitive line
\emph{sċēofan}; present system also leveled \emph{scūfan} & scēofan /
scūfan & necessary background, but not the selected comparison for
\emph{sċēaf} \\
1/3 sg. preterite & *skáub & documented trace output: \emph{sċēaf} &
sċēaf & exact match between selected input and Old English target \\
preterite plural & *skúbun & later leveled plural \emph{scufon} beside
expected \emph{sċufun} under the corrected cascade & scufon & poorer
comparison for the singular-preterite target \\
past participle & *skúbanaz & attested participial line \emph{scofen} &
scofen & valid clean cell, but not the chosen one \\
\end{longtable}
}

\subsubsection{span --- OE spanne}\label{span-oe-spanne}

Derivation: citation reconstruction \emph{*spannō}; selected input
\emph{*spánnai} \textgreater{} \emph{spanne} (late analogy).

\paragraph{Derivation trace}\label{derivation-trace-55}

Proto input: \emph{*spánnai}

\begingroup
\setlength{\fboxsep}{6pt}

\noindent\fbox{%
\begin{minipage}{0.97\linewidth}
\small
\begin{tabularx}{\linewidth}{@{}>{\raggedright\arraybackslash}p{0.300\linewidth}>{\centering\arraybackslash}X>{\raggedright\arraybackslash}p{0.540\linewidth}@{}}
\begin{minipage}[t]{\linewidth}
\centering\textbf{Earlier Germanic changes}\par
\vspace{0.35em}
\raggedright
\centering\textbf{West Germanic}\par
\raggedright
\vspace{0.2em}
\begin{tabular}{@{}>{\raggedright\arraybackslash}p{0.70\linewidth}@{\hspace{0.55em}}>{\raggedright\arraybackslash}p{0.16\linewidth}@{\hspace{0.25em}}}
PWGmc Ai Monophthongization & \emph{*spánnē} \\
\end{tabular}
\vspace{0.6em}
\centering\textbf{Northwest Germanic}\par
\raggedright
\vspace{0.2em}
\raggedright [no change]\par
\end{minipage}
&
&
\begin{minipage}[t]{\linewidth}
\centering\textbf{Old English changes}\par
\vspace{0.35em}
\raggedright
\begin{tabular}{@{}>{\raggedright\arraybackslash}p{0.74\linewidth}@{\hspace{0.55em}}>{\raggedright\arraybackslash}p{0.16\linewidth}@{\hspace{0.25em}}}
OE Unstressed Long Vowel Shortening & \emph{*spánne} \\
\end{tabular}
\end{minipage}
\\
\end{tabularx}
\end{minipage}%
} \endgroup

Outcome: \emph{spanne}

\paragraph{Reconstruction and comparative
evidence}\label{reconstruction-and-comparative-evidence-55}

Orel reconstructs the noun as \emph{*spannō}, and Seebold likewise gives
Old English \emph{spann} under the same noun family
(\citeproc{ref-Orel2003}{Orel 2003}; \citeproc{ref-Seebold1970}{Seebold
1970, 450}). The selected input \emph{*spánnai} is therefore not a rival
headword, but a specific dative singular cell of the feminine ō-stem
paradigm (\citeproc{ref-SieversBrunner1965}{Sievers 1965}).

\paragraph{Old English evidence}\label{old-english-evidence-55}

The reviewed lexicographic evidence more directly supports the citation
noun \emph{spann} than the exact form \emph{spanne}. Clark Hall gives
\emph{spann}, and \emph{spanne} is accordingly treated as the selected
regular dative singular comparison form rather than as a dictionary
headword (\citeproc{ref-ClarkHall1960}{Clark Hall 1960}).

\paragraph{Development to Old
English}\label{development-to-old-english-55}

Citation \emph{*spannō} yields \emph{span}. The oblique cell
\emph{*spánnai} therefore supplies the conservative comparison form: it
preserves the medial geminate and yields \emph{spanne}, while citation
\emph{*spannō} gives the nominative background form.

\paragraph{Paradigm comparison}\label{paradigm-comparison-16}

The comparison below is manual. It shows the contrast between the
citation form and the selected dative singular cell.

{\def\LTcaptype{none} % do not increment counter
\begin{longtable}[]{@{}
  >{\raggedright\arraybackslash}p{(\linewidth - 8\tabcolsep) * \real{0.2000}}
  >{\raggedright\arraybackslash}p{(\linewidth - 8\tabcolsep) * \real{0.2000}}
  >{\raggedright\arraybackslash}p{(\linewidth - 8\tabcolsep) * \real{0.2000}}
  >{\raggedright\arraybackslash}p{(\linewidth - 8\tabcolsep) * \real{0.2000}}
  >{\raggedright\arraybackslash}p{(\linewidth - 8\tabcolsep) * \real{0.2000}}@{}}
\toprule\noalign{}
\begin{minipage}[b]{\linewidth}\raggedright
PGmc cell / interpretation
\end{minipage} & \begin{minipage}[b]{\linewidth}\raggedright
Candidate input
\end{minipage} & \begin{minipage}[b]{\linewidth}\raggedright
OE output or comparison
\end{minipage} & \begin{minipage}[b]{\linewidth}\raggedright
OE comparison form
\end{minipage} & \begin{minipage}[b]{\linewidth}\raggedright
Result
\end{minipage} \\
\midrule\noalign{}
\endhead
\bottomrule\noalign{}
\endlastfoot
citation nominative singular & *spannō & compact-trace output:
\emph{span} & spann & useful citation-form background, but not the
selected target \\
selected dative singular & *spánnai & compact-trace output:
\emph{spanne} & spanne & exact match for the chosen conservative cell \\
\end{longtable}
}

\subsubsection{thistle --- OE þistles}\label{thistle-oe-uxfeistles}

Derivation: citation reconstruction \emph{*θéstilaz}; selected input
\emph{*θístilas} \textgreater{} \emph{þistles} (late analogy).

\paragraph{Derivation trace}\label{derivation-trace-56}

Proto input: \emph{*θístilas}

\begingroup
\setlength{\fboxsep}{6pt}

\noindent\fbox{%
\begin{minipage}{0.97\linewidth}
\small
\begin{tabularx}{\linewidth}{@{}>{\raggedright\arraybackslash}p{0.300\linewidth}>{\centering\arraybackslash}X>{\raggedright\arraybackslash}p{0.540\linewidth}@{}}
\begin{minipage}[t]{\linewidth}
\centering\textbf{Earlier Germanic changes}\par
\vspace{0.35em}
\raggedright
\centering\textbf{West Germanic}\par
\raggedright
\vspace{0.2em}
\raggedright [no change]\par
\vspace{0.6em}
\centering\textbf{Northwest Germanic}\par
\raggedright
\vspace{0.2em}
\raggedright [no change]\par
\end{minipage}
&
&
\begin{minipage}[t]{\linewidth}
\centering\textbf{Old English changes}\par
\vspace{0.35em}
\raggedright
\begin{tabular}{@{}>{\raggedright\arraybackslash}p{0.66\linewidth}@{\hspace{0.55em}}>{\raggedright\arraybackslash}p{0.22\linewidth}@{\hspace{0.25em}}}
Anglo Frisian Brightening & \emph{*θístilæs} \\
OE L Adjacent Syncope & \emph{*θístlæs} \\
OE Unstressed AE Merger & \emph{*θístles} \\
\end{tabular}
\end{minipage}
\\
\end{tabularx}
\end{minipage}%
} \endgroup

Outcome: \emph{þistles}

\paragraph{Reconstruction and comparative
evidence}\label{reconstruction-and-comparative-evidence-56}

The comparative tradition is divided. Orel prints \emph{*þe(x)stilaz},
while Kluge-Seebold gives \emph{*þistila-} (\citeproc{ref-Orel2003}{Orel
2003, 458}; \citeproc{ref-KlugeSeebold2011}{Kluge 2011}). The
comparative label \emph{*θéstilaz} therefore remains in view as the
lexeme-level headword, while the selected input \emph{*θístilas} is a
specific genitive singular cell.

\paragraph{Old English evidence}\label{old-english-evidence-56}

The ordinary simplex headword tradition is broken \emph{þistel} /
\emph{ðistel}. Clark Hall gives \emph{ðistel} as the noun headword
(\citeproc{ref-ClarkHall1960}{Clark Hall 1960, 326}). The selected
target here is the genitive singular \emph{þistles}, which preserves the
same stem in an oblique form where the cluster is medial.

\paragraph{Development to Old
English}\label{development-to-old-english-56}

Campbell's discussion of cluster nouns shows the contrast clearly.
Simplex forms often develop a parasite vowel in word-final obstruent +
sonorant clusters, while comparable medial clusters remain unbroken; his
examples include \emph{hrefn}, \emph{tacn}, \emph{wépn}, and \emph{botm}
beside forms with parasitic vowels elsewhere in the same lexical class
(\citeproc{ref-Campbell1959}{Campbell 1959, 151}). The selected genitive
singular \emph{*θístilas} therefore supplies the conservative comparison
form: the cluster is medial and the regular development yields
\emph{þistles}, while the simplex nominative belongs to the broken
headword tradition \emph{þistel}.

\paragraph{Paradigm comparison}\label{paradigm-comparison-17}

The comparison below is manual. It shows the contrast between the
citation form and the selected genitive singular cell.

{\def\LTcaptype{none} % do not increment counter
\begin{longtable}[]{@{}
  >{\raggedright\arraybackslash}p{(\linewidth - 8\tabcolsep) * \real{0.2000}}
  >{\raggedright\arraybackslash}p{(\linewidth - 8\tabcolsep) * \real{0.2000}}
  >{\raggedright\arraybackslash}p{(\linewidth - 8\tabcolsep) * \real{0.2000}}
  >{\raggedright\arraybackslash}p{(\linewidth - 8\tabcolsep) * \real{0.2000}}
  >{\raggedright\arraybackslash}p{(\linewidth - 8\tabcolsep) * \real{0.2000}}@{}}
\toprule\noalign{}
\begin{minipage}[b]{\linewidth}\raggedright
PGmc cell / interpretation
\end{minipage} & \begin{minipage}[b]{\linewidth}\raggedright
Candidate input
\end{minipage} & \begin{minipage}[b]{\linewidth}\raggedright
OE output or comparison
\end{minipage} & \begin{minipage}[b]{\linewidth}\raggedright
OE comparison form
\end{minipage} & \begin{minipage}[b]{\linewidth}\raggedright
Result
\end{minipage} \\
\midrule\noalign{}
\endhead
\bottomrule\noalign{}
\endlastfoot
citation nominative singular & *θéstilaz & computed output: \emph{þistl}
& þistel & useful citation-form background, but not the selected
target \\
selected genitive singular & *θístilas & computed output: \emph{þistles}
& þistles & exact match for the chosen conservative cell \\
\end{longtable}
}

\subsubsection{make (iptv.2sg) --- OE maca}\label{make-iptv.2sg-oe-maca}

Derivation: citation reconstruction \emph{*makōną}; selected input
\emph{*mákô} \textgreater{} \emph{maca} (late analogy).

\paragraph{Derivation trace}\label{derivation-trace-57}

Proto input: \emph{*mákô}

\begingroup
\setlength{\fboxsep}{6pt}

\noindent\fbox{%
\begin{minipage}{0.97\linewidth}
\small
\begin{tabularx}{\linewidth}{@{}>{\raggedright\arraybackslash}p{0.300\linewidth}>{\centering\arraybackslash}X>{\raggedright\arraybackslash}p{0.540\linewidth}@{}}
\begin{minipage}[t]{\linewidth}
\centering\textbf{Earlier Germanic changes}\par
\vspace{0.35em}
\raggedright
\centering\textbf{West Germanic}\par
\raggedright
\vspace{0.2em}
\raggedright [no change]\par
\vspace{0.6em}
\centering\textbf{Northwest Germanic}\par
\raggedright
\vspace{0.2em}
\raggedright [no change]\par
\end{minipage}
&
&
\begin{minipage}[t]{\linewidth}
\centering\textbf{Old English changes}\par
\vspace{0.35em}
\raggedright
\begin{tabular}{@{}>{\raggedright\arraybackslash}p{0.74\linewidth}@{\hspace{0.55em}}>{\raggedright\arraybackslash}p{0.16\linewidth}@{\hspace{0.25em}}}
Anglo Frisian Brightening & \emph{*mækô} \\
OE A Restoration & \emph{*makô} \\
OE Unstressed Long Vowel Shortening & \emph{*maka} \\
\end{tabular}
\end{minipage}
\\
\end{tabularx}
\end{minipage}%
} \endgroup

Outcome: \emph{maca}

\paragraph{Reconstruction and comparative
evidence}\label{reconstruction-and-comparative-evidence-57}

The make-family belongs to the Old English class-II weak verbs. Campbell
cites \emph{lapian, macian} among verbs with restored \emph{a}
(\citeproc{ref-Campbell1959}{Campbell 1959, sect. 159}). Ringe and
Taylor place the Germanic verb in the same class, comparing West
Germanic continuants such as Old Frisian \emph{makia}, Old Saxon
\emph{makon}, and Old High German \emph{mahhon}
(\citeproc{ref-RingeTaylor2014}{Ringe and Taylor 2014, 191}).

The selected input \emph{*mákô} is not the citation form of the lexeme
but a finite paradigm cell. Ringe and Taylor give the class-II weak
imperative singular as -a \textless{} \emph{*-ō}, which makes this cell
the relevant comparison point for the Old English form treated here
(\citeproc{ref-RingeTaylor2014}{Ringe and Taylor 2014, 314}).

\paragraph{Old English evidence}\label{old-english-evidence-57}

The dictionary headword is \emph{macian}
(\citeproc{ref-ClarkHall1960}{Clark Hall 1960, 193}). The selected form
in this entry is therefore not the lemma but the imperative singular
\emph{maca}, chosen as a paradigm form beside the headword \emph{macian}
and the related finite form \emph{macaþ}.

That distinction matters for the comparison. The lexical history of the
verb is still the history of \emph{macian}, but the finite cell isolates
the regular outcome of trimoric \emph{*ō} more cleanly than the citation
form does.

\paragraph{Development to Old
English}\label{development-to-old-english-57}

From \emph{*mákô}, Anglo-Frisian brightening first gives \emph{*mækô}.
Campbell cites \emph{macian} among the class-II verbs with restored
\emph{a} (\citeproc{ref-Campbell1959}{Campbell 1959, sect. 159}). Ringe
and Taylor's class-II weak imperative singular -a \textless{} \emph{*-ō}
then supports the later finite ending
(\citeproc{ref-RingeTaylor2014}{Ringe and Taylor 2014, 314}).

The same development explains why earlier fronted forms of the
\emph{mæċa} type do not control the entry. Once trimoric \emph{*ô} is
treated as a back-vocalic trigger for restoration, the imperative
singular falls into line with the broader \emph{macian} family.

\paragraph{Paradigm comparison}\label{paradigm-comparison-18}

A paradigm comparison identifies which finite cell of the make-family
matches the Old English form chosen here. The comparison below is
manual.

{\def\LTcaptype{none} % do not increment counter
\begin{longtable}[]{@{}
  >{\raggedright\arraybackslash}p{(\linewidth - 8\tabcolsep) * \real{0.2000}}
  >{\raggedright\arraybackslash}p{(\linewidth - 8\tabcolsep) * \real{0.2000}}
  >{\raggedright\arraybackslash}p{(\linewidth - 8\tabcolsep) * \real{0.2000}}
  >{\raggedright\arraybackslash}p{(\linewidth - 8\tabcolsep) * \real{0.2000}}
  >{\raggedright\arraybackslash}p{(\linewidth - 8\tabcolsep) * \real{0.2000}}@{}}
\toprule\noalign{}
\begin{minipage}[b]{\linewidth}\raggedright
PGmc cell / interpretation
\end{minipage} & \begin{minipage}[b]{\linewidth}\raggedright
Candidate input
\end{minipage} & \begin{minipage}[b]{\linewidth}\raggedright
Old English outcome or comparison
\end{minipage} & \begin{minipage}[b]{\linewidth}\raggedright
OE comparison form
\end{minipage} & \begin{minipage}[b]{\linewidth}\raggedright
Result
\end{minipage} \\
\midrule\noalign{}
\endhead
\bottomrule\noalign{}
\endlastfoot
lexeme-level infinitive & \emph{*mákōjaną} & comparative continuation
\emph{macian} & macian & ordinary headword of the verb, but not the
selected finite cell \\
selected imperative singular & \emph{*mákô} & compact-trace output:
\emph{maca} & maca & exact match between input, output, and selected
paradigm form \\
present third singular companion & \emph{*mákōθi} & comparative
companion \emph{macaþ} & macaþ & useful family control, but not the
target of this entry \\
\end{longtable}
}

\subsubsection{make (3sg) --- OE macaþ}\label{make-3sg-oe-macauxfe}

Derivation: citation reconstruction \emph{*makōną}; selected input
\emph{*mákōθi} \textgreater{} \emph{macaþ} (late analogy).

\paragraph{Derivation trace}\label{derivation-trace-58}

Proto input: \emph{*mákōθi}

\begingroup
\setlength{\fboxsep}{6pt}

\noindent\fbox{%
\begin{minipage}{0.97\linewidth}
\small
\begin{tabularx}{\linewidth}{@{}>{\raggedright\arraybackslash}p{0.300\linewidth}>{\centering\arraybackslash}X>{\raggedright\arraybackslash}p{0.540\linewidth}@{}}
\begin{minipage}[t]{\linewidth}
\centering\textbf{Earlier Germanic changes}\par
\vspace{0.35em}
\raggedright
\centering\textbf{West Germanic}\par
\raggedright
\vspace{0.2em}
\begin{tabular}{@{}>{\raggedright\arraybackslash}p{0.64\linewidth}@{\hspace{0.55em}}>{\raggedright\arraybackslash}p{0.16\linewidth}@{\hspace{0.25em}}}
PWGmc Early I Apocope & \emph{*mákōθ} \\
\end{tabular}
\vspace{0.6em}
\centering\textbf{Northwest Germanic}\par
\raggedright
\vspace{0.2em}
\raggedright [no change]\par
\end{minipage}
&
&
\begin{minipage}[t]{\linewidth}
\centering\textbf{Old English changes}\par
\vspace{0.35em}
\raggedright
\begin{tabular}{@{}>{\raggedright\arraybackslash}p{0.70\linewidth}@{\hspace{0.55em}}>{\raggedright\arraybackslash}p{0.16\linewidth}@{\hspace{0.25em}}}
Anglo Frisian Brightening & \emph{*mækōθ} \\
OE A Restoration & \emph{*makōθ} \\
OE Late O Shortening & \emph{*makaθ} \\
\end{tabular}
\end{minipage}
\\
\end{tabularx}
\end{minipage}%
} \endgroup

Outcome: \emph{macaþ}

\paragraph{Reconstruction and comparative
evidence}\label{reconstruction-and-comparative-evidence-58}

Kroonen derives the Old English verb from \emph{*makōjan-} on the
make-family base \emph{*maka-} (\citeproc{ref-Kroonen2013}{Kroonen
2013}). Ringe and Taylor likewise derive Old English \emph{macian} from
PWGmc \emph{*makon} through \emph{*mekojan}
(\citeproc{ref-RingeTaylor2014}{Ringe and Taylor 2014, 191}).

The selected input \emph{*mákōθi} is therefore a finite 3sg cell of the
same family, not the citation form of the verb.

\paragraph{Old English evidence}\label{old-english-evidence-58}

Clark Hall lemmatizes the verb as \emph{macian}
(\citeproc{ref-ClarkHall1960}{Clark Hall 1960}). The relevant comparison
form here is the normalized present-third-singular \emph{macaþ}, set
beside the dictionary headword and the related imperative singular
\emph{maca}.

Campbell's class-II paradigm makes the ordinary 3sg ending \emph{-aþ},
while his dialect survey allows secondary \emph{-e-} spellings in some
traditions (\citeproc{ref-Campbell1959}{Campbell 1959, sect. 356.4};
\citeproc{ref-Campbell1959}{1959, sect. 757}). \emph{Macaþ} is thus the
regular comparison form for the non-\emph{j} 3sg cell.

\paragraph{Development to Old
English}\label{development-to-old-english-58}

After early loss of final \emph{-i}, \emph{*mákōθi} yields
\emph{*mákōθ}. Anglo-Frisian brightening gives \emph{*mækōθ}, but
Campbell lists \emph{macian} among the class-II verbs with restored
\emph{a}, so the stem returns to \emph{mak-} before the ending is
reduced (\citeproc{ref-Campbell1959}{Campbell 1959, sect. 159}).

The ending then follows the ordinary class-II 3sg development.
Campbell's lufas, \emph{-aþ} (\textless{} \emph{-ōsi}, \emph{-ōþi)} and
Ringe and Taylor's discussion of stable \emph{a} in the finite
non-\emph{j} cells point to \emph{*makōθ} \textgreater{} \emph{*makaθ}
\textgreater{} \emph{macaþ} (\citeproc{ref-Campbell1959}{Campbell 1959,
sect. 356.4}; \citeproc{ref-RingeTaylor2014}{Ringe and Taylor 2014,
80}).

\paragraph{Paradigm comparison}\label{paradigm-comparison-19}

The comparison below is manual. It distinguishes the selected 3sg cell
from the make-family lemma and from the companion imperative form.

{\def\LTcaptype{none} % do not increment counter
\begin{longtable}[]{@{}
  >{\raggedright\arraybackslash}p{(\linewidth - 8\tabcolsep) * \real{0.2000}}
  >{\raggedright\arraybackslash}p{(\linewidth - 8\tabcolsep) * \real{0.2000}}
  >{\raggedright\arraybackslash}p{(\linewidth - 8\tabcolsep) * \real{0.2000}}
  >{\raggedright\arraybackslash}p{(\linewidth - 8\tabcolsep) * \real{0.2000}}
  >{\raggedright\arraybackslash}p{(\linewidth - 8\tabcolsep) * \real{0.2000}}@{}}
\toprule\noalign{}
\begin{minipage}[b]{\linewidth}\raggedright
PGmc cell / interpretation
\end{minipage} & \begin{minipage}[b]{\linewidth}\raggedright
Candidate input
\end{minipage} & \begin{minipage}[b]{\linewidth}\raggedright
OE outcome or comparison
\end{minipage} & \begin{minipage}[b]{\linewidth}\raggedright
OE comparison form
\end{minipage} & \begin{minipage}[b]{\linewidth}\raggedright
Result
\end{minipage} \\
\midrule\noalign{}
\endhead
\bottomrule\noalign{}
\endlastfoot
lexeme-level infinitive & \emph{*makōjaną} & dictionary headword
\emph{macian} & macian & family background, but not the selected cell \\
selected 3sg present & \emph{*mákōθi} & trace output \emph{macaþ} &
macaþ & exact match \\
imperative singular companion & \emph{*mákô} & related finite form
\emph{maca} & maca & useful control, but not the target \\
\end{longtable}
}

\subsubsection{bore (iptv.2sg) --- OE bora}\label{bore-iptv.2sg-oe-bora}

Derivation: citation reconstruction \emph{*burōną}; selected input
\emph{*búrô} \textgreater{} \emph{bora} (late analogy).

\paragraph{Derivation trace}\label{derivation-trace-59}

Proto input: \emph{*búrô}

\begingroup
\setlength{\fboxsep}{6pt}

\noindent\fbox{%
\begin{minipage}{0.97\linewidth}
\small
\begin{tabularx}{\linewidth}{@{}>{\raggedright\arraybackslash}p{0.460\linewidth}>{\centering\arraybackslash}X>{\raggedright\arraybackslash}p{0.460\linewidth}@{}}
\begin{minipage}[t]{\linewidth}
\centering\textbf{Earlier Germanic changes}\par
\vspace{0.35em}
\raggedright
\centering\textbf{West Germanic}\par
\raggedright
\vspace{0.2em}
\raggedright [no change]\par
\vspace{0.6em}
\centering\textbf{Northwest Germanic}\par
\raggedright
\vspace{0.2em}
\begin{tabular}{@{}>{\raggedright\arraybackslash}p{0.74\linewidth}@{\hspace{0.55em}}>{\raggedright\arraybackslash}p{0.16\linewidth}@{\hspace{0.25em}}}
\mbox{NWGmc U Lowering} & \emph{*bórô} \\
\end{tabular}
\end{minipage}
&
&
\begin{minipage}[t]{\linewidth}
\centering\textbf{Old English changes}\par
\vspace{0.35em}
\raggedright
\begin{tabular}{@{}>{\raggedright\arraybackslash}p{0.74\linewidth}@{\hspace{0.55em}}>{\raggedright\arraybackslash}p{0.16\linewidth}@{\hspace{0.25em}}}
OE Unstressed Long Vowel Shortening & \emph{*bóra} \\
\end{tabular}
\end{minipage}
\\
\end{tabularx}
\end{minipage}%
} \endgroup

Outcome: \emph{bora}

\paragraph{Reconstruction and comparative
evidence}\label{reconstruction-and-comparative-evidence-59}

Kroonen reconstructs the bore-family verb as \emph{*burojan-} and cites
Old English \emph{borian} among its continuants
(\citeproc{ref-Kroonen2013}{Kroonen 2013}). Ringe and Taylor give the
class-II weak imperative singular as -a \textless{} \emph{*-ō}
(\citeproc{ref-RingeTaylor2014}{Ringe and Taylor 2014, 314}).

The selected input \emph{*búrô} is therefore an imperative cell of the
same family, not the citation form of the verb.

\paragraph{Old English evidence}\label{old-english-evidence-59}

Clark Hall lemmatizes the verb as \emph{borian}
(\citeproc{ref-ClarkHall1960}{Clark Hall 1960, 48}). The comparison form
here is the normalized imperative singular \emph{bora}, used beside the
headword and the related 3sg form \emph{boraþ}.

The imperative is thus a paradigm form rather than a replacement for the
dictionary lemma. It is the most direct Old English comparator for the
non-\emph{j} finite cell represented by \emph{*búrô}.

\paragraph{Development to Old
English}\label{development-to-old-english-59}

Northwest Germanic lowering first gives \emph{*bórô} from \emph{*búrô},
and late shortening of the unstressed long vowel then yields
\emph{*bóra}, whence \emph{bora}.

Ringe and Taylor's class-II imperative singular -a \textless{}
\emph{*-ō} points to exactly this type of outcome
(\citeproc{ref-RingeTaylor2014}{Ringe and Taylor 2014, 314}). The
selected form therefore isolates the regular finite-cell development
more cleanly than the remodelled infinitive does.

\paragraph{Paradigm comparison}\label{paradigm-comparison-20}

The comparison below is manual. It distinguishes the selected imperative
cell from the bore-family lemma and from the companion 3sg form.

{\def\LTcaptype{none} % do not increment counter
\begin{longtable}[]{@{}
  >{\raggedright\arraybackslash}p{(\linewidth - 8\tabcolsep) * \real{0.2000}}
  >{\raggedright\arraybackslash}p{(\linewidth - 8\tabcolsep) * \real{0.2000}}
  >{\raggedright\arraybackslash}p{(\linewidth - 8\tabcolsep) * \real{0.2000}}
  >{\raggedright\arraybackslash}p{(\linewidth - 8\tabcolsep) * \real{0.2000}}
  >{\raggedright\arraybackslash}p{(\linewidth - 8\tabcolsep) * \real{0.2000}}@{}}
\toprule\noalign{}
\begin{minipage}[b]{\linewidth}\raggedright
PGmc cell / interpretation
\end{minipage} & \begin{minipage}[b]{\linewidth}\raggedright
Candidate input
\end{minipage} & \begin{minipage}[b]{\linewidth}\raggedright
OE outcome or comparison
\end{minipage} & \begin{minipage}[b]{\linewidth}\raggedright
OE comparison form
\end{minipage} & \begin{minipage}[b]{\linewidth}\raggedright
Result
\end{minipage} \\
\midrule\noalign{}
\endhead
\bottomrule\noalign{}
\endlastfoot
lexeme-level infinitive & \emph{*burōjaną} & dictionary headword
\emph{borian} & borian & family background, but not the selected cell \\
selected imperative singular & \emph{*búrô} & trace output \emph{bora} &
bora & exact match \\
3sg present companion & \emph{*búrōθi} & related finite form
\emph{boraþ} & boraþ & useful control, but not the target \\
\end{longtable}
}

\subsubsection{bore (3sg) --- OE boraþ}\label{bore-3sg-oe-borauxfe}

Derivation: citation reconstruction \emph{*burōną}; selected input
\emph{*búrōθi} \textgreater{} \emph{boraþ} (late analogy).

\paragraph{Derivation trace}\label{derivation-trace-60}

Proto input: \emph{*búrōθi}

\begingroup
\setlength{\fboxsep}{6pt}

\noindent\fbox{%
\begin{minipage}{0.97\linewidth}
\small
\begin{tabularx}{\linewidth}{@{}>{\raggedright\arraybackslash}p{0.540\linewidth}>{\centering\arraybackslash}X>{\raggedright\arraybackslash}p{0.300\linewidth}@{}}
\begin{minipage}[t]{\linewidth}
\centering\textbf{Earlier Germanic changes}\par
\vspace{0.35em}
\raggedright
\centering\textbf{West Germanic}\par
\raggedright
\vspace{0.2em}
\begin{tabular}{@{}>{\raggedright\arraybackslash}p{0.74\linewidth}@{\hspace{0.55em}}>{\raggedright\arraybackslash}p{0.16\linewidth}@{\hspace{0.25em}}}
PWGmc Early I Apocope & \emph{*búrōθ} \\
\end{tabular}
\vspace{0.6em}
\centering\textbf{Northwest Germanic}\par
\raggedright
\vspace{0.2em}
\begin{tabular}{@{}>{\raggedright\arraybackslash}p{0.74\linewidth}@{\hspace{0.55em}}>{\raggedright\arraybackslash}p{0.16\linewidth}@{\hspace{0.25em}}}
\mbox{NWGmc U Lowering} & \emph{*bórōθ} \\
\end{tabular}
\end{minipage}
&
&
\begin{minipage}[t]{\linewidth}
\centering\textbf{Old English changes}\par
\vspace{0.35em}
\raggedright
\begin{tabular}{@{}>{\raggedright\arraybackslash}p{0.64\linewidth}@{\hspace{0.55em}}>{\raggedright\arraybackslash}p{0.16\linewidth}@{\hspace{0.25em}}}
OE Late O Shortening & \emph{*bóraθ} \\
\end{tabular}
\end{minipage}
\\
\end{tabularx}
\end{minipage}%
} \endgroup

Outcome: \emph{boraþ}

\paragraph{Reconstruction and comparative
evidence}\label{reconstruction-and-comparative-evidence-60}

Kroonen reconstructs the bore-family verb as \emph{*burojan-} and cites
Old English \emph{borian} among its reflexes
(\citeproc{ref-Kroonen2013}{Kroonen 2013}). The selected form isolates
the finite 3sg cell \emph{*búrōθi} rather than the infinitive.

Campbell's class-II pattern lufas, \emph{-aþ} (\textless{} \emph{-ōsi},
\emph{-ōþi)} and Ringe and Taylor's account of stable \emph{a} in the
class-II 2sg and 3sg make this finite cell the relevant comparison form
for the ending (\citeproc{ref-Campbell1959}{Campbell 1959, sect. 356.4};
\citeproc{ref-RingeTaylor2014}{Ringe and Taylor 2014, 80}).

\paragraph{Old English evidence}\label{old-english-evidence-60}

Clark Hall lemmatizes the verb as \emph{borian}
(\citeproc{ref-ClarkHall1960}{Clark Hall 1960}). The relevant comparison
form here is the normalized present-third-singular \emph{boraþ}, used
beside the headword and the imperative singular \emph{bora}.

Campbell's dialect survey allows secondary \emph{-e-} and \emph{-o-}
spellings in 2sg and 3sg class-II forms, but the basic ending remains
\emph{-aþ} (\citeproc{ref-Campbell1959}{Campbell 1959, sect. 757}).
\emph{Boraþ} is therefore the regular comparison form for this
non-\emph{j} 3sg cell.

\paragraph{Development to Old
English}\label{development-to-old-english-60}

Early loss of final \emph{-i} first gives \emph{*búrōθ} from
\emph{*búrōθi}. Northwest Germanic lowering then produces \emph{*bórōθ},
and late shortening of unstressed \emph{ō} yields \emph{*bóraθ}, whence
\emph{boraþ}.

Campbell's class-II ending evidence and Ringe and Taylor's discussion of
stable \emph{a} in the finite non-\emph{j} cells support exactly this
sequence (\citeproc{ref-Campbell1959}{Campbell 1959, sect. 356.4};
\citeproc{ref-RingeTaylor2014}{Ringe and Taylor 2014, 80}).

\paragraph{Paradigm comparison}\label{paradigm-comparison-21}

The comparison below is manual. It distinguishes the selected 3sg cell
from the bore-family lemma and from the companion imperative form.

{\def\LTcaptype{none} % do not increment counter
\begin{longtable}[]{@{}
  >{\raggedright\arraybackslash}p{(\linewidth - 8\tabcolsep) * \real{0.2000}}
  >{\raggedright\arraybackslash}p{(\linewidth - 8\tabcolsep) * \real{0.2000}}
  >{\raggedright\arraybackslash}p{(\linewidth - 8\tabcolsep) * \real{0.2000}}
  >{\raggedright\arraybackslash}p{(\linewidth - 8\tabcolsep) * \real{0.2000}}
  >{\raggedright\arraybackslash}p{(\linewidth - 8\tabcolsep) * \real{0.2000}}@{}}
\toprule\noalign{}
\begin{minipage}[b]{\linewidth}\raggedright
PGmc cell / interpretation
\end{minipage} & \begin{minipage}[b]{\linewidth}\raggedright
Candidate input
\end{minipage} & \begin{minipage}[b]{\linewidth}\raggedright
OE outcome or comparison
\end{minipage} & \begin{minipage}[b]{\linewidth}\raggedright
OE comparison form
\end{minipage} & \begin{minipage}[b]{\linewidth}\raggedright
Result
\end{minipage} \\
\midrule\noalign{}
\endhead
\bottomrule\noalign{}
\endlastfoot
lexeme-level infinitive & \emph{*burōjaną} & dictionary headword
\emph{borian} & borian & family background, but not the selected cell \\
selected 3sg present & \emph{*búrōθi} & trace output \emph{boraþ} &
boraþ & exact match \\
imperative singular companion & \emph{*búrô} & related finite form
\emph{bora} & bora & useful control, but not the target \\
\end{longtable}
}

\subsubsection{learn (iptv.2sg) --- OE
liorna}\label{learn-iptv.2sg-oe-liorna}

Derivation: citation reconstruction \emph{*liznōjaną}; selected input
\emph{*líznô} \textgreater{} \emph{liorna} (late analogy).

\paragraph{Derivation trace}\label{derivation-trace-61}

Proto input: \emph{*líznô}

\begingroup
\setlength{\fboxsep}{6pt}

\noindent\fbox{%
\begin{minipage}{0.97\linewidth}
\small
\begin{tabularx}{\linewidth}{@{}>{\raggedright\arraybackslash}p{0.300\linewidth}>{\centering\arraybackslash}X>{\raggedright\arraybackslash}p{0.540\linewidth}@{}}
\begin{minipage}[t]{\linewidth}
\centering\textbf{Earlier Germanic changes}\par
\vspace{0.35em}
\raggedright
\centering\textbf{West Germanic}\par
\raggedright
\vspace{0.2em}
\raggedright [no change]\par
\vspace{0.6em}
\centering\textbf{Northwest Germanic}\par
\raggedright
\vspace{0.2em}
\raggedright [no change]\par
\end{minipage}
&
&
\begin{minipage}[t]{\linewidth}
\centering\textbf{Old English changes}\par
\vspace{0.35em}
\raggedright
\begin{tabular}{@{}>{\raggedright\arraybackslash}p{0.74\linewidth}@{\hspace{0.55em}}>{\raggedright\arraybackslash}p{0.16\linewidth}@{\hspace{0.25em}}}
OE Breaking & \emph{*líornô} \\
OE Unstressed Long Vowel Shortening & \emph{*líorna} \\
\end{tabular}
\end{minipage}
\\
\end{tabularx}
\end{minipage}%
} \endgroup

Outcome: \emph{liorna}

\paragraph{Reconstruction and comparative
evidence}\label{reconstruction-and-comparative-evidence-61}

Ringe and Taylor give Old English \emph{liornian} \textasciitilde{}
\emph{leornian} from a learn-family base of the \emph{*lizn-} type
(\citeproc{ref-RingeTaylor2014}{Ringe and Taylor 2014, 38}), and Kroonen
likewise keeps the weak verb as \emph{*liznōn-}
(\citeproc{ref-Kroonen2013}{Kroonen 2013, 380}). Fulk cites the same Old
English family from \emph{*liznō-} (\citeproc{ref-Fulk2018}{Fulk 2018}).

The selected input \emph{*líznô} is a finite imperative cell of that
family, not the citation form of the verb.

\paragraph{Old English evidence}\label{old-english-evidence-61}

Clark Hall gives the ordinary headword as \emph{leornian}
(\citeproc{ref-ClarkHall1960}{Clark Hall 1960}). Brunner, however,
explicitly records \emph{leornian, nordh. auch liorna}, and Campbell
notes that beside \emph{leornian} Northumbrian forms with \emph{io}
occur where original \emph{eo} and \emph{io} remain distinct
(\citeproc{ref-SieversBrunner1965}{Sievers 1965, sect. 417} Anm. 10;
\citeproc{ref-Campbell1959}{Campbell 1959, sect. 123} n.~2).

\emph{Liorna} can therefore be treated as an attested Northumbrian
finite form, while \emph{leornian} remains the better-known dictionary
headword.

\paragraph{Development to Old
English}\label{development-to-old-english-61}

The selected form develops regularly as \emph{*líznô} \textgreater{}
\emph{*lírnô} by rhotacism, then \emph{*líornô} by breaking before
\emph{rn}, and finally \emph{*líorna} by late shortening of the
unstressed long vowel.

Campbell's Northumbrian \emph{io} evidence and Ringe and Taylor's
explicit statement that no form of \emph{liornian} stood in an
i-umlauting environment support this stem shape
(\citeproc{ref-Campbell1959}{Campbell 1959, sect. 123} n.~2;
\citeproc{ref-RingeTaylor2014}{Ringe and Taylor 2014, 247}). The
West-Saxon-looking \emph{eo} forms belong to a different dialectal
presentation of the same family.

\paragraph{Paradigm comparison}\label{paradigm-comparison-22}

The comparison below is manual. It distinguishes the selected imperative
cell from the learn-family infinitive and from the companion 3sg form.

{\def\LTcaptype{none} % do not increment counter
\begin{longtable}[]{@{}
  >{\raggedright\arraybackslash}p{(\linewidth - 8\tabcolsep) * \real{0.2000}}
  >{\raggedright\arraybackslash}p{(\linewidth - 8\tabcolsep) * \real{0.2000}}
  >{\raggedright\arraybackslash}p{(\linewidth - 8\tabcolsep) * \real{0.2000}}
  >{\raggedright\arraybackslash}p{(\linewidth - 8\tabcolsep) * \real{0.2000}}
  >{\raggedright\arraybackslash}p{(\linewidth - 8\tabcolsep) * \real{0.2000}}@{}}
\toprule\noalign{}
\begin{minipage}[b]{\linewidth}\raggedright
PGmc cell / interpretation
\end{minipage} & \begin{minipage}[b]{\linewidth}\raggedright
Candidate input
\end{minipage} & \begin{minipage}[b]{\linewidth}\raggedright
OE outcome or comparison
\end{minipage} & \begin{minipage}[b]{\linewidth}\raggedright
OE comparison form
\end{minipage} & \begin{minipage}[b]{\linewidth}\raggedright
Result
\end{minipage} \\
\midrule\noalign{}
\endhead
\bottomrule\noalign{}
\endlastfoot
lexeme-level infinitive & \emph{*liznōjaną} & Northumbrian
\emph{liornian}; dictionary headword often \emph{leornian} & liornian /
leornian & family background, but not the selected cell \\
selected imperative singular & \emph{*líznô} & trace output and
Brunner's Northumbrian \emph{liorna} & liorna & exact match \\
3sg present companion & \emph{*líznōθi} & related finite form
\emph{liornaþ} & liornaþ & useful control, but not the target \\
\end{longtable}
}

\subsubsection{learn (3sg) --- OE
liornaþ}\label{learn-3sg-oe-liornauxfe}

Derivation: citation reconstruction \emph{*liznōjaną}; selected input
\emph{*líznōθi} \textgreater{} \emph{liornaþ} (late analogy).

\paragraph{Derivation trace}\label{derivation-trace-62}

Proto input: \emph{*líznōθi}

\begingroup
\setlength{\fboxsep}{6pt}

\noindent\fbox{%
\begin{minipage}{0.97\linewidth}
\small
\begin{tabularx}{\linewidth}{@{}>{\raggedright\arraybackslash}p{0.440\linewidth}>{\centering\arraybackslash}X>{\raggedright\arraybackslash}p{0.440\linewidth}@{}}
\begin{minipage}[t]{\linewidth}
\centering\textbf{Earlier Germanic changes}\par
\vspace{0.35em}
\raggedright
\centering\textbf{West Germanic}\par
\raggedright
\vspace{0.2em}
\begin{tabular}{@{}>{\raggedright\arraybackslash}p{0.68\linewidth}@{\hspace{0.55em}}>{\raggedright\arraybackslash}p{0.16\linewidth}@{\hspace{0.25em}}}
PWGmc Early I Apocope & \emph{*lírnōθ} \\
\end{tabular}
\vspace{0.6em}
\centering\textbf{Northwest Germanic}\par
\raggedright
\vspace{0.2em}
\raggedright [no change]\par
\end{minipage}
&
&
\begin{minipage}[t]{\linewidth}
\centering\textbf{Old English changes}\par
\vspace{0.35em}
\raggedright
\begin{tabular}{@{}>{\raggedright\arraybackslash}p{0.68\linewidth}@{\hspace{0.55em}}>{\raggedright\arraybackslash}p{0.20\linewidth}@{\hspace{0.25em}}}
OE Breaking & \emph{*líornōθ} \\
OE Late O Shortening & \emph{*líornaθ} \\
\end{tabular}
\end{minipage}
\\
\end{tabularx}
\end{minipage}%
} \endgroup

Outcome: \emph{liornaþ}

\paragraph{Reconstruction and comparative
evidence}\label{reconstruction-and-comparative-evidence-62}

Ringe and Taylor give Old English \emph{liornian} \textasciitilde{}
\emph{leornian} from a learn-family base of the \emph{*lizn-} type
(\citeproc{ref-RingeTaylor2014}{Ringe and Taylor 2014, 38}), and Kroonen
likewise keeps the weak verb as \emph{*liznōn-}
(\citeproc{ref-Kroonen2013}{Kroonen 2013, 380}). The selected input
\emph{*líznōθi} is the finite 3sg cell of that family, not the citation
form of the verb.

For the ending, Campbell's lufas, \emph{-aþ} (\textless{} \emph{-ōsi},
\emph{-ōþi)} and Ringe and Taylor's discussion of stable \emph{a} in the
class-II 2sg and 3sg make the non-\emph{j} 3sg cell the relevant
comparison point (\citeproc{ref-Campbell1959}{Campbell 1959, sect.
356.4}; \citeproc{ref-RingeTaylor2014}{Ringe and Taylor 2014, 80}).

\paragraph{Old English evidence}\label{old-english-evidence-62}

Clark Hall gives the ordinary headword as \emph{leornian}
(\citeproc{ref-ClarkHall1960}{Clark Hall 1960}). Brunner records
Northumbrian finite forms in \emph{liorn-}, including \emph{liorna} and
the 3sg \emph{liornes}, beside the West-Saxon-looking \emph{leornian}
tradition (\citeproc{ref-SieversBrunner1965}{Sievers 1965, sect. 417}
Anm. 10). Campbell likewise notes Northumbrian forms with \emph{io}
beside \emph{leornian} (\citeproc{ref-Campbell1959}{Campbell 1959, sect.
123} n.~2).

The relevant comparison form here is the normalized 3sg \emph{liornaþ}.
The directly cited Old English evidence supports the finite stem
\emph{liorn-}; the exact \emph{-aþ} ending follows the regular class-II
3sg pattern.

\paragraph{Development to Old
English}\label{development-to-old-english-62}

The selected form develops as \emph{*líznōθi} \textgreater{}
\emph{*lírnōθi} by rhotacism, then \emph{*lírnōθ} after early apocope of
final \emph{-i}, then \emph{*líornōθ} by breaking before \emph{rn}, and
finally \emph{*líornaθ} \textgreater{} \emph{liornaþ} by late shortening
of the unstressed long vowel.

Campbell's Northumbrian \emph{io} evidence and Ringe and Taylor's
statement that no form of \emph{liornian} stood in an i-umlauting
environment support the stem, while Campbell's class-II ending evidence
supports the final \emph{-aþ} (\citeproc{ref-Campbell1959}{Campbell
1959, sect. 123} n.~2; \citeproc{ref-Campbell1959}{Campbell 1959, sect.
356.4}; \citeproc{ref-RingeTaylor2014}{Ringe and Taylor 2014, 247}).

\paragraph{Paradigm comparison}\label{paradigm-comparison-23}

The comparison below is manual. It distinguishes the selected 3sg cell
from the learn-family infinitive and from the companion imperative form.

{\def\LTcaptype{none} % do not increment counter
\begin{longtable}[]{@{}
  >{\raggedright\arraybackslash}p{(\linewidth - 8\tabcolsep) * \real{0.2000}}
  >{\raggedright\arraybackslash}p{(\linewidth - 8\tabcolsep) * \real{0.2000}}
  >{\raggedright\arraybackslash}p{(\linewidth - 8\tabcolsep) * \real{0.2000}}
  >{\raggedright\arraybackslash}p{(\linewidth - 8\tabcolsep) * \real{0.2000}}
  >{\raggedright\arraybackslash}p{(\linewidth - 8\tabcolsep) * \real{0.2000}}@{}}
\toprule\noalign{}
\begin{minipage}[b]{\linewidth}\raggedright
PGmc cell / interpretation
\end{minipage} & \begin{minipage}[b]{\linewidth}\raggedright
Candidate input
\end{minipage} & \begin{minipage}[b]{\linewidth}\raggedright
OE outcome or comparison
\end{minipage} & \begin{minipage}[b]{\linewidth}\raggedright
OE comparison form
\end{minipage} & \begin{minipage}[b]{\linewidth}\raggedright
Result
\end{minipage} \\
\midrule\noalign{}
\endhead
\bottomrule\noalign{}
\endlastfoot
lexeme-level infinitive & \emph{*liznōjaną} & Northumbrian
\emph{liornian}; dictionary headword often \emph{leornian} & liornian /
leornian & family background, but not the selected cell \\
selected 3sg present & \emph{*líznōθi} & trace output \emph{liornaþ} &
liornaþ & exact match \\
imperative singular companion & \emph{*líznô} & trace output and
Brunner's Northumbrian \emph{liorna} & liorna & useful control, but not
the target \\
\end{longtable}
}

\subsubsection{lick (iptv.2sg) --- OE
licca}\label{lick-iptv.2sg-oe-licca}

Derivation: citation reconstruction \emph{*likkōną}; selected input
\emph{*líkkô} \textgreater{} \emph{licca} (late analogy).

\paragraph{Derivation trace}\label{derivation-trace-63}

Proto input: \emph{*líkkô}

\begingroup
\setlength{\fboxsep}{6pt}

\noindent\fbox{%
\begin{minipage}{0.97\linewidth}
\small
\begin{tabularx}{\linewidth}{@{}>{\raggedright\arraybackslash}p{0.300\linewidth}>{\centering\arraybackslash}X>{\raggedright\arraybackslash}p{0.540\linewidth}@{}}
\begin{minipage}[t]{\linewidth}
\centering\textbf{Earlier Germanic changes}\par
\vspace{0.35em}
\raggedright
\centering\textbf{West Germanic}\par
\raggedright
\vspace{0.2em}
\raggedright [no change]\par
\vspace{0.6em}
\centering\textbf{Northwest Germanic}\par
\raggedright
\vspace{0.2em}
\raggedright [no change]\par
\end{minipage}
&
&
\begin{minipage}[t]{\linewidth}
\centering\textbf{Old English changes}\par
\vspace{0.35em}
\raggedright
\begin{tabular}{@{}>{\raggedright\arraybackslash}p{0.74\linewidth}@{\hspace{0.55em}}>{\raggedright\arraybackslash}p{0.16\linewidth}@{\hspace{0.25em}}}
OE Unstressed Long Vowel Shortening & \emph{*líkka} \\
\end{tabular}
\end{minipage}
\\
\end{tabularx}
\end{minipage}%
} \endgroup

Outcome: \emph{licca}

\paragraph{Reconstruction and comparative
evidence}\label{reconstruction-and-comparative-evidence-63}

Ringe and Taylor place the verb among the West Germanic class-II weak
verbs, with PWGmc \emph{*li}/\emph{ekkōn} continuing as Old English
\emph{liccian}, Old Saxon \emph{likkon}, and Old High German
\emph{lecchon} (\citeproc{ref-RingeTaylor2014}{Ringe and Taylor 2014}).
Orel gives the fuller weak-verb reconstruction \emph{*likkōjanan} with
the same Old English continuation (\citeproc{ref-Orel2003}{Orel 2003,
285}).

Campbell's weak class-II discussion gives present forms such as lufas,
\emph{-aþ} (\textless{} \emph{-ōsi}, \emph{-ōþi)}
(\citeproc{ref-Campbell1959}{Campbell 1959, sect. 356.4}). Ringe and
Taylor likewise note that class-II weak present 2sg. \emph{-as(t)} and
3sg. \emph{-aþ} have stable \emph{a}
(\citeproc{ref-RingeTaylor2014}{Ringe and Taylor 2014, 80}). The form
treated here is therefore not that remodeled infinitive but a finite
cell in bare trimoric \emph{*-ō}.

\paragraph{Old English evidence}\label{old-english-evidence-63}

Bosworth-Toller lemmatizes the verb as \emph{liccian}
(\citeproc{ref-BosworthToller1898}{Bosworth and Toller 1898}). Campbell
cites \emph{liccian} among Old English forms with preserved geminate
\emph{cc} (\citeproc{ref-Campbell1959}{Campbell 1959, sect. 398.1}).
Brunner likewise cites \emph{liccian}
(\citeproc{ref-SieversBrunner1965}{Sievers 1965, sect. 45} Anm. 3). The
Old English evidence therefore establishes the verbal headword and its
consonantal frame securely.

The selected target in this entry is the imperative singular
\emph{licca}. It is a paradigm form chosen beside the headword
\emph{liccian} and the related present \emph{liccaþ}, not a separately
lemmatized citation word.

\paragraph{Development to Old
English}\label{development-to-old-english-63}

With the stem \emph{licc-} established, the remaining development is
brief. Campbell's class-II present endings lufas, \emph{-aþ}
(\textless{} \emph{-ōsi}, \emph{-ōþi)} support late \emph{-a} in this
finite cell (\citeproc{ref-Campbell1959}{Campbell 1959, sect. 356.4}).
Ringe and Taylor likewise note stable \emph{a} in the class-II 2sg and
3sg (\citeproc{ref-RingeTaylor2014}{Ringe and Taylor 2014, 80}). The
same stem consonantism that appears in \emph{liccian} is preserved here,
giving \emph{cc} throughout the finite form.

\paragraph{Paradigm comparison}\label{paradigm-comparison-24}

The comparison below is manual.

{\def\LTcaptype{none} % do not increment counter
\begin{longtable}[]{@{}
  >{\raggedright\arraybackslash}p{(\linewidth - 8\tabcolsep) * \real{0.2000}}
  >{\raggedright\arraybackslash}p{(\linewidth - 8\tabcolsep) * \real{0.2000}}
  >{\raggedright\arraybackslash}p{(\linewidth - 8\tabcolsep) * \real{0.2000}}
  >{\raggedright\arraybackslash}p{(\linewidth - 8\tabcolsep) * \real{0.2000}}
  >{\raggedright\arraybackslash}p{(\linewidth - 8\tabcolsep) * \real{0.2000}}@{}}
\toprule\noalign{}
\begin{minipage}[b]{\linewidth}\raggedright
PGmc cell / interpretation
\end{minipage} & \begin{minipage}[b]{\linewidth}\raggedright
Candidate input
\end{minipage} & \begin{minipage}[b]{\linewidth}\raggedright
Old English outcome or comparison
\end{minipage} & \begin{minipage}[b]{\linewidth}\raggedright
OE comparison form
\end{minipage} & \begin{minipage}[b]{\linewidth}\raggedright
Result
\end{minipage} \\
\midrule\noalign{}
\endhead
\bottomrule\noalign{}
\endlastfoot
lexeme-level infinitive & \emph{*líkkōjaną} & manual probe output
\emph{liccian} & liccian & ordinary dictionary headword of the verb, but
not the selected finite cell \\
selected imperative singular & \emph{*líkkô} & manual probe output
\emph{licca} & licca & exact match between the chosen input and the
selected target \\
present third singular companion & \emph{*líkkōθi} & manual probe output
\emph{liccaþ} & liccaþ & useful family control, but not the target of
this entry \\
\end{longtable}
}

\subsubsection{lick (3sg) --- OE liccaþ}\label{lick-3sg-oe-liccauxfe}

Derivation: citation reconstruction \emph{*likkōną}; selected input
\emph{*líkkōθi} \textgreater{} \emph{liccaþ} (late analogy).

\paragraph{Derivation trace}\label{derivation-trace-64}

Proto input: \emph{*líkkōθi}

\begingroup
\setlength{\fboxsep}{6pt}

\noindent\fbox{%
\begin{minipage}{0.97\linewidth}
\small
\begin{tabularx}{\linewidth}{@{}>{\raggedright\arraybackslash}p{0.440\linewidth}>{\centering\arraybackslash}X>{\raggedright\arraybackslash}p{0.440\linewidth}@{}}
\begin{minipage}[t]{\linewidth}
\centering\textbf{Earlier Germanic changes}\par
\vspace{0.35em}
\raggedright
\centering\textbf{West Germanic}\par
\raggedright
\vspace{0.2em}
\begin{tabular}{@{}>{\raggedright\arraybackslash}p{0.68\linewidth}@{\hspace{0.55em}}>{\raggedright\arraybackslash}p{0.16\linewidth}@{\hspace{0.25em}}}
PWGmc Early I Apocope & \emph{*líkkōθ} \\
\end{tabular}
\vspace{0.6em}
\centering\textbf{Northwest Germanic}\par
\raggedright
\vspace{0.2em}
\raggedright [no change]\par
\end{minipage}
&
&
\begin{minipage}[t]{\linewidth}
\centering\textbf{Old English changes}\par
\vspace{0.35em}
\raggedright
\begin{tabular}{@{}>{\raggedright\arraybackslash}p{0.68\linewidth}@{\hspace{0.55em}}>{\raggedright\arraybackslash}p{0.16\linewidth}@{\hspace{0.25em}}}
OE Late O Shortening & \emph{*líkkaθ} \\
\end{tabular}
\end{minipage}
\\
\end{tabularx}
\end{minipage}%
} \endgroup

Outcome: \emph{liccaþ}

\paragraph{Reconstruction and comparative
evidence}\label{reconstruction-and-comparative-evidence-64}

Ringe and Taylor place the verb among the West Germanic class-II weak
verbs, with PWGmc \emph{*li}/\emph{ekkōn} continuing as Old English
\emph{liccian}, Old Saxon \emph{likkon}, and Old High German
\emph{lecchon} (\citeproc{ref-RingeTaylor2014}{Ringe and Taylor 2014}).
Orel gives the fuller weak-verb reconstruction \emph{*likkōjanan} with
the same Old English continuation (\citeproc{ref-Orel2003}{Orel 2003,
285}).

The selected form in this entry is the non-\emph{j} present third
singular \emph{*líkkōθi}, not the remodeled infinitive. Campbell states
the class-II present endings as lufas, \emph{-aþ} (\textless{}
\emph{-ōsi}, \emph{-ōþi)} (\citeproc{ref-Campbell1959}{Campbell 1959,
sect. 356.4}). Ringe and Taylor likewise note that class-II weak present
2sg. \emph{-as(t)} and 3sg. \emph{-aþ} have stable \emph{a}
(\citeproc{ref-RingeTaylor2014}{Ringe and Taylor 2014, 80}).

\paragraph{Old English evidence}\label{old-english-evidence-64}

Bosworth-Toller lemmatizes the verb as \emph{liccian}
(\citeproc{ref-BosworthToller1898}{Bosworth and Toller 1898}). The same
consonantal frame appears in Campbell's and Brunner's grammatical
citations of \emph{liccian} (\citeproc{ref-Campbell1959}{Campbell 1959,
sect. 398.1}; \citeproc{ref-SieversBrunner1965}{Sievers 1965, sect. 45}
Anm. 3). The Old English headword is therefore clear even though the
entry here is not about the citation form.

The form treated here is the present third singular \emph{liccaþ}. It is
a selected paradigm form beside the lemma \emph{liccian} and the related
imperative \emph{licca}, not a separately lemmatized headword.

\paragraph{Development to Old
English}\label{development-to-old-english-64}

\emph{*líkkōθi} first loses final \emph{-i}, giving \emph{*líkkōθ}.
Campbell's class-II present endings lufas, \emph{-aþ} (\textless{}
\emph{-ōsi}, \emph{-ōþi)} support the regular 3sg outcome \emph{-aþ}
(\citeproc{ref-Campbell1959}{Campbell 1959, sect. 356.4}). Ringe and
Taylor likewise note stable \emph{a} in the class-II 2sg and 3sg
(\citeproc{ref-RingeTaylor2014}{Ringe and Taylor 2014, 80}). Because
this ending never contains \emph{-j-}, the form does not pass through an
i-umlauted \emph{-eþ} stage.

\paragraph{Paradigm comparison}\label{paradigm-comparison-25}

The comparison below is manual.

{\def\LTcaptype{none} % do not increment counter
\begin{longtable}[]{@{}
  >{\raggedright\arraybackslash}p{(\linewidth - 8\tabcolsep) * \real{0.2000}}
  >{\raggedright\arraybackslash}p{(\linewidth - 8\tabcolsep) * \real{0.2000}}
  >{\raggedright\arraybackslash}p{(\linewidth - 8\tabcolsep) * \real{0.2000}}
  >{\raggedright\arraybackslash}p{(\linewidth - 8\tabcolsep) * \real{0.2000}}
  >{\raggedright\arraybackslash}p{(\linewidth - 8\tabcolsep) * \real{0.2000}}@{}}
\toprule\noalign{}
\begin{minipage}[b]{\linewidth}\raggedright
PGmc cell / interpretation
\end{minipage} & \begin{minipage}[b]{\linewidth}\raggedright
Candidate input
\end{minipage} & \begin{minipage}[b]{\linewidth}\raggedright
Old English outcome or comparison
\end{minipage} & \begin{minipage}[b]{\linewidth}\raggedright
OE comparison form
\end{minipage} & \begin{minipage}[b]{\linewidth}\raggedright
Result
\end{minipage} \\
\midrule\noalign{}
\endhead
\bottomrule\noalign{}
\endlastfoot
lexeme-level infinitive & \emph{*líkkōjaną} & manual probe output
\emph{liccian} & liccian & ordinary dictionary headword of the verb, but
not the selected finite cell \\
imperative singular companion & \emph{*líkkô} & manual probe output
\emph{licca} & licca & useful family control, but not the target of this
entry \\
selected present third singular & \emph{*líkkōθi} & manual probe output
\emph{liccaþ} & liccaþ & exact match between the chosen input and the
selected target \\
\end{longtable}
}

\subsubsection{show (iptv.2sg) --- OE
sċēawa}\label{show-iptv.2sg-oe-sux10bux113awa}

Derivation: citation reconstruction \emph{*skawōną}; selected input
\emph{*skáwô} \textgreater{} \emph{sċēawa} (late analogy).

\paragraph{Derivation trace}\label{derivation-trace-65}

Proto input: \emph{*skáwô}

\begingroup
\setlength{\fboxsep}{6pt}

\noindent\fbox{%
\begin{minipage}{0.97\linewidth}
\small
\begin{tabularx}{\linewidth}{@{}>{\raggedright\arraybackslash}p{0.300\linewidth}>{\centering\arraybackslash}X>{\raggedright\arraybackslash}p{0.540\linewidth}@{}}
\begin{minipage}[t]{\linewidth}
\centering\textbf{Earlier Germanic changes}\par
\vspace{0.35em}
\raggedright
\centering\textbf{West Germanic}\par
\raggedright
\vspace{0.2em}
\raggedright [no change]\par
\vspace{0.6em}
\centering\textbf{Northwest Germanic}\par
\raggedright
\vspace{0.2em}
\raggedright [no change]\par
\end{minipage}
&
&
\begin{minipage}[t]{\linewidth}
\centering\textbf{Old English changes}\par
\vspace{0.35em}
\raggedright
\begin{tabular}{@{}>{\raggedright\arraybackslash}p{0.74\linewidth}@{\hspace{0.55em}}>{\raggedright\arraybackslash}p{0.16\linewidth}@{\hspace{0.25em}}}
OE Aw Long Diphthong & \emph{*skḗawô} \\
OE Sk Palatalization & \emph{*ʃḗawô} \\
OE Unstressed Long Vowel Shortening & \emph{*ʃḗawa} \\
\end{tabular}
\end{minipage}
\\
\end{tabularx}
\end{minipage}%
} \endgroup

Outcome: \emph{sċēawa}

\paragraph{Reconstruction and comparative
evidence}\label{reconstruction-and-comparative-evidence-65}

Orel reconstructs the verb as \emph{*skawōjanan} and cites Old English
\emph{sceáwian} beside Old Frisian \emph{skawia}, Old Saxon
\emph{skawōn}, and Old High German \emph{scouwōn}
(\citeproc{ref-Orel2003}{Orel 2003, 337}). The selected input in this
entry is not that infinitive but the imperative singular \emph{*skáwô},
a finite class-II cell with imperative -a \textless{} \emph{*-ō}
(\citeproc{ref-RingeTaylor2014}{Ringe and Taylor 2014, 314}).

That distinction matters because the imperative singular provides the
direct comparison for the Old English form treated here. The lexical
history still belongs to the \emph{sceáwian} verb, but the selected cell
isolates the finite \emph{-a} outcome more clearly than the citation
form does.

\paragraph{Old English evidence}\label{old-english-evidence-65}

Bright lists \emph{scēawian} and explicitly gives the imperative
singular \emph{scēawa} under that headword
(\citeproc{ref-BrightCassidyRingler1971}{Bright 1971, 346}). The form
treated here is therefore an attested finite paradigm form, not a
reconstructed convenience form.

The spelling used in this entry is normalized \emph{sċēawa}, while
Bright's glossary gives source spelling \emph{scēawa}. The ordinary Old
English headword remains \emph{scēawian}; \emph{sċēawa} is the
imperative singular chosen beside it.

\paragraph{Development to Old
English}\label{development-to-old-english-65}

Campbell lists \emph{scéawian} under the West Germanic \emph{*auw}
developments (\citeproc{ref-Campbell1959}{Campbell 1959, sect. 120}).
Ringe and Taylor's class-II weak imperative singular -a \textless{}
\emph{*-ō} supports the late finite ending that yields \emph{sċēawa}
(\citeproc{ref-RingeTaylor2014}{Ringe and Taylor 2014, 314}). The result
is therefore the expected finite singular form of the \emph{scēawian}
family rather than an analogical replacement of the headword.

\paragraph{Paradigm comparison}\label{paradigm-comparison-26}

The comparison below is manual.

{\def\LTcaptype{none} % do not increment counter
\begin{longtable}[]{@{}
  >{\raggedright\arraybackslash}p{(\linewidth - 8\tabcolsep) * \real{0.2000}}
  >{\raggedright\arraybackslash}p{(\linewidth - 8\tabcolsep) * \real{0.2000}}
  >{\raggedright\arraybackslash}p{(\linewidth - 8\tabcolsep) * \real{0.2000}}
  >{\raggedright\arraybackslash}p{(\linewidth - 8\tabcolsep) * \real{0.2000}}
  >{\raggedright\arraybackslash}p{(\linewidth - 8\tabcolsep) * \real{0.2000}}@{}}
\toprule\noalign{}
\begin{minipage}[b]{\linewidth}\raggedright
PGmc cell / interpretation
\end{minipage} & \begin{minipage}[b]{\linewidth}\raggedright
Candidate input
\end{minipage} & \begin{minipage}[b]{\linewidth}\raggedright
Old English outcome or comparison
\end{minipage} & \begin{minipage}[b]{\linewidth}\raggedright
OE comparison form
\end{minipage} & \begin{minipage}[b]{\linewidth}\raggedright
Result
\end{minipage} \\
\midrule\noalign{}
\endhead
\bottomrule\noalign{}
\endlastfoot
lexeme-level infinitive & \emph{*skáwōjaną} & manual probe output
\emph{sċēawian} & scēawian & ordinary dictionary headword of the verb,
but not the selected finite cell \\
selected imperative singular & \emph{*skáwô} & manual probe output
\emph{sċēawa} & scēawa / normalized sċēawa & exact match between the
chosen input and the selected target \\
present third singular companion & \emph{*skáwōθi} & manual probe output
\emph{sċēawaþ} & sċēawaþ & useful family control, but not the target of
this entry \\
\end{longtable}
}

\subsubsection{show (3sg) --- OE
sċēawaþ}\label{show-3sg-oe-sux10bux113awauxfe}

Derivation: citation reconstruction \emph{*skawōną}; selected input
\emph{*skáwōθi} \textgreater{} \emph{sċēawaþ} (late analogy).

\paragraph{Derivation trace}\label{derivation-trace-66}

Proto input: \emph{*skáwōθi}

\begingroup
\setlength{\fboxsep}{6pt}

\noindent\fbox{%
\begin{minipage}{0.97\linewidth}
\small
\begin{tabularx}{\linewidth}{@{}>{\raggedright\arraybackslash}p{0.320\linewidth}>{\centering\arraybackslash}X>{\raggedright\arraybackslash}p{0.520\linewidth}@{}}
\begin{minipage}[t]{\linewidth}
\centering\textbf{Earlier Germanic changes}\par
\vspace{0.35em}
\raggedright
\centering\textbf{West Germanic}\par
\raggedright
\vspace{0.2em}
\begin{tabular}{@{}>{\raggedright\arraybackslash}p{0.64\linewidth}@{\hspace{0.55em}}>{\raggedright\arraybackslash}p{0.16\linewidth}@{\hspace{0.25em}}}
PWGmc Early I Apocope & \emph{*skáwōθ} \\
\end{tabular}
\vspace{0.6em}
\centering\textbf{Northwest Germanic}\par
\raggedright
\vspace{0.2em}
\raggedright [no change]\par
\end{minipage}
&
&
\begin{minipage}[t]{\linewidth}
\centering\textbf{Old English changes}\par
\vspace{0.35em}
\raggedright
\begin{tabular}{@{}>{\raggedright\arraybackslash}p{0.68\linewidth}@{\hspace{0.55em}}>{\raggedright\arraybackslash}p{0.20\linewidth}@{\hspace{0.25em}}}
OE Aw Long Diphthong & \emph{*skḗawōθ} \\
OE Sk Palatalization & \emph{*ʃḗawōθ} \\
OE Late O Shortening & \emph{*ʃḗawaθ} \\
\end{tabular}
\end{minipage}
\\
\end{tabularx}
\end{minipage}%
} \endgroup

Outcome: \emph{sċēawaþ}

\paragraph{Reconstruction and comparative
evidence}\label{reconstruction-and-comparative-evidence-66}

Orel reconstructs the verb as \emph{*skawōjanan} and cites Old English
\emph{sceáwian} beside Old Frisian \emph{skawia}, Old Saxon
\emph{skawōn}, and Old High German \emph{scouwōn}
(\citeproc{ref-Orel2003}{Orel 2003, 337}). The selected input in this
entry is the present third singular \emph{*skáwōθi}, a finite class-II
cell with stable \emph{a} in the 3sg ending
(\citeproc{ref-RingeTaylor2014}{Ringe and Taylor 2014, 80}).

Campbell states the class-II present endings as lufas, \emph{-aþ}
(\textless{} \emph{-ōsi}, \emph{-ōþi)}
(\citeproc{ref-Campbell1959}{Campbell 1959, sect. 356.4}). Ringe and
Taylor likewise note that class-II weak present 2sg. \emph{-as(t)} and
3sg. \emph{-aþ} have stable \emph{a}
(\citeproc{ref-RingeTaylor2014}{Ringe and Taylor 2014, 80}). The
relevant comparison is therefore the 3sg cell itself, not an i-umlauted
alternative.

\paragraph{Old English evidence}\label{old-english-evidence-66}

Bright lists the simplex headword \emph{scēawian} and the imperative
\emph{scēawa}, and under \emph{geond-scēawian} also records a third
singular \emph{-sceawað} (\citeproc{ref-BrightCassidyRingler1971}{Bright
1971}). The evidence thus establishes the \emph{scēaw-} / \emph{-awað}
finite-cell pattern directly.

The form written here as \emph{sċēawaþ} is the normalized simplex
comparison form for that weak class-II pattern. It is therefore not a
dictionary headword but a selected finite form aligned with the attested
\emph{scēaw-} evidence and the directly cited \emph{-sceawað} ending
pattern.

\paragraph{Development to Old
English}\label{development-to-old-english-66}

Campbell lists \emph{scéawian} under the same West Germanic \emph{*auw}
development (\citeproc{ref-Campbell1959}{Campbell 1959, sect. 120}).
\emph{*skáwōθi} therefore belongs to the \emph{scēaw-} family before the
class-II 3sg ending is applied. Campbell's chronology and Ringe and
Taylor's stable-\emph{a} discussion show that the class-II 3sg ending
gives \emph{-aþ}, not \emph{-eþ} (\citeproc{ref-Campbell1959}{Campbell
1959, sect. 356.4}; \citeproc{ref-RingeTaylor2014}{Ringe and Taylor
2014, 80}). Because the ending never contains \emph{-j-}, no i-umlaut
applies.

\paragraph{Paradigm comparison}\label{paradigm-comparison-27}

The comparison below is manual.

{\def\LTcaptype{none} % do not increment counter
\begin{longtable}[]{@{}
  >{\raggedright\arraybackslash}p{(\linewidth - 8\tabcolsep) * \real{0.2000}}
  >{\raggedright\arraybackslash}p{(\linewidth - 8\tabcolsep) * \real{0.2000}}
  >{\raggedright\arraybackslash}p{(\linewidth - 8\tabcolsep) * \real{0.2000}}
  >{\raggedright\arraybackslash}p{(\linewidth - 8\tabcolsep) * \real{0.2000}}
  >{\raggedright\arraybackslash}p{(\linewidth - 8\tabcolsep) * \real{0.2000}}@{}}
\toprule\noalign{}
\begin{minipage}[b]{\linewidth}\raggedright
PGmc cell / interpretation
\end{minipage} & \begin{minipage}[b]{\linewidth}\raggedright
Candidate input
\end{minipage} & \begin{minipage}[b]{\linewidth}\raggedright
Old English outcome or comparison
\end{minipage} & \begin{minipage}[b]{\linewidth}\raggedright
OE comparison form
\end{minipage} & \begin{minipage}[b]{\linewidth}\raggedright
Result
\end{minipage} \\
\midrule\noalign{}
\endhead
\bottomrule\noalign{}
\endlastfoot
lexeme-level infinitive & \emph{*skáwōjaną} & manual probe output
\emph{sċēawian} & scēawian & ordinary dictionary headword of the verb,
but not the selected finite cell \\
imperative singular companion & \emph{*skáwô} & manual probe output
\emph{sċēawa} & scēawa & useful family control, but not the target of
this entry \\
selected present third singular & \emph{*skáwōθi} & manual probe output
\emph{sċēawaþ} & normalized sċēawaþ; source-side pattern \emph{-sceawað}
& exact match for the selected cell \\
\end{longtable}
}

\clearpage

\subsection{Part V. Reconstructed Old English
comparators}\label{part-v.-reconstructed-old-english-comparators}

These entries use an explicitly reconstructed Old English-stage
comparator for the branch being modelled. The relevant comparison is
therefore later than the Proto-Germanic citation form but still belongs
to the lexical derivation layer.

\subsubsection{knob --- OE cnobba}\label{knob-oe-cnobba}

Derivation: citation reconstruction \emph{*knúppaz}; selected input
\emph{*knúbbô} \textgreater{} \emph{cnobba} (reconstructed Old English
comparator).

\paragraph{Derivation trace}\label{derivation-trace-67}

Proto input: \emph{*knúbbô}

\begingroup
\setlength{\fboxsep}{6pt}

\noindent\fbox{%
\begin{minipage}{0.97\linewidth}
\small
\begin{tabularx}{\linewidth}{@{}>{\raggedright\arraybackslash}p{0.460\linewidth}>{\centering\arraybackslash}X>{\raggedright\arraybackslash}p{0.460\linewidth}@{}}
\begin{minipage}[t]{\linewidth}
\centering\textbf{Earlier Germanic changes}\par
\vspace{0.35em}
\raggedright
\centering\textbf{West Germanic}\par
\raggedright
\vspace{0.2em}
\raggedright [no change]\par
\vspace{0.6em}
\centering\textbf{Northwest Germanic}\par
\raggedright
\vspace{0.2em}
\begin{tabular}{@{}>{\raggedright\arraybackslash}p{0.74\linewidth}@{\hspace{0.55em}}>{\raggedright\arraybackslash}p{0.16\linewidth}@{\hspace{0.25em}}}
\mbox{NWGmc U Lowering} & \emph{*knóbbô} \\
\end{tabular}
\end{minipage}
&
&
\begin{minipage}[t]{\linewidth}
\centering\textbf{Old English changes}\par
\vspace{0.35em}
\raggedright
\begin{tabular}{@{}>{\raggedright\arraybackslash}p{0.74\linewidth}@{\hspace{0.55em}}>{\raggedright\arraybackslash}p{0.16\linewidth}@{\hspace{0.25em}}}
OE Unstressed Long Vowel Shortening & \emph{*knóbba} \\
\end{tabular}
\end{minipage}
\\
\end{tabularx}
\end{minipage}%
} \endgroup

Outcome: \emph{cnobba}

\paragraph{Reconstruction and comparative
evidence}\label{reconstruction-and-comparative-evidence-67}

The wider knob-family is not uniform. Kroonen's discussion of the
Germanic n-stems points to related voiced and voiceless branches within
this group (\citeproc{ref-Kroonen2011}{Kroonen 2011}). The citation
reconstruction \emph{*knúppaz} represents the broader cognate-set
headword, while the selected input \emph{*knúbbô} represents the voiced
weak-noun branch treated here.

That distinction matters because the Old English record is uneven. The
better attested OE material belongs to the voiceless branch, but the
present entry represents the reconstructed OE form that would continue
the voiced branch behind later English \textbf{knob}.

\paragraph{Old English evidence}\label{old-english-evidence-67}

Bosworth-Toller and Clark Hall preserve Old English evidence of the
\emph{cnopp} / \emph{cnoppa} type
(\citeproc{ref-BosworthToller1898}{Bosworth and Toller 1898};
\citeproc{ref-ClarkHall1960}{Clark Hall 1960}). Those forms are genuine
Old English evidence, but they belong to the voiceless branch of the
family.

The target \emph{cnobba} is different in status. It is a
\textbf{reconstructed Old English form}, not a directly attested one.
The point of using it here is to give the voiced branch an explicit
OE-stage representation instead of allowing the attested \emph{cnoppa}
branch to stand in for a different prehistory. The choice of
\emph{cnobba} is therefore a modeling and comparative decision rather
than a settled point of Old English philology.

\paragraph{Development to Old
English}\label{development-to-old-english-67}

From the selected weak-noun input \emph{*knúbbô}, the regular Old
English outcome is \emph{cnobba}, with Proto-Germanic \emph{kn-}
represented in Old English as \emph{cn-} and with the expected weak-noun
ending.

The entry therefore does not claim that \emph{cnobba} is attested. Its
claim is different: if the voiced weak-noun branch is the one to be
represented, then \emph{cnobba} is the regular Old English form
corresponding to that branch.

\paragraph{Reconstruction status}\label{reconstruction-status}

The comparison below is manual. It keeps apart the reconstructed target
and the better-attested neighboring forms.

{\def\LTcaptype{none} % do not increment counter
\begin{longtable}[]{@{}
  >{\raggedright\arraybackslash}p{(\linewidth - 4\tabcolsep) * \real{0.3333}}
  >{\raggedright\arraybackslash}p{(\linewidth - 4\tabcolsep) * \real{0.3333}}
  >{\raggedright\arraybackslash}p{(\linewidth - 4\tabcolsep) * \real{0.3333}}@{}}
\toprule\noalign{}
\begin{minipage}[b]{\linewidth}\raggedright
Form
\end{minipage} & \begin{minipage}[b]{\linewidth}\raggedright
Status
\end{minipage} & \begin{minipage}[b]{\linewidth}\raggedright
Relevance to this entry
\end{minipage} \\
\midrule\noalign{}
\endhead
\bottomrule\noalign{}
\endlastfoot
\emph{*knúbbô} -\textgreater{} \emph{cnobba} & reconstructed OE form;
trace-supported & selected target \\
\emph{cnopp} / \emph{cnoppa} & attested OE branch & important control
form, but belongs to the voiceless branch \\
\emph{cnæp} & attested OE form from another family & not part of the
present lexeme line \\
\end{longtable}
}

This remains the most review-sensitive item in the present pilot batch,
because the choice between reconstructed \emph{cnobba} and attested
\emph{cnoppa} is still a comparator-policy question rather than a
settled point of OE attestation.

\subsubsection{reek --- OE rēac}\label{reek-oe-rux113ac}

Derivation: citation reconstruction \emph{*ráukiz}; selected input
\emph{*ráukaz} \textgreater{} \emph{rēac} (reconstructed Old English
comparator).

\paragraph{Derivation trace}\label{derivation-trace-68}

Proto input: \emph{*ráukaz}

\begingroup
\setlength{\fboxsep}{6pt}

\noindent\fbox{%
\begin{minipage}{0.97\linewidth}
\small
\begin{tabularx}{\linewidth}{@{}>{\raggedright\arraybackslash}p{0.320\linewidth}>{\centering\arraybackslash}X>{\raggedright\arraybackslash}p{0.520\linewidth}@{}}
\begin{minipage}[t]{\linewidth}
\centering\textbf{Earlier Germanic changes}\par
\vspace{0.35em}
\raggedright
\centering\textbf{West Germanic}\par
\raggedright
\vspace{0.2em}
\raggedright [no change]\par
\vspace{0.6em}
\centering\textbf{Northwest Germanic}\par
\raggedright
\vspace{0.2em}
\begin{tabular}{@{}>{\raggedright\arraybackslash}p{0.74\linewidth}@{\hspace{0.55em}}>{\raggedright\arraybackslash}p{0.16\linewidth}@{\hspace{0.25em}}}
\mbox{PGmc Final Z Deletion} & \emph{*ráuka} \\
\end{tabular}
\end{minipage}
&
&
\begin{minipage}[t]{\linewidth}
\centering\textbf{Old English changes}\par
\vspace{0.35em}
\raggedright
\begin{tabular}{@{}>{\raggedright\arraybackslash}p{0.70\linewidth}@{\hspace{0.55em}}>{\raggedright\arraybackslash}p{0.16\linewidth}@{\hspace{0.25em}}}
OE Au Fronting & \emph{*ráeuka} \\
OE Diphthong Leveling & \emph{*rēaka} \\
PWGmc Final Bare A Loss & \emph{*rēak} \\
\end{tabular}
\end{minipage}
\\
\end{tabularx}
\end{minipage}%
} \endgroup

Outcome: \emph{rēac}

\paragraph{Reconstruction and comparative
evidence}\label{reconstruction-and-comparative-evidence-68}

The wider noun family is represented by \emph{*ráukiz} / \emph{*rauki-},
with Old English \emph{rēc} as the attested noun reflex in the
comparative dictionaries (\citeproc{ref-Kroonen2013}{Kroonen 2013, 446};
\citeproc{ref-Orel2003}{Orel 2003, 338}). The selected input
\emph{*ráukaz} is therefore not the lexeme-level headword, but the form
used here for the Old English derivation.

\paragraph{Old English evidence}\label{old-english-evidence-68}

The attested noun is \emph{rēc}, not \emph{rēac}. Clark Hall records
\emph{rēc} as the noun and also preserves related forms such as
\emph{rēcels}; Kroonen likewise gives OE \emph{rēc} under the noun
family (\citeproc{ref-ClarkHall1960}{Clark Hall 1960, 255};
\citeproc{ref-Kroonen2013}{Kroonen 2013, 446}). Clark Hall and Seebold
also record verbal \emph{rēac} as the preterite of \emph{rēocan}, but
that verbal form is separate from the noun treated here
(\citeproc{ref-ClarkHall1960}{Clark Hall 1960, 254};
\citeproc{ref-Seebold1970}{Seebold 1970, 380}).

The selected target \emph{rēac} is therefore a reconstructed West Saxon
noun form, not a directly attested manuscript headword.

\paragraph{Development to Old
English}\label{development-to-old-english-68}

From \emph{*ráukaz}, the regular West Saxon development gives
\emph{rēac}. The attested noun \emph{rēc} belongs to the same lexical
family, but reflects a later smoothed surface form rather than the
regular noun target represented here.

\paragraph{Form note}\label{form-note}

The distinction here is between an attested noun headword \emph{rēc} and
a reconstructed regular West Saxon target \emph{rēac}. The latter is
treated as the modelling target, while the former remains philological
background.

\subsubsection{strew --- OE strīeġan}\label{strew-oe-strux12beux121an}

Derivation: \emph{*stráwjaną} \textgreater{} \emph{strīeġan}
(reconstructed Old English comparator).

\paragraph{Derivation trace}\label{derivation-trace-69}

Proto input: \emph{*stráwjaną}

\begingroup
\setlength{\fboxsep}{6pt}

\noindent\fbox{%
\begin{minipage}{0.97\linewidth}
\footnotesize
\begin{tabularx}{\linewidth}{@{}>{\raggedright\arraybackslash}p{0.280\linewidth}>{\centering\arraybackslash}X>{\raggedright\arraybackslash}p{0.560\linewidth}@{}}
\begin{minipage}[t]{\linewidth}
\centering\textbf{Earlier Germanic changes}\par
\vspace{0.35em}
\raggedright
\centering\textbf{West Germanic}\par
\raggedright
\vspace{0.2em}
\raggedright [no change]\par
\vspace{0.6em}
\centering\textbf{Northwest Germanic}\par
\raggedright
\vspace{0.2em}
\raggedright [no change]\par
\end{minipage}
&
&
\begin{minipage}[t]{\linewidth}
\centering\textbf{Old English changes}\par
\vspace{0.35em}
\raggedright
\begin{tabular}{@{}>{\raggedright\arraybackslash}p{0.62\linewidth}@{\hspace{0.55em}}>{\raggedright\arraybackslash}p{0.28\linewidth}@{\hspace{0.25em}}}
OE Awj Glide Formation & \emph{*stráujaną} \\
OE Au Fronting & \emph{*stráeujaną} \\
OE Diphthong Leveling & \emph{*strēajaną} \\
OE Heavy Syllable Nasal Apocope & \emph{*strēajan} \\
OE Secondary Nasalization & \emph{*strēająn} \\
OE I Umlaut & \emph{*strīejąn} \\
OE Weak Tail Reduction & \emph{*strīejan} \\
OE J Strengthening After Front Diphthong & \emph{*strīeʒan} \\
\end{tabular}
\end{minipage}
\\
\end{tabularx}
\end{minipage}%
} \endgroup

Outcome: \emph{strīeġan}

\paragraph{Reconstruction and comparative
evidence}\label{reconstruction-and-comparative-evidence-69}

Kroonen cites the inherited weak verb as \emph{*straujan-} and gives Old
English \emph{streowian} as its dictionary continuation
(\citeproc{ref-Kroonen2013}{Kroonen 2013, 483}). Ringe and Taylor make
the split within Old English explicit: the inherited class-I verb is
continued by Anglian \emph{strēgan}, while West Saxon \emph{streowian}
is a remodelled class-II verb (\citeproc{ref-RingeTaylor2014}{Ringe and
Taylor 2014, sect. 6.1} n.~27).

The aw-series comparison is important here. Luick groups
\emph{*strauwjan} with the same set as \emph{*hauwja-} and
\emph{*kauwjan}, yielding Anglian \emph{strēzan} beside West Saxon forms
of the \emph{hīez}, ciezan type (\citeproc{ref-Luick1914}{Luick
1914-\/-1940, sect. 98}). Fulk likewise allows an early West Saxon
\emph{*striegan} directly from Proto-Germanic \emph{*straujana}
(\citeproc{ref-Fulk2018}{Fulk 2018, sect. 4.10} n.~1).

\paragraph{Old English evidence}\label{old-english-evidence-69}

The attested inherited Old English form is \emph{strēgan} in Anglian.
The attested West Saxon citation forms are \emph{strewian},
\emph{streowian}, and \emph{strēawian}, which belong to the remodelled
class-II branch (\citeproc{ref-RingeTaylor2014}{Ringe and Taylor 2014,
sect. 6.1} n.~27; \citeproc{ref-Campbell1959}{Campbell 1959, sect.
753.7}).

The target \emph{strīeġan} is therefore a \textbf{reconstructed Old
English form}, not an attested manuscript lemma. It is the reconstructed
West Saxon reflex of the inherited class-I verb, chosen to keep the
inherited branch distinct from the better-attested remodelled West Saxon
lemma.

\paragraph{Development to Old
English}\label{development-to-old-english-69}

From \emph{*stráwjaną}, the inherited West Saxon development passes
through \emph{*straujaną}, fronting and leveling to a \emph{*strēajan-}
stage, i-umlaut to \emph{*strīejan}, and retention or strengthening of
the glide after the front diphthong, written here as \emph{ġ}. The
resulting form is \emph{strīeġan}.

This differs from Anglian \emph{strēgan}, where smoothing removes the
diphthongal sequence, and from West Saxon \emph{strewian} /
\emph{streowian} / \emph{strēawian}, where the verb has already been
remodelled into class II (\citeproc{ref-Fulk2018}{Fulk 2018, sect. 4.10}
n.~1; \citeproc{ref-Campbell1959}{Campbell 1959, sect. 753.7}).

\paragraph{Reconstruction status}\label{reconstruction-status-1}

The comparison below is manual. It keeps apart the attested inherited
branch, the attested remodelled branch, and the reconstructed West Saxon
comparator.

{\def\LTcaptype{none} % do not increment counter
\begin{longtable}[]{@{}
  >{\raggedright\arraybackslash}p{(\linewidth - 4\tabcolsep) * \real{0.3333}}
  >{\raggedright\arraybackslash}p{(\linewidth - 4\tabcolsep) * \real{0.3333}}
  >{\raggedright\arraybackslash}p{(\linewidth - 4\tabcolsep) * \real{0.3333}}@{}}
\toprule\noalign{}
\begin{minipage}[b]{\linewidth}\raggedright
Form or branch
\end{minipage} & \begin{minipage}[b]{\linewidth}\raggedright
Status
\end{minipage} & \begin{minipage}[b]{\linewidth}\raggedright
Relevance to this entry
\end{minipage} \\
\midrule\noalign{}
\endhead
\bottomrule\noalign{}
\endlastfoot
\emph{strēgan} & attested Anglian inherited class-I form & proves that
the inherited verb survived into Old English \\
\emph{strīeġan} & reconstructed West Saxon inherited class-I form;
trace-supported & selected target \\
\emph{strewian} / \emph{streowian} / \emph{strēawian} & attested
remodelled West Saxon class-II forms & genuine OE evidence, but not the
inherited branch modeled here \\
\end{longtable}
}

\clearpage

\subsection{Part VI. Known but unmodelled
remodellings}\label{part-vi.-known-but-unmodelled-remodellings}

These entries preserve cases where the historical remodelling is broadly
understood, but the current deterministic transducer does not model that
later reshaping directly.

\subsubsection{fire --- OE fȳre}\label{fire-oe-fux233re}

Derivation: \emph{*fūri} yields regular \emph{fȳr}; the selected target
is \emph{fȳre} (known but unmodelled remodelling).

\paragraph{Derivation trace}\label{derivation-trace-70}

Proto input: \emph{*fūri}

\begingroup
\setlength{\fboxsep}{6pt}

\noindent\fbox{%
\begin{minipage}{0.97\linewidth}
\small
\begin{tabularx}{\linewidth}{@{}>{\raggedright\arraybackslash}p{0.300\linewidth}>{\centering\arraybackslash}X>{\raggedright\arraybackslash}p{0.540\linewidth}@{}}
\begin{minipage}[t]{\linewidth}
\centering\textbf{Earlier Germanic changes}\par
\vspace{0.35em}
\raggedright
\centering\textbf{West Germanic}\par
\raggedright
\vspace{0.2em}
\raggedright [no change]\par
\vspace{0.6em}
\centering\textbf{Northwest Germanic}\par
\raggedright
\vspace{0.2em}
\raggedright [no change]\par
\end{minipage}
&
&
\begin{minipage}[t]{\linewidth}
\centering\textbf{Old English changes}\par
\vspace{0.35em}
\raggedright
\begin{tabular}{@{}>{\raggedright\arraybackslash}p{0.74\linewidth}@{\hspace{0.55em}}>{\raggedright\arraybackslash}p{0.16\linewidth}@{\hspace{0.25em}}}
OE I Umlaut & \emph{*fȳri} \\
\mbox{OE High Vowel Apocope} & \emph{*fȳr} \\
\end{tabular}
\end{minipage}
\\
\end{tabularx}
\end{minipage}%
} \endgroup

Transducer outcome: \emph{fȳr}

Selected target: \emph{fȳre}

\paragraph{Reconstruction and comparative
evidence}\label{reconstruction-and-comparative-evidence-70}

Kroonen places the lexeme in a heteroclitic family \emph{*fōr}
\textasciitilde{} \emph{*fun-} and explains the front-mutated West
Germanic forms from an oblique form of the \emph{*fu(w)eri} type
(\citeproc{ref-Kroonen2013}{Kroonen 2013}). The selected input
\emph{*fūri} therefore does not function as an arbitrary substitute for
the headword: it represents the specific inherited cell that supplies
the \emph{i} needed for i-umlaut.

That distinction matters because the Old English target combines a
regular inherited form \emph{fȳr} with an attested analogical surface
form \emph{fȳre}.

\paragraph{Old English evidence}\label{old-english-evidence-70}

Bosworth-Toller records \emph{fyr} as the noun `fire' and also preserves
oblique \emph{fyre} in the Old English record
(\citeproc{ref-BosworthToller1898}{Bosworth and Toller 1898, 288}). The
first is the regular inherited outcome of the phonological development
from the selected input; the second shows the later restoration of a
final \emph{-e} within the paradigm.

The entry therefore concerns the relation between a regular inherited
oblique input and an attested Old English surface form that has
undergone later morphological remodeling.

\paragraph{Development to Old
English}\label{development-to-old-english-70}

From \emph{*fūri}, i-umlaut changes \emph{ū} to \emph{ȳ}, and subsequent
loss of the final high vowel after a heavy syllable yields \emph{fȳr}
(\citeproc{ref-RingeTaylor2014}{Ringe and Taylor 2014};
\citeproc{ref-Hogg1992}{Hogg 1992}; \citeproc{ref-Campbell1959}{Campbell
1959}). The inherited phonology is complete at that point.

\emph{fȳre} is later than that inherited output. Its final \emph{-e}
belongs to analogical restoration in the Old English paradigm rather
than to the original Proto-Germanic ending. The form therefore remains
\emph{known\_unmodelled}: the deterministic phonology is regular, but
the attested surface form includes later morphological rebuilding.

\paragraph{Paradigm comparison}\label{paradigm-comparison-28}

The comparison below is manual. It distinguishes the lexeme-level
etymological background from the inherited cell that actually produces
the front-mutated form and from the later analogical surface result.

{\def\LTcaptype{none} % do not increment counter
\begin{longtable}[]{@{}
  >{\raggedright\arraybackslash}p{(\linewidth - 8\tabcolsep) * \real{0.2000}}
  >{\raggedright\arraybackslash}p{(\linewidth - 8\tabcolsep) * \real{0.2000}}
  >{\raggedright\arraybackslash}p{(\linewidth - 8\tabcolsep) * \real{0.2000}}
  >{\raggedright\arraybackslash}p{(\linewidth - 8\tabcolsep) * \real{0.2000}}
  >{\raggedright\arraybackslash}p{(\linewidth - 8\tabcolsep) * \real{0.2000}}@{}}
\toprule\noalign{}
\begin{minipage}[b]{\linewidth}\raggedright
PGmc cell / interpretation
\end{minipage} & \begin{minipage}[b]{\linewidth}\raggedright
Candidate input
\end{minipage} & \begin{minipage}[b]{\linewidth}\raggedright
OE output or comparison
\end{minipage} & \begin{minipage}[b]{\linewidth}\raggedright
OE comparison form
\end{minipage} & \begin{minipage}[b]{\linewidth}\raggedright
Result
\end{minipage} \\
\midrule\noalign{}
\endhead
\bottomrule\noalign{}
\endlastfoot
lexeme-level heteroclitic headword & \emph{*fōr} \textasciitilde{}
\emph{*fun-} & comparative background only & fire family & explains the
wider lexeme, but not the selected oblique input \\
inherited oblique cell & *fūri & compact-trace output: \emph{fȳr} & fȳr
& regular inherited output from the selected input \\
later analogical surface form & --- & attested \emph{fȳre} with restored
\emph{-e} & fȳre & genuine OE target, but not the direct phonological
output \\
\end{longtable}
}

\subsubsection{tap --- OE tæppa}\label{tap-oe-tuxe6ppa}

Derivation: \emph{*táppô} yields regular \emph{tappa}; the selected
target is \emph{tæppa} (known but unmodelled remodelling).

\paragraph{Derivation trace}\label{derivation-trace-71}

Proto input: \emph{*táppô}

\begingroup
\setlength{\fboxsep}{6pt}

\noindent\fbox{%
\begin{minipage}{0.97\linewidth}
\small
\begin{tabularx}{\linewidth}{@{}>{\raggedright\arraybackslash}p{0.300\linewidth}>{\centering\arraybackslash}X>{\raggedright\arraybackslash}p{0.540\linewidth}@{}}
\begin{minipage}[t]{\linewidth}
\centering\textbf{Earlier Germanic changes}\par
\vspace{0.35em}
\raggedright
\centering\textbf{West Germanic}\par
\raggedright
\vspace{0.2em}
\raggedright [no change]\par
\vspace{0.6em}
\centering\textbf{Northwest Germanic}\par
\raggedright
\vspace{0.2em}
\raggedright [no change]\par
\end{minipage}
&
&
\begin{minipage}[t]{\linewidth}
\centering\textbf{Old English changes}\par
\vspace{0.35em}
\raggedright
\begin{tabular}{@{}>{\raggedright\arraybackslash}p{0.74\linewidth}@{\hspace{0.55em}}>{\raggedright\arraybackslash}p{0.16\linewidth}@{\hspace{0.25em}}}
Anglo Frisian Brightening & \emph{*tæppô} \\
OE A Restoration & \emph{*tappô} \\
OE Unstressed Long Vowel Shortening & \emph{*tappa} \\
\end{tabular}
\end{minipage}
\\
\end{tabularx}
\end{minipage}%
} \endgroup

Transducer outcome: \emph{tappa}

Selected target: \emph{tæppa}

\paragraph{Reconstruction and comparative
evidence}\label{reconstruction-and-comparative-evidence-71}

Orel gives the noun under \emph{*tappòn} and already connects it with
Old English \emph{tæppa} (\citeproc{ref-Orel2003}{Orel 2003}). The
selected input is therefore the inherited noun itself; the entry does
not depend on a different lexeme-level proto or a different inherited
noun cell.

\paragraph{Old English evidence}\label{old-english-evidence-71}

The Old English noun family is well attested. Orel gives \emph{tæppa},
and Clark Hall records \emph{tæppa} together with derivatives
\emph{tæppere} and \emph{tæppestre} (\citeproc{ref-Orel2003}{Orel 2003};
\citeproc{ref-ClarkHall1960}{Clark Hall 1960, 305}). The target is
therefore a real Old English noun form, not a reconstructed convenience
spelling.

\paragraph{Development to Old
English}\label{development-to-old-english-71}

From \emph{*táppô}, the regular inherited noun path gives \emph{tappa}.
The attested target \emph{tæppa} therefore stands outside that regular
phonological development.

The mismatch is historically intelligible, but it is not solved here by
a new inherited input. A related j-verb pathway would give
\emph{teppan}, not the noun target \emph{tæppa}. The entry accordingly
remains \emph{known\_unmodelled}.

\paragraph{Form comparison}\label{form-comparison-2}

{\def\LTcaptype{none} % do not increment counter
\begin{longtable}[]{@{}
  >{\raggedright\arraybackslash}p{(\linewidth - 6\tabcolsep) * \real{0.2500}}
  >{\raggedright\arraybackslash}p{(\linewidth - 6\tabcolsep) * \real{0.2500}}
  >{\raggedright\arraybackslash}p{(\linewidth - 6\tabcolsep) * \real{0.2500}}
  >{\raggedright\arraybackslash}p{(\linewidth - 6\tabcolsep) * \real{0.2500}}@{}}
\toprule\noalign{}
\begin{minipage}[b]{\linewidth}\raggedright
Form type
\end{minipage} & \begin{minipage}[b]{\linewidth}\raggedright
Input or form
\end{minipage} & \begin{minipage}[b]{\linewidth}\raggedright
OE output or comparison
\end{minipage} & \begin{minipage}[b]{\linewidth}\raggedright
Result
\end{minipage} \\
\midrule\noalign{}
\endhead
\bottomrule\noalign{}
\endlastfoot
regular inherited noun path & *táppô & compact-trace output:
\emph{tappa} & regular output, but not the target \\
attested OE target & --- & \emph{tæppa} & genuine target form, but
analogically remodelled in the present classification \\
related j-verb background & *táppjaną & \emph{teppan} & related
formation, but not the noun target \\
\end{longtable}
}

\clearpage

\subsection{Part VII. Unexplained or deliberately unmodelled
exceptions}\label{part-vii.-unexplained-or-deliberately-unmodelled-exceptions}

These entries preserve a mismatch between the regular transducer output
and the selected Old English target. They are retained as documented
lexical exceptions rather than treated as evidence for further
sound-change repair.

\subsubsection{buck --- OE bucc}\label{buck-oe-bucc}

Derivation: \emph{*búkkaz} yields regular \emph{bocc}; the selected
target is \emph{bucc} (unexplained exception).

\paragraph{Derivation trace}\label{derivation-trace-72}

Proto input: \emph{*búkkaz}

\begingroup
\setlength{\fboxsep}{6pt}

\noindent\fbox{%
\begin{minipage}{0.97\linewidth}
\small
\begin{tabularx}{\linewidth}{@{}>{\raggedright\arraybackslash}p{0.540\linewidth}>{\centering\arraybackslash}X>{\raggedright\arraybackslash}p{0.300\linewidth}@{}}
\begin{minipage}[t]{\linewidth}
\centering\textbf{Earlier Germanic changes}\par
\vspace{0.35em}
\raggedright
\centering\textbf{West Germanic}\par
\raggedright
\vspace{0.2em}
\raggedright [no change]\par
\vspace{0.6em}
\centering\textbf{Northwest Germanic}\par
\raggedright
\vspace{0.2em}
\begin{tabular}{@{}>{\raggedright\arraybackslash}p{0.74\linewidth}@{\hspace{0.55em}}>{\raggedright\arraybackslash}p{0.16\linewidth}@{\hspace{0.25em}}}
\mbox{NWGmc U Lowering} & \emph{*bókkaz} \\
\mbox{PGmc Final Z Deletion} & \emph{*bókka} \\
\end{tabular}
\end{minipage}
&
&
\begin{minipage}[t]{\linewidth}
\centering\textbf{Old English changes}\par
\vspace{0.35em}
\raggedright
\begin{tabular}{@{}>{\raggedright\arraybackslash}p{0.70\linewidth}@{\hspace{0.55em}}>{\raggedright\arraybackslash}p{0.16\linewidth}@{\hspace{0.25em}}}
PWGmc Final Bare A Loss & \emph{*bókk} \\
\end{tabular}
\end{minipage}
\\
\end{tabularx}
\end{minipage}%
} \endgroup

Transducer outcome: \emph{bocc}

Selected target: \emph{bucc}

\paragraph{Reconstruction and comparative
evidence}\label{reconstruction-and-comparative-evidence-72}

Kroonen and Orel both reconstruct the word with a geminate stop,
\emph{*bukkaz}, and both also preserve parallel n-stem material behind
Old English \emph{bucca} (\citeproc{ref-Kroonen2013}{Kroonen 2013, 121};
\citeproc{ref-Orel2003}{Orel 2003}). The selected input therefore
remains identical with the lexeme label: no alternative inherited cell
accounts for the form.

\paragraph{Old English evidence}\label{old-english-evidence-72}

Old English preserves a mixed lexical picture. Campbell cites
\emph{bucca} in the exception set for this phonological environment,
while Clark Hall and Bosworth-Toller show that Old English has both
\emph{bucca} and \emph{bucc} (\citeproc{ref-Campbell1959}{Campbell
1959}; \citeproc{ref-ClarkHall1960}{Clark Hall 1960};
\citeproc{ref-BosworthToller1898}{Bosworth and Toller 1898, 122}). The
a-stem citation form \emph{bucc} is the target treated here, with
\emph{bucca} kept as genuine philological background from the same
lexical family.

\paragraph{Development to Old
English}\label{development-to-old-english-72}

From \emph{*búkkaz}, the regular inherited path gives \emph{bocc}. That
is the form expected under the ordinary lowering pattern in this
environment. \emph{bucc} therefore remains outside the deterministic
phonology.

No accepted inherited cell repairs the mismatch. A high-vowel
alternative would introduce i-umlaut and produce a \emph{byċċ}-type form
rather than the target. \emph{bucc} is therefore best treated as a
documented exception, not as a regular paradigm-cell survival.

\paragraph{Form comparison}\label{form-comparison-3}

{\def\LTcaptype{none} % do not increment counter
\begin{longtable}[]{@{}
  >{\raggedright\arraybackslash}p{(\linewidth - 6\tabcolsep) * \real{0.2500}}
  >{\raggedright\arraybackslash}p{(\linewidth - 6\tabcolsep) * \real{0.2500}}
  >{\raggedright\arraybackslash}p{(\linewidth - 6\tabcolsep) * \real{0.2500}}
  >{\raggedright\arraybackslash}p{(\linewidth - 6\tabcolsep) * \real{0.2500}}@{}}
\toprule\noalign{}
\begin{minipage}[b]{\linewidth}\raggedright
Form type
\end{minipage} & \begin{minipage}[b]{\linewidth}\raggedright
Input or form
\end{minipage} & \begin{minipage}[b]{\linewidth}\raggedright
OE output or comparison
\end{minipage} & \begin{minipage}[b]{\linewidth}\raggedright
Result
\end{minipage} \\
\midrule\noalign{}
\endhead
\bottomrule\noalign{}
\endlastfoot
regular inherited noun path & *búkkaz & compact-trace output:
\emph{bocc} & regular output, but not the target \\
attested OE target & --- & \emph{bucc} & genuine target form, but
unexplained in the present classification \\
parallel OE lexical background & --- & \emph{bucca} & related n-stem
form, not the present target \\
\end{longtable}
}

\subsubsection{fowl --- OE fugol}\label{fowl-oe-fugol}

Derivation: \emph{*fúglaz} yields regular \emph{fogol}; the selected
target is \emph{fugol} (unexplained exception).

\paragraph{Derivation trace}\label{derivation-trace-73}

Proto input: \emph{*fúglaz}

\begingroup
\setlength{\fboxsep}{6pt}

\noindent\fbox{%
\begin{minipage}{0.97\linewidth}
\small
\begin{tabularx}{\linewidth}{@{}>{\raggedright\arraybackslash}p{0.440\linewidth}>{\centering\arraybackslash}X>{\raggedright\arraybackslash}p{0.440\linewidth}@{}}
\begin{minipage}[t]{\linewidth}
\centering\textbf{Earlier Germanic changes}\par
\vspace{0.35em}
\raggedright
\centering\textbf{West Germanic}\par
\raggedright
\vspace{0.2em}
\raggedright [no change]\par
\vspace{0.6em}
\centering\textbf{Northwest Germanic}\par
\raggedright
\vspace{0.2em}
\begin{tabular}{@{}>{\raggedright\arraybackslash}p{0.74\linewidth}@{\hspace{0.55em}}>{\raggedright\arraybackslash}p{0.16\linewidth}@{\hspace{0.25em}}}
\mbox{NWGmc U Lowering} & \emph{*fóglaz} \\
\mbox{PGmc Final Z Deletion} & \emph{*fógla} \\
\end{tabular}
\end{minipage}
&
&
\begin{minipage}[t]{\linewidth}
\centering\textbf{Old English changes}\par
\vspace{0.35em}
\raggedright
\begin{tabular}{@{}>{\raggedright\arraybackslash}p{0.70\linewidth}@{\hspace{0.55em}}>{\raggedright\arraybackslash}p{0.16\linewidth}@{\hspace{0.25em}}}
PWGmc Final Bare A Loss & \emph{*fógl} \\
OE Epenthetic Vowel & \emph{*fógol} \\
\end{tabular}
\end{minipage}
\\
\end{tabularx}
\end{minipage}%
} \endgroup

Transducer outcome: \emph{fogol}

Selected target: \emph{fugol}

\paragraph{Reconstruction and comparative
evidence}\label{reconstruction-and-comparative-evidence-73}

The noun is the ordinary Germanic a-stem \emph{*fúglaz}, continued by
forms such as Old Norse \emph{fugl} and Old High German \emph{fogal}
(\citeproc{ref-Kroonen2013}{Kroonen 2013, 197};
\citeproc{ref-Orel2003}{Orel 2003, 155}). There is no stem-class or
paradigm-cell dispute behind this entry. The comparative headword and
the selected input are the same.

The difficulty lies instead in the stressed root vowel. Under the
regular West Germanic and Old English development, that \emph{u} should
lower before the following non-high vowel, yielding an \emph{o}-vocalism
(\citeproc{ref-RingeTaylor2014}{Ringe and Taylor 2014, 42--43};
\citeproc{ref-Campbell1959}{Campbell 1959, 43}).

\paragraph{Old English evidence}\label{old-english-evidence-73}

Old English dictionaries record the noun as \emph{fugol}, with variant
spelling \emph{fugel} (\citeproc{ref-BosworthToller1898}{Bosworth and
Toller 1898, 282}; \citeproc{ref-ClarkHall1960}{Clark Hall 1960, 138}).
The target is therefore an attested ordinary Old English noun, not a
reconstructed or selectively chosen paradigm form.

The attested word already contains the crucial problem. Its medial
\emph{-o-} is the ordinary parasite vowel of Old English cluster
phonology, but the root \emph{fu-} retains \emph{u} where the regular
history predicts \emph{fo-} (\citeproc{ref-Campbell1959}{Campbell 1959,
150}).

\paragraph{Development to Old
English}\label{development-to-old-english-73}

From \emph{*fúglaz}, the regular cascade yields \emph{fogol}: the root
vowel lowers before the following non-high vowel
(\citeproc{ref-RingeTaylor2014}{Ringe and Taylor 2014, 42--43};
\citeproc{ref-Campbell1959}{Campbell 1959, 43}), final \emph{z} is lost,
and the cluster is resolved by the usual medial vowel
(\citeproc{ref-RingeTaylor2014}{Ringe and Taylor 2014, 345};
\citeproc{ref-Campbell1959}{Campbell 1959, 150}). That is the expected
inherited outcome.

The attested Old English noun is \emph{fugol}, not \emph{fogol}. Luick
and later handbooks treat this preservation of \emph{u} as a small
inherited residue, not as a categorical sound law
(\citeproc{ref-Luick1914}{Luick 1914-\/-1940, 148};
\citeproc{ref-RingeTaylor2014}{Ringe and Taylor 2014, 47}). The item
therefore remains a genuine lexical exception rather than the output of
a recoverable regular mechanism.

\paragraph{Expected and attested
forms}\label{expected-and-attested-forms}

The comparison below is manual. It distinguishes the regular prediction
from the attested Old English noun.

{\def\LTcaptype{none} % do not increment counter
\begin{longtable}[]{@{}
  >{\raggedright\arraybackslash}p{(\linewidth - 4\tabcolsep) * \real{0.3333}}
  >{\raggedright\arraybackslash}p{(\linewidth - 4\tabcolsep) * \real{0.3333}}
  >{\raggedright\arraybackslash}p{(\linewidth - 4\tabcolsep) * \real{0.3333}}@{}}
\toprule\noalign{}
\begin{minipage}[b]{\linewidth}\raggedright
Form
\end{minipage} & \begin{minipage}[b]{\linewidth}\raggedright
Status
\end{minipage} & \begin{minipage}[b]{\linewidth}\raggedright
Relevance to this entry
\end{minipage} \\
\midrule\noalign{}
\endhead
\bottomrule\noalign{}
\endlastfoot
\emph{fogol} & computed regular output from \emph{*fúglaz} & establishes
the expected inherited outcome \\
\emph{fugol} & attested Old English form & selected target; preserves
unexplained root \emph{u} \\
\emph{fugel} & attested variant spelling & secondary spelling variant of
the attested noun \\
\end{longtable}
}

The unresolved point lies only in the root vowel. The medial \emph{-o-}
is regular, but no accepted lautgesetzlich pathway has been found from
\emph{*fúglaz} to attested \emph{fugol}.

\subsubsection{rust --- OE rust}\label{rust-oe-rust}

Derivation: \emph{*rústō} yields regular \emph{rost}; the selected
target is \emph{rust} (unexplained exception).

\paragraph{Derivation trace}\label{derivation-trace-74}

Proto input: \emph{*rústō}

\begingroup
\setlength{\fboxsep}{6pt}

\noindent\fbox{%
\begin{minipage}{0.97\linewidth}
\small
\begin{tabularx}{\linewidth}{@{}>{\raggedright\arraybackslash}p{0.460\linewidth}>{\centering\arraybackslash}X>{\raggedright\arraybackslash}p{0.460\linewidth}@{}}
\begin{minipage}[t]{\linewidth}
\centering\textbf{Earlier Germanic changes}\par
\vspace{0.35em}
\raggedright
\centering\textbf{West Germanic}\par
\raggedright
\vspace{0.2em}
\raggedright [no change]\par
\vspace{0.6em}
\centering\textbf{Northwest Germanic}\par
\raggedright
\vspace{0.2em}
\begin{tabular}{@{}>{\raggedright\arraybackslash}p{0.74\linewidth}@{\hspace{0.55em}}>{\raggedright\arraybackslash}p{0.16\linewidth}@{\hspace{0.25em}}}
\mbox{NWGmc U Lowering} & \emph{*róstō} \\
\mbox{NWGmc Final Long O Raising} & \emph{*róstu} \\
\end{tabular}
\end{minipage}
&
&
\begin{minipage}[t]{\linewidth}
\centering\textbf{Old English changes}\par
\vspace{0.35em}
\raggedright
\begin{tabular}{@{}>{\raggedright\arraybackslash}p{0.74\linewidth}@{\hspace{0.55em}}>{\raggedright\arraybackslash}p{0.16\linewidth}@{\hspace{0.25em}}}
\mbox{OE High Vowel Apocope} & \emph{*róst} \\
\end{tabular}
\end{minipage}
\\
\end{tabularx}
\end{minipage}%
} \endgroup

Transducer outcome: \emph{rost}

Selected target: \emph{rust}

\paragraph{Reconstruction and comparative
evidence}\label{reconstruction-and-comparative-evidence-74}

The comparative dictionaries do not support a single citation
reconstruction uniformly. Orel cites \emph{*rustaz} sb.m./f.~with Old
English \emph{rust} and Old Saxon and Old High German \emph{rost}
(\citeproc{ref-Orel2003}{Orel 2003}). The form \emph{*rústō} therefore
stands here as a competing citation reconstruction rather than as the
best-supported inherited headword.

That disagreement does not remove the central problem. Whether one
starts from \emph{*rústō} or from source-supported \emph{*rustaz}, the
regular citation-form history points toward \emph{rost}, not toward the
attested Old English noun.

\paragraph{Old English evidence}\label{old-english-evidence-74}

The Old English noun is attested, not reconstructed. Clark Hall gives
\emph{rūst} m. (\citeproc{ref-ClarkHall1960}{Clark Hall 1960, 245}), and
Bosworth-Toller records \emph{rúst} (? and rust)
(\citeproc{ref-BosworthToller1898}{Bosworth and Toller 1898, 677}). The
form is normalized here as \emph{rust} from that attested record.

Those dictionary entries also matter morphologically. They support a
masculine noun, which aligns better with Orel's \emph{*rustaz} than with
the competing \emph{*rústō} preserved in the header.

\paragraph{Development to Old
English}\label{development-to-old-english-74}

Under the regular lowering of stressed \emph{u} before a following
non-high vowel, the citation-form input gives \emph{rost}, not
\emph{rust} (\citeproc{ref-Campbell1959}{Campbell 1959};
\citeproc{ref-RingeTaylor2014}{Ringe and Taylor 2014}). The same regular
result follows from the comparative citation-form reconstruction
\emph{*rustaz}.

A high-vowel comparator such as instrumental-type \emph{*rústu} would
yield \emph{rust} regularly, but that does not explain the attested
citation form of the noun. No accepted regular pathway from the citation
form to attested \emph{rust} has been established.

\paragraph{Expected and attested
forms}\label{expected-and-attested-forms-1}

The comparison below is manual. It separates the regular inherited
outcomes from the attested Old English noun.

{\def\LTcaptype{none} % do not increment counter
\begin{longtable}[]{@{}
  >{\raggedright\arraybackslash}p{(\linewidth - 4\tabcolsep) * \real{0.3333}}
  >{\raggedright\arraybackslash}p{(\linewidth - 4\tabcolsep) * \real{0.3333}}
  >{\raggedright\arraybackslash}p{(\linewidth - 4\tabcolsep) * \real{0.3333}}@{}}
\toprule\noalign{}
\begin{minipage}[b]{\linewidth}\raggedright
Form / interpretation
\end{minipage} & \begin{minipage}[b]{\linewidth}\raggedright
Status
\end{minipage} & \begin{minipage}[b]{\linewidth}\raggedright
Relevance to this entry
\end{minipage} \\
\midrule\noalign{}
\endhead
\bottomrule\noalign{}
\endlastfoot
\emph{*rústō} -\textgreater{} \emph{rost} & computed regular output from
one competing citation reconstruction & shows that this citation
reconstruction does not reach attested \emph{rust} \\
\emph{*rustaz} -\textgreater{} \emph{rost} & expected regular output
from the source-supported citation reconstruction & shows that
correcting the stem class does not solve the vowel problem \\
\emph{*rústu} -\textgreater{} \emph{rust} & regular high-vowel
comparator & useful negative control, but not a defensible citation-form
solution \\
\emph{rust} & attested Old English noun, normalized from \emph{rūst} /
\emph{rúst} / \emph{rust} & selected target; the citation-form
development remains unexplained \\
\end{longtable}
}

\subsubsection{wolf --- OE wulf}\label{wolf-oe-wulf}

Derivation: \emph{*wúlfaz} yields regular \emph{wolf}; the selected
target is \emph{wulf} (unexplained exception).

\paragraph{Derivation trace}\label{derivation-trace-75}

Proto input: \emph{*wúlfaz}

\begingroup
\setlength{\fboxsep}{6pt}

\noindent\fbox{%
\begin{minipage}{0.97\linewidth}
\small
\begin{tabularx}{\linewidth}{@{}>{\raggedright\arraybackslash}p{0.540\linewidth}>{\centering\arraybackslash}X>{\raggedright\arraybackslash}p{0.300\linewidth}@{}}
\begin{minipage}[t]{\linewidth}
\centering\textbf{Earlier Germanic changes}\par
\vspace{0.35em}
\raggedright
\centering\textbf{West Germanic}\par
\raggedright
\vspace{0.2em}
\raggedright [no change]\par
\vspace{0.6em}
\centering\textbf{Northwest Germanic}\par
\raggedright
\vspace{0.2em}
\begin{tabular}{@{}>{\raggedright\arraybackslash}p{0.74\linewidth}@{\hspace{0.55em}}>{\raggedright\arraybackslash}p{0.16\linewidth}@{\hspace{0.25em}}}
\mbox{NWGmc U Lowering} & \emph{*wólfaz} \\
\mbox{PGmc Final Z Deletion} & \emph{*wólfa} \\
\end{tabular}
\end{minipage}
&
&
\begin{minipage}[t]{\linewidth}
\centering\textbf{Old English changes}\par
\vspace{0.35em}
\raggedright
\begin{tabular}{@{}>{\raggedright\arraybackslash}p{0.70\linewidth}@{\hspace{0.55em}}>{\raggedright\arraybackslash}p{0.16\linewidth}@{\hspace{0.25em}}}
PWGmc Final Bare A Loss & \emph{*wólf} \\
\end{tabular}
\end{minipage}
\\
\end{tabularx}
\end{minipage}%
} \endgroup

Transducer outcome: \emph{wolf}

Selected target: \emph{wulf}

\paragraph{Reconstruction and comparative
evidence}\label{reconstruction-and-comparative-evidence-75}

The inherited noun is an a-stem: Kroonen gives \emph{*wulfa-}, and the
regular Old English development of that citation form is the same one
reflected by Old High German \emph{wolf}
(\citeproc{ref-Kroonen2013}{Kroonen 2013};
\citeproc{ref-RingeTaylor2014}{Ringe and Taylor 2014}). Campbell
accordingly names Old English \emph{wulf} as an exception to the regular
lowering of stressed \emph{u} before a following non-high vowel
(\citeproc{ref-Campbell1959}{Campbell 1959}).

The older literature often notices that the exceptional words cluster
near labials. Bülbring lists \emph{full}, \emph{wulle}, and \emph{wulf}
together, but he also says that the ordinary rule still gives \emph{o}
in comparable forms such as \emph{folc} and \emph{bolt}
(\citeproc{ref-Bulbring1902}{Bülbring 1902}). Luick rejects a
categorical labial blocker on the same grounds and prefers a lexical or
analogical account instead (\citeproc{ref-Luick1914}{Luick 1914-\/-1940,
148}).

\paragraph{Old English evidence}\label{old-english-evidence-75}

Old English \emph{wulf} is an attested noun, and the handbook tradition
treats it as the ordinary lexical form while simultaneously recognizing
its exceptional vowel (\citeproc{ref-Campbell1959}{Campbell 1959};
\citeproc{ref-SieversBrunner1965}{Sievers 1965}). The oblique form
\emph{wulfe} also belongs to the record, but Sievers-Brunner warns that
this type continues \emph{wulfi} only with i-umlaut later abandoned
across the paradigm (\citeproc{ref-SieversBrunner1965}{Sievers 1965}).

That warning matters because the surviving oblique forms do not supply a
clean regular route back to bare \emph{wulf}. They belong to the same
lexeme, but they do not remove the explanatory problem presented by the
citation form.

\paragraph{Development to Old
English}\label{development-to-old-english-75}

Under the regular lowering of stressed \emph{u}, the citation-form input
gives \emph{wolf}, not \emph{wulf} (\citeproc{ref-Campbell1959}{Campbell
1959}; \citeproc{ref-RingeTaylor2014}{Ringe and Taylor 2014}). The
compact trace shows exactly that path: \emph{*wúlfaz} \textgreater{}
\emph{*wólfaz} \textgreater{} \emph{*wólfa} \textgreater{} \emph{wolf}.

A high-vowel oblique input would behave differently. The following high
vowel would block the lowering of \emph{u}, but the same environment
would trigger i-umlaut, so the regular result would be \emph{wylf} or
\emph{wylfe}, not bare \emph{wulf}
(\citeproc{ref-SieversBrunner1965}{Sievers 1965}). The attested noun
therefore remains unexplained at the citation-form level.

\paragraph{Expected and attested
forms}\label{expected-and-attested-forms-2}

The comparison below is manual. It separates the regular inherited
outcomes from the attested Old English noun.

{\def\LTcaptype{none} % do not increment counter
\begin{longtable}[]{@{}
  >{\raggedright\arraybackslash}p{(\linewidth - 4\tabcolsep) * \real{0.3333}}
  >{\raggedright\arraybackslash}p{(\linewidth - 4\tabcolsep) * \real{0.3333}}
  >{\raggedright\arraybackslash}p{(\linewidth - 4\tabcolsep) * \real{0.3333}}@{}}
\toprule\noalign{}
\begin{minipage}[b]{\linewidth}\raggedright
Form / interpretation
\end{minipage} & \begin{minipage}[b]{\linewidth}\raggedright
Status
\end{minipage} & \begin{minipage}[b]{\linewidth}\raggedright
Relevance to this entry
\end{minipage} \\
\midrule\noalign{}
\endhead
\bottomrule\noalign{}
\endlastfoot
\emph{*wúlfaz} -\textgreater{} \emph{wolf} & computed regular output
from the citation form & shows the regular development expected from the
inherited a-stem \\
\emph{OHG wolf} & comparative regular cognate & confirms that the
\emph{o}-vocalism is the ordinary outcome \\
\emph{*wúlfi} / \emph{*wúlfis} -\textgreater{} \emph{wylf} /
\emph{wylfe} & expected high-vowel control forms & shows why oblique
high-vowel cells do not solve the noun's vowel history \\
\emph{wulf} & attested Old English noun & selected target; the
preservation of \emph{u} remains unexplained \\
\end{longtable}
}

\subsubsection{wool --- OE wull}\label{wool-oe-wull}

Derivation: \emph{*wúllō} yields regular \emph{woll}; the selected
target is \emph{wull} (unexplained exception).

\paragraph{Derivation trace}\label{derivation-trace-76}

Proto input: \emph{*wúllō}

\begingroup
\setlength{\fboxsep}{6pt}

\noindent\fbox{%
\begin{minipage}{0.97\linewidth}
\small
\begin{tabularx}{\linewidth}{@{}>{\raggedright\arraybackslash}p{0.460\linewidth}>{\centering\arraybackslash}X>{\raggedright\arraybackslash}p{0.460\linewidth}@{}}
\begin{minipage}[t]{\linewidth}
\centering\textbf{Earlier Germanic changes}\par
\vspace{0.35em}
\raggedright
\centering\textbf{West Germanic}\par
\raggedright
\vspace{0.2em}
\raggedright [no change]\par
\vspace{0.6em}
\centering\textbf{Northwest Germanic}\par
\raggedright
\vspace{0.2em}
\begin{tabular}{@{}>{\raggedright\arraybackslash}p{0.74\linewidth}@{\hspace{0.55em}}>{\raggedright\arraybackslash}p{0.16\linewidth}@{\hspace{0.25em}}}
\mbox{NWGmc U Lowering} & \emph{*wóllō} \\
\mbox{NWGmc Final Long O Raising} & \emph{*wóllu} \\
\end{tabular}
\end{minipage}
&
&
\begin{minipage}[t]{\linewidth}
\centering\textbf{Old English changes}\par
\vspace{0.35em}
\raggedright
\begin{tabular}{@{}>{\raggedright\arraybackslash}p{0.74\linewidth}@{\hspace{0.55em}}>{\raggedright\arraybackslash}p{0.16\linewidth}@{\hspace{0.25em}}}
\mbox{OE High Vowel Apocope} & \emph{*wóll} \\
\end{tabular}
\end{minipage}
\\
\end{tabularx}
\end{minipage}%
} \endgroup

Transducer outcome: \emph{woll}

Selected target: \emph{wull}

\paragraph{Reconstruction and comparative
evidence}\label{reconstruction-and-comparative-evidence-76}

The inherited form is a feminine ō-stem \emph{*wúllō}. In the ordinary
phonological history of West Germanic, stressed \emph{u} lowers before a
following non-high vowel, so the regular Old English outcome is an
\emph{o}-form. Campbell's discussion of the parallel adjective
\emph{full}, with OHG \emph{foll} as the regular comparator, shows that
the handbooks treat this as a genuine exception cluster rather than as a
place where the rule itself is doubtful
(\citeproc{ref-Campbell1959}{Campbell 1959, sect. 115}).

Bülbring likewise lists \emph{wulle} among the traditional
\emph{u}-preserving exceptions (\citeproc{ref-Bulbring1902}{Bülbring
1902, sect. 116}). The comparative evidence therefore establishes two
things at once: the regular result should be \emph{woll}, and Old
English still has a lexical exception of the \emph{wull} / \emph{wulle}
type.

\paragraph{Old English evidence}\label{old-english-evidence-76}

The Old English target is given here as \emph{wull}, a normalized lexeme
form. Handbook discussion often cites \emph{wulle}, the feminine weak
form of the noun (\citeproc{ref-Bulbring1902}{Bülbring 1902, sect.
116}). Both point to the same lexical item and to the same exceptional
preservation of root \emph{u}.

The OE evidence therefore does not remove the problem. It confirms that
the language has a \emph{u}-form where the regular phonology would have
produced \emph{o}.

\paragraph{Development to Old
English}\label{development-to-old-english-76}

From \emph{*wúllō}, the regular sequence is lowering of stressed
\emph{u} before a non-high vowel, followed by the ordinary later
reductions of the ending. The regular outcome is therefore \emph{woll}.

That regular derivation is not the attested Old English form. Luick
rejects a simple phonological rule that would protect this word alone,
and Ringe and Taylor state the larger problem plainly: we do not really
know why \emph{*u} failed to lower in forms of this sort
(\citeproc{ref-Luick1914}{Luick 1914-\/-1940, 148};
\citeproc{ref-RingeTaylor2014}{Ringe and Taylor 2014, sect. 2.3.1}).

\paragraph{What remains unexplained}\label{what-remains-unexplained}

The comparison below is manual. It separates the regular trace result
from the attested lexical exception.

{\def\LTcaptype{none} % do not increment counter
\begin{longtable}[]{@{}
  >{\raggedright\arraybackslash}p{(\linewidth - 4\tabcolsep) * \real{0.3333}}
  >{\raggedright\arraybackslash}p{(\linewidth - 4\tabcolsep) * \real{0.3333}}
  >{\raggedright\arraybackslash}p{(\linewidth - 4\tabcolsep) * \real{0.3333}}@{}}
\toprule\noalign{}
\begin{minipage}[b]{\linewidth}\raggedright
Form
\end{minipage} & \begin{minipage}[b]{\linewidth}\raggedright
Status
\end{minipage} & \begin{minipage}[b]{\linewidth}\raggedright
Relevance to this entry
\end{minipage} \\
\midrule\noalign{}
\endhead
\bottomrule\noalign{}
\endlastfoot
\emph{*wúllō} -\textgreater{} \emph{woll} & regular trace output & shows
what the deterministic sound laws produce \\
\emph{wull} / \emph{wulle} & attested OE exception & selected target and
philological fact to be recorded \\
high-vowel escape from another paradigm cell & unsupported for this noun
& rejected because the feminine ō-stem paradigm supplies no suitable
escape cell \\
\end{longtable}
}

\clearpage

\subsection*{References}\label{references}
\addcontentsline{toc}{subsection}{References}

\protect\phantomsection\label{refs}
\begin{CSLReferences}{1}{1}
\bibitem[\citeproctext]{ref-Adamczyk2001}
Adamczyk, Elżbieta. 2001. {`Old {English} Reflexes of {Sievers'} Law'}.
\emph{Studia Anglica Posnaniensia: An International Review of English
Studies} 36: 61--72. \url{http://hdl.handle.net/10593/18335}.

\bibitem[\citeproctext]{ref-Bammesberger1997}
Bammesberger, Alfred. 1997. {`Die Vorform von Altenglisch {h{æ}rfest}'}.
\emph{Anglia: Zeitschrift f{ü}r Englische Philologie} 115 (2): 223--30.
\url{https://doi.org/10.1515/angl.1997.115.2.223}.

\bibitem[\citeproctext]{ref-BosworthToller1898}
Bosworth, Joseph, and T. Northcote Toller. 1898. \emph{An {Anglo-Saxon}
Dictionary}. Clarendon Press.

\bibitem[\citeproctext]{ref-BrightCassidyRingler1971}
Bright, James W. 1971. \emph{Bright's {Old English} Grammar and Reader}.
3rd edn. Edited by Frederic G. Cassidy and Richard N. Ringler. Holt,
Rinehart; Winston.

\bibitem[\citeproctext]{ref-Bulbring1902}
Bülbring, Karl D. 1902. \emph{Altenglisches Elementarbuch}. Winter.

\bibitem[\citeproctext]{ref-Campbell1959}
Campbell, Alistair. 1959. \emph{Old {English} Grammar}. Clarendon Press.

\bibitem[\citeproctext]{ref-ClarkHall1960}
Clark Hall, J. R. 1960. \emph{A Concise {Anglo-Saxon} Dictionary}. 4th
edn. University of Toronto Press.

\bibitem[\citeproctext]{ref-Crist2002}
Crist, Sean J. 2002. {`An Analysis of *{z} Loss in {West Germanic}'}.
Unpublished manuscript.

\bibitem[\citeproctext]{ref-Fulk2018}
Fulk, R. D. 2018. \emph{A Comparative Grammar of the Early {Germanic}
Languages}. Vol. 3. Studies in Germanic Linguistics. John Benjamins.

\bibitem[\citeproctext]{ref-Hogg1992}
Hogg, Richard M. 1992. \emph{A Grammar of Old English. Volume 1:
Phonology}. Blackwell.

\bibitem[\citeproctext]{ref-Kilday2024}
Kilday, Douglas G. 2024. {`{Crist's} Law, {Smith's} Law, and {English}
\emph{Wizen}'}. Unpublished manuscript.

\bibitem[\citeproctext]{ref-KlugeSeebold2011}
Kluge, Friedrich. 2011. \emph{Etymologisches w{ö}rterbuch Der Deutschen
Sprache}. 25th edn. Edited by Elmar Seebold. De Gruyter.

\bibitem[\citeproctext]{ref-Kroonen2011}
Kroonen, Guus. 2011. \emph{The {Proto-Germanic} n-Stems: A Study in
Diachronic Morphophonology}. Vol. 18. Leiden Studies in Indo-European.
Rodopi.

\bibitem[\citeproctext]{ref-Kroonen2013}
Kroonen, Guus. 2013. \emph{Etymological Dictionary of {Proto-Germanic}}.
Vol. 11. Leiden Indo-European Etymological Dictionary Series. Brill.

\bibitem[\citeproctext]{ref-Lloyd1966}
Lloyd, Albert L. 1966. {`Is There an {a-}Umlaut of *{i} in {Germanic}?'}
\emph{Language} 42 (4): 738--45. \url{https://doi.org/10.2307/411829}.

\bibitem[\citeproctext]{ref-Luick1914}
Luick, Karl. 1914-\/-1940. \emph{Historische Grammatik Der Englischen
Sprache}. Tauchnitz.

\bibitem[\citeproctext]{ref-Mayrhofer1992}
Mayrhofer, Manfred. 1992-\/-2001. \emph{Etymologisches w{ö}rterbuch Des
Altindoarischen}. Carl Winter.

\bibitem[\citeproctext]{ref-Orel2003}
Orel, Vladimir. 2003. \emph{A Handbook of {Germanic} Etymology}. Brill.

\bibitem[\citeproctext]{ref-Ringe2006}
Ringe, Don. 2006. \emph{From {Proto-Indo-European} to {Proto-Germanic}}.
1st edn. Vol. 1. A Linguistic History of English. Oxford University
Press.

\bibitem[\citeproctext]{ref-RingeTaylor2014}
Ringe, Don, and Ann Taylor. 2014. \emph{The Development of {Old
English}}. Vol. 2. A Linguistic History of English. Oxford University
Press.

\bibitem[\citeproctext]{ref-Seebold1970}
Seebold, Elmar. 1970. \emph{Vergleichendes Und Etymologisches
w{ö}rterbuch Der Germanischen Starken Verben}. Mouton.

\bibitem[\citeproctext]{ref-SieversBrunner1965}
Sievers, Eduard. 1965. \emph{Altenglische Grammatik}. 3rd edn. Edited by
Karl Brunner. Niemeyer.

\bibitem[\citeproctext]{ref-Streitberg1896}
Streitberg, Wilhelm. 1896. \emph{Urgermanische Grammatik}. Carl Winter.

\bibitem[\citeproctext]{ref-Sweet1953}
Sweet, Henry. 1953. \emph{An {Anglo-Saxon} Primer}. 9th edn. Clarendon
Press.

\bibitem[\citeproctext]{ref-GermanicSlavicBaltic2025}
Viredaz, Rémy. 2025. \emph{{Germanic}, {Slavic} and {Baltic}
{`Thousand'} Once More}. Unpublished manuscript, Geneva, 20 October
2025; full text (unfinished).
\url{https://www.academia.edu/144462167/2025_Germanic_Slavic_and_Baltic_thousand_once_more_full_text_unfinished_}.

\end{CSLReferences}

\end{document}
